%!TEX TS-program = xelatex
%!TEX TS-options = -shell-escape
%!TEX root = ../OpAlg2_SS16/operatoralgebren2.tex
\RequirePackage{fix-cm} 
\documentclass[a4paper, twoside, headsepline, index=totoc,toc=listof,toc=bibliography,toc=index, fontsize=10pt, cleardoublepage=empty, headinclude, DIV=12, BCOR=5mm, titlepage,draft]{scrartcl}
%!TEX root = ../AnaTopGeo_SS14/ana_top_geo.tex
\usepackage{scrtime} % KOMA, Uhrzeit ermoeglicht

%--Pakete zum "Programmieren"
% ======================================================================================
\usepackage{etoolbox}
\usepackage{letltxmacro}
\usepackage{ifthen}
% ======================================================================================

%--Farbdefinitionen und Grafiken (muss vor tikz geladen werden)
% ======================================================================================
\usepackage[usenames, table, x11names]{xcolor}
\definecolor{dark_gray}{gray}{0.45}
\definecolor{light_gray}{gray}{0.6}
\definecolor{fb10_blue}{cmyk}{0.8,0.4,0.13,0.07}
\usepackage[final]{graphicx}
\usepackage{adjustbox}
\newcommand{\cfbox}[2]{% coloured frame box
	\ifmmode
	\mathchoice{\adjustbox{cfbox=#1}{$\displaystyle#2$}}{\adjustbox{cfbox=#1}{$\textstyle#2$}}{\adjustbox{cfbox=#1}{$\scriptstyle#2$}}{\adjustbox{cfbox=#1}{$\scriptscriptstyle#2$}}
	\else
	\adjustbox{cfbox=#1}{#2}
	\fi
}
% ======================================================================================

%--Zum Zeichnen/ TikZ-Kram (vor polyglossia bzw. babel geladen werden)
% ======================================================================================
\usepackage{tikz}
\usepackage{tikz-cd}
\usetikzlibrary{external}
\tikzset{>=latex}
\usetikzlibrary{%
	shapes,
	arrows.meta,
	intersections,
	calc,
	3d,
	decorations.pathreplacing,decorations.markings,decorations.pathmorphing,
	angles,
	quotes,
}
\tikzexternalize[prefix=tikz/,up to date check=diff]
\pgfkeys{/pgf/images/include external/.code=\includegraphics{#1}}
\tikzset{external/system call={lualatex \tikzexternalcheckshellescape -halt-on-error -interaction=batchmode --shell-escape -jobname "\image" "\texsource"}}
\AtBeginEnvironment{tikzcd}{\tikzexternaldisable} % tikzexternalize fuer tikzcd deaktivieren, da inkompatibel
\AtEndEnvironment{tikzcd}{\tikzexternalenable}
\tikzset{% um Inkompatibilitaeten von quotes und polyglossia bzw. babel zu vermeiden
  every picture/.append style={
    execute at begin picture={\shorthandoff{"}},
    execute at end picture={\shorthandon{"}}
  }
}
\usepackage{pgfplots}
\usepgfplotslibrary{colormaps}
\newcommand*\circled[1]{\tikzexternaldisable\tikz[baseline=(char.base)]{\node[shape=circle,draw,inner sep=2pt] (char) {#1};}\tikzexternalenable}
% ======================================================================================



%-- Mathepakete etc.
% ======================================================================================
\usepackage[T1]{fontenc}
\renewcommand{\rmdefault}{zpltlf}
\usepackage{mathtools} % beinhaltet amsmath
\mathtoolsset{showonlyrefs,centercolon,showmanualtags}
\newtagform{brackets}[\textbf]{[}{]}
\usetagform{brackets}
\usepackage{fix-cm}
\usepackage[bbgreekl]{mathbbol}
\usepackage{amssymb,marvosym} 
\usepackage{nicefrac} % schräge Brüche
\usepackage{faktor}
\newcommand{\Faktor}[1]{\faktor[\textstyle]{#1}}
\usepackage{xfrac}
\usepackage{cancel}
\usepackage{mathdots} % Verbesserung von Punkten wie zB \ldots
\usepackage[bb=px]{mathalfa} % \mathbb als px font
\usepackage{centernot}
\usepackage{stackrel}
\DeclareSymbolFont{bbold}{U}{bbold}{m}{n}
\DeclareSymbolFontAlphabet{\mathbbold}{bbold}
\newcommand{\ind}{\mathbbold{1}} % charakteristische-Funktion-Eins
\def\mathul#1#2{\color{#1}\underline{{\color{black}#2}}\color{black}} %farbiges Untersteichen im Mathe-Modus
\renewcommand{\le}{\leqslant}
\renewcommand{\ge}{\geqslant}
% ======================================================================================


%-- Von xfrac erzeuge font warnings ignorieren
% ======================================================================================
\usepackage{silence}
\WarningFilter{latexfont}{Size substitutions with differences}
\WarningFilter{latexfont}{Font shape `U/bbold/m/n' in size}
% ======================================================================================


%-- Typographie/Polyglossia
% ======================================================================================
\usepackage[euler-digits]{eulervm} % vor fontspec laden!
\usepackage[no-math]{fontspec}
\usepackage{polyglossia} % moderner babel-ersatz
\setmainlanguage[spelling=new,babelshorthands=true]{german}
\shorthandoff{"}
\setotherlanguage{english}
\defaultfontfeatures{Mapping=tex-text, WordSpace={1.2}, Ligatures={Required,Common,Contextual},Extension=.otf} %


\setmainfont{TeXGyrePagellaX}[UprightFont=*-Regular,BoldFont=*-Bold,ItalicFont=*-Italic,BoldItalicFont=*-BoldItalic,ItalicFeatures={Style=Historic},Ligatures={Required,Common,Contextual,Historic}]
\setsansfont{texgyreadventor}[Scale=MatchUppercase, UprightFont=*-regular, BoldFont=*-bold, ItalicFont=*-italic, BoldItalicFont=*-bolditalic]
\setmonofont{SourceCodePro}[Scale=0.9,UprightFont=*-Regular, BoldFont=*-Semibold, ItalicFont=*-Light]
\usepackage{xltxtra}
\usepackage{fontawesome}
\usepackage[final]{microtype}
\usepackage[draft=false]{scrlayer-scrpage} 
\flushbottom
% ======================================================================================


%-- Aufzählungen
% ======================================================================================
\usepackage[shortlabels,inline]{enumitem}
\setlist[itemize,1]{label=\faCaretRight}
\setlist[enumerate]{font=\bfseries}
\setlist[description]{font=\normalfont\bfseries}
\usepackage{multicol}
% ======================================================================================


%-- Floats/Figures/Tabellen
% ======================================================================================
\usepackage{wrapfig}
\usepackage{float}
\usepackage[margin=10pt, font=small, labelfont={sf, bf}, format=plain, indention=1em]{caption}
\captionsetup[wrapfigure]{name=Abb. }
\usepackage{booktabs}
% ======================================================================================


%-- korrekte Anführungszeichen und Zitierbefehle
% ======================================================================================
\usepackage[autostyle,german=quotes,english=british]{csquotes}
% ======================================================================================


%--Indexverarbeitung
% ======================================================================================
\usepackage{makeidx}
\newcommand{\bet}[1]{\textbf{\emph{#1}}}
\newcommand{\Index}[1]{\bet{#1}\index{#1}}
\makeindex
\setindexpreamble{{\noindent\sffamily\small Die \emph{Seitenzahlen} sind mit Hyperlinks versehen und somit anklickbar} \par \bigskip}
\renewcommand{\indexpagestyle}{scrheadings}
% ======================================================================================


%-- Marginnotes/Todonotes/Footnotes
% ======================================================================================
\deffootnote[1.5em]{1.5em}{1.5em}{\textsuperscript{\thefootnotemark}\ }
\usepackage[fulladjust]{marginnote}
\renewcommand*{\marginfont}{\itshape\footnotesize}
\usepackage[textsize=small]{todonotes}
\usepackage{ragged2e}
\renewcommand*{\raggedleftmarginnote}{\RaggedLeft}
\renewcommand*{\raggedrightmarginnote}{\RaggedRight}
\LetLtxMacro{\oldtodo}{\todo}
\renewcommand{\todo}[2][]{\tikzexternaldisable\oldtodo[#1]{#2}\tikzexternalenable}
\LetLtxMacro{\oldmissingfigure}{\missingfigure}
\renewcommand{\missingfigure}[2][]{\tikzexternaldisable\oldmissingfigure[{#1}]{#2}\tikzexternalenable}
% ======================================================================================


% -- BibLaTeX
% ======================================================================================
\usepackage[%
	backend=biber,
	sortlocale=auto,
	natbib,
	hyperref,
	backref,
	style=alphabetic
	]%
{biblatex}
\renewcommand*{\mkbibnamelast}[1]{%
  \ifmknamesc{\textsc{#1}}{#1}}
\renewcommand*{\mkbibnameprefix}[1]{%
  \ifboolexpr{ test {\ifmknamesc} and test {\ifuseprefix} }
    {\textsc{#1}}
    {#1}}
\def\ifmknamesc{%
  \ifboolexpr{ test {\ifcurrentname{labelname}}
               or test {\ifcurrentname{author}}
               or ( test {\ifnameundef{author}} and test {\ifcurrentname{editor}} ) }}
\addbibresource{../!config/quellen.bib}
% ======================================================================================

%--Konfiguration von Hyperref und Cleveref
% ======================================================================================
\usepackage[hidelinks, pdfpagelabels,  bookmarksopen=true, bookmarksnumbered=true, linkcolor=black, urlcolor=SkyBlue2, plainpages=false,pagebackref, citecolor=black, hypertexnames=true, pdfauthor={Jannes Bantje}, pdfborderstyle={/S/U}, linkbordercolor=SkyBlue2, colorlinks=false,final,backref=false]{hyperref}
\usepackage[nameinlink,noabbrev]{cleveref}
\newcommand{\appendLink}[1]{#1\,\faExternalLink}
\newcommand{\hrefsym}[2]{\href{#1}{\texttt{\appendLink{#2}}}}
\newcommand{\hrefsymX}[2]{\href{#1}{\appendLink{#2}}}
\newcommand{\hrefsymmail}[2]{\href{#1}{\texttt{\faEnvelopeO\,#2}}}
\renewcommand{\url}[1]{\hrefsym{#1}{\nolinkurl{#1}}}
% ======================================================================================


% -- QR-Codes (hinter hyperref laden!)
% ======================================================================================
\usepackage{qrcode}
% ======================================================================================

%--Römische Zahlen
% ======================================================================================
\newcommand{\RM}[1]{\MakeUppercase{\romannumeral #1{}}}
% ======================================================================================

%-- Definition von diversen Mathe-Befehlen
% ======================================================================================
%!TEX root = mitschrift_main.tex

% -- Zum Finetuning von Befehlen
% ======================================================================================
\makeatletter
\newcommand{\raisemath}[1]{\mathpalette{\raisem@th{#1}}}
\newcommand{\raisem@th}[3]{\raisebox{#1}{$#2#3$}}
\makeatother
\makeatletter
\newcommand{\killDescendersM}[1]{\mathpalette{\killD@scendersM{#1}}}
\newcommand{\killD@scendersM}[2]{\raisebox{0pt}[\height][0pt]{$#2#1$}}
\makeatother
\DeclareRobustCommand{\minwidthbox}[2]{%
  \ifmmode
    \expandafter\mathmakebox
  \else
    \expandafter\makebox
  \fi
  [\ifdim#2<\width\width\else#2\fi]{#1}%
}
% ======================================================================================


%-- Klammerbefehle
% ======================================================================================
\DeclarePairedDelimiter{\abs}{\lvert}{\rvert}
\DeclarePairedDelimiter{\floor}{\lfloor}{\rfloor}
\DeclarePairedDelimiter{\ceil}{\lceil}{\rceil}
\DeclarePairedDelimiter\norm{\Vert}{\Vert}
\DeclarePairedDelimiter\enbrace{(}{)}
\DeclarePairedDelimiter\benbrace{[}{]}
\DeclarePairedDelimiter\bbenbrace{[\![}{]\!]}
\DeclarePairedDelimiter\lenbrace{<}{>}
\DeclarePairedDelimiter\angbrace{\langle}{\rangle}
\newcommand{\ssbrace}[1]{{\scriptscriptstyle\enbrace{#1}}}
\newcommand{\ssbbrace}[1]{{\scriptscriptstyle\benbrace{#1}}}
% ======================================================================================

%-- Mengen
% ======================================================================================
\newcommand\SetSymbol[1][]{\nonscript\:#1\vert\allowbreak\nonscript\:\mathopen{}}
\providecommand\given{} % to make it exist
\DeclarePairedDelimiterX\set[1]\{\}{\renewcommand\given{\SetSymbol[\delimsize]}#1}
% ======================================================================================

%-- Skalarprodukt (3 Varianten) 
% ======================================================================================
\DeclarePairedDelimiterX\sprod[2]{\langle}{\rangle}{#1\,\delimsize\vert\,#2}
\DeclarePairedDelimiterX\skal[2]{\langle}{\rangle}{#1\,,\,#2}
\makeatletter
\DeclareFontFamily{OMX}{MnSymbolE}{}
\DeclareSymbolFont{MnLargeSymbols}{OMX}{MnSymbolE}{m}{n}
\SetSymbolFont{MnLargeSymbols}{bold}{OMX}{MnSymbolE}{b}{n}
\DeclareFontShape{OMX}{MnSymbolE}{m}{n}{
    <-6>  MnSymbolE5
   <6-7>  MnSymbolE6
   <7-8>  MnSymbolE7
   <8-9>  MnSymbolE8
   <9-10> MnSymbolE9
  <10-12> MnSymbolE10
  <12->   MnSymbolE12
}{}
\DeclareFontShape{OMX}{MnSymbolE}{b}{n}{
    <-6>  MnSymbolE-Bold5
   <6-7>  MnSymbolE-Bold6
   <7-8>  MnSymbolE-Bold7
   <8-9>  MnSymbolE-Bold8
   <9-10> MnSymbolE-Bold9
  <10-12> MnSymbolE-Bold10
  <12->   MnSymbolE-Bold12
}{}
\let\llangle\@undefined
\let\rrangle\@undefined
\DeclareMathDelimiter{\llangle}{\mathopen}%
                     {MnLargeSymbols}{'164}{MnLargeSymbols}{'164}
\DeclareMathDelimiter{\rrangle}{\mathclose}%
                     {MnLargeSymbols}{'171}{MnLargeSymbols}{'171}
\makeatother
\DeclarePairedDelimiterX\sskal[2]{\llangle}{\rrangle}{#1\,,\,#2}
% ======================================================================================

%-- Abbildungsdefinition
% ======================================================================================
\newcommand{\mapdef}[5]{%
	\[
		\begin{array}{rcl}
			\textstyle #1 &\xrightarrow{\minwidthbox{#5}{2em}} & \textstyle #2 \\[0.5ex]
			\textstyle #3 &\xmapsto{\minwidthbox{\mbox{ }}{2em}} & \textstyle #4
		\end{array}
	\]
}
% ======================================================================================

%-- modifiziertes Stackrel 
% ======================================================================================
\newcommand{\StackText}[2]{\stackrel{\mbox{\scriptsize #1}}{#2}}
\newcommand{\StackTextClap}[2]{\stackrel{\mathclap{\mbox{\scriptsize #1}}}{#2}}
% ======================================================================================

%-- Blitz
% ======================================================================================
\newcommand{\light}{\text{\raisebox{-.3ex}{\Large\Lightning}}}
% ======================================================================================


%-- Underbrace u.Ä. als Befehl in LaTeX-Syntax (und ohne Spacingprobleme mit nachfolgenden Operatoren...)
% ======================================================================================
\newcommand{\Underbrace}[2]{{\underbrace{#1}_{#2}}}
\newcommand{\Underbracket}[2]{{\underbracket[0.7pt][2pt]{#1}_{#2}}}
\newcommand{\Overbracket}[2]{{\overbracket[0.7pt][2pt]{#1}^{#2}}}
% ======================================================================================


%-- Deklaration weiterer Operatoren (allgemein)
% ======================================================================================
\DeclareMathOperator{\re}{Re} % Realteil
\let\Re\relax
\DeclareMathOperator{\Re}{Re} % Realteil
\DeclareMathOperator{\im}{im} % Bild
\let\Im\relax
\DeclareMathOperator{\Im}{Im} % Bild
\DeclareMathOperator{\id}{id} % identische Abbildung
\DeclareMathOperator{\conj}{conj} % Konjugation
\DeclareMathOperator{\sgn}{sgn} % Signum
\DeclareMathOperator{\End}{End} % Endomorphismen
\DeclareMathOperator{\Hom}{Hom} % Homomorphismen
\DeclareMathOperator{\Iso}{Iso} % Isomorphismen
\DeclareMathOperator{\Aut}{Aut} % Automorphismen
\DeclareMathOperator{\Span}{span} % Span
\DeclareMathOperator{\coker}{coker} % Kokern
\DeclareMathOperator{\Tr}{Tr} % Spur,Trace
\DeclareMathOperator{\pr}{pr} % Projektion
\DeclareMathOperator{\diag}{diag} % Diagonalmatrix
\DeclareMathOperator{\Rg}{Rg} % Rang
\DeclareMathOperator{\const}{const} % konstante Abbildung
\DeclareMathOperator{\Spur}{Spur} % Spur
\DeclareMathOperator{\Arg}{Arg} % Argument
\DeclareMathOperator{\dist}{dist} % Distanz
\DeclareMathOperator{\supp}{supp} % Träger
\DeclareMathOperator{\Char}{char} % Charakteristik
% ======================================================================================


%-- Deklaration weiterer Operatoren (Differentiale etc.)
% ======================================================================================
\DeclareMathOperator{\grad}{grad} % Gradient
\DeclareMathOperator{\dive}{div} % Gradient
\DeclareMathOperator{\rot}{rot} % Rotation
\newcommand{\D}{\ensuremath{\mathrm{D}\mkern-1.0mu}} % Differential
\newcommand{\mathd}{\ensuremath{\mathrm{d}\mkern-1.0mu}} % äußere Ableitung
\newcommand{\Tmap}{\ensuremath{\mathrm{T}\mkern-0.85mu}} % Tangentialraum
\let\Tang\Tmap
\DeclareMathOperator{\Diff}{Diff}
\newcommand{\diff}[2]{\ensuremath{\frac{{\partial #1}}{{\partial #2}} }}
\newcommand{\diffd}[2]{\ensuremath{\frac{\mathd #1}{\mathd #2} }}
% ======================================================================================


%-- Deklaration weiterer Operatoren (Topologie)
% ======================================================================================
\newcommand*\interior[1]{\overset{\smash{\raisebox{-0.18ex}{$\scriptstyle\circ$}}}{#1}}
\newcommand{\sing}{{\raisemath{1.1pt}{\scriptscriptstyle\mathrm{sing}}}}
\newcommand{\pt}{\mathrm{pt}}
\DeclareMathOperator{\Zyl}{Zyl}
\DeclareMathOperator{\Tel}{Tel}
\newcommand{\op}{\mathrm{op}}
\DeclareMathOperator{\Sp}{Sp}
\DeclareMathOperator{\Keg}{Keg}
\newcommand{\slashedi}{i\hspace{-3.5pt}/}
\newcommand{\cupp}{\smallsmile}
\newcommand{\capp}{\smallfrown}
\DeclareMathOperator*{\colim}{colim}
\DeclareMathOperator{\PD}{PD}
\newcommand{\lf}{\mathrm{lf}}
\DeclareMathOperator{\sig}{sig}
\DeclareMathOperator{\Tor}{Tor}
\DeclareMathOperator{\Ext}{Ext}
\DeclareMathOperator{\AW}{AW}
\DeclareMathOperator{\Proj}{Proj}
\DeclareMathOperator{\Gr}{Gr}
\DeclareMathOperator{\res}{res}
\DeclareMathOperator{\Spec}{Spec}
\DeclareMathOperator{\co}{co}
\DeclareMathOperator{\ch}{ch}
\DeclareMathOperator{\wOp}{w}
\DeclareMathOperator{\Ar}{Ar}
\newcommand{\actson}{\mathrel{\curvearrowright}}
\let\acts\actson
\let\action\actson
\DeclareMathSymbol{\bbDelta}{\mathord}{bbold}{"01}
\newcommand{\DDelta}{\bbDelta}
\DeclareMathOperator{\Star}{Star}
\DeclareMathOperator{\Link}{Link}
\DeclareMathOperator{\EPK}{EPK}
\DeclareMathOperator{\Vol}{Vol}
\newcommand{\cell}{{\raisemath{1.1pt}{\scriptscriptstyle\mathrm{cell}}}}
\DeclarePairedDelimiter{\homologieklasse}{\llbracket}{\rrbracket}
\newcommand{\rand}[1]{\ensuremath{\partial^{\scriptscriptstyle #1}}}
\DeclareMathOperator{\ab}{ab}
% ======================================================================================


%-- Deklaration von Operatoren (Liegruppen)
% ======================================================================================
\DeclareMathOperator{\GL}{GL}
\DeclareMathOperator{\SO}{SO}
\DeclareMathOperator{\Ad}{Ad}
\DeclareMathOperator{\ad}{ad}
\DeclareMathOperator{\On}{O}
\DeclareMathOperator{\Un}{U}
\DeclareMathOperator{\SU}{SU}
\DeclareMathOperator{\Mat}{Mat}
\DeclareRobustCommand{\Der}{\mathop{\mathfrak{der}}}
\DeclareMathOperator{\SL}{SL}
\DeclareMathOperator{\Graph}{Graph}
\DeclareMathOperator{\Int}{Int}
\DeclareRobustCommand{\intAlg}{\mathop{\mathfrak{int}}}
\DeclareMathOperator{\aut}{aut}
\DeclareMathOperator{\Rad}{Rad}
\DeclareMathOperator{\Nil}{Nil}
\DeclareMathOperator{\rad}{rad}
\DeclareMathOperator{\nil}{nil}
\DeclareMathOperator{\Ric}{Ric}
\DeclareMathOperator{\ric}{ric}
\newcommand{\bi}{\mathrm{bi}}
\DeclareMathOperator{\Isom}{Isom}
\DeclareMathOperator{\Sym}{Sym}
\newcommand{\opL}{\ensuremath{\mathrm{L}\mkern-0.6mu}}
% ======================================================================================

%-- Deklaration von Operatoren (Funktionalanalysis)
% ======================================================================================
\DeclareMathOperator{\tr}{tr}
\newcommand{\w}{\mkern1mu\mathrm{w}}
\newcommand{\sa}{\mathrm{sa}}
\newcommand{\vb}{\mathrm{v\mkern-2.5mu.b\mkern-1.5mu.}} % vollständig beschränkt
\newcommand{\so}{\mathrm{\mkern.3mu s\mkern-1.4mu.\mkern-.6mu o\mkern-1.7mu.}} % \newcommand{\so}{\mathrm{s.o.}}
\newcommand{\solim}{\so\text{-}\mkern-0.8mu\lim}
\newcommand{\wo}{\mathrm{w\mkern-3mu.\mkern-.4mu o\mkern-1.7mu.}}
\newcommand{\Top}[1]{\mathcal{T}_{\mkern-2.3mu #1}}
\newcommand{\weakT}[1]{\ensuremath{\mathcal{T}_{#1}^{\mkern+1.0mu\text{\raisebox{0.4ex}{$\mathrm{w}$}}}}}
\newcommand{\weakTstar}[1]{\ensuremath{\mathcal{T}_{#1}^{\mkern+1.0mu\text{\raisebox{0.4ex}{$\mathrm{w}$}}^*}}}
\newcommand{\TWeakStar}{\Top{\w^*}}
\newcommand{\TWeakOp}{\Top{\wo}}
\newcommand{\Tso}{\Top{\so}}
\newcommand{\finSub}{\subset\mkern-0.7mu \subset}
\DeclareMathOperator{\Inv}{Inv}
\newcommand{\simm}{{\hspace{-1.6pt}\raisemath{0.5pt}{\sim}}}
\newcommand{\plus}{{\hspace{-1.6pt}+}}
\DeclareMathOperator{\ev}{ev}
\DeclareMathOperator{\Alg}{Alg}
\DeclareMathOperator{\her}{her}
\newcommand{\subher}{\subset_{\her}}
\newcommand{\grenzw}[1]{\xrightarrow{\minwidthbox{#1}{1.4em}}}
\newcommand{\grenzwl}[1]{\xleftarrow{\minwidthbox{#1}{1.4em}}}
\newcommand{\grenzwIn}[1]{\grenzw{\raisemath{-2pt}{#1}}}
\newcommand{\MyTo}[1]{\tikzexternaldisable\mathbin{\tikz[baseline] \draw[-to,line width=.4pt] (0ex,0.94ex) -- (#1,0.94ex);}\tikzexternalenable}
\newcommand{\dlim}{%
    \mathchoice
      {\lim\limits_{\MyTo{4.2ex}}}% \displaystyle
      {\lim\limits_{\MyTo{2.8ex}}}% \textstyle
      {\lim\limits_{\MyTo{2.3ex}}}% \scriptstyle
      {\lim\limits_{\MyTo{2.3ex}}}% \scriptscriptstyle
}
\newcommand{\Dlim}{\killDescendersM{\dlim}}
\DeclareMathOperator{\sep}{sep}
\DeclareMathOperator{\diam}{diam}
\DeclareMathOperator{\conv}{conv}
\DeclareMathOperator{\Prim}{Prim}
\DeclareMathOperator{\hull}{hull}
\DeclareMathOperator{\red}{red}
\DeclarePairedDelimiterX\bra[1]{\langle}{\rvert}{#1\,}
\DeclarePairedDelimiterX\ket[1]{\lvert}{\rangle}{\,#1}
\DeclarePairedDelimiterX\bracket[2]{\langle}{\rangle}{#1\,\delimsize\vert\,#2}
\newcommand{\tensormax}{\mathbin{\otimes_{\max}}}
\newcommand{\tensormin}{\mathbin{\otimes_{\min}}}
\DeclareMathOperator{\Ped}{Ped}
\newcommand{\alg}{\mathrm{alg}}
\DeclareMathOperator{\CPC}{CPC}
\DeclareMathOperator{\CP}{CP}
\DeclareMathOperator{\UPC}{UPC}
\newcommand{\DeltaOp}{\mathbin{\Delta}}
\newcommand{\kernedP}{\mathcal{P}\mkern-2mu}
\newcommand{\Pinfty}{\kernedP_{\infty}}
\DeclareMathOperator{\Groth}{Groth}
\DeclareMathOperator{\rk}{rk}
\newcommand{\MvN}{\mathrm{MvN}}
% ======================================================================================

%-- Kategorien
% ======================================================================================
\DeclareMathOperator{\Mor}{Mor}
\DeclareMathOperator{\Obj}{Obj}
\newcommand{\TOP}{\textsc{Top}}
\newcommand{\HTOP}{\textsc{HTop}}
\newcommand{\VR}{\textsc{VR}}
\newcommand{\MOD}{\textsc{Mod}}
\newcommand{\Mod}[1]{#1\text{-}\MOD}
\newcommand{\MONOIDE}{\textsc{Monoide}}
\newcommand{\SET}{\textsc{Set}}
\newcommand{\MAN}{\textsc{Man}}
\newcommand{\GRUPPEN}{\textsc{Gruppen}}
\newcommand{\ABELGRUPPEN}{\textsc{Abel.Gruppen}}
\newcommand{\ABEL}{\textsc{Abel}}
\newcommand{\KAT}{\textsc{Kat}}
\newcommand{\FUN}{\textsc{Fun}}
\newcommand{\SIMP}{\textsc{Simp}}
\newcommand{\VEKT}{\textsc{Vekt}}
\newcommand{\CH}{\textsc{Ch}}
\newcommand{\CSTARUN}{C^*\text{-}\textsc{Alg}^{\raisemath{-2.5pt}{1}}}
\newcommand{\CSTAR}{C^*\text{-}\textsc{Alg}}
\newcommand{\AB}{\textsc{Ab}}
% ======================================================================================
% ======================================================================================



% -- theorem packages
% ======================================================================================
\usepackage{amsthm}
\usepackage{thmtools,thm-restate}
\usepackage{mdframed}
\renewcommand{\listtheoremname}{Übersicht aller Aussagen}
\usepackage{bookmark}
\bookmarksetup{open,numbered}
\makeatletter
\newcommand*{\theorembookmark}{%
  \bookmark[
    dest=\@currentHref,
    rellevel=1,
    keeplevel,
  ]{%
    \thmt@thmname\space\csname the\thmt@envname\endcsname
    \ifx\thmt@shortoptarg\@empty
    \else
      \space(\thmt@shortoptarg)%
    \fi
  }%
}   
\makeatother
% ======================================================================================

% -- Definition der einzelnen Theorem-Umgebungen
% ======================================================================================
\declaretheoremstyle[%
	headfont=\sffamily\bfseries,
	notefont=\normalfont\sffamily\scshape,
	bodyfont=\normalfont,
	headformat=\NUMBER\ \NAME\NOTE,
	headpunct=.,
	postheadspace=1em,
	spaceabove=15pt,spacebelow=10pt,
	shaded={bgcolor=gray!20},
	postheadhook=\theorembookmark]%
{mainstyle}
\declaretheoremstyle[%
	headfont=\sffamily\bfseries,
	notefont=\normalfont\sffamily\scshape,
	bodyfont=\normalfont,
	headformat=\NUMBER\ \NAME\NOTE,
	headpunct=.,
	postheadspace=1em,
	spaceabove=15pt,spacebelow=10pt,
	shaded={bgcolor=fb10_blue!20},
	postheadhook=\theorembookmark]%
{mainstyle_blue}
\declaretheoremstyle[%
	headfont=\sffamily\bfseries,
	notefont=\normalfont\sffamily\scshape,
	bodyfont=\normalfont,
	headformat=\NUMBER\ \NAME\NOTE,
	headpunct=.,
	postheadspace=1em,
	spaceabove=15pt,spacebelow=10pt,
	postheadhook=\theorembookmark]%
{mainstyle_unshaded}
\declaretheoremstyle[%
	headfont=\sffamily\bfseries,
	notefont=\normalfont\sffamily\scshape,
	bodyfont=\normalfont,
	headformat=\NUMBER\NAME\NOTE,
	headpunct=.,
	postheadspace=1em,
	spaceabove=15pt,spacebelow=10pt,
	% shaded={bgcolor=gray!20},
	postheadhook=\theorembookmark]%
{mainstyle_unnumbered}
\declaretheoremstyle[%
	headfont=\sffamily\bfseries,
	notefont=\normalfont\sffamily\scshape,
	bodyfont=\normalfont,
	headformat=swapnumber,
	headpunct=.,
	postheadspace=1em,
	spaceabove=15pt,spacebelow=10pt,
	shaded={bgcolor=gray!20},
	postheadhook=\theorembookmark,
	qed=\qedsymbol]%
{mainstyleB}
\declaretheoremstyle[%
	headfont=\bfseries\scshape,
	bodyfont=\normalfont,
	headpunct=:,
	postheadspace=1em,
	spacebelow=12pt,spaceabove=2pt,
	qed=\qedsymbol]%
{beweise}
\declaretheoremstyle[%
	headfont=\bfseries\scshape,
	bodyfont=\normalfont,
	headpunct=:,
	postheadspace=1em,
	spacebelow=12pt,spaceabove=2pt]%
{beweisskizze}
\declaretheoremstyle[%
	headfont=\sffamily\bfseries,
	bodyfont=\normalfont,
	headpunct=:,
	postheadspace=1em,
	spacebelow=10pt,spaceabove=10pt]%
{bemerkungen}
\declaretheorem[name=Definition,parent=section,style=mainstyle_blue]{definition}
\declaretheorem[name=Definition \& Proposition,refname=Proposition,sharenumber=definition,style=mainstyle_blue]{definitionP}
\declaretheorem[name=Definition,numbered=no,style=mainstyle_unnumbered]{definition*}
\declaretheorem[name=Theorem,sharenumber=definition,style=mainstyle]{theorem}
\declaretheorem[name=Theorem,numbered=no,style=mainstyle_unnumbered]{theorem*}
\declaretheorem[name=Proposition,sharenumber=definition,style=mainstyle,refname=Proposition]{proposition}
\declaretheorem[name=Lemma,sharenumber=definition,style=mainstyle]{lemma}
\declaretheorem[name=Satz,sharenumber=definition,style=mainstyle,refname=Satz]{satz}
\declaretheorem[name=Satz,sharenumber=definition,style=mainstyle_unshaded]{satzUnshaded}
\declaretheorem[name=Definition,sharenumber=definition,style=mainstyle_unshaded]{definitionUnshaded}
\declaretheorem[name=Satz,numbered=no,style=mainstyle_unnumbered]{satz*}
\declaretheorem[name=Korollar,sharenumber=definition,style=mainstyle,refname=Korollar]{korollar}
\declaretheorem[name=Korollar,sharenumber=definition,style=mainstyleB,refname=Korollar]{korollarB}
\declaretheorem[name=Frage,numbered=no,style=mainstyle_unnumbered]{frage}
\declaretheorem[name=Frage,sharenumber=definition,style=mainstyle_unshaded]{frageA}
\declaretheorem[name=Erinnerung,sharenumber=definition,style=mainstyle_unshaded]{erinnerungA}
\declaretheorem[name=Ausblick,sharenumber=definition,style=mainstyle_unshaded]{ausblick}
\declaretheorem[name=Konvention,sharenumber=definition,style=mainstyle]{konvention}
\declaretheorem[name=Notation,sharenumber=definition,style=mainstyle_unshaded]{notation}
\declaretheorem[name=Bemerkung,sharenumber=definition,style=mainstyle_unshaded,refname=Bemerkung]{bemerkung}
\declaretheorem[name=Bemerkung,numbered=no,style=mainstyle_unnumbered]{bemerkung*}
\declaretheorem[name=Beispiel,sharenumber=definition,style=mainstyle_unshaded,refname=Beispiel]{beispiel}
\declaretheorem[name=Beispiel,numbered=no,style=mainstyle_unnumbered]{beispiel*}
\declaretheorem[name=Exkurs,numbered=no,style=mainstyle_unnumbered]{exkurs*}
\declaretheorem[name=Beweis,numbered=no,style=beweise]{beweis}
\declaretheorem[name=Übung,numbered=no,style=bemerkungen]{uebung}
\declaretheorem[name=Erinnerung,numbered=no,style=bemerkungen]{erinnerung}

% english versions
\declaretheorem[name=Remark,sharenumber=definition,style=mainstyle_unshaded]{remark}
\declaretheorem[name=Remark,numbered=no,style=mainstyle_unnumbered]{remark*}
\declaretheorem[name=Example,sharenumber=definition,style=mainstyle_unshaded]{example}
\declaretheorem[name=Corollary,sharenumber=definition,style=mainstyle]{corollary}
\let\proof\relax
\declaretheorem[name=Proof,numbered=no,style=beweise]{proof}
\declaretheorem[name=Sketch of Proof,numbered=no,style=beweisskizze]{sketch}
% ======================================================================================

%--Inhaltsverzeichnis
% ======================================================================================
\usepackage[tocindentauto]{tocstyle}
\usetocstyle{KOMAlike}
% ======================================================================================

%-- Dinge, die erst am Ende getan werden dürfen
% ======================================================================================
\shorthandon{"}
\usepackage{ellipsis}
% ======================================================================================


\newcommand{\fach}{Topologie \RM{1}.}
\newcommand{\semester}{WiSe 2014/2015}
\newcommand{\homepage}{https://wwwmath.uni-muenster.de/reine/u/topos/lehre/WS2014-2015/Topologie1/Topologie1.html}
\publishers{\scalebox{11}{\Huge$H_*$}}
\newcommand{\prof}{Prof.\ Dr.\ Arthur Bartels}
\providecommand{\verfasser}{Jannes Bantje}

%--Konfiguration von scrheadings
\setheadsepline{1pt}[\color{light_gray}]
\pagestyle{scrheadings}
%\fancyhf[H, F]{}
\clearscrheadfoot

\providecommand{\shortFach}{\fach}
\lehead{\includegraphics[height=0.6 cm,keepaspectratio]{../!config/Bilder/Logo_WWU_Muenster_light_gray.pdf}}
\rehead{\rule{0cm}{0.6cm}\footnotesize\sffamily\color{light_gray}\verfasser{} -- Mitschrift \shortFach}
\lohead{\rule{0cm}{0.6cm}\footnotesize\sffamily\color{light_gray}Stand: \today \; \thistime[:]}
\rohead{\includegraphics[height=0.6 cm,keepaspectratio]{../!config/Bilder/fb10logo_gray.pdf}}


\ofoot[{ \color{dark_gray} \LARGE \sffamily \thepage}]{{ \color{dark_gray} \LARGE \sffamily \thepage}} %hier wir auch der plain Stil bearbeitet!
\automark{section}
\ifoot{ \color{dark_gray} \small \leftmark}

%--Metadaten
\providecommand{\mail}{j.bantje@wwu.de}
\author{\verfasser}
\titlehead{\includegraphics[height=1.5cm, keepaspectratio]{../!config/Bilder/Logo_WWU_Muenster.pdf}%
\hfill \includegraphics[height=1.3cm, keepaspectratio]{../!config/Bilder/fb10logo.pdf}}
\title{Skript \fach}
\subtitle{Mitschrift der Vorlesung  \enquote{\fach} von \prof}

\KOMAoptions{DIV=last}


\begin{document}
\pagenumbering{Roman}
\maketitle
\begin{abstract}
\section*{Aktuelle Version verfügbar bei}
\newcommand{\dieBreite}{11cm}
\begin{minipage}{4cm}
	\qrcode[height=3.3cm, version=6]{https://gitlab.com/JaMeZ-B/LaTeX-WWU}
\end{minipage}
\hfill
\begin{minipage}{\dieBreite}
	% \includegraphics[height=0.6cm, keepaspectratio]{../!config/Bilder/wm_no_bg.pdf}
	\includegraphics[height=0.8cm, keepaspectratio]{../!config/Bilder/wm_no_bg.pdf}\\
	\url{https://gitlab.com/JaMeZ-B/LaTeX-WWU} \smallskip\\
	Das zentrale Repository des \enquote{\LaTeX-WWU}-Projekts befindet sich auf der Plattform GitLab.com.
	Neben der Koordination aller Beteiligten werden über diesen Dienst auch die PDFs gebaut, die in der Readme verlinkt sind.
\end{minipage}\\[1cm]
\begin{minipage}{4cm}
	\qrcode[height=3.3cm, version=6]{https://github.com/JaMeZ-B/latex-wwu}
\end{minipage}
\hfill
\begin{minipage}{\dieBreite}
	\includegraphics[height=0.6cm, keepaspectratio]{../!config/Bilder/github_octo.pdf}
	\includegraphics[height=0.6cm, keepaspectratio]{../!config/Bilder/GitHub_Logo.pdf}\\
	\url{https://github.com/JaMeZ-B/latex-wwu} \smallskip\\
	Die Entwicklung des \enquote{\LaTeX-WWU}-Projekts hat ursprünglich auf GitHub stattgefunden, ist mittlerweile aber zu GitLab gewechselt.
	Das GitHub-Repository wird stündlich automatisch aktualisiert, Merge-Requests werden aber nicht mehr entgegengenommen.
\end{minipage}\\[1cm]
% \begin{minipage}{4cm}
% 	\qrcode[height=3.3cm, version=6]{https://uni-muenster.sciebo.de/public.php?service=files&t=965ae79080a473eb5b6d927d7d8b0462}
% \end{minipage}
% \hfill
% \begin{minipage}{\dieBreite}
% 	\raisebox{-2pt}{\includegraphics[height=0.6cm, keepaspectratio]{../!config/Bilder/sciebo_logo.pdf}}
% 	\resizebox{!}{0.5cm}{\large \sffamily\textbf{sciebo}} {\sffamily\large die Campuscloud} \\
% 	\resizebox{\dieBreite}{!}{\footnotesize\url{https://uni-muenster.sciebo.de/public.php?service=files&t=965ae79080a473eb5b6d927d7d8b0462}}\smallskip\\
% 	Sciebo ist ein Dropbox-Ersatz der Hochschulen in NRW, der von der Uni Münster in leitender Position auf Basis der OpenSource-Software Owncloud aufgebaut wurde.
% \end{minipage}\\[1cm]
\hrule \mbox{ }\\[0.7cm]
\begin{minipage}{4cm}
	\qrcode[height=3.3cm, version=6]{\homepage}
\end{minipage}
\hfill
\begin{minipage}{\dieBreite}
	\resizebox{!}{0.5cm}{\large\sffamily\textbf{Vorlesungshomepage}}\\
	\resizebox{\dieBreite}{!}{\footnotesize\url{\homepage}}\smallskip\\
	Hier ist ein Link zur offiziellen Vorlesungshomepage.
\end{minipage}
\newpage
\section*{Vorwort --- Mitarbeit am Skript}
Dieses Dokument ist eine Mitschrift aus der Vorlesung \enquote{\fach, \semester}, gelesen von \prof. 
Der Inhalt entspricht weitestgehend dem Tafelanschrieb. 
Für die Korrektheit des Inhalts übernehme ich keinerlei Garantie! 
Für Bemerkungen und Korrekturen -- und seien es nur Rechtschreibfehler -- bin ich sehr dankbar. 
Korrekturen lassen sich prinzipiell auf drei Wegen einreichen: 
\begin{itemize}
	\item Persönliches Ansprechen in der Uni, Mails an \hrefsymmail{mailto:\mail}{\mail} (gerne auch mit annotieren PDFs) oder Kommentare auf \url{https://gitlab.com/JaMeZ-B/LaTeX-WWU}.
	\item \emph{Direktes} Mitarbeiten am Skript: Den Quellcode poste ich auf GitLab (siehe oben), also stehen vielfältige Möglichkeiten der Zusammenarbeit zur Verfügung:
	Zum Beispiel durch Kommentare am Code über die Website und die Kombination Fork und Merge-Request. 
	Wer sich verdient macht oder ein Skript zu einer Vorlesung, die ich nicht besuche, beisteuern will, dem gewähre ich gerne auch Schreibzugriff.
	
	Beachten sollte man dabei, dass dazu ein Account bei \url{gitlab.com} notwendig ist, der allerdings ohne Angabe von persönlichen Daten angelegt werden kann. 
	Wer bei GitLab (bzw. dem zugrunde liegenden Open-Source-Programm \enquote{\texttt{git}}) -- verständlicherweise -- Hilfe beim Einstieg braucht, dem helfe ich gerne weiter. 
	Es gibt aber auch zahlreiche empfehlenswerte Tutorials im Internet.\footnote{zB. \url{https://try.github.io/levels/1/challenges/1}, ist auf Englisch, aber dafür interaktiv}
	\item \emph{Indirektes} Mitarbeiten: \TeX-Dateien per Mail verschicken. 
	
	Dies ist nur dann sinnvoll, wenn man einen ganzen Abschnitt ändern möchte (zB. einen alternativen Beweis geben), da ich die Änderungen dann per Hand einbauen muss! Ich freue mich aber auch über solche Beiträge!
\end{itemize}
\section*{Inhalt der Vorlesung}
In der Vorlesung werden wir uns mit algebraischer Topologie beschäftigen. 
In der algebraischen Topologie werden topologische Fragen, wie zum Beispiel wann $\mathbb{R}^n$ und $\mathbb{R}^m$ homöomorph sind, in algebraische Fragen übersetzt. 
Ein wichtiges Hilfsmittel für solche Übersetzungen sind Homologiegruppen. 
Wir werden in der Vorlesung diese Homologiegruppen definieren und Hilfsmittel entwickeln, um Homologiegruppen zu berechnen. 
Es ergeben sich dann schöne topologische Anwendungen, wie zum Beispiel die topologische Invarianz der Dimension. 
\section*{Literatur}
\begin{itemize}
	\item \textsc{Bredon, G.E.}: \emph{Topology and Geometry}, Graduate Texts in Mathematics, Springer Verlag 1993
	\item \textsc{Hatcher, A.}: \href{http://www.math.cornell.edu/~hatcher/AT/ATpage.html}{\appendLink{\emph{Algebraic Topology}}}, Cambridge University Press, 2002
	\item \textsc{Lück, W.}: \emph{Algebraische Topologie - Homologie und Mannigfaltigkeiten}, Vieweg, 2005
	\item \textsc{tom Dieck, T.}: \emph{Topologie}, 2. Auflage, de Gruyter 2000
	\item \textsc{Waldhausen, F.}: \href{https://www.math.uni-bielefeld.de/~fw/ein.pdf}{\appendLink{\emph{Skript zur Einführung in die Topologie}}}, 
	\href{https://www.math.uni-bielefeld.de/~fw/at.pdf}{\appendLink{\emph{Skript zur Topologie}}}
\end{itemize}
\end{abstract}

\tableofcontents
\cleardoubleoddemptypage

\pagenumbering{arabic}
\setcounter{page}{1}
\setcounter{footnote}{0}

\section{Kategorien, Funktoren und natürliche Transformationen} % (fold)
\label{sec:1}

\begin{definition}[{name=[Kategorie]}]
	Eine \Index{Kategorie} $\mathcal{C}$ besteht aus:
	\begin{enumerate}[i),itemsep=0pt]
		\item Einer Klasse $\Obj(\mathcal{C})$. Die Elemente von $\Obj(\mathcal{C})$ heißen die \Index{Objekte} von $\mathcal{C}$.
		\item Zu je zwei Objekten $A,B$ aus $\mathcal{C}$ einer Menge $\Mor_{\mathcal{C}}(A,B)$. Die Elemente von $\Mor_\mathcal{C}(A,B)$ heißen \Index{Morphismen} von $\mathcal{C}$.
		\item Zu je drei Objekten $A,B,C \in \Obj(\mathcal{C})$ aus einer Abbildung: 
		\[
			\Mor_\mathcal{C}(B,C) \times \Mor_\mathcal{C}(A,B) \to \Mor_\mathcal{C}(A,C) \enspace, \quad (f,g) \mapsto f \circ g
		\]
		genannt die \Index{Komposition} in $\mathcal{C}$. 
	\end{enumerate}
	Dabei müssen folgende Axiome  erfüllt sein:
	\begin{enumerate}[(i)]
		\item Die Komposition ist \Index{assoziativ}: Für Objekte $A,B,C,D$ von $\mathcal{C}$ und $f \in \Mor_\mathcal{C}(C,D)$, $ g \in \Mor_\mathcal{C}(B,C)$ und $ h \in \Mor_\mathcal{C}(A,B)$ gilt immer
		\[
			f \circ (g \circ h) = (f \circ g) \circ h
		\]
		\item Die Komposition ist \Index{unital}: Für jedes Objekt $A$ von $\mathcal{C}$ gibt es einen Morphismus $\id_A \in \Mor_\mathcal{C}(A,A)$ so, dass
		\[
			\forall f \in \Mor_\mathcal{C}(B,A) : \id_A \circ f = f  \qquad  \text{und} \qquad  \forall f \in \Mor_\mathcal{C}(A,B) : f \circ \id_A = f
		\]
	\end{enumerate}
\end{definition}

\begin{beispiel}[{name=[Kategorien]}]
	\leavevmode
	\begin{enumerate}[(1),itemsep=0pt]
		\item Die Kategorie der Mengen $\SET$: Objekte sind Mengen und Morphismen sind Abbildungen.
		\item Die Kategorie der Gruppen $\GRUPPEN$: Objekte sind Gruppen und Morphismen sind Gruppenhomomorphismen.
		\item Die Kategorie der $K$-Vektorräume $K\text{-}\VR$: Objekte sind $K$-Vektorräume und die Morphismen sind $K$-lineare Abbildungen.
		\item Die Kategorie der $R$-Moduln $R\text{-}\MOD$: Objekte sind $R$-Moduln und Morphismen sind $R$-lineare Abbildungen.
		\item Die Kategorie der $C^\infty$-Mannigfaltigkeiten $C^\infty\text{-}\MAN$: Objekte sind $C^\infty$-Mannigfaltigkeiten und Morphismen sind $C^\infty$-Abbildungen.
		\item Die Kategorie der topologischen Räume $\TOP$: Objekte sind topologische Räume und Morphismen sind stetige Abbildungen.
		\item Die Kategorie der punktierten topologischen Räume $\TOP_\bullet$: Objekte sind punktiert"=topologische Räume und Morphismen sind punktiert-stetige Abbildungen.
		\item Die Kategorie $\HTOP$: Objekte sind topologische Räume und Morphismen sind Homotopieklassen von stetigen Abbildungen.
		\item Die Kategorie $\HTOP_\bullet$: Objekte sind punktiert-topologische Räume und Morphismen sind Homotopieklassen von punktiert-stetigen Abbildungen.
		\item Sei $G$ eine Gruppe. Wir erhalten eine Kategorie $\mathcal{C}_G$ mit genau einem Objekt $*$ und Morphismen $\Mor_{\mathcal{C}_G}(*,*) = G$. Die Komposition wird durch die
		Verknüpfung in der Gruppe festgelegt.
		\item Ist $\mathcal{C}$ eine Kategorie, so ist $\mathcal{C}^{\op}$ eine Kategorie, wobei $\Obj(\mathcal{C}^\op) = \Obj(\mathcal{C})$ und 
		\(
			\Mor_{\mathcal{C}^\op}(A,B) := \Mor_\mathcal{C}(B,A)
		\).
		Die Komposition ist gegeben durch $f \circ^\op g := g \circ f$.
	\end{enumerate}
\end{beispiel}

\begin{bemerkung}[{name=[kleine Kategorie]}]
	\leavevmode
	\begin{enumerate}[(i)]
		\item Eine Kategorie heißt \bet{klein}\index{Kategorie!klein}, wenn ihre Objekte eine Menge bilden.
		\item Statt $f \in \Mor_\mathcal{C}(A,B)$ schreiben wir oft $f \colon A \to B$ oder $A \xrightarrow{f} B$.
	\end{enumerate}
\end{bemerkung}

\begin{definition}[{name=[Funktor]}]
	Seien $\mathcal{C}$ und $\mathcal{D}$ Kategorien: Ein \Index{Funktor} $F \colon \mathcal{C} \to \mathcal{D}$ ordnet jedem Objekt $C$ von $\mathcal{C}$ ein Objekt $F(C)$ von 
	$\mathcal{D}$ und ordnet jedem Morphismus $f : C \to C'$ in $\mathcal{C}$ einen Morphismus $F(f) \colon F(C) \to F(C')$ in $\mathcal{D}$ zu. Dabei muss gelten:
	\[
		F(f \circ g) = F(f) \circ F(g) \quad \text{und} \quad F(\id_C) = \id_{F(C)}
	\]
\end{definition}

\begin{beispiel}[{name=[Funktoren]}]
	\leavevmode
	\begin{enumerate}[(1)]
		\item Es gibt offensichtliche \enquote{Vergiss}-Funktoren:
		\begin{gather}
			K\text{-}\VR \longrightarrow  \SET\qquad  R\text{-}\MOD \longrightarrow  \SET \qquad \TOP \longrightarrow  \SET\\
			\TOP_\bullet \longrightarrow  \SET  \qquad 
			\TOP_\bullet \longrightarrow   \TOP
		\end{gather}
		\item Sei $\mathcal{C}$ eine Kategorie und $C$ ein Objekt von $\mathcal{C}$. 
		Der durch $C$ \bet{dargestellte Funktor}\index{Funktor!dargestellter} $F_C \colon \mathcal{C} \to \SET$ ist definiert durch 
		\[
			\begin{array}{ll}
				F_C(A) = \Mor_\mathcal{C}(C,A) &\text{ für } A \in  \Obj(\mathcal{C}) \\
				F_C(f) \colon \Mor_\mathcal{C}(C,A) \to \Mor_\mathcal{C}(C,A'), g \mapsto f \circ  g &\text{ für } f \colon A \to A' \text{ in } \Mor(A,A')
			\end{array}
		\]
		\item Die Fundamentalgruppe definiert einen Funktor
		\[
			\pi_1 \colon \HTOP_\bullet \longrightarrow \GRUPPEN
		\]
	\end{enumerate}
	Ist $F \colon \mathcal{C} \to \mathcal{D}$ ein Funktor so schreiben wir oft kürzer und ungenauer $f_* := F(f)$.
	Kleine Kategorien und Funktoren bilden die Kategorie $\KAT$.
\end{beispiel}

\begin{definition}[{name=[natürliche Transformation]},label=def:16]
	Seien $F,G \colon \mathcal{C} \to \mathcal{D}$ zwei Funktoren. 
	Eine \Index{natürliche Transformation} $\tau \colon F \to G$ ordnet jedem $C \in \Obj(\mathcal{C})$ einen Morphismus $\tau_C \colon F(C) \to G(C)$ in $\mathcal{D}$ zu, sodass für jedes $f \colon C \to C'$ in $\mathcal{C}$
	\[
		\begin{tikzcd}[sep=large]
			F(C) \rar["F(f)"] \dar["\tau_C"]& F(C') \dar["\tau_{C'}"] \\
			G(C) \rar["G(f)"] & G(C')
		\end{tikzcd}
	\]
	kommutiert.
\end{definition}

Natürliche Transformationen lassen sich komponieren. 
Für $\tau \colon F \to G$, $\eta \colon G \to H$ ist $\eta \circ \tau \colon F \to H$ gegeben durch
\marginnote{$F,G,H \colon \mathcal{C} \to \mathcal{D}$}
\[
	(\eta \circ \tau)_C :=  \eta_C \circ \tau_C \colon F(C) \to H(C)
\]
Genauer: Für feste kleine Kategorien $\mathcal{C}$ und $\mathcal{D}$ bilden die Funktoren $\mathcal{C} \to \mathcal{D}$ mit den natürlichen Transformationen eine Kategorie $\FUN(\mathcal{C}, \mathcal{D})$.

In der algebraischen Topologie werden topologische Fragen, wie zum Beispiel, wann $\mathbb{R}^n$ und $\mathbb{R}^m$ homöomorph sind, in algebraische Fragen übersetzt. 
Eine Möglichkeit für eine solche Übersetzung sind Funktoren von einer Kategorie von topologischen Räumen, z.B. $\TOP, \TOP_\bullet, \HTOP_\bullet$, in eine algebraisch Kategorie, z.B. $\GRUPPEN, K\text{-}\VR, \ABELGRUPPEN, R\text{-}\MOD$. 
Ein Beispiel für einen solchen Funktor ist die Fundamentalgruppe $\pi_1$. 
Ein Nachteil der Fundamentalgruppe ist, dass diese oft schwierig zu berechnen ist. 
Wir werden in dieser Vorlesung weitere Funktoren und Methoden für ihre Berechnung kennenlernen.

\newpage

\section{Kombinatorische Beschreibungen topologischer Räume} % (fold)
\label{sec:2}

\begin{beispiel}[{name=[kombinatorische Beschreibung eines Raumes]}]
	\Index{Graphen} sind einerseits topologische Räume und andererseits kombinatorische Objekte:
	\begin{align*}
		\vcenter{\hbox{\begin{tikzpicture}[scale=1.2]
			\coordinate (A) at (0,1);
			\coordinate (B) at (1,1);
			\coordinate (C) at (2,1);
			\coordinate (1) at (0,0);
			\coordinate (2) at (1,0);
			\coordinate (3) at (2,0);
			\foreach \x in {1,2,3} {
				\foreach \y in {A,B,C}{
					\draw[thick] (\x) -- (\y);
				};
			}
			\foreach \x in {A,B,C} {
				\foreach \y in {1,2,3}{
					\draw[thick] (\x) -- (\y);
				};
				\draw[fill=DodgerBlue3, thick] (\x) circle[radius=0.05] node[above]{$\x$};
			}
			\foreach \x in {1,2,3} {
				\draw[fill=DodgerBlue3,thick] (\x) circle[radius=0.05] node[below]{$\x$};
			}
		\end{tikzpicture} }}\qquad 
		&\longleftrightarrow \quad \begin{array}{l}
			V =  \set{A,B,C,1,2,3} \\
			E =  \set{ \rule{0cm}{0.3cm} { \scriptstyle \set{A,1}, \set{A,2}, \set{A,3} , \set{B,1}, \set{B,2}, \set{B,3}, \set{C,1} , \set{C,2}, \set{C,3}}}  
		\end{array}\\
		\vcenter{\hbox{\begin{tikzpicture}[scale=1.2]
			\coordinate (1) at (0,1);
			\coordinate (2) at (0.5,1.5);
			\coordinate (3) at (1,1);
			\coordinate (4) at (1,0);
			\coordinate (5) at (0,0);
			\coordinate (6) at (0.5,0.5);
			\foreach \x in {2,3,5,6} {
				\draw[thick] (1) -- (\x);
			}
			\foreach \x/\y in {2/3,3/4,3/6,4/5,4/6,5/6}{
				\draw[thick] (\x) -- (\y);
			}
			\foreach \x/\y in {1/left,2/above,3/right,4/right,5/left,6/right}{
				\draw[fill=DodgerBlue3, thick] (\x) circle[radius=0.05] node[\y]{$\x$};
			}
		\end{tikzpicture} }}\qquad
		&\longleftrightarrow \quad \begin{array}{l}
			V = \set{1,2,3,4,5,6} \\
			E = \set{ \rule{0cm}{0.3cm} {\scriptstyle \set{1,2}, \set{1,3}, \set{1,5} , \set{1,6}, \set{2,3}, \set{3,4}, \set{3,6} , \set{4,5}, \set{4,6}, \set{5,6}} }  
		\end{array}
	\end{align*}
\end{beispiel}

\begin{definition}[{name=[simplizialer Komplex]}]
	Ein \Index{simplizialer Komplex} $K=(V, \Sigma)$ besteht aus einer Menge $V$ und einer Menge $\Sigma$ von nichtleeren, endlichen Teilmengen von $V$, sodass gilt
	\begin{enumerate}[(i)]
		\item $\set{v} \in \Sigma$ für alle $v \in V$
		\item $\sigma \in \Sigma$, $\emptyset \not= \tau \subseteq \sigma \Longrightarrow \tau \in\Sigma.$ \hfill {\footnotesize(Abgeschlossen bzgl. Teilmengenbildung)}
	\end{enumerate}
	Die Elemente von $V$ heißen die \Index{Ecken} oder \Index{Vertices} von $K$. 
	Die Elemente von $\Sigma$ heißen die \Index{Simplizes} von $K$. 
	Enthält $\sigma \in \Sigma$ genau $n+1$ Elemente, so heißt $\sigma$ ein \bet{$\mathbf{n}$-Simplex}\index{n-Simplex@$n$-Simplex}. 
	Ist $\tau \subseteq \sigma$, wobei $\sigma$ ein $n$-Simplex und $\tau$ ein $n-1$-Simplex ist, so heißt $\tau$ eine \Index{Seite} von $\sigma$. 
	Ist $V$ geordnet, so heißt $K$ geordnet. Ist $V$ endlich, so heißt $K$ endlich.
\end{definition}

\begin{beispiel}
	\begin{align*}
		\vcenter{\hbox{\begin{tikzpicture}[scale=1.]
			\coordinate (0) at (0,0);
			\coordinate (1) at (1,0.8);
			\coordinate (2) at (1,-0.7);
			\coordinate (3) at (2,0.2);
			\filldraw[DodgerBlue2!40] (0) -- (1) -- (2) -- cycle ;
			\foreach \x/\y in {0/1,0/2,1/2,1/3,2/3} {
				\draw[thick] (\x) -- (\y);
			};
			\foreach \x/\pos in {0/left,1/above,2/below,3/right} {
				\draw[fill=DodgerBlue3,thick] (\x) circle[radius=0.06] node[\pos]{$\x$};
			}
		\end{tikzpicture}}} \qquad 
		\begin{array}{l}
			\minwidthbox{V}{1em} =  \set{0,1,2,3} \\
			\minwidthbox{\Sigma}{1em} =  \set{ \rule{0cm}{0.3cm} {\set{0}, \set{1}, \set{2}, \set{3}, \set{0,1}, \set{0,2},\set{1,2},\set{1,3},\set{2,3},\set{0,1,2}}}  
		\end{array}
	\end{align*}
	Ist $\sigma$ eine endliche Menge, so heißt $\Delta^\sigma := \enbrace[\big]{\sigma, \mathcal{P}(\sigma) \setminus \set{\emptyset}}$ der $\sigma$-Simplex. Für 
	$\sigma = \set{0, \ldots , n}$ schreiben wir $\Delta^n := \Delta^{\set{0, \ldots ,n}}$. Es ist
	\[
		\Delta^0  = \vcenter{\hbox{\begin{tikzpicture}
			\draw[fill=DodgerBlue3,thick] (0,0) circle[radius=0.06];
		\end{tikzpicture}}} , \qquad 
		\Delta^1 = \vcenter{\hbox{\begin{tikzpicture}
			\coordinate (0) at (0,0);
			\coordinate (1) at (0.4,0.3);
			\draw[thick] (1) -- (0);
			\foreach \x in {0,1} {
				\draw[fill=DodgerBlue3,thick] (\x) circle[radius=0.06];
			}
		\end{tikzpicture}}} , \qquad 
		\Delta^2 = \vcenter{\hbox{\begin{tikzpicture}
			\coordinate (0) at (0,0);
			\coordinate (1) at (0.6,0);
			\coordinate (2) at (60:0.6);
			\draw[thick, fill=DodgerBlue2!40] (1) -- (0) -- (2) -- cycle;
			\foreach \x in {0,1,2} {
				\draw[fill=DodgerBlue3,thick] (\x) circle[radius=0.06];
			}
		\end{tikzpicture}}}, \qquad 
		\Delta^3 = 
		\vcenter{\hbox{\begin{tikzpicture}
				\coordinate (0) at (0,0);
				\coordinate (1) at (1,0);
				\coordinate (2) at (0.5,0,0.866);
				\coordinate (3) at (0.5,0.866,0.433);
				\foreach \x in {0,1,2,3}{
					\foreach \y in {0,1,2,3}{
						\foreach \z/\col in {0/DodgerBlue2,1/DodgerBlue1,2/DodgerBlue3,3/DodgerBlue2}{
						\draw[thick, fill=\col!40, fill opacity=0.4] (\x) -- (\y) -- (\z) -- cycle;
						};
					};
				};
				\foreach \x in {0,1,2,3} {
					\draw[fill=DodgerBlue3,thick] (\x) circle[radius=0.06];
				}
			\end{tikzpicture}}}
	\]
\end{beispiel}

\begin{definition}[{name=[Unterkomplex]}]
	Sei $K=(V, \Sigma)$ ein simplizialer Komplex. 
	Ein \Index{Unterkomplex} von $K$ ist ein simplizialer Komplex $K_0 = (V_0, \Sigma_0)$ mit $V_0 \subseteq V$ und $\Sigma_0 \subseteq \Sigma$.
	
	Sei $K=(V,\Sigma)$ ein endlicher simplizialer Komplex. 
	Dann ist $K$ ein Unterkomplex von $\Delta^V$.
\end{definition}

\begin{definition}[{name=[simpliziale Abbildung]},label=def:25]
	Eine \Index{simpliziale Abbildung} $f \colon K_1=(V_1, \Sigma_1) \to K_2=(V_2, \Sigma_2)$ zwischen simplizialen Komplexen ist eine Abbildung $f \colon V_1 \to V_2$, sodass $f(\sigma_1) \in \Sigma_2$ für alle $\sigma_1 \in \Sigma_1$. 
	Simpliziale Komplexe zusammen mit simplizialen Abbildungen bilden eine Kategorie $\SIMP$.
\end{definition}

\begin{definition}[{name=[geometrische Realisierung]}]
	Sei $\sigma$ eine endliche Menge. 
	Sei $\mathbb{R}^\sigma = \prod_{v \in \sigma} \mathbb{R}$ mit der Produkttopologie. 
	Wir setzen nun 
	\[
		\abs*{\Delta^\sigma} := \set*{x=(x_v)_{v \in \sigma} \in \mathbb{R}^\sigma \given \sum_{v \in \sigma}x_v = 1 , x_v \in [0,1] \,\, \forall v \in \sigma}  
	\]
	Ist $\tau \subseteq \sigma$, so erhalten wir eine Abbildung $\iota_\tau^\sigma \colon \abs*{\Delta^\tau} \to \abs*{\Delta^\sigma}$ indem wir $(x_v)_{v \in \tau} \in \abs*{\Delta^\tau}$ durch $x_v=0$ für $v \in \sigma \setminus \tau$ zu $(x_v)_{v \in \sigma}$ auffüllen. 
	Ist $\eta \subseteq \tau \subseteq \sigma$ so gilt offensichtlich $\iota^\sigma_\eta = \iota_\tau^\sigma \circ  \iota_\eta^\tau$.
	Sei nun $K=(V, \Sigma)$ ein simplizialer Komplex. 
	Die \Index{geometrische Realisierung} $\abs*{K}$ von $K$ ist definiert als 
	\[
		\abs*{K} := \coprod_{\sigma \in \Sigma} \set{\sigma} \times \abs*{\Delta^\sigma} \Big/ \!\sim
	\] 
	wobei $\sim$ die durch $(\tau, x) \sim \enbrace[\big]{\sigma, \iota_\tau^\sigma(x)}$ für $\tau \subseteq \sigma$ erzeugte Äquivalenzrelation ist.\marginnote{Anschaulich \enquote{klebt} $\sim$ die geometrischen Realisierungen aller Einzelteile passend zusammen} 
	Versehen mit der Quotiententopologie ist $\abs*{K}$ ein topologischer Raum.
\end{definition}

\begin{itemize}
	\item In Aufgabe 4 von Blatt 1 haben wir gezeigt, dass $\abs*{K}$ homöomorph zu einem Unterraum $X$ von $\mathbb{R}^n$ mit $n:= \# V$ ist via $\benbrace[\big]{\enbrace*{\set{\sigma}, (x_{v_i})_{v_i \in \sigma} } } \mapsto \sum_{v_i \in \sigma} x_{v_i} \cdot e_i$.
	\item In Aufgabe 2 von Blatt 2 haben wir gezeigt, dass $\abs*{\,\cdot\,} $ einen Funktor $\SIMP \to \TOP$ definiert.
	Dazu definiert man eine Abbildung zunächst auf den zugrunde liegenden Räumen durch
	\[
		\enbrace*{\sigma, \sum\nolimits_{v \in \sigma} x_v \cdot v} \longmapsto \enbrace*{f(\sigma), \sum\nolimits_{v \in \sigma} x_v \cdot f(v)} 
	\]
	und wendet diese unter Benutzung der universellen Eigenschaft der Quotiententopologie auf Repräsentanten an.\footnote{hier wird eine formale Summe anstatt eines Tupels benutzt; da $\sigma$ stets eine endliche Menge ist, macht dies aber keinen Unterschied.}
\end{itemize}

\begin{beispiel}[{name=[Rand des $n$-Simplizes und der Torus]}]
	\leavevmode
	\begin{itemize}
		\item Sei $\partial \Delta^n := \enbrace[\big]{\set{0, \ldots ,n}, \mathcal{P}\enbrace[\big]{\set{0,\ldots ,n}} \setminus \set{\emptyset, \set{0, \ldots , n}}} = $ 
		\enquote{$\Delta^n \setminus \set{0,\ldots ,n}$}. $\partial \Delta^n$ heißt der \Index{Rand des $n$-Simplizes}. Es gilt
		\begin{align*}
			\abs*{\partial \Delta^1} &= \enspace\vcenter{\hbox{\begin{tikzpicture}
			\draw[fill] (0,0) circle[radius=0.04];
			\draw[fill] (0.155,0) circle[radius=0.04];
		\end{tikzpicture}}} \enspace\cong S^0\\ 
			\abs*{\partial \Delta^2} &=  \vcenter{\hbox{\begin{tikzpicture}
			\coordinate (0) at (0,0);
			\coordinate (1) at (0.6,0);
			\coordinate (2) at (60:0.6);
			\draw[thick] (1) -- (0) -- (2) -- cycle;
		\end{tikzpicture}}} \cong S^1 \\
			\abs*{\partial \Delta^3} &= \vcenter{\hbox{
			\begin{tikzpicture}
				\coordinate (0) at (0,0);
				\coordinate (1) at (1,0);
				\coordinate (2) at (0.5,0,0.866);
				\coordinate (3) at (0.5,0.866,0.433);
				\foreach \x in {0,1,2,3}{
					\foreach \y in {0,1,2,3}{
						\foreach \z/\col in {0/DodgerBlue2,1/DodgerBlue1,2/DodgerBlue3,3/DodgerBlue2}{
						\draw[thick, fill=\col!40, fill opacity=0.4] (\x) -- (\y) -- (\z) -- cycle;
						};
					};
				};
			\end{tikzpicture}}} \cong S^2 
		\end{align*}
		Allgemein gilt $\abs*{\partial \Delta^n} \cong S^{n-1}$, siehe Aufgabe 1 von Blatt 3. 
		\item $T^2 = S^1 \times S^1 \cong 
		\vcenter{\hbox{\begin{tikzpicture}[scale=0.3]
			\draw (-3.5,0) .. controls (-3.5,2) and (-1.5,2.5) .. (0,2.5);
			\draw[xscale=-1] (-3.5,0) .. controls (-3.5,2) and (-1.5,2.5) .. (0,2.5);
			\draw[rotate=180] (-3.5,0) .. controls (-3.5,2) and (-1.5,2.5) .. (0,2.5);
			\draw[yscale=-1] (-3.5,0) .. controls (-3.5,2) and (-1.5,2.5) .. (0,2.5);

			\draw (-2,.2) .. controls (-1.5,-0.3) and (-1,-0.5) .. (0,-.5) .. controls (1,-0.5) and (1.5,-0.3) .. (2,0.2);

			\draw (-1.75,0) .. controls (-1.5,0.3) and (-1,0.5) .. (0,.5) .. controls (1,0.5) and (1.5,0.3) .. (1.75,0);
		\end{tikzpicture}}} \cong 
		\abs*{\vcenter{\hbox{\begin{tikzpicture}[>=to, scale=0.8]
			\filldraw[color=DodgerBlue2!20,] (0,0) rectangle (3,3);
			\foreach \x in {0,1,2}{
				\foreach \y in {0,1,2}{
					\draw (\x,\y) -- ++ (1,1);
				};
			};
			\draw (0,0) grid (3,3);
			\draw[very thick] (0,0) -- (0,3) -- (3,3) -- (3,0) -- cycle;
			\draw[->>, thick] (0,1) -- (0,1.5);
			\draw[->, thick] (1,3) -- (1.5,3);
			\draw[->, thick] (1,0) -- (1.5,0);
			\draw[->>, thick] (3,1) -- (3,1.5);
			\clip (-0.2,-0.2) rectangle (3.2,3.2);
		\end{tikzpicture}}}} $
	\end{itemize}
\end{beispiel}

\begin{definition}[{name=[Polyeder]}]
	Ein topologischer Raum $X$ heißt ein \Index{Polyeder}, falls er homöomorph zur geometrischen Realisierung eines simplizialen Komplexes ist. 
\end{definition}

\begin{definition}[{name=[Euler-Charakteristik]}]
	Sei $K=(V,\Sigma)$ ein endlicher simplizialer Komplex. Sei 
	\[
		a_n := \# \set[\big]{\sigma \in \Sigma \given \sigma \text{ ist ein $n$-Simplex von $K$}}
	\] 
	Dann heißt $\chi(K) := \sum_{n \in \mathbb{N}} (-1)^n \cdot a_n $ die \Index{Euler-Charakteristik} von $K$.
\end{definition}

Die Eulercharakteristik lässt sich einfach durch Abzählen bestimmen: 
\begin{align*}
	\chi \enbrace*{\vcenter{\hbox{\begin{tikzpicture}
		\coordinate (0) at (0,0);
		\coordinate (1) at (0.6,0);
		\coordinate (2) at (60:0.6);
		\draw[thick] (1) -- (0) -- (2) -- cycle;
	\end{tikzpicture}}} } &= 3 -3 = 0 \\
	\chi \enbrace[\bigg]{\vcenter{\hbox{\begin{tikzpicture}
		\coordinate (0) at (0,0);
		\coordinate (1) at (0.6,0);
		\coordinate (2) at (0.6,0.6);
		\coordinate (3) at (0,0.6);
		\draw[thick] (0) -- (1) -- (2) -- (3) -- cycle;
	\end{tikzpicture}}} } &= 4-4 =0 \\
	\chi \enbrace[\big]{\benbrace*{n\text{-Eck}}} &= n-n = 0
\end{align*}
Selbstverständlich funktioniert dies auch bei dreidimensionalen Komplexen:
\begin{align*}
	\chi \enbrace*{\partial \Delta^3 = \vcenter{\hbox{\begin{tikzpicture}
			\coordinate (0) at (0,0);
			\coordinate (1) at (1,0);
			\coordinate (2) at (0.5,0,0.866);
			\coordinate (3) at (0.5,0.866,0.433);
			\draw[sharp corners, dashed] (0) -- (1);
			\draw[sharp corners] (1) -- (2) -- (0) -- (3) -- (2) -- (1) -- (3) -- (0);
		\end{tikzpicture}}} } &= 4 - 6 +4 = 2 \\
	\chi\enbrace[\Bigg]{ \vcenter{\hbox{\begin{tikzpicture}
		\coordinate (0) at (0,0,-0.6);
		\coordinate (1) at (0.6,0,-0.6);
		\coordinate (2) at (0.6,0.6,-0.6);
		\coordinate (3) at (0,0.6,-0.6);
		\draw[thick, dashed] (3) -- (0) -- (1);
		\draw[thick] (1) -- (2) -- (3);
		\draw[thick, dashed] (0,0) -- (0);
		\draw[thick] (0.6,0) -- (0.6,0,-0.6);
		\draw[thick] (0.6,0.6) -- (0.6,0.6,-0.6);
		\draw[thick] (0,0.6) -- (0,0.6,-0.6) -- ++ (0.2,0);
		\draw[thick] (0,0) rectangle (0.6,0.6);
	\end{tikzpicture}}} } &= 8 -12 +6 = 2 \\
	\chi \enbrace*{ \vcenter{\hbox{\begin{tikzpicture}
		\coordinate (A1) at (0,0);
		\coordinate (A2) at (0.6,0.2);
		\coordinate (A3) at (1,0);
		\coordinate (A4) at (0.4,-0.2);
		\coordinate (B1) at (0.5,0.5);
		\coordinate (B2) at (0.5,-0.5);
		\begin{scope}[thick,dashed,]
		\draw (A1) -- (A2) -- (A3);
		\draw (B1) -- (A2) -- (B2);
		\end{scope}
		\begin{scope}[thick]
		\draw[] (A1) -- (A4) -- (B1);
		\draw[] (A1) -- (A4) -- (B2);
		\draw[] (A3) -- (A4) -- (B1);
		\draw[] (A3) -- (A4) -- (B2);
		\draw (B1) -- (A1) -- (B2) -- (A3) --cycle;
		\end{scope}
	\end{tikzpicture}}} } &= 6- 12  +8 = 2 
\end{align*}

\begin{satz}[{name=[die Eulercharakteristik ist eine topologische Invariante]}]
	Seinen $K$ und $K'$ endliche simpliziale Komplexe. 
	Gilt $\abs*{K} \cong \abs*{K'}$, so gilt 
	\[
		\chi(K)= \chi(K').
	\]
\end{satz}
\begin{beweis}
	\emph{Siehe \ref{kor:129}.}
\end{beweis}

Die Euler-Charakteristik ist also eine topologische Invariante von simplizialen Komplexen. 
Sie hängt nur von der topologischen Struktur von $\abs*{K}$ und nicht von der kombinatorischen Struktur von $K$ ab.

\newpage

\section{Simpliziale Homologie} % (fold)
\label{sec:3}

\begin{definition}[{name=[freier Modul]}]
	Sei $S$ eine Menge. Sei $R$ ein Ring.\marginnote{Ring mit Eins} 
	Der \bet{freie $R$-Modul mit Basis $S$}\index{freier Modul}, $R[S]$, besteht aus allen endlichen formalen $R$-Linearkombinationen
	\[
		\sum_{s \in S} r_s \cdot s \quad , r_s \in R \text{ und } r_s \not=0 \text{ für endliche viele }s \in S
	\]
	Die $R$-Modulstruktur auf $R[S]$ ist definiert durch:
	\begin{align*}
		r \cdot \enbrace*{\sum_{s\in S} r_s \cdot s} := \sum_{s \in S} (r \cdot r_s) \cdot s \qquad \qquad 
		\enbrace*{\sum_{s \in S} r_s \cdot s} + \enbrace*{\sum_{s \in S} r_s' s} := \sum_{s \in S} (r_s + r_s') \cdot s
	\end{align*}
	Mittels $s= \sum_{s' \in S} \delta_{s,s'} \cdot s'$ fassen wir $S$ als Teilmenge von $R[S]$ auf.
\end{definition}

Die freien Moduln sind durch die folgende universelle Eigenschaft charakterisiert:

\begin{lemma}[{name=[Universelle Eigenschaft von freien Moduln]},label=univ_frei_mod]
	Sei $S$ eine Menge, $R$ ein Ring. Es gilt\\
	\begin{minipage}[c]{0.7\textwidth}
		\begin{enumerate}[1)]
			\item Ist $M$ ein $R$-Modul und $i \colon S \to M$ eine Abbildung, so gibt es genau eine $R$-lineare Abbildung $\varphi \colon R[S] \to M$, die $i$ fortsetzt.
			\item Sei $f\colon R[S]\to M$ $R$-linear und $p\colon N\twoheadrightarrow M$ $R$-linear und surjektiv. Dann gibt es $\hat{f}\colon R[S]\to N$ mit $p\circ \hat{f}=f$.
		
			Moduln mit dieser Eigenschaft bezeichnet man als \bet{projektiv}\index{projektiver Modul}. Freie Moduln sind also projektiv.
		\end{enumerate}
	\end{minipage}\hfill
	\begin{minipage}[c]{0.25\textwidth}
		\[
			\begin{tikzcd}
				\mbox{ } & N \dar["p", two heads] \\
				R[S] \rar["f"] \urar["\hat{f}", dashed] & M
			\end{tikzcd}
		\]
	\end{minipage}
\end{lemma}
\begin{beweis}
	\leavevmode
	\begin{enumerate}[1)]
		\item Eine solche Abbildung ist gegeben durch
		\(
			\varphi \enbrace*{\sum_{s \in S} r_s \cdot s} = \sum_{s \in S} r_s \cdot i(s)
		\).
		Ist $\varphi'$ eine zweite, so gilt 
		\begin{align*}
			\hat{\varphi} \enbrace*{\sum_{s \in S} r_s \cdot s}  = \sum_{s \in S} r_s \cdot \hat{\varphi}(s) = \sum_{s  \in S} r_s \cdot i(s)
		\end{align*}
		\item Wähle für jedes $s \in S$ ein Urbild $j(s)\in N$ für $f(s) \in M$ unter $p\colon N \twoheadrightarrow M$. 
		Nun wende 1) auf $j \colon S \to N$ an, um $\hat{f} \colon R[S] \to N$ mit $\hat{f}(s) = j(s)$ zu erhalten. 
		Nun ist $ p \circ \hat{f}(s) = f(s)$ für alle $s \in S$. 
		Mit der Eindeutigkeit aus 1) folgt $p \circ \hat{f} = f$.\qedhere
	\end{enumerate}
\end{beweis}

\begin{beispiel}[{name=[freie Moduln]}]
	Viele uns bereits bekannte Strukturen sind freie Moduln:
	\begin{enumerate}[(i)]
		\item Sei $K$ ein Körper und $V,M,N$ seien $K$-Vektorräume, $f\colon V \to M$, $p\colon N \to M$ seien $K$-linear und $p$ surjektiv. 
		Dann gibt es $\hat{f}\colon V\to N$ mit $p \circ \hat{f} = f$. 
		Wir können Lemma \ref{univ_frei_mod} benutzen, da $V$ eine Basis $B$ hat, also $V \cong K[B]$.
		\item Sei $R=\mathbb{Z}$, sei $V= \nicefrac{\mathbb{Z}}{2 \mathbb{Z}}$ und $M = \nicefrac{\mathbb{Z}}{2 \mathbb{Z}}$, sowie $f= \id \colon V \to M$, $N= \mathbb{Z}$ und $p \colon \mathbb{Z} \twoheadrightarrow \nicefrac{\mathbb{Z}}{2 \mathbb{Z}}$ die Projektion mit $p(n)= n + 2 \mathbb{Z}$. 
		Dann gibt es \emph{keine} $\mathbb{Z}$-lineare Abbildung $\hat{f} \colon \nicefrac{\mathbb{Z}}{2 \mathbb{Z}} \to \mathbb{Z}$ mit $p \circ \hat{f} = f$, denn falls $\hat{f} (1+ 2 \mathbb{Z}) = n \in \mathbb{Z}$ ist, so folgt 
		\[
			2 n = 2 \cdot \hat{f}(1+ 2 \mathbb{Z}) = \hat{f} \enbrace[\big]{2 \cdot (1+ 2 \mathbb{Z})} = \hat{f} (0+ 2 \mathbb{Z}) =0
		\]
		also $n=0$ und somit $p \circ \hat{f}=0$. 
		$\nicefrac{\mathbb{Z}}{2 \mathbb{Z}}$ ist also kein freier Modul.
	\end{enumerate}
\end{beispiel}

\begin{definition}[{name=[{$n$-ter Kettenmodul}]}]
	Sei $K=(V, \Sigma)$ ein simplizialer Komplex. Sei $\Sigma_n := \set{\sigma \in \Sigma \given \sigma \text{ ist ein $n$-Simplex}}$. Sei $R$ ein Ring. Der 
	\bet{$\mathbf{n}$-te Kettenmodul}\index{n-ter Kettenmodul@$n$-ter Kettenmodul!simplizial} von $K$ über $R$ ist definiert als 
	\[
		C_n(K;R) := R[\Sigma_n]
	\]
	Ist $R=\mathbb{Z}$, so schreiben wir auch kurz $C_n(K) := C_n(K;\mathbb{Z})$.
\end{definition}

\begin{definition}[{name=[{Orientierung eines $(n-1)$-Simplizes als Teilmenge eines $n$-Simplizes}]}]
	Sei $K=(V,\Sigma)$ ein geordneter simplizialer Komplex. 
	Sei $\sigma = \set{v_0, \ldots ,v_n}$ ein $n$-Simplex von $K$. 
	Bezüglich der Ordnung von $K$ sei dabei $v_0 < v_1 < \ldots < v_n$. 
	Sei $\tau$ ein $(n-1)$-Simplex von $K$. Definiere
	\[
		\varepsilon_\sigma^\tau := \begin{cases}
			(-1)^i, &\text{ falls }\tau = \set{v_0, \ldots , v_{i-1}, v_{i+1}, \ldots , v_n} \\
			0 , &\text{ sonst}
		\end{cases}
	\]
\end{definition}

Es gilt offensichtlich $\varepsilon^\tau_\sigma \neq 0$ genau dann, wenn $\tau$ eine Seite von $\sigma$ ist.
Für den folgenden geordneten simplizialen Komplex gilt
\[
	\vcenter{\hbox{\begin{tikzpicture}
		\coordinate (0) at (0,0);
		\coordinate (1) at (0.5,0.866);
		\coordinate (2) at (1,0);
		\draw[very thick,
			fill=DodgerBlue2!40,
			postaction={decorate},
			decoration={markings,
				mark=at position 1/6 with {\arrow{to}},
				mark=at position 0.5 with {\arrow{to}},
				mark=at position 5/6 with {\arrow{to reversed}}}] 
		(0) -- (1) -- (2) -- cycle;
		\foreach \x/\pos in{0/left,1/above,2/right}{
			\node[\pos] at (\x) {$\x$};
		}
	\end{tikzpicture}}} \qquad 
	\varepsilon_{\set{0,1,2}}^{\set{0,1}} = 1 , \quad \varepsilon_{\set{0,1,2}}^{\set{0,2}}=-1 , \quad \varepsilon_{\set{0,1,2}}^{\set{1,2}} =1    
\]

\begin{definition}[{name=[{$n$-te Randabbildung}]}]
	Sei $K=(V,\Sigma)$ ein geordneter simplizialer Komplex. 
	Die \bet{$\mathbf{n}$-te Randabbildung}\index{n-te Randabbildung@$n$-te Randabbildung!simplizial} 
	\[
		\partial_n \colon \underset{=R[\Sigma_n]}{C_n(K;R)} \longrightarrow \underset{=R[\Sigma_{n-1}]}{C_{n-1}(K;R)}
	\]
	ist für ein Basiselement $\sigma\in \Sigma_n$ definiert durch
	\[
		\partial_n(\sigma) = \sum_{\mathclap{\tau \in \Sigma_{n-1}}} \varepsilon_\sigma^\tau \cdot \tau
	\]
\end{definition}

\begin{proposition}[{name=[{Für die Randabbildung gilt $\partial_{n-1} \circ \partial_n = 0$}]}]
	Es gilt
	\[
		\partial_{n-1} \circ \partial_n = 0
	\]
	Damit folgt $\im \partial_n \subseteq \ker \partial_{n-1}$.
\end{proposition}
\begin{beweis}
	Sei $\sigma$ ein $n$-Simplex von $K$. Dann gilt
	\begin{align*}
		\partial_{n-1} \enbrace[\big]{\partial_n (\sigma)} = \partial_{n-1} \enbrace*{ \smashoperator[r]{\sum_{\tau \in \Sigma_{n-1}}} \varepsilon_\sigma^\tau \cdot \tau} = 
		\sum_{\mathclap{\tau \in \Sigma_{n-1}}} \varepsilon_{\sigma}^\tau \cdot \partial_{n-1} (\tau) 
		 &= \smashoperator[r]{\sum_{\tau \in \Sigma_{n-1}}} \varepsilon_\sigma^\tau \enbrace*{ \sum_{\eta \in \Sigma_{n-2}} \varepsilon_\tau^\eta \cdot \eta} \\
		 &= \sum_{\eta \in \Sigma_{n-2}} \enbrace*{ \sum_{\tau \in \Sigma_{n-1}, \eta \subseteq \tau \subseteq \sigma} \!\!\!\!\! \! \! \varepsilon_\tau^\eta \cdot \varepsilon_\sigma^\tau}
		 \eta  \\
		 &\stackrel{\mathclap{\text{Lemma \ref{lem:38}}}}{=} \quad \enspace 0\qedhere
	\end{align*}
\end{beweis}

\begin{lemma}[label=lem:38,{name=[{Ein $(n-2)$-Simplex ist Seite von genau zwei $(n-1)$-Simplizes}]}]
	Sei $\sigma$ ein $n$-Simplex. Sei $\eta \subseteq \sigma$ ein $(n-2)$-Simplex. Dann gibt es genau zwei $(n-1)$-Simplizes $\tau, \tau'$ von $K$, die eine Seite von $\sigma$
	sind und $\eta$ als Seite enthalten. Es gilt
	\[
		\varepsilon_\tau^\eta \cdot \varepsilon^\tau_\sigma = - \varepsilon_{\tau'}^\eta \cdot \varepsilon^{\tau'}_\sigma
	\]
\end{lemma}
\begin{beweis}
	Sei $\sigma= \set{v_0, \ldots ,v_n}$ mit $v_0 < v_1 < \ldots < v_n$. 
	Dann ist $\eta = \set{v_0, \ldots , v_n} \setminus \set{v_i, v_j}$ mit $i<j$. 
	Dann sind $\tau = \set{v_0, \ldots , v_{i-1}, v_{i+1}, \ldots , v_n}$ und $\tau' = \set{v_0, \ldots , v_{j-1}, v_{j+1}, \ldots , v_n}$ die gesuchten $(n-1)$-Simplizes und es gilt
	\[
		\varepsilon_\sigma^\tau = (-1)^i \enspace, \enspace \varepsilon_\sigma^{\tau'} = (-1)^j \enspace, \enspace \varepsilon^\eta_\tau = (-1)^{j-1} \enspace, \enspace 
		\varepsilon_{\tau'}^\eta = (-1)^i \qedhere
	\]
\end{beweis}

Wir vereinbaren die folgende Notation: Für $\sigma \in \Sigma_n$, $ \sigma= \set{v_0, \ldots ,v_n} $ mit $v_0 < \ldots < v_n$ schreiben wir 
\[
	\delta^j \sigma := \set{v_0 , \ldots , v_{j-1}, v_{j+1}, \ldots , v_n}  \in \Sigma_{n-1}
\]
Dann ist $\partial_n(\sigma) = \sum_{j=0}^{n} (-1)^j \cdot \delta^j \sigma$.

\begin{definition}[{name=[{$n$-ter Homologiemodul}]}]
	Sei $K$ ein geordneter simplizialer Komplex. Der \bet{$\mathbf{n}$-te Homologiemodul}\index{n-ter Homologiemodul@$n$-ter Homologiemodul!simplizial} von 
	$K$ über $R$ ist definiert als\marginnote{In Aufgabe 3 von Blatt 3 haben wir gezeigt, dass die Homologiemoduln für verschiedene Ordnungen auf der gleichen Menge isomorph sind.}
	\[
		H_n(K;R) := \frac{\ker \partial_n \colon C_n(K;R) \to C_{n-1}(K;R)}{\im \partial_{n+1} \colon C_{n+1}(K;R) \to C_n(K;R)} 
	\]
	Für $n=0$ interpretieren wir $\partial_0$ als die Nullabbildung. Daher gilt 
	\[
		H_0(K;R) = \frac{C_0(K;R)}{\im \partial_1 \colon C_1(K;R) \to C_0(K;R)} 
	\]
\end{definition}

\begin{beispiel}
	Für ausreichend \enquote{einfache} Komplexe können wir die Homologie direkt ausrechnen:
	\begin{enumerate}[(1)]
		\item Sei $K= \bullet$. Dann ist $C_0(K;R) \cong R$ und $C_i(K;R) = 0$ für $i>0$. Weiter ist $H_0(K;R) = C_0(K;R) \cong R$ und $H_i(K;R) =0$ für $i>0$.
		\item Sein nun $K=K_n$ das $n$-Eck, wobei $n \ge 3$.  Also
		\begin{align*}
			\Sigma_0 &= \set[\big]{\tau_1=\set{1}, \tau_2=\set{2}, \ldots , \tau_n = \set{n} \rule{0cm}{0.9em}} \\
	 		\Sigma_1 &= \set[\big]{\sigma_1 = \set{1,2}, \ldots , \sigma_{n-1} = \set{n-1, n}, \sigma_n=\set{1,n} \rule{0cm}{0.9em}} 
		\end{align*}
		Dann gilt
		\begin{align*}
			C_0(K;R) &= R[\Sigma_0] \cong R^{n} \\
			C_1(K;R) &= R[\Sigma_1] \cong R^{n} \\
			C_i(K;R) &= 0 \text{ für $i>1$, insbesondere } H_i(K;R) = 0 \text{ für } i>1
		\end{align*}
		Es ist dann $\partial_2 = \partial_3 = \ldots  = 0$. Interessant ist $\partial_1 \colon C_1(K;R) \to C_0(K;R)$, denn es gilt
		\[
			\partial_1(\sigma_i) = \begin{cases}
				\tau_{i+1} - \tau_i, &\text{ falls }i=1, \ldots ,n-1\\
				\tau_n - \tau_1 , &\text{ falls } i=n
			\end{cases}
		\]
		Mit $\sigma_1' := \sigma_1, \sigma_2' := \sigma_2, \ldots , \sigma_{n-1}' := \sigma_{n-1}$ und $ \sigma_n' := - \sigma_n$ gilt dann
		\[
			\partial_1 (\sigma_i') = \tau_{i+1} - \tau_i
		\]
		mit der Konvention $\tau_{n+1} = \tau_1$. Also gilt
		\begin{align*}
			\partial_1 \enbrace*{\sum_{i=1}^{n} r_i \sigma_i' } = \sum_{i=1}^{n} r_i \enbrace*{\tau_{i+1}- \tau_i} = \sum_{i=1}^{n} (r_{i-1} - r_i) \tau_i  
		\end{align*}
		Es folgt
		\begin{align*}
			\ker \partial_1 &= \set*{\sum_{i=1}^{n} r_i \cdot \sigma_i' \given r_1=r_2= \ldots =r_n} \subseteq C_1(K;R) \\
			\im \partial_1 &= \set*{\sum_{i=1}^{n} s_i \cdot \tau_i \given \sum_{i=1}^{n} s_i =0} \subseteq C_0(K;R)
		\end{align*}
		Ist $\sum_{i=1}^{n} s_i =0$, so folgt für $r_n=0, r_1 = -s_1, r_2=-s_1 -s_2, \ldots , r_{n-1} = -s_1 - s_2 - \ldots - s_{n-1}$
		\[
			\partial_1 \enbrace*{\sum_{i=1}^{n} r_i \sigma_i'} = \sum_{i=1}^{n} s_i \tau_i 
		\]
		Es folgt $ %{\color{Firebrick1}[K] \in}\, 
		H_1(K;R) = \nicefrac{\ker \partial_1}{\im \partial_2} = \ker \partial_1 \cong R % \,{\color{Firebrick1}\ni 1}
		$ und 
		\[
			% { \color{Gold3} [\tau_i] =[p] \in }\,
			H_0(K;R) = \nicefrac{\ker \partial_0}{\im \partial_1} = \nicefrac{R[\Sigma_0]}{ \set*{\sum_{i=1}^{n} s_i \tau_i \given \sum_{i=1}^{n} s_i =0} } \cong R 
			% \, { \color{Gold3} \ni 1 }
		\]
		(via $\sum_{i=1}^{n} s_i \tau_i \longmapsto \sum_{i=1}^{n} s_i$ und Homomorphiesatz)
	\end{enumerate}
\end{beispiel}

\begin{definition}[{name=[{$R$-Kettenkomplex, $n$-te Homologie eines $R$-Kettenkomplexes}]}]
	Sei $R$ ein Ring. Ein \bet{$\mathbf{R}$-Kettenkomplex}\index{R-Kettenkomplex@$R$-Kettenkomplex} $(C_*, d_*)$ ist eine Folge von $R$-Moduln $(C_n)_{n \in \mathbb{N}}$ zusammen mit $R$-linearen Abbildungen $(d_n \colon C_n \to C_{n-1})_{n \ge 1}$, so dass $d_n \circ d_{n+1} = 0$. 
	Die $d_n$ heißen die \Index{Randabbildungen} von $(C_*,d_*)$, die $C_n$ die \Index{Kettenmoduln}.
	Die \bet{$\mathbf{n}$-te Homologie}\index{n-te Homologie@$n$-te Homologie} eines Kettenkomplexes $(C_*,d_*)$ ist definiert als
	\[
		H_n(C_*, d_*) := \frac{\ker d_n \colon C_n \to C_{n-1}}{\im d_{n+1} \colon C_{n+1} \to C_n} 
	\]
\end{definition}

\begin{bemerkung}
	\leavevmode
	\begin{enumerate}[(i)]
		\item Oft werden auch Kettenkomplexe betrachtet, die $\mathbb{Z}$ statt $\mathbb{N}$ verwenden.
		\item Ein Kettenkomplex $(C_*,d_*)$ heißt endlich erzeugt, wenn alle $C_n$ endlich erzeugte $R$-Moduln sind und $C_n \not= 0$ nur für endlich viele $n$ ist.
	\end{enumerate}
\end{bemerkung}

\begin{definition}[{name=[Simplizialer Kettenkomplex]}]
	Zu einem geordneten simplizialen Komplex $K$ heißt $\enbrace[\big]{C_*(K;R), \partial_*}$ der \bet{simpliziale Kettenkomplex}\index{simplizialer Kettenkomplex} von $K$ über $R$.
\end{definition}

Ist $K$ endlich, so ist $\enbrace[\big]{C_*(K;R), \partial_*}$ endlich erzeugt.


\newpage

\section{Euler-Charakteristik von Kettenkomplexen} % (fold)
\label{sec:4}

\begin{erinnerungA}[{name=[Torsionsgruppe]}]
	Sei $A$ eine abelsche Gruppe. Dann ist 
	\[
		T A := \set*{a \in A \given \exists n \ge 1 : n \cdot a = \underbracket[1pt][2pt]{a + \ldots + a}_{\text{$n$-mal}} = 0} 
	\]
	die \Index{Torsionsgruppe} von $A$. Ist $A$ endlich erzeugt, so ist auch $TA$ endlich erzeugt und es gibt Primzahlpotenzen $p_1^{n_1}, \ldots , p_k^{n_k}$ mit
	\marginnote{abelsche Gruppen sind $\mathbb{Z}$-Moduln}
	\[
		TA \cong \nicefrac{\mathbb{Z}}{(p_1^{n_1})} \oplus \ldots \oplus \nicefrac{\mathbb{Z}}{(p_k^{n_k})}
	\]
	Weiter gibt es dann $n$ mit $A \cong \mathbb{Z}^n \oplus TA$. $\Rg A := n$ ist der \bet{Rang} von $A$\index{Rang einer abelschen Gruppe}.
\end{erinnerungA}

\begin{lemma}[{name=[Rang einer endlich erzeugten abelschen Gruppe]}]
	Sei $A$ eine endlich erzeugte abelsche Gruppe. Dann gilt 
	\[
		\Rg A = \max \set[\big]{m \given \exists C \le A, C \cong \mathbb{Z}^m}.
	\]
\end{lemma}
\begin{beweis}
	Da $A \cong \mathbb{Z}^n \oplus T A$ ist, genügt zu zeigen: Ist $\varphi \colon \mathbb{Z}^m \to \mathbb{Z}^n \oplus T A$ ein injektiver Gruppenhomomorphismus, so ist 
	$m \le n$. Schreibe $\varphi = \varphi_0 \oplus \varphi_1$, also $\varphi(a) = \varphi_0(a) + \varphi_1(a)$ mit $\varphi_0 \colon \mathbb{Z}^m \to \mathbb{Z}^n$, 
	$\varphi_1 \colon \mathbb{Z}^m \to T A$. 

	Behauptung: $\varphi_0$ ist injektiv. Zu $v \in \mathbb{Z}^m$ wähle $k\ge 1$ mit $k \cdot \varphi_1(v) = 0 \in T A$. Dann ist $\varphi_1(k v) = k \varphi_1(v)=0$.
	Ist $v \not= 0$, so ist $k \cdot v \not= 0 \in \mathbb{Z}^m$, also ist $\varphi(k v) \not= 0$ und damit $\varphi_0(k v) \not= 0$. Es folgt $\varphi_0(v) \not= 0$, da sonst
	$\varphi_0(k \cdot v) = k \cdot \varphi_0(v)= 0$. Damit ist die Behauptung gezeigt.

	Gruppenhomomorphismen $\varphi_0 \colon \mathbb{Z}^m \to \mathbb{Z}^n$ werden durch $n \times m$-Matrizen beschrieben: Es gibt $A= (a_{ij}) \in \mathbb{Z}^{n \times m}$ mit
	\[
		\varphi_0 \begin{pmatrix}
			z_1 \\ \vdots \\ z_m
		\end{pmatrix} = \begin{pmatrix}
			\sum_{j=1}^{m} a_{1j} \cdot z_j \\
			\vdots \\
			\sum_{j=1}^{m} a_{nj} \cdot z_j
		\end{pmatrix}
	\]
	Ist $m >n$, so hat $A$, aufgefasst als Matrix über $\mathbb{Q}$, einen Kern. Es gibt also einen Vektor $w = \begin{psmallmatrix}a_1/b_1 \\ : \\ a_m/b_m\end{psmallmatrix}$
	mit $A w = 0$, $w \not=0$. Dann ist $(b_1, \ldots , b_m) \cdot w \in \mathbb{Z}^m$ und 
	\[
		\varphi_0 \enbrace[\big]{(b_1, \ldots , b_m)\cdot w} = A \cdot \enbrace[\big]{(b_1, \ldots , b_m) \cdot w} = (b_1, \ldots , b_m) \cdot A \cdot w = 0
	\]
	Dies ist ein Widerspruch zur Injektivität von $\varphi_0$. Also gilt $m \le n$.
\end{beweis}

\begin{definition}[{name=[Kurze exakte Sequenz]}]
	Seien $M_0, M_1, M_2$ $R$-Moduln und $f_0 \colon M_0 \to M_1$, $f_1 \colon M_1 \to M_2$ $R$-lineare Abbildungen. Dann heißt 
	\begin{equation*}
		\begin{tikzcd}
			M_0 \rar["f_0", hook] & M_1 \rar["f_1", two heads]& M_2 
		\end{tikzcd} \label{eq:43} \tag{*}
	\end{equation*}
	eine \Index{kurze exakte Sequenz}, wenn gilt:
	\begin{multicols}{3}
		\begin{enumerate}[(i)]
			\item $f_0$ ist injektiv,
			\item $\im f_0 = \ker f_1$,
			\item $f_1$ ist surjektiv.
		\end{enumerate}
	\end{multicols}
\end{definition}

Oft sagt man $M_0 \xrightarrow{f_0} M_1 \xrightarrow{f_1} M_2$ ist exakt in $M_1$, wenn $\im f_0 = \ker f_1$ ist. Dann ist \eqref{eq:43} eine kurze exakte Folge, wenn
\[
	\begin{tikzcd}
		0 \rar & M_0 \rar["f_0"] & M_1 \rar["f_1"]& M_2 \rar & 0
	\end{tikzcd}
\]
exakt in $M_0, M_1$ und $M_2$ ist.

\begin{beispiel}
	Es gibt zahlreiche Beispiele für kurze exakte Sequenzen:
	\begin{enumerate}[1)]
		\item \[
			\begin{tikzcd}[row sep=0.3em]
				M_0 \rar["i",hook] & M_0 \oplus M_1 \rar["p",two heads] & M_1 \\
				v_0 \rar["i", mapsto] & (v_0,0) \enspace (v_0,v_1) \rar["p",mapsto] & v_1
			\end{tikzcd}
		\]
		ist eine kurze exakte Folge.
		\item $\begin{tikzcd}
			\mathbb{Z} \rar[hook,"\cdot n"] & \mathbb{Z} \rar[two heads] &  \nicefrac{\mathbb{Z}}{n \mathbb{Z}}
		\end{tikzcd}$
		ist eine kurze exakte Folge.
	\end{enumerate}
\end{beispiel}

\begin{bemerkung}[{name=[Spaltung einer kurzen exakten Sequenz]}]
	Eine \Index{Spaltung} für eine kurze exakte Folge $M_0 \xrightarrow{f_0} M_1 \xrightarrow{f_1} M_2$ ist eine $R$-lineare Abbildung $s \colon M_2 \to M_1$ mit $f_1 \circ s = \id_{M_2}$. In diesem Fall erhalten wir einen Isomorphismus $M_0 \oplus M_2 \to M_1$, $(v_0, v_2) \mapsto f_0(v_0) + s(v_2)$.
	\begin{description}
		\item[Injektivität:] Sei $(v_0,v_2) \in M_0 \oplus M_2$ mit $f_0(v_0) + s(v_2) =0$. Dann gilt
		\[
			0 = f_1 \enbrace[\big]{f_0(v_0) + s(v_2)} = v_2 \implies f_0(v_0)=0 \implies v_0 =0
		\]
		\item[Surjektivität:] Sei $v_1 \in M_1$. Betrachte $v_2 := f_1(v_1)$. Dann ist $v_1 -s(v_2) \in \ker f_1 = \im f_0$, also gibt es $v_0$ in $M_0$ mit
		$f_0(v_0) = v_1 - s(v_1)$. Damit ist $v_1= f_0(v_0)+ s(v_1)$.
	\end{description}
\end{bemerkung}

\begin{bemerkung}[{name=[{Existenz einer Spaltung, wenn $M_2$ frei ist}]}]
	\begin{enumerate}[1)]
		\item Die kurze exakte Sequenz $\begin{tikzcd}
			\mathbb{Z} \rar[hook,"\cdot n"] & \mathbb{Z} \rar[two heads] &  \nicefrac{\mathbb{Z}}{n \mathbb{Z}}
		\end{tikzcd}$
		spaltet nicht.
		\item Ist $\begin{tikzcd}
			M_0 \rar["f_0"] & M_1 \rar["f_1"] & M_2
		\end{tikzcd}$
		eine kurze exakte Folge mit $M_2$ frei, also $M_2 \cong R[S]$, so spaltet die Folge.
		Insbesondere ist $M_1 \cong M_0 \oplus M_2$.
		Dazu wenden wir Lemma \ref{univ_frei_mod} an und erhalten eine $R$-lineare Abbildung $s \colon M_2 \to M_1$, sodass folgendes Diagramm kommutiert
		\[
			\begin{tikzcd}
				\mbox{ } & M_1 \dar["f_1",two heads] \\
				R[S] \urar["s", dashed] \rar["\id"] &  R[S]
			\end{tikzcd}
		\]
		$s$ ist also die gesuchte Spaltung.
		Dieses Vorgehen funktioniert bei jedem projektiven Modul.\index{projektiver Modul}
	\end{enumerate}
\end{bemerkung}

\begin{bemerkung}[{name=[Der Rang ist additiv]},label=bem:47]
	Der Rang für alle endlich erzeugten abelschen Gruppen ist additiv:
	\[
		\Rg(A \oplus B) =\Rg A + \Rg B,
	\]
	da $A \cong \mathbb{Z}^{\Rg A} \oplus T A$, $B \cong \mathbb{Z}^{\Rg B} \oplus TB$, also 
	\[
		A \oplus B \cong \mathbb{Z}^{\Rg A + \Rg B} \oplus \underbrace{T A \oplus T B}_{=T(A \oplus B)}
	\]
\end{bemerkung}

\begin{proposition}[{name=[{Der Rang ist für kurze exakte Folgen endl. erz., abelscher Gruppen additiv}]},label=prop:48]
	Der Rang von endlich erzeugten abelschen Gruppen ist additiv für kurze exakte Folgen: Ist
	\begin{equation*}
		\begin{tikzcd}
			A \rar["i",hook] & B \rar["p", two heads] & C
		\end{tikzcd} \label{eq:48}\tag{$\star$}
	\end{equation*}
	eine kurze exakte Folge von endlich 
	erzeugten abelschen Gruppen, so gilt $\Rg B = \Rg A + \Rg C$.
\end{proposition}
\begin{beweis}
	Spaltet \eqref{eq:48}, so ist $B \cong A \oplus C$ und die Behauptung folgt aus der \autoref{bem:47}.
	Wir können annehmen, dass $C = \mathbb{Z}^n \oplus T C$ gilt. Wir erhalten folgendes Diagramm: 
	\[
		\begin{tikzcd}
			A \rar["i",hook] \dar[equal] & p ^{-1}(\mathbb{Z}^n) \rar["p", two heads] \dar[hook,"i''"]& \mathbb{Z}^n \dar[hook,"i'"]\\
			A \rar["i",hook] & B \rar["p", two heads] \dar[two heads,"p''"]& C \dar[two heads,"p'"] & \text{\eqref{eq:48}}\\
			& \nicefrac{B}{p ^{-1}(\mathbb{Z}^n)} \rar[two heads, hook,"p' \circ p"]& T C
		\end{tikzcd} 
	\]
	Ist $p' \circ p(b)=0$, so gibt es $v \in \mathbb{Z}^n$ mit $i'(v)= p(b)$. Nun gibt es $v' \in p ^{-1}(\mathbb{Z}^n)$ mit $p(v')=v$. Es folgt 
	\[
		p \enbrace[\big]{i''(v')} = p(b) 
	\]
	also $i''(v') -b \in \ker p = i(A)$. Da $i(A) \subseteq i'' \enbrace*{p ^{-1} (\mathbb{Z}^n)}$, folgt $b \in i'' \enbrace*{p ^{-1}(\mathbb{Z}^n)}$.
	In dem Diagramm sind die Spalten und Zeilen exakt. Da $\mathbb{Z}^n$ frei ist, spaltet 
	$\begin{tikzcd}
		A \rar[hook,"i"] & p^{-1}(\mathbb{Z}^n) \rar["p", two heads] & \mathbb{Z}^n
	\end{tikzcd}$
	und es gilt 
	\[
		\Rg A + \Rg C = \Rg A + n = \Rg \enbrace*{p ^{-1}(\mathbb{Z}^n)} 
	\]
	Es bleibt zu zeigen: $\Rg B = \Rg  \enbrace*{p ^{-1}(\mathbb{Z}^n)}$. \medskip \\
	Nebenbei: Ist $A' \subseteq B'$, so gilt $\Rg A' \le \Rg B'$, denn 
	\[
		\Rg A' = \max \set*{n' \given \mathbb{Z}^{n'} \cong C' \le A'} \le \max \set*{n' \given \mathbb{Z}^{n'} \cong C' \le B'} = \Rg B'.
	\]
	Betrachte die kurze exakte Folge
	\[
		\begin{tikzcd}
			p^{-1}(\mathbb{Z}^n) \rar & B \rar & \nicefrac{B}{p ^{-1}(\mathbb{Z}^n)}
		\end{tikzcd}
	\]
	Wegen $\nicefrac{B}{p ^{-1}(\mathbb{Z}^n)} \cong T C$ gibt es ein
	$k$ mit $k \cdot \enbrace*{b+ p ^{-1}(\mathbb{Z}^n)} =0$ für alle $b \in B$. Also $k \cdot b \in p ^{-1}(\mathbb{Z}^n)$. Es folgt 
	$kB \subseteq  p ^{-1}(\mathbb{Z}^n)$. Ist  $B \cong \mathbb{Z}^{\Rg B}+ TB$, so ist $k B \cong (k \mathbb{Z})^{\Rg B}+ T(kB)$ und $\Rg kB = n = \Rg B$. 
	Mit der Nebenbemerkung folgt $\Rg p ^{-1}(\mathbb{Z}^n) = \Rg B$ aus $kB \subseteq  p ^{-1}(\mathbb{Z}^n) \subseteq B$.
\end{beweis}

Man kann die Proposition auch mittels Tensorprodukten beweisen.
Dazu geht man wie folgt vor:
\begin{itemize}
	\item $\mathbb{Q} \otimes_\mathbb{Z} \colon \mathbb{Z}\text{-}\MOD \to \mathbb{Q}\text{-}\VR $ ist ein Funktor.
	\item $\Rg A= \dim_\mathbb{Q} (\mathbb{Q} \otimes_\mathbb{Z} A)$
	\begin{enumerate}[(a)]
		\item $\mathbb{Q} \otimes_\mathbb{Z} (X \oplus Y) = \mathbb{Q} \otimes_\mathbb{Z} X \oplus  \mathbb{Q} \otimes_\mathbb{Z} Y$
		\item $\mathbb{Q} \otimes_\mathbb{Z} \mathbb{Z}^n \cong \mathbb{Q}^n$
		\item $\mathbb{Q} \otimes_\mathbb{Z} T \cong 0$ für $T$ ein Torsionsmodul.
	\end{enumerate}
	\item Ist $\begin{tikzcd}[cramped,sep=small]
		A \rar & B \rar & C
	\end{tikzcd}$ eine kurze exakte Folge, so ist auch 
	\(
		\begin{tikzcd}[cramped,sep=small]
				\mathbb{Q} \otimes_\mathbb{Z} A \rar & \mathbb{Q} \otimes_\mathbb{Z} B \rar & \mathbb{Q} \otimes_\mathbb{Z} C
			\end{tikzcd}
	\) 
	eine kurze exakte Folge.
	\item Jede kurze exakte Folge $\begin{tikzcd}
		V_0 \rar & V_1 \rar & V_2
	\end{tikzcd}$ von $\mathbb{Q}$-Vektorräumen spaltet und daher gilt $V_1 \cong V_0 \oplus V_2$ und $\dim_\mathbb{Q}(V \oplus W) = \dim_\mathbb{Q}(V)+ \dim_\mathbb{Q} (W)$.
\end{itemize}

\begin{definition}[{name=[{Eulercharakteristik von Z-Kettenkomplexen}]}]
	Sei $(C_*,d_*)$ ein endlich erzeugter $\mathbb{Z}$-Kettenkomplex.
	\[
		\chi(C_*) := \sum_{i=0}^{\infty} \Rg(C_i) = \sum_{i=0}^{n_0} \Rg(C_i) 
	\]
	heißt die \bet{Eulercharakteristik}\index{Euler-Charakteristik!von $\mathbb{Z}$-Kettenkomplexen} von $(C_*,d_*)$. 
	Dabei ist $n_0$ so gewählt, dass $C_n =0$ für alle $n >n_0$.
\end{definition}

\begin{satz}[{name=[Eulercharakteristik eines endlich erzeugten Z-Kettenkomplexes]},label=satz:410]
	Sei $(C_*,d_*)$ ein endlich erzeugter $\mathbb{Z}$-Kettenkomplex. Dann gilt 
	\[
		\chi(C_*) = \sum_{i=1}^{\infty} (-1)^i \Rg H_i(C_*,d_*).
	\]
\end{satz}
\begin{beweis}
	Sei $B_{n+1} := \im d_{n+1} \colon C_{n+1} \to C_n$ und $Z_n := \ker d_n \colon C_n \to C_{n-1}$. Also $H_n(C_*,d_*) = \nicefrac{Z_n}{B_{n+1}}$.
	Als Untermoduln von $C_n$ sind $Z_n$ und $B_n$ endlich erzeugt (LA2). Insbesondere ist auch $H_n(C_*,d_*)$ endlich erzeugt und der Rang somit definiert. Auch $B_{n+1}$
	ist Untermodul des endlich erzeugten $\mathbb{Z}$-Moduls $C_n$ und somit endlich erzeugt.
	Wir erhalten kurze exakte Folgen:
	\[
		\begin{tikzcd}[row sep=0.4em]
			B_{n+1} \rar[hook] & Z_n \rar[two heads] & H_n(C_*, d_*) \\
			Z_n \rar[hook] & C_n \rar[two heads] & B_n \quad 
		\end{tikzcd}
	\]
	Nach \autoref{prop:48} gilt $\Rg Z_n = \Rg B_{n+1} + \Rg H_n(C_*,d_*)$ und $\Rg C_n = \Rg Z_n + \Rg B_n$. Also gilt
	\begin{align*}
		\chi(C_*) = \sum_{i=0}^{\infty} (-1)^i \Rg C_i &= \sum_{i=0}^{\infty} (-1)^i \enbrace[\big]{\Rg Z_i + \Rg B_i} \\  
		&= \sum_{i=0}^{\infty} (-1)^i \enbrace[\Big]{\Rg H_i(C_*,d_*) + \Rg B_{i+1} + \Rg B_i} \\
		&= \sum_{i=0}^{\infty} (-1)^i \Rg H_i(C_*,d_*) + \sum_{i=0}^{\infty} (-1)^i \Rg B_{i+1} + \sum_{i=0}^{\infty}(-1)^i \Rg B_i \\
		&= \sum_{i=0}^{\infty} (-1)^i \Rg H_i(C_*,d_*)  \qedhere  
	\end{align*}
\end{beweis}

\begin{korollar}[{name=[Die Eulercharakteristik eines simplizialen Komplexes]}]
	Sei $K$ ein endlicher, geordneter simplizialer Kettenkomplex. 
	Dann gilt\marginnote{Eine vereinfachte Variante dieser Aussage kam in Aufgabe 2 von Blatt3 vor: Dort haben wirden Homologiemodul über einen Körper betrachtet. Deshalb konnte wir da dann die Dimensionsformeln anwenden, um die beiden Vorraussetzungen für die Rechnung in \ref{satz:410} 
	zu erhalten.}
	\[
		\chi(K) = \sum_{i=1}^{\dim K} (-1)^i \Rg H_i(K;\mathbb{Z})
	\]
\end{korollar}
\begin{beweis}
	Wende den \autoref{satz:410} auf den simplizialen Kettenkomplex von $K$ an. Da 
	\[
		\Rg C_n(K;\mathbb{Z}) = \Rg \mathbb{Z}[\Sigma_n] = \# \text{ $n$-Simplizes in $K$}
	\]
	ist
	\[
		\chi(K) = \chi \enbrace[\big]{C_*(K;\mathbb{Z}), \partial_*} = \sum_{i=0}^{\infty} (-1)^i \Rg H_i \enbrace[\big]{C_*(K;\mathbb{Z}), \partial_*}  =
		\sum_{i=0}^{\infty} (-1)^i \Rg H_i \enbrace*{K;\mathbb{Z}} \qedhere
	\]
\end{beweis}

\newpage

\section{Singuläre Homologie} % (fold)
\label{sec:5}

\begin{definition}[{name=[{Singuläre Simplizes und $n$-ter singulärer Kettenmodul}]}]
	Sei $X$ ein topologischer Raum. 
	Sei $S_n(X)$ die Menge aller stetigen Abbildungen $\sigma \colon \abs*{\Delta^n} \to X$. 
	Elemente von $S_n(X)$ heißen \Index{singuläre Simplizes} in $X$. 
	Sei $R$ ein Ring. 
	Der \bet{$\mathbf{n}$-te (singuläre) Kettenmodul}\index{n-ter Kettenmodul@$n$-ter Kettenmodul!singulär} von $X$ über $R$ ist
	\[
		C_n(X;R) := R \benbrace[\big]{S_n(X)} 
	\]
\end{definition}

\begin{definition}[{name=[{Einschränkung eines singulären Simplizes auf eine Seite}]},label=def:52]
	Für $\sigma \in S_n(X)$ und $j \in \set{0,\ldots ,n}$ sei $\delta_j \sigma \in S_{n-1}(X)$ die Einschränkung von $\sigma$ auf die $j$-te Seite von $\abs*{\Delta^n}$, also $\delta_j(\sigma) = \sigma \circ \iota_{n,j}$, wobei $\iota_{n,j} \colon \abs{\Delta^{n-1}} \to \abs*{\Delta^n}$ die Inklusion der $j$-ten Seite ist:
	\[
		\iota_{n,j}(x_0, \ldots ,x_{n-1}) = \enbrace[\big]{x_0, \ldots , x_{j-1}, 0, x_j, \ldots , x_{n-1}}
	\]
\end{definition}

\begin{bemerkung}[{name=[{Inklusion der $j$-ten Seite}]}]
	Es gilt $\iota_{n,j} = \abs*{i_{n,j}}$ wobei $i_{n,j} \colon \Delta^{n-1} \to \Delta^n$ gegeben ist durch \marginnote{vergleiche Blatt 2, Aufgabe 2}
	\[
		i_{n,j}(k) = \begin{cases}
		k, &\text{ falls }k < j\\
		k+1, &\text{ falls } k\ge j
	\end{cases}
	\]
\end{bemerkung}

\begin{definition}[{name=[{$n$-te singuläre Randabbildung}]}]
	Die \bet{$\mathbf{n}$-te singuläre Randabbildung}\index{n-te Randabbildung@$n$-te Randabbildung!singulär} $\partial_n \colon C_n(X;R) \to C_{n-1}(X;R)$ ist 
	definiert durch 
	\[
		\partial_n(\sigma) = \sum_{j=0}^{n} (-1)^j \cdot \delta_j(\sigma).
	\]
\end{definition}

\begin{proposition}[{name=[{Für die Randabbildungen gilt $\partial_{n-1} \circ \partial_n = 0$}]},label=prop:55]
	Es gilt
	\[
		\partial_{n-1} \circ \partial_n = 0
	\]
\end{proposition}
\begin{beweis}
	Wir rechnen die Behauptung einfach für $\sigma \colon \abs*{\Delta^n}\to X$ nach:
	\begin{align*}
		\partial_{n-1} \circ \partial_n(\sigma) = \partial_{n-1} \enbrace*{\sum_{j=0}^{n} (-1)^j \delta_j(\sigma)} &= 
		\sum_{k=0}^{n-1} \sum_{j=0}^{n} (-1)^{j+k} \delta_k \delta_j(\sigma)   \\
		&= \,\sum_{0 \le k < j \le n} \!\!\! (-1)^{k+j} \delta_k \delta_j(\sigma) +\!\!\!\smashoperator[r]{\sum_{0 \le j \le k \le n-1}} \enspace(-1)^{k+j} \delta_k\delta_j(\sigma) \\
		&\StackText{\ref{lem:56}}{=} \, \sum_{0 \le k < j \le n} \!\!\!(-1)^{k+j} \delta_k \delta_j(\sigma) + \!\!\!\smashoperator[r]{\sum_{0 \le j \le k \le n-1}} \enspace(-1)^{k+j} \delta_j \delta_{k+1}(\sigma) \\
		&= \,\sum_{0 \le k < j \le n} \!\!\! (-1)^{k+j} \delta_k \delta_j(\sigma) + \smashoperator[r]{\sum_{0 \le j < k \le n}} \enspace(-1)^{k-1+j} \delta_j \delta_k(\sigma) \\
		&= 0 \qedhere
	\end{align*}
\end{beweis}

\begin{lemma}[{name=[Hilfslemma für Proposition \ref{prop:55}]},label=lem:56]
	Für $0 \le j \le k \le n-1$ und $\sigma \in S_n(X)$ ist $\delta_k \delta_j(\sigma) = \delta_j \delta_{k+1}(\sigma)$.
\end{lemma}
\begin{beweis}
	Es gilt
	\begin{align*}
		\delta_k \enbrace[\big]{\delta_j(\sigma)} \enbrace*{x_0, \ldots , x_{n-2}} &= \delta_j(\sigma) \enbrace[\big]{x_0, \ldots , x_{k-1}, 0, x_k , \ldots ,x_{n-2}}  \\
		&= \enbrace[\big]{x_0, \ldots , x_{j-1},0,x_j, \ldots , x_{k-1},0,x_k, \ldots, x_{n-2}}  \\
		\intertext{und}
		\delta_j \delta_{k+1}(\sigma) \enbrace*{x_0, \ldots , x_{n-2}} &= \delta_{k+1}(\sigma) \enbrace[\big]{x_0, \ldots , x_{j-1},0,x_j,\ldots , x_{n-2}} \\
		&= \enbrace[\big]{x_0, \ldots , x_{j-1},0,x_j, \ldots , x_{k-1},0, x_k, \ldots , x_{n-2}} \qedhere
	\end{align*}
\end{beweis}

\begin{definition}[{name=[{Singuläre Homologie von $X$}]}]
	Sei $X$ ein topologischer Raum. 
	Die Homologie des \bet{singulären Kettenkomplex}\index{singulärer Kettenkomplex} über $R$, $\enbrace[\big]{C_*(X;R), \partial_*}$, heißt die \Index{singuläre Homologie von $X$} mit Koeffizienten in $R$:
	\[
		H_n(X;R) := H_n \enbrace[\big]{C_*(X;R), \partial_*} 
	\]  
	Für $R=\mathbb{Z}$ schreiben wir kürzer $C_*(X) := C_*(X;\mathbb{Z})$ und $H_n(X) := H_n(X;\mathbb{Z})$.
	\begin{itemize}
		\item Die Elemente von $C_n(X;R)$ heißen \bet{$\mathbf{n}$-Ketten}\index{n-Ketten@$n$-Ketten}.
		\item Die Elemente von $\im \partial_{n+1} \subseteq C_n(X;R)$ heißen \bet{$\mathbf{n}$-Ränder}\index{n-Ränder@$n$-Ränder}.
		\item Die Elemente von $\ker \partial_{n} \subseteq C_n(X;R)$ heißen \bet{$\mathbf{n}$-Zykel}\index{n-Zykel@$n$-Zykel}.
	\end{itemize}
	Jeder $n$-Zykel $\sigma \in C_n(X;R)$ bestimmt eine \Index{Homologieklasse} $[\sigma] := \sigma + \im \partial_{n+1} \in H_n(X;R)$.
\end{definition}

\begin{beispiel}[{name=[Die Homologie des Ein-Punkt-Raumes]},label=bsp:58]
	Ist $X= \set{x_0}$ der Ein-Punkt-Raum, so ist 
	\[
		H_* \enbrace[\big]{\set{x_0};R} \cong \begin{cases}
		R, &\text{ falls }*=0\\
		0 , &\text{ sonst } 
	\end{cases} 
	\]
\end{beispiel}
\begin{beweis}
	Es gibt für jedes $n$ genau eine Abbildung $\sigma_n \colon \abs*{\Delta^n} \to \set{x_0}$ nämlich die konstante. 
	Diese ist offensichtlich stetig. 
	Also ist $C_n \enbrace*{\set{x_0};R}=R[\sigma_n]$. 
	Für alle $j$ ist $\delta_j \sigma_n = \sigma_{n-1}$. Daher ist
	\[
		\partial_n(\sigma_n) = \begin{cases}
			0, &\text{ falls }n \text{ ungerade}\\
			\sigma_{n-1} , &\text{ falls } n \text{ gerade}
		\end{cases}
	\]
	Der singuläre Kettenkomplex von $\set{x_0}$ hat also folgende Gestalt:
	\[
		\begin{tikzcd}[row sep=0.4em]
			C_0 \enbrace[\big]{\set{x_0};R} \arrow[draw=none]{d}[sloped,auto=false]{\cong} 
			& C_1 \enbrace[\big]{\set{x_0};R} \lar["0", swap] \arrow[draw=none]{d}[sloped,auto=false]{\cong} & 
			C_2 \enbrace[\big]{\set{x_0};R} \lar["\cong", swap] \arrow[draw=none]{d}[sloped,auto=false]{\cong} 
			& C_3 \enbrace[\big]{\set{x_0};R} \lar["0", swap] \arrow[draw=none]{d}[sloped,auto=false]{\cong}
			& \cdots \lar["\cong",swap] \arrow[draw=none]{d}[sloped,auto=false]{\cong} \\
			R & R \lar["0"] & R \lar["\id"] & R \lar["0"] & R \lar["\id"]
		\end{tikzcd}
	\]
	Es folgt
	\[
		H_n \enbrace[\big]{\set{x_0};R } \cong \begin{cases}
			R, &\text{ falls }n=0\\
			0 , &\text{ sonst}
		\end{cases} \qedhere
	\]
\end{beweis}

\begin{proposition}[{name=[{Eigenschaften von $H_0(X)$ für $X \not= \emptyset$ und $X$ wegzusammenhängend}]}]
	\leavevmode\begin{enumerate}[1),itemsep=0pt]
		\item Ist $X \not= \emptyset$, so ist $H_0(X) \not=0$.\marginnote{Eine analoge Aussage für simpliziale Komplexe haben wir in Aufgabe 3 von Blatt 2 bewiesen}
		\item Ist $X$ wegzusammenhängend, so gilt $H_0(X) \cong \mathbb{Z}$.\footnote{Tatsächlich gilt auch die Umkehrung, wie in Aufgabe 4 von Blatt 3 gezeigt.}
	\end{enumerate}
\end{proposition}
\begin{beweis}
	\begin{enumerate}[1)]
		\item Sei die sogenannte \Index{Augmentierung} $\epsilon \colon C_0(X) \to \mathbb{Z}$ definiert durch 
		\[
			\epsilon \enbrace*{\sum_{\sigma \in S_0(X)} \hspace{-1ex} r_\sigma \cdot \sigma } := \sum_\sigma r_\sigma.
		\]
		Dann ist $\epsilon \circ \partial_1 \colon C_1(X) \to \mathbb{Z}$ trivial, denn für $\sigma \in S_1(X)$ ist $\epsilon \circ \partial_1(\sigma)= \epsilon(\delta_0 \sigma - \delta_1 \sigma) = 0$.
		Daher induziert $\epsilon$ eine Abbildung $\overline{\epsilon} \colon H_0(X) \to \mathbb{Z}$. 
		Da $X \not= \emptyset$ ist, gibt es einen singulären $0$-Simplex $\sigma \colon \abs*{\Delta^0} \to X$. 
		Für $\sigma$ gilt $\overline{\epsilon}( [\sigma])=1$ und daher ist $\overline{\epsilon}$ surjektiv und $H_0(X) \not= 0$.
		\item Wir zeigen, dass $\overline{\epsilon} \colon H_0(X) \to \mathbb{Z}$ ein Isomorphismus ist, falls $X$ wegzusammenhängend ist. 
		Dazu zeigen wir $\ker \epsilon = \im \partial_1$. 
		Dazu definieren wir $s \colon C_0(X) \to C_1(X)$ wie folgt: Sei $\sigma_0$ ein fest gewählter singulärer $0$-Simplex.
		Zu $\sigma \in S_0(X)$ gibt es, da $X$ wegzusammenhängend ist, $s(\sigma) \in S_1(X)$ mit $\delta_0 s(\sigma)= \sigma$, $\delta_1 s(\sigma)= \sigma_0$.
		Es gilt nun 
		\begin{align*}
			\partial_1 \circ s(\sigma) = \sigma - \sigma_0 = \enbrace*{\id_{C_0(X)} - i \circ \epsilon }(\sigma)
		\end{align*}
		wobei $i \colon \mathbb{Z} \to C_0(X)$, $n \mapsto n \cdot \sigma_0$. 
		Ist nun $v \in \ker \epsilon$, so folgt 
		\[
			\partial_1 \circ s(v)= v- i \circ \underbracket[1pt][2pt]{\epsilon(v)}_{=0} = v
		\]
		also ist $v \in \im \partial_1$ und es folgt die Behauptung. \qedhere
	\end{enumerate}
\end{beweis}

\begin{bemerkung}[{name=[{Vorgehen, um später zu zeigen: $H_*(K;R) \cong H_*(\abs*{K};R)$}]}]
	Sei $K=(V,\Sigma)$ ein geordneter simplizialer Komplex. 
	Sei $\sigma= \set{v_o, \ldots , v_n} \in \Sigma_n $ ein $n$-Simplex von $K$ mit $v_0 < v_1 < \ldots < v_n$. 
	Wir ordnen $\sigma$ den singulären $n$-Simplex $f_\sigma \colon \abs*{\Delta^n} \to \abs*{\Delta^\sigma} \subseteq \abs*{K}$ zu, wobei
	\begin{align*}
		f_\sigma \enbrace*{\sum_{i=0}^{n} t_i \cdot e_i} = \sum_{i=0}^{n} t_i \cdot v_i. 
	\end{align*}
	Nun erhalten wir eine Abbildung 
	\[
		C_*(K;R) \to C_*(\abs*{K},R) \quad , \qquad   \sum_{\sigma \in \Sigma_n} r_\sigma \cdot \sigma \longmapsto \sum_{\sigma \in \Sigma_n} r_\sigma \cdot f_\sigma.
	\]
	Diese Abbildung induziert(!) eine weitere Abbildung $H_*(K;R) \xrightarrow{\cong} H_*(\abs*{K};R )$ von der wir später zeigen werden, dass sie ein Isomorphismus ist.
\end{bemerkung}

\newpage

\section{Funktorialität} % (fold)
\label{sec:6}

\begin{definition}[{name=[{$R$-Kettenabbildung}]},label=def:61]
	Seien $(C_*,d_*)$ und $(C_*',d_*')$ zwei $R$-Kettenkomplexe. 
	Dann ist eine \bet{$\mathbf{R}$-Kettenabbildung}\index{R-Kettenabbildung@$R$-Kettenabbildung} $f_* \colon (C_*,d_*) \to (C_*',d_*')$ eine Folge von $R$-linearen Abbildungen $f_n \colon C_n \to C_n'$, sodass
	\[
		d_n' \circ f_n = f_{n-1} \circ d_n
	\]
	für alle $n\ge 1$ gilt. Also kommutiert folgendes Diagramm
	\[
		\begin{tikzcd}[sep=normal]
			C_0 \dar["f_0", Firebrick3] & C_1 \lar["d_1"'] \dar["f_1", Firebrick3] & C_2 \lar["d_2"'] \dar["f_2", Firebrick3]& C_3 \lar["d_3"'] \dar["f_3", Firebrick3]& \ldots \lar \\
			C_0' & C_1' \lar["d_1'"]  & C_2' \lar["d_2'"]& C_3' \lar["d_3'"]& \ldots \lar
		\end{tikzcd}
	\]
\end{definition}

\begin{bemerkung}[{name=[{Induzierte Abbildung einer Kettenabbildung auf Homologie}]}]
	Ist $f_* \colon (C_*,d_*) \to (C_*',d_*')$ eine Kettenabbildung, so erhalten wir eine induzierte Abbildung auf Homologie $H_n(f_*) \colon H_n(C_*,d_*) \to H_n(C_*',d_*')$ durch
	\[
		H_n(f_*) \enbrace*{[v]} := \benbrace*{f_n(v)}  
	\]
	für $v \in \ker d_n$. Dies ist wohldefiniert, denn:
	\begin{itemize}
		\item $d_n' \enbrace[\big]{f_n(v)} = f_{n-1} \enbrace[\big]{d_n(v)} = 0 $, also $f_n(v) \in \ker d_n'$ für $v \in \ker d_n$.
		\item Ist $[v]=[w]$, mit $v,w \in\ker d_n$, so gibt es $x \in C_{n+1}$ mit $d_{n+1}(x)=v-w$. Dann ist 
		\[
			d_{n+1}'\enbrace[\big]{f_{n+1}(x)} = f_n \enbrace[\big]{d_{n+1}(x)} = f_n(v)- f_n(w) 
		\]
		also $\benbrace*{f_n(v)} = \benbrace*{f_n(w)} \in H_n(C_*',d_*')$.
	\end{itemize}
\end{bemerkung}

\begin{bemerkung}[{name=[Homologie definiert einen Funktor]},label=bem:63]
	\leavevmode
	\begin{enumerate}[a)]
		\item $R$-Kettenkomplexe mit $R$-Kettenabbildungen bilden die Kategorie $R\text{-}\textsc{Ketten}$.
		\item Homologie definiert nun einen Funktor: $H_n \colon R\text{-}\textsc{Ketten} \to R\text{-}\MOD$.
	
		Es gilt $H_n(\id_{(C_*,d_*)}) = \id_{H_n(C_*,d_*)}$ und $H_n(f \circ g) = H_n(f) \circ H_n(g)$, da
		\[
			H_n(f \circ g) \enbrace*{[v]} = \benbrace*{f_n \circ g_n(v)} = \benbrace*{f_n \enbrace[\big]{g_n(v)} } = H_n(f) \enbrace[\big]{\benbrace*{g_n(v)}} = H_n(f) \enbrace[\big]{H_n(g)(v)}     
		\]
		\item Definiert man die Kategorie der \Index{graduierten} $R$-Moduln als die Kategorie, deren Objekte Folgen $(V_n)_n$ von $\mathbb{R}$-Moduln sind und deren
		Morphismen Folgen von $R$-linearen Abbildungen $(f_n)_{n}$ sind, so kann man die $H_n$, $n \in \mathbb{N}$ zu einem Funktor
		\[
			H_* \colon R\text{-}\textsc{Ketten} \longrightarrow \textsc{Gr}\text{-}R\text{-}\MOD \quad , \qquad (C_*,d_*) \longmapsto \enbrace[\big]{H_n(C_*,d_*)}_n 
		\]
		zusammensetzen.
	\end{enumerate}
\end{bemerkung}

\begin{definition}[{name=[Induzierte Abbildung auf Kettenkomplexen]}]
	Sei $f \colon X \to Y$ stetig. Wir definieren $C_n(f;R) \colon C_n(X;R) \to C_n(Y;R)$ durch
	\[
		C_n(f;R) \enbrace*{\smashoperator[r]{\sum_{\sigma \in S_n(X)}} r_\sigma \cdot  \sigma} := \smashoperator[r]{\sum_{\sigma \in S_n(X)}} r_\sigma \cdot \enbrace*{f \circ \sigma} 
	\]
	Dies ist wohldefiniert, denn für $\sigma \in S_n(X)$ ist $f \circ \sigma \in S_n(Y)$.
\end{definition}

\begin{proposition}[{name=[{Der singuläre Kettenkomplex über $R$ bildet einen Funktor}]},label=prop:65]
	Mit dieser Definition von $C_*(f;R)$ wird der singuläre Kettenkomplex über $R$ zu einem Funktor
	\[
		C_*(-;R) \colon \textsc{Top} \to R\text{-}\textsc{Ketten}
	\]
\end{proposition}
\begin{beweis}
	\leavevmode
	\begin{enumerate}[(i),itemsep=0pt]
		\item $C_*(f;R)$ ist eine $R$-Kettenabbildung: Es gilt
		\begin{align*}
			C_{n-1}(f;R) \circ \partial_n(\sigma) = C_{n-1} (f;R) \enbrace*{\sum_{l=0}^{n}  (-1)^l \cdot \sigma \circ \iota_{n,l}} &=
			\sum_{l=0}^{n} (-1)^l \cdot (f \circ \sigma) \circ \iota_{n,l} \\
			&= \partial_n(f \circ \sigma) = \partial \enbrace[\big]{C_n(f;R)(\sigma)}.
		\end{align*}
		\item Zu zeigen: $C_*(f \circ g;R) = C_*(f;R) \circ C_*(g;R)$. Es gilt
		\begin{align*}
			C_n(f \circ g;R)(\sigma) = f \circ (g \circ \sigma) = C_n(f;R)(g \circ \sigma) = C_n(f;R) \enbrace[\big]{C_n(g;R)(\sigma)}.
		\end{align*}
		\item Zu zeigen: $C_*(\id_X;R) = \id_{C_*(X;R)}$. Es gilt
		\[
			C_n({\id_X};R)(\sigma) = {\id_{X}} \circ {\sigma} = \sigma \qedhere
		\]
	\end{enumerate}
\end{beweis}

\begin{korollar}[{name=[{Homologie ist ein Funktor $\TOP \to R\text{-}\textsc{Mod}$}]}]
	Mit $H_n(f;R) := H_n \enbrace*{C_*(f;R)}$ wird $H_n(-;R)$ zu einem Funktor: $H_n(-;R) \colon \TOP \to R\text{-}\textsc{Mod}$.
\end{korollar}
\begin{beweis}
	$H_n(-;R)$ ist die Komposition des Funktors $C_*(-;R) \colon \TOP \to R\text{-}\textsc{Ketten}$ aus \autoref{prop:65} mit dem Funktor $H_n \colon R\text{-}\textsc{Ketten} \to R\text{-}\textsc{Mod}$ aus \autoref{bem:63}.
\end{beweis}

\begin{bemerkung}[{name=[Notation für die auf Homologie induzierte Abbildung]}]
	Oft schreiben wir kurz $f_* = H_n(f;R)$.
	Für $\benbrace*{\sum_{\sigma \in S_n(X)} r_\sigma \sigma} \in H_n(X;R)$ ist
	\[
		f_* \enbrace*{\benbrace*{\sum\limits_{\sigma \in S_n(X)} r_\sigma \sigma}} =  \benbrace*{\sum\limits_{\sigma \in S_n(X)} r_\sigma
		f\circ \sigma} \in H_n(Y;R)
	\]
\end{bemerkung}

\begin{definition}[{name=[Summe oder Koprodukt von topologischen Räumen]}]
	Seien $X_i$, $i \in I$ topologische Räume. Mit 
	\[
		X := \coprod_{i \in I} X_i 
	\]
	bezeichnen wir die \bet{Summe}\index{Summe!topologischer Räume} (oder auch das \bet{Koprodukt}\index{Koprodukt!topologischer Räume}) der $X_i$. 
	Als Menge ist $X$ die disjunkte Vereinigung der $X_i$. 
	$U \subseteq X$ ist offen genau dann, wenn $X_i \cap U \subseteq X_i$ offen ist für jedes $i \in I$. 
	Für jedes $i_0$ erhalten wir eine stetige Inklusion $j_{i_0} \colon X_{i_0} \to \coprod_{i} X_i$.
	Für jedes $i_0$ ist $X_{i_0} \subseteq \coprod_i X_i$ offen und abgeschlossen.
\end{definition}

\begin{definition}[{name=[{Summe von $R$-Moduln}]}]
	Seien $V_i$, $i \in I$ $R$-Moduln. Mit 
	\[
		V := \bigoplus_{i \in I} V_i
	\]
	bezeichnen wir die \bet{Summe}\index{Summe!von $R$-Moduln} (auch das \bet{Koprodukt}\index{Koprodukt!von $R$-Moduln}) der $V_i$. 
	Elemente von $V$ sind $I$-Folgen $(v_i)_{i \in I}$ mit $v_i \in V_i$ und $v_i=0$ füralle bis auf endlich viele $i$. 
	Die $R$-Modulstruktur ist erklärt durch: $(v_i)_{i \in I} + (w_i)_{i \in I} := (v_i +w_i)_{i \in I}$ und $r \cdot (v_i)_{i \in I} := (r \cdot v_i)_{i \in I}$
	% \[
	% 	(v_i)_{i \in I} + (w_i)_{i \in I} := (v_i +w_i)_{i \in I} \qquad ,\qquad r \cdot (v_i)_{i \in I} := (r \cdot v_i)_{i \in I}
	% \]
	Für jedes $i_0 \in I$ erhalten wir eine $R$-lineare Abbildung $j_{i_0} \colon V_{i_0} \to V$ mit
	\[
		\enbrace*{j_{i_0} (v)}_{i} = \begin{cases}
			v, &\text{ falls }i=i_0\\
			0 ,&\text{ sonst}
		\end{cases} 
	\]
\end{definition}

Seien $V_i$ für $i \in I$ $R$-Moduln. Sei $W$ ein weiterer $R$-Modul. Dann gibt es zu jeder Folge $f_i \colon V_i \to W$ von $R$-linearen Abbildungen eine
$R$-lineare Abbildung 
\(
	\oplus_{i \in I} f_i \colon \bigoplus_{i \in I} V_i \to W  
\)
mit
\[
	\enbrace*{\oplus_{i \in I} f_i} \enbrace[\big]{(v_i)_{i \in I}} = \sum_{i \in I} f_i(v_i)  
\]
Ist umgekehrt $f \colon \bigoplus_{i \in I} V_i \to W$ eine $R$-lineare Abbildung, so ist $f_i := f \circ j_{i}$ eine Folge von $R$-linearen Abbildungen mit 
$f = \bigoplus_{i \in I} f_i$.

\begin{satz}[{name=[{Homologie des Koproduktes topologischer Räume}]}]
	Sei $X = \coprod_{i \in I} X_i$. Dann induzieren die Inklusionen $j_i \colon X_i \hookrightarrow X$ einen Isomorphismus
	\[
		\bigoplus_{i \in I} H_n(X_i;R) \xrightarrow{\bigoplus_{i \in I} (j_i)_*}  H_n(X;R)
	\]
\end{satz}
\begin{beweis}
	Da die $X_i \subseteq X$ offen und abgeschlossen sind und $\abs*{\Delta^n}$ zusammenhängend ist (sogar wegzusammenhängend), gibt es für jedes 
	$\sigma \colon \abs*{\Delta^n} \to X $ ein eindeutiges $i$ mit $\im \sigma \subseteq  X_i$. Es gilt also $S_n(X) = {\dot{\bigcup}} S_n(X_i)$. Daher induzieren die $j_i$
	für jedes $n$ einen Isomorphismus
	\[
		{\textstyle\bigoplus\limits_{i \in I}} C_n(j_i;R) \colon \bigoplus_{i \in I}C_n(X_i;R) \xrightarrow{\enspace\cong \enspace} C_n(X;R) 
	\]
	Da diese Isomorphismen mit den Randabbildungen vertauschen, erhalten wir einen Isomorphismus von $R$-Kettenkomplexen
	\[
		\varphi := {\textstyle\bigoplus\limits_{i \in I}} C_*(j_i;R) \colon \bigoplus C_*(X_i;R) \xrightarrow{\enspace\cong \enspace} C_*(X;R)
	\]
	Da $\bigoplus_i H_n \enbrace[\big]{C_*(X_i;R)} \cong H_n \enbrace[\big]{\bigoplus_i C_*(X_i;R)}$ induzieren dann auch die $(j_i)_*$ einen Isomorphismus
	\[
		{\textstyle\bigoplus\limits_{i \in I}} (j_i)_* \colon \bigoplus_{i \in I} H_n(X_i;R) \xrightarrow{\enspace\cong \enspace} H_n(X;R) \qedhere
	\]
\end{beweis}

\newpage

\section{Homotopieinvarianz} % (fold)
\label{sec:7}

\begin{bemerkung}[{name=[Die induzierte Abbildung eines Homöomorphismus ist ein Isomorphismus]}]
	Sei $f \colon X \to Y$ ein Homöomorphismus. Dann ist 
	\(
		f_* \colon H_n(X;R) \longrightarrow H_n(Y;R)
	\)
	ein Isomorphismus.
\end{bemerkung}
\begin{beweis}
	Da Homologie ein Funktor ist, gilt
	\begin{align*}
		f_* \circ (f^{-1})_* &= (f \circ f ^{-1})_* = (\id_Y)_* = \id_{H_n(Y;R)} \qquad \text{ und} \\
		(f ^{-1})_* \circ f_* &= (f ^{-1} \circ f)_* = (\id_X)_* = \id_{H_n(X;R)}
	\end{align*}
	Also ist $(f_*)^{-1}=(f ^{-1})_*$. Insbesondere ist $f_*$ ein Isomorphismus.
\end{beweis}

\begin{definition}[{name=[Homotopieäquivalenz und Homotopieinverse]}]
	Eine stetige Abbildung $f \colon X \to Y$ heißt eine \Index{Homotopieäquivalenz}, falls es eine stetige Abbildung $g \colon Y \to X$ gibt, so dass $g \circ f$ homotop zu $\id_X$ ist und $f \circ g$ homotop zu $\id_Y$ ist. 
	$g$ heißt dann eine \Index{Homotopieinverse} zu $f$.
\end{definition}
\begin{enumerate}[(i)]
	\item $f$ ist genau dann eine Homotopieäquivalenz, wenn $[f]$ in $\HTOP$ invertierbar ist. In $\HTOP$ ist dann $[f]^{-1} = [g]$.
	\item Die Homotopieinverse ist eindeutig bis auf Homotopie.
\end{enumerate}

\begin{satz}[{name=[Homotopieinvarianz der Homologie]},label=satz:73]
	Seien $f,g \colon X \to Y$ stetige Abbildungen. 
	Sind $f$ und $g$ homotop, so gilt für die induzierten Abbildungen $H_n(f;R) = H_n(g;R)$ für alle $n \in \mathbb{N}$.
\end{satz}
\begin{beweis}
	\emph{Siehe \autoref{lem:78}, sowie \autoref{bem:79}.}
\end{beweis}

\begin{korollar}[{name=[Die induzierte Abbildung einer Homotopieäquivalenz ist ein Isomorphismus]},label=kor:74]
	Ist $f \colon X \to Y$ eine Homotopieäquivalenz, so ist $f_* \colon H_n(X;R) \to H_n(Y;R)$ ein Isomorphismus.
\end{korollar}
\begin{beweis}
	Sei $g \colon Y \to X$ ein Homotopieinverses zu $f$. Es folgt 
	\begin{align*}
		f_* \circ g_* &\stackrel{\text{Funktor}}{=} (f \circ g)_* \stackrel{\text{Homotopieinvarianz}}{=} (\id_Y)_* \stackrel{\text{Funktor}}{=} \id_{H_n(Y;R)} \\
		g_* \circ f_* &\stackrel{\text{Funktor}}{=} (g \circ f)_* \stackrel{\text{Homotopieinvarianz}}{=} (\id_X)_* \stackrel{\text{Funktor}}{=} \id_{H_n(X;R)}\qedhere
	\end{align*}
\end{beweis}

\begin{korollar}[{name=[Homologie eines kontrahierbaren Raumes]},label=kor:75]
	Ist $X$ \Index{kontrahierbar}, d.h. es gibt eine Homotopie $H \colon X \times [0,1] \to X$ mit $H_0 =\id_X$ und $H_1$ konstant, so gilt 
	\[
		H_n(X;R) \cong \begin{cases}
			R, &\text{ falls }n=0\\
			0, &\text{ sonst}
		\end{cases}
	\]
\end{korollar}
\begin{beweis}
	Sei $\set{x_0} =\im H_1$. Dann ist die Inklusion $\set{x_0} \to X$ eine Homotopieäquivalenz. Also gilt nach \autoref{kor:74}
	\[
		H_n(X;R) \cong H_n \enbrace[\big]{\set{x_0};R} \StackText{\ref{bsp:58}}{=} \begin{cases}
			R, &\text{ falls }n=0\\
			0, &\text{ sonst}
		\end{cases} \qedhere
	\]
\end{beweis}

Es gilt beispielsweise
\[
	H_n(\mathbb{R}^k;R) \cong \begin{cases}
		R, &\text{ falls }n=0\\
		0, &\text{ sonst}
	\end{cases} \qquad \text{ und } \qquad 
	H_n(D^k;R) \cong \begin{cases}
		R, &\text{ falls }n=0\\
		0, &\text{ sonst}
	\end{cases}
\]

\begin{definition}[{name=[Kettenhomotopie]},label=def:76]
	Seien $(C_*,d_*)$ und $(C_*',d_*')$ $R$-Kettenkomplexe. 
	Seien $f_*,g_* \colon (C_*,d_*) \to (C_*',d_*')$ $R$-Kettenabbildungen.\marginnote{Kettenhomotopie definiert eine Äquivalenzrelation wie in Aufgabe 2 von Blatt 5 gezeigt.}
	Eine \Index{Kettenhomotopie} von $f_*$ nach $g_*$ ist eine Folge von $R$-linearen Abbildungen $h_n \colon C_n \to C_{n+1}'$, $n \in \mathbb{N}$, sodass für alle $n$ gilt:
	\[
		d_{n+1}' \circ h_n + h_{n-1} \circ d_n = f_n - g_n
	\] 
	In diesem fall heißen $f_*$ und $g_*$ \Index{kettenhomotop}.
\end{definition}
\begin{figure}[th]
	\centering{
	\begin{tikzcd}[row sep=2.5em, column sep=3em]
		C_{n-1} \drar["h_{n-1}"'] & C_n \lar["d_n"'] \dar[shift left,"f_n"] \dar[shift right,"g_n"'] \drar["h_n"] \\
		& C_n' & C_{n+1}' \lar["d_{n+1}'"]
	\end{tikzcd}
	\caption{Diagramm zu \autoref{def:76} (Kettenhomotopie)}}
\end{figure}

\begin{proposition}[{name=[{Induzierte Abbildungen von kettenhomotopen Kettenabbildungen sind gleich}]}]
	Seien $f_*,g_* \colon (C_*,d_*) \to (C_*',d_*')$ Kettenabbildungen. Sind $f_*$ und $g_*$ kettenhomotop, so gilt für alle $n$
	\[
		H_n(f_*) = H_n(g_*)
	\]
\end{proposition}
\begin{beweis}
	Sei $x \in H_n(C_*,d_*)$. Also $x=[v]$ mit $v \in \ker d_n \colon C_n \to C_{n-1}$. Dann gilt, da $v \in \ker d_n$
	\begin{align*}
		H_n(f_*)(x) = H_n(f_*) \enbrace*{[v]} = \benbrace*{f_n(v)}  &= \benbrace[\Big]{g_n(v) + {\underbrace{ d_{n+1}' \circ h_n(v)}_{\in \im d_{n+1}'}} + {\underbrace{h_{n-1} \circ d_n(v)}_{=0}}} \\[-1pt]
		&= \benbrace*{g_n(v)} = H_n(g_*) \enbrace*{[v]} = H_n(g_*)(x) \qedhere
	\end{align*}
\end{beweis}

\begin{lemma}[{name=[{Reduktion von \autoref{satz:73} auf Beweis eines Spezialfalles}]},label=lem:78]
	Die Homotopieinvarianz von $H_n(-;R)$, also \autoref{satz:73}, folgt aus folgenden Spezialfall:
	\begin{equation*}
		\text{Für $i_0,i_1 \colon X \hookrightarrow X \times [0,1]$ gilt $i_0(x)= (x,0)$, $i_1(x)=(x,1)$ $\implies$ $(i_0)_* = (i_1)_*$} \label{eq:78} \tag{\#}
	\end{equation*}
\end{lemma}
\begin{beweis}
	Sei $H \colon X \times [0,1] \to Y$ eine Homotopie zwischen $f,g \colon X \to Y$, also $f = H \circ i_0$ und $g = H \circ i_1$. Dann folgt mittels Funktorialität 
	\[
		f_* = (H \circ i_0)_* = H_* \circ (i_0)_* = H_* \circ (i_1)_* = (H \circ i_1)_* = g_* \qedhere
	\]
\end{beweis}

\begin{bemerkung}[{name=[{Bemerkung zum Vorgehen beim Beweis von \textbf{[\#]}}]},label=bem:79]
	Zum Beweis von \eqref{eq:78} werden wir eine explizite Kettenhomotopie $h_*$ zwischen $C_*(i_0;R)$ und $C_*(i_1;R)$ konstruieren. 
	Wir brauchen also $h_n \colon C_n(X;R) \to C_{n+1}(X \times [0,1];R)$ für alle $n \in \mathbb{N}$ mit 
	\begin{equation*}
		\partial_{n+1} \circ h_n + h_{n-1} \circ \partial_n = C_n(i_0;R)- C_n(i_1;R). \label{eq:79} \tag{\#\#}
	\end{equation*}
\end{bemerkung}

Wir besprechen  zunächst eine möglichen -- aber formal nicht korrekten -- Ansatz:
Für $\sigma \colon \abs*{\Delta^n} \to X $ setze $h_n(\sigma) = \sigma \times \id_{[0,1]} \colon \abs*{\Delta^n} \times [0,1] \to X \times [0,1]$. 
Dann ist \enquote{$h_{n-1}(\partial_n \sigma) = \partial_n \sigma \times \id_{[0,1]}$} und 
\[
	\partial_{n+1}(h_n \sigma) = \partial_{n+1}(\sigma \times \id_{[0,1]}) = (\sigma \times \id)\Big|_{\partial (\abs*{\Delta^n} \times [0,1] )} = \partial_n \sigma \times \id_{[0,1]} + i_0 \sigma + i_1 \sigma
\]
da $\partial \enbrace*{\abs*{\Delta^n} \times [0,1]} = \partial \abs*{\Delta^n} \times [0,1] \,\dot{\cup}|, \abs*{\Delta^n} \times \set{0,1}$. \medskip

\noindent Um daraus Sinn zu machen, zerlegen wir $\abs*{\Delta^n} \times [0,1]$ in eine Vereinigung von $(n+1)$-Simplizes. $\Delta^n \times [0,1]$ hat die folgende Form für
$n \in \set{0,1,2}$:
\begin{figure}[H]
	\centering{
	\scalebox{0.9}{\begin{tikzpicture}[scale=2.3]
		\coordinate (1) at (0,1);
		\coordinate (2) at (0,0);
		
		\coordinate (3) at (1.7,1);
		\path (3) ++ (1,0) coordinate (4) ++ (0,-1) coordinate (6) ++ (-1,0) coordinate (5);
		% \coordinate (4) at (3,1);
		% \coordinate (5) at (2,0);
		% \coordinate (6) at (3,0);
		
		\coordinate (7) at (4.5,1);
		\path (7) ++ (0.9,-0.15) coordinate (8);
		\path (7) ++ (0.5,0.18) coordinate (9);
		\path (7) ++ (0,-1) coordinate (10);
		\path (8) ++ (0,-1) coordinate (11);
		\path (9) ++ (0,-1) coordinate (12);
		
		\draw[thick] (2) -- (1) node[above, yshift=0.4cm]{$\Delta^0$} node[midway, left, xshift=-0.3cm]{$[0,1]$};
		
		\draw[thick] (3) node[left]{$(0,0)$}-- (4) node[midway, above, yshift=0.4cm]{$\Delta^1$} node[right]{$(1,0)$} -- (6) node[right]{$(1,1)$} 
		-- (5) node[left]{$(0,1)$} -- cycle;
		
		\draw[thick] (7) node[left]{$(0,0)$} node[above, yshift=0.4cm, xshift=0.9cm]{$\Delta^2$} -- (8) node[right]{$(2,0)$} -- (9) node[right,yshift=2pt]{$(1,0)$} -- cycle;
		
		\draw[thick] (7) -- (10) node[left]{$(0,1)$} -- (11) node[right]{$(2,1)$} -- (8);
		\draw[thick, dashed] (10) -- (12) node[above left]{$(1,1)$} -- (9);
		\draw[thick, dashed] (12) -- (11);
		
		\foreach \x in {1,...,12}{
			\filldraw (\x) circle(0.03);
		}
	\end{tikzpicture}}
	\caption{$\Delta^n \times [0,1]$ für $n=0,1,2$}}
\end{figure}
\noindent Die Zerlegung in $3$-Simplizes ist für $\Delta^2 \times [0,1]$ gegeben durch
\begin{figure}[H]
	\centering{
	\[
		\vcenter{\hbox{
		\begin{tikzpicture}[scale=2.3]
			\coordinate (7) at (0,1);
			\path (7) ++ (0.9,-0.15) coordinate (8);
			\path (7) ++ (0.5,0.18) coordinate (9);
			\path (7) ++ (0,-1) coordinate (10);
			\path (8) ++ (0,-1) coordinate (11);
			\path (9) ++ (0,-1) coordinate (12);
		
			\draw[thick] (7)  -- (8) -- (9) -- cycle;
			\draw[thick] (7) -- (10)  -- (11) -- (8);
			\draw[thick, dashed] (10) -- (12) -- (9);
			\draw[thick, dashed] (12) -- (11);
		
			\foreach \x in {7,...,12}{
				\filldraw (\x) circle(0.02);
			}
		\end{tikzpicture}
		}}
		\enspace =\enspace
		\vcenter{\hbox{
		\begin{tikzpicture}[scale=2.3]
			\coordinate (7) at (0,1);
			\path (7) ++ (0.9,-0.15) coordinate (8);
			\path (7) ++ (0.5,0.18) coordinate (9);
			\path (7) ++ (0,-1) coordinate (10);
			\path (8) ++ (0,-1) coordinate (11);
			\path (9) ++ (0,-1) coordinate (12);
		
			\draw[thick, gray] (7)  -- (8) -- (9) -- cycle;
			\draw[thick, gray] (7) -- (10)  -- (11) -- (8);
			\draw[thick, dashed, gray] (10) -- (12) -- (9);
			\draw[thick, dashed, gray] (12) -- (11);
			\draw[very thick, Firebrick3] (7) -- (10) -- (11) -- (12) -- (7) -- (11) -- (12) -- (10) -- cycle;
		
			\foreach \x in {7,...,12}{
				\filldraw (\x) circle(0.02);
			}
		\end{tikzpicture}
		}} 
		\enspace\cup \enspace
		\vcenter{\hbox{
		\begin{tikzpicture}[scale=2.3]
			\coordinate (7) at (0,1);
			\path (7) ++ (0.9,-0.15) coordinate (8);
			\path (7) ++ (0.5,0.18) coordinate (9);
			\path (7) ++ (0,-1) coordinate (10);
			\path (8) ++ (0,-1) coordinate (11);
			\path (9) ++ (0,-1) coordinate (12);
		
			\draw[thick, gray] (7)  -- (8) -- (9) -- cycle;
			\draw[thick, gray] (7) -- (10)  -- (11) -- (8);
			\draw[thick, dashed, gray] (10) -- (12) -- (9);
			\draw[thick, dashed, gray] (12) -- (11);
			
			\draw[very thick, Goldenrod3] (7) -- (9) -- (12) -- (7) -- (11) -- (12) -- (9) -- (11) -- cycle;
		
			\foreach \x in {7,...,12}{
				\filldraw (\x) circle(0.02);
			}
		\end{tikzpicture}
		}} 
		\enspace\cup \enspace
		\vcenter{\hbox{
		\begin{tikzpicture}[scale=2.3]
			\coordinate (7) at (0,1);
			\path (7) ++ (0.9,-0.15) coordinate (8);
			\path (7) ++ (0.5,0.18) coordinate (9);
			\path (7) ++ (0,-1) coordinate (10);
			\path (8) ++ (0,-1) coordinate (11);
			\path (9) ++ (0,-1) coordinate (12);
		
			\draw[thick, gray] (7)  -- (8) -- (9) -- cycle;
			\draw[thick, gray] (7) -- (10)  -- (11) -- (8);
			\draw[thick, dashed, gray] (10) -- (12) -- (9);
			\draw[thick, dashed, gray] (12) -- (11);
			
			\draw[very thick,RoyalBlue3] (7) -- (8) -- (9) -- (11) -- (7) -- (8) -- (11)-- (9) -- cycle;
		
			\foreach \x in {7,...,12}{
				\filldraw (\x) circle(0.02);
			}
		\end{tikzpicture}
		}}
	\]
	\caption{Zerlegung von $\Delta^2 \times [0,1]$ in $3$-Simplizes}}
\end{figure}
\noindent Dieses Vorgehen verallgemeinern wir jetzt für beliebiges $n \in \mathbb{N}$:

\begin{definition}[{name=[Affin lineare Abbildungen zur Konstruktion der Kettenhomotopie]},label=def:710]
	Für $j=0, \ldots ,n$ seien $k_{n,j} \colon \abs{\Delta^{n+1}} \to \abs{\Delta^n} \times [0,1]$ und $\iota_{n,j} \colon \abs*{\Delta^{n-1}} \to \abs*{\Delta^n}$\marginnote{$e_1, \ldots , e_n$ sind wieder die Einheitsvektoren von $\mathbb{R}^n$} die eindeutigen affin linearen Abbildungen, für die gilt:
	\begin{align*}
		k_{n,j}(e_0)&=(e_0,0) & \iota_{n,j}(e_0) &= e_0\\
		k_{n,j}(e_1)&= (e_1,0) \enspace\ldots & \iota_{n,j}(e_1)&=e_1  \enspace \ldots \\
		k_{n,j}(e_j)&= (e_j,0) & \iota_{n,j}(e_{j-1})&= e_{j-1}\\
		k_{n,j}(e_{j+1})&= (e_j,1) \enspace\ldots & \iota_{n,j}(e_j)&= e_{j+1} \enspace \ldots \\
		k_{n,j}(e_{n+1})&= (e_n,1)& \iota_{n,j}(e_{n-1}) &= e_n
	\end{align*}
\end{definition}

\begin{lemma}[{name=[{Gleichungen für die Abbildungen aus \autoref{def:710}}]},label=lem:711]
	Es gelten folgende fünf Gleichungen für die Abbildungen $k_{n,j}$ und $\iota_{n,j}$:
	\begin{enumerate}[(i)]
		\item Für $0 \le l < j \le n$ gilt $k_{n,j} \circ \iota_{n+1,l} = \enbrace*{\iota_{n,l} \times \id_{[0,1]}} \circ k_{n-1,j-1}$
		\item Für $1 \le j+1 < l \le n+1$ gilt $k_{n,j} \circ \iota_{n+1,l} = \enbrace*{\iota_{n,l-1} \times \id_{[0,1]}} \circ k_{n-1,j} $
		\item Für $1 \le j+1 = l \le n$ gilt $k_{n,j} \circ \iota_{n+1,l} = k_{n,j+1} \circ \iota_{n+1,l}$
		\item Für $l=0,j=0$ ist $k_{n,j} \circ \iota_{n+1,l} = i_1 :  \Delta^n \to \Delta^n \times [0,1]$
		\item Für $l=n,j=n+1$ ist $k_{n,j} \circ \iota_{n+1,l} = i_0 : \Delta^n \to \Delta^n \times [0,1]$
	\end{enumerate}
\end{lemma}
\begin{beweis}
	Wir zeigen hier nur exemplarisch die erste Gleichung:
	Die linke Seite der Gleichung entspricht
		\[
			\begin{tikzcd}[row sep=2.8em, column sep=0.5em, font=\small]
				\abs*{\Delta^n} \dar["\iota_{n+1,l}" description] && e_0 \dar[mapsto] & e_{l-1} \dar[mapsto] & e_l \drar[mapsto]& \cdots & e_j \drar[mapsto]& \cdots 
				& e_n \drar[mapsto]& \\
				\abs*{\Delta^{n+1}} \dar["k_{n,j}" description] && e_0 \dar[mapsto] & e_{l-1} \dar[mapsto] & e_l \dar[mapsto] & e_{l+1} \dar[mapsto]& e_{j} \dar[mapsto] 
				& e_{j+1} \arrow[ddl,mapsto, bend left] & e_n & e_{n+1} \arrow[ddl,mapsto, bend left] \\
				\abs*{\Delta^{n}} \times [0,1] && (e_0,0) & (e_{l-1},0) & (e_l,0) & (e_{l+1},0)& (e_j,0) & & \\
				 &&  &  &  & & (e_j,1) & \cdots & (e_n,1)
			\end{tikzcd}
		\]
		und die rechte Seite entspricht
		\[
			\begin{tikzcd}[row sep=2.8em, column sep=0.2em, font=\small]
				\abs*{\Delta^n} \dar["k_{n-1,j-1}" description] && e_0 \dar[mapsto] & e_{l-1} \dar[mapsto] & e_l \dar[mapsto]& \cdots 
				& e_{j-1} \dar[mapsto]& e_j \arrow[ddl,bend left,mapsto]& e_{j+1} \arrow[ddl,bend left,mapsto] & \cdots & e_n \arrow[ddl,mapsto,bend left] \\
				\abs*{\Delta^{n-1}} \times [0,1] \arrow[dd,"\iota_{n,l}\times \id_{[0,1]}" description] && (e_0,0) \arrow[dd,mapsto] & (e_{l-1},0) \arrow[dd,mapsto] 
				& (e_l,0) \arrow[ddr, mapsto] & \cdots & (e_{j-1},0) \arrow[dd, controls={+(225:1.2) and +(135:1.2)}, mapsto] &   &  & \cdots \\
				 &&&&&& (e_{j-1},1) \arrow[dd, controls={+(-45:1) and +(45:1)}, end anchor={[xshift=2ex]north}, mapsto] & (e_{j},1) \arrow[dd, mapsto] & \cdots 
				 & (e_{n-1},1) \arrow[dd,mapsto]\\
				 \abs*{\Delta^n} \times [0,1] && (e_0,0)  & (e_{l-1},0) & (e_l,0) & (e_{l+1},0) &(e_j,0) & & & \\
				 &&&&&& (e_j,1) & (e_{j+1},1) & \cdots & (e_n,1)
			\end{tikzcd}
		\]
		Also folgt insgesamt die Gleichheit. Die anderen Gleichungen folgen genauso.
\end{beweis}

Damit können wir nun die im vorigen schon vereinfachte zentrale Aussage dieses Abschnitt beweisen:

\begin{beweis}[{name={von \eqref{eq:79}}}]
	Sei $h_n \colon C_n(X;R) \to C_{n+1}(X \times [0,1];R)$ definiert durch
	\[
		h_n(\sigma) := \sum_{j=0}^{n} (-1)^j \cdot \enbrace*{\sigma \times \id} \circ k_{n,j}
	\]
	$h_*$ ist nun die gesuchte Kettenhomotopie, da für jedes $n \in \mathbb{N}$ unter Anwendung der Gleichungen aus \autoref{lem:711} gilt:
	\begin{align*}
		\partial_{n+1} \enbrace[\big]{h_n(\sigma)} &= \partial_{n+1} \enbrace*{\sum_{j=0}^{n} (-1)^j \cdot (\sigma \times \id) \circ k_{n,j}} \\
		&=  \sum_{l=0}^{n+1} \sum_{j=0}^{n} (-1)^{j+l} \cdot \enbrace*{ \sigma \times \id} \circ k_{n,j} \circ \iota_{n+1,l} \tag*{(nach Definition \ref{def:52})}\\
		&= \smashoperator[r]{\sum_{0 \le l < j \le n}} \, (-1)^{j+l} \enbrace[\big]{(\sigma \circ \iota_{n,l}) \times \id} \circ k_{n-1,j-1} 
		+ \enspace\smashoperator{\sum_{1 \le j+1 < l \le n+1}}\, (-1)^{j+l} \enbrace[\big]{(\sigma \circ \iota_{n,l-1}) \times \id} \circ k_{n-1,j} \\
		& \quad + \smashoperator[r]{\sum_{1 \le j+1=l \le n}} \,(-1)^{j+l} (\sigma \times \id) \circ k_{n,j+1} \circ \iota_{n+1,l} 
		+ \smashoperator[r]{\sum_{1 \le j=l \le n}} \, (-1)^{j+l} (\sigma\times\id) \circ k_{n,j}\circ \iota_{n+1,l} \\
		& \quad + \sum_{0=l=j}  (-1)^{j+l} (\sigma \times \id) \circ i_1 
		+ \smashoperator[r]{\sum_{n+1=j+1=l}}\, (-1)^{j+l} (\sigma \times \id) \circ i_0 \\
		&= \smashoperator[r]{\sum_{0 \le l \le j \le n-1}}\, (-1)^{j+l+1} \enbrace[\big]{(\sigma \circ \iota_{n,l}) \times \id} \circ k_{n-1,j} 
		+ \enspace\smashoperator{\sum_{1 \le j+1 \le l \le n}} \,(-1)^{j+l+1} \enbrace[\big]{(\sigma \circ \iota_{n,l}) \times \id} \circ k_{n-1,j} \\
		& \quad + \underbrace{\smashoperator[r]{\sum_{1 \le j=l \le n}}\, (-1)^{j+l-1} (\sigma \times \id) \circ k_{n,j} \circ \iota_{n+1,l} 
		+ \smashoperator[r]{\sum_{1 \le j=l \le n}} \,(-1)^{j+l} (\sigma\times \id) \circ k_{n,j} \circ \iota_{n+1,l}}_{=0} \\
		& \quad + \sum_{0=l=j}  (-1)^{j+l} (\sigma \times \id) \circ i_1 + \smashoperator[r]{\sum_{n+1=j+1=l}} (-1)^{j+l} (\sigma\times \id) \circ i_0 \\
		&= - \smashoperator[r]{\sum_{\substack{0 \le l\le n \\ 0 \le j \le n-1}}} \,(-1)^{j+l} \enbrace[\big]{ \enbrace*{\sigma \circ \iota_{n,l}} \times \id } \circ k_{n-1,j} 
		+ (\sigma\times \id) \circ i_1  - (\sigma\times \id) \circ i_0 \\
		&= - \sum_{j=0}^{n-1} (-1)^j \sum_{l=0}^{n} (-1)^{l} \enbrace[\big]{ \enbrace*{\sigma \circ \iota_{n,l}} \times \id } \circ k_{n-1,j}
		+ {\underbrace{(\sigma\times \id) \circ i_1}_{= i_1 \circ \sigma}}  - \underbrace{(\sigma\times \id) \circ i_0}_{= i_0 \circ \sigma} \\
		&= - h_{n-1}(\partial_n \sigma) + C_n(i_1;R)(\sigma) - C_n(i_0;R)(\sigma) \qedhere  
	\end{align*}
\end{beweis}

\begin{bemerkung}[{name=[Konstruktion einer Kettenhomotopie aus einer Homotopie]}]
	Ist $H \colon X \times [0,1] \to Y$ eine Homotopie zwischen $f$ und $g$, so erhalten wir eine Kettenhomotopie $k$ zwischen $C_*(f;R)$ und $C_*(g;R)$ durch
	\[
		k_n(\sigma) := C_{n+1}(H;R) \circ h_n = \sum_{j=0}^{n} (-1)^j H \circ (\sigma \times \id) \circ k_{n,j}
	\]
\end{bemerkung}

\newpage

\section{Homologie von Paaren} % (fold)
\label{sec:8}

\begin{definition}[{name=[Paar von topologischen Räumen]}]
	Sei $A$ ein Teilraum von $X$, dann heißt $(X,A)$ ein \Index{Paar von topologischen Räumen}. 
	Eine Abbildung von Paaren $f \colon (X,A) \to (Y,B)$ ist eine stetige Abbildung $f \colon X \to Y$ mit $f(A) \subseteq B$. 
	Manchmal schreiben wir $(f,f|_{A}) \colon (X,A) \to (Y,B)$. 
	Die Kategorie von Paaren von topologischen Räumen bezeichnen wir mit $\TOP^2$.
\end{definition}

\begin{definition}[{name=[Singulärer Kettenkomplex und Homologiemodul von Paaren]}]
	Sei $(X,A)$ ein Paar. Dann definieren wir den \bet{singulären Kettenkomplex}\index{singulärer Kettenkomplex!von Paaren} von $(X,A)$ über $R$ durch
	\begin{align*}
		C_n(X,A;R) &:= \nicefrac{C_n(X;R)}{C_n(A;R)} \\
		\partial_n \enbrace[\big]{\sigma+ C_n(A;R)} &:= \partial_n \sigma + C_{n-1}(A;R)
	\end{align*}
	Dies ist wohldefiniert, da $\partial_n \enbrace[\big]{C_n(A;R)} \subseteq C_{n-1}(A;R)$. Weiter heißt
	\[
		H_n(X,A;R) := H_n \enbrace[\big]{C_*(X,A;R),\partial_*} 
	\]
	der \bet{$\mathbf{n}$-te singuläre Homologiemodul}\index{n-ter Homologiemodul@$n$-ter Homologiemodul!singulär, von Paaren} von $(X,A)$ mit Koeffizienten in $R$.
\end{definition}

\begin{bemerkung}[{name=[Der singuläre Kettenkomplex von Paaren definiert einen Funktor]}]
	Der singuläre Kettenkomplex von Paaren definiert offensichtlich einen Funktor $\TOP^2 \to R\text{-}\textsc{Ketten}$.\marginnote{Homologie von Paaren ist auch wieder homotopieinvariant, wie in Aufgabe 1 von Blatt 7 gezeigt.} 
	Durch Komposition mit Homologie (als Funktor $R\text{-}\textsc{Ketten} \to R\text{-}\MOD$) erhalten wir einen Funktor
	\[
		H_n(-,-;R) \colon \TOP^2 \longrightarrow R\text{-}\MOD
	\]
	Via $X \mapsto (X,\emptyset)$ können wir jeden Raum auch als Paar auffassen.
	Somit sind Paare eine Verallgemeinerung von Räumen, passend dazu gilt auch $H_n(X,\emptyset;R) = H_n(X;R)$.
\end{bemerkung}

Es stellt sich nun die folgende Frage:
Können wir $H_n(X,A;R)$ durch $H_n(X;R)$ und $H_n(A;R)$ ausdrücken? 
\begin{description}
	\item[Ansatz 1:] 
	\[
		H_n(X,A;R) \cong \nicefrac{H_n(X;R)}{H_n(A;R)}
	\]
	\faWarning \, Problem: $H_n(A;R)$ ist kein Untermodul von $H_n(X;R)$.
	\item[Ansatz 2:] Ist $H_n(X;R) \to H_n(X,A;R)$ surjektiv?
	
	Sei $x \in H_n(X,A;R)$. Dann gibt es $\sigma + C_n(A;R) \in C_n(X,A;R)$ mit $\partial_n \enbrace[\big]{\sigma + C_n(A;R)} = 0$ und $x= [\sigma]$. Es ist 
	$\sigma \in C_n(X;R)$, aber wir wissen nur $\partial_n \sigma \in C_{n-1}(A;R)$, nicht $\partial_n(\sigma)= 0$. \hfill {\large \faMehO}
\end{description}

\begin{beispiel}[{name=[{Homologie des Paares $(\abs*{\Delta^n}, \abs*{\partial \Delta^n})$}]}]
	In $C_n \enbrace[\big]{\abs*{\Delta^n}, \abs*{\partial \Delta^n};R}$ gilt für $\sigma := \id \colon \abs*{\Delta^n} \to \abs*{\Delta^n}$, dass $\partial_n([\sigma])=0$ ist, da 
	$\partial_n \sigma \in C_{n-1} \enbrace*{\abs*{\partial \Delta^n};R}$. Wir werden später sehen, dass $[\sigma]$ den Homologiemodul
	\[
		H_n \enbrace[\big]{\abs*{\Delta^n}, \abs*{\partial \Delta^n};R  } \cong R 
	\]
	erzeugt.
\end{beispiel}

\begin{definition}[{name=[{Lange exakte Sequenz von $R$-Moduln}]}]
	Eine Folge von $R$-linearen Abbildungen zwischen $R$-Moduln $M_i$
	\[
		\begin{tikzcd}
			\cdots \rar & M_n \rar["f_n"]  & M_{n-1} \rar["f_{n-1}"] & M_{n-2} \rar & \cdots \rar & M_1 \rar & M_0
		\end{tikzcd}
	\]
	heißt \Index{lange exakte Folge}, wenn sie exakt an jeder Stelle $M_i$ ist, d.h. für alle $i$ gilt $\im f_i = \ker f_{i-1}$.
\end{definition}

\begin{satz}[{name=[Existenz der Paarsequenz]},label=satz:86]
	Es gibt eine natürliche Transformation $\partial_n$ von $(X,A) \mapsto H_n(X,A;R)$ nach $(X,A) \mapsto H_{n-1}(A;R)$, sodass für jedes Paar $(X,A)$
	\begin{equation*}
		\begin{tikzcd}[row sep=1em, column sep=2.4em]
			\cdots  \rar["\partial_{n+1}"] & H_n(A;R) \rar["i_*"] & H_n(X;R) \rar["j_*"] & H_n(X,A;R) \rar["\partial_n"] 
			& H_{n-1}(A;R) \rar & \cdots \\
			\cdots \rar["\partial_1"] & H_0(A;R) \rar & H_0(X;R) \rar & H_0(X,A;R) \rar & 0
		\end{tikzcd} \label{eq:88}\tag{*}
	\end{equation*}
	eine lange exakte Folge ist. 
	Dabei sind $i \colon A \hookrightarrow X$ und $j \colon (X,\emptyset) \hookrightarrow (X,A)$ die Inklusionen. Das bedeutet: 
	
	Für jedes Paar $(X,A)$ haben wir eine $R$-lineare Abbildung $\partial_n \colon H_n(X,A;R) \to H_{n-1}(A;R)$, sodass für jede Abbildung $(f,f|_{A}) \colon (X,A) \to (Y,B)$ von Paaren folgendes Diagramm kommutiert
	\[
		\begin{tikzcd}
			H_{n}(X,A;R) \rar["\partial_n"] \dar["{(f,f|_A)_*}"] & H_{n-1}(A;R) \dar["{(f|_A)_*}"] \\
			H_n(Y,B;R) \rar["\partial_n"] & H_{n-1}(B;R)
		\end{tikzcd}
	\]
\end{satz}
\begin{beweis}
	\emph{siehe Seite~\pageref{bew:86}.}
\end{beweis}

\begin{korollar}[{name=[{Homologiemodul wenn $A$ kontraktibel ist}]}]
	\leavevmode
	\begin{enumerate}[(1)]
		\item Ist die Inklusion $i \colon A \to X$ eine Homotopieäquivalenz, so ist $H_n(X,A;R)= 0$ für alle $n$.
		\item Sei $A$ kontraktibel. Dann gilt für alle $n \ge 1$. $H_n(X;R) \cong H_n(X,A;R)$
	\end{enumerate}
\end{korollar}
\begin{beweis}
	\begin{enumerate}[(1)]
		\item Betrachte den folgenden Ausschnitt aus der langen exakten Folge:
		\[
			\begin{tikzcd}
				H_n(A;R) \rar["i_*"] & H_n(X;R) \rar["j_*"] & H_n(X,A;R) \rar["\partial_n"] & H_{n-1}(A;R) \rar["i_*"] & H_n(X;R)
			\end{tikzcd}
		\]
		Nach \autoref{kor:74} ist $i_* \colon H_n(A;R) \to H_n(X;R)$ ein Isomorphismus. Wegen der Exaktheit gilt
		$\ker j_* = \im i_* = H_n(X;R)$, also folgt $j_* = 0$. Da $\im \partial_n = \ker i_* = 0$, folgt $\partial_n=0$. Nun ist 
		\[
			H_n(X,A;R) = \ker \partial_n = \im j_* = 0
		\]
		\item Folgt leicht aus \autoref{kor:75} und der Exaktheit der Paarsequenz \eqref{eq:88}. Für $n=1$ brauchen wir außerdem, dass $\partial_1$ für $A$ wegzusammenhängend
		trivial ist. Dies folgt aus Aufgabe 1a) von Blatt 6, denn dort wurde gezeigt, dass $i_* \colon H_0(A) \to H_0(X)$ injektiv ist. \qedhere
	\end{enumerate}
\end{beweis}

\begin{definition}[{name=[Kurze exakte Folge von Kettenkomplexen]}]
	Seien 
	\begin{equation*}
		\begin{tikzcd}
			(C_*,d_*) \rar["i_*"] & (C_*',d_*') \rar["p_*"] & (C_*'',d_*'') 
		\end{tikzcd} \label{eq:810}\tag{\#}
	\end{equation*}
	\hyperref[def:61]{$R$-Kettenabbildungen}. Ist für jedes $n$ die Folge
	\(
		\begin{tikzcd}
			C_n \rar["i_n"] & C_n' \rar["p_n"] & C_n''
		\end{tikzcd}
	\)
	kurz exakt, so heißt \eqref{eq:810} eine \Index{kurze exakte Folge von Kettenkomplexen}.
\end{definition}

\begin{beispiel}
	Für jedes Paar $(X,A)$ ist 
	\[
		\begin{tikzcd}
			\enbrace*{C_*(A;R),\partial_*} \rar["i_*"] & \enbrace*{C_*(X;R),\partial_*} \rar["j_*"] & \enbrace*{C_*(X,A;R),\partial_*}   
		\end{tikzcd}
	\]
	eine kurze exakte Folge von $R$-Kettenkomplexen.
\end{beispiel}

\begin{satz}[{name={Schlangenlemma}},label=satz:810]
	Sei 
	\[
		\begin{tikzcd}
			(C_*,d_*) \rar["i_*"] & (C_*',d_*') \rar["p_*"] & (C_*'',d_*'')
		\end{tikzcd}
	\]
	eine kurze exakte Folge von $R$-Kettenkomplexen.
	\begin{enumerate}[a)]
		\item Für jedes $n$ gibt es eine eindeutige wohldefinierte $R$-lineare Abbildung $\partial_n \colon H_n(C_*'',d_*'') \to H_{n-1}(C_*,d_*)$ mit:
		Für $v' \in C_n'$ mit $d_n'' \circ p_n(v')=0$ ist 
		\[
			\partial_n \enbrace[\Big]{\underbrace{\benbrace*{p_n(v')}}_{\mathclap{\in H_n(C_*'',d_*'')}} } = [v] \in H_{n-1}(C_*,d_*)
		\]
		wobei $v \in C_{n-1}$ bestimmt ist durch $i_{n-1}(v)= d_n'(v')$.
		\item 
		\[
			\begin{tikzcd}
				\cdots \rar["\partial_{n+1}"] & H_n(C_*,d_*) \rar["i_*"] & 
					H_n(C_*',d_*') \arrow[d, phantom, ""{coordinate, name=Z}] \rar["p_*"] 
						& H_n(C_*'',d_*'') \arrow[dll, "\partial_n" very near start, rounded corners, to path={
																		-- ([xshift=2ex]\tikztostart.east)
																		|- (Z) [near end]\tikztonodes
																		-| ([xshift=-2ex]\tikztotarget.west)
																		-- (\tikztotarget)
																		}] & \\
				& H_{n-1}(C_*,d_*) \rar["i_*"] 
					& H_{n-1}(C_*',d_*') \rar["p_*"] \arrow[d, phantom, ""{coordinate, name=R}]
						& \quad \cdots  \quad \arrow[dll, "\partial_1" very near start, rounded corners, to path={
																		-- ([xshift=2ex]\tikztostart.east)
																		|- (R) [near end]\tikztonodes
																		-| ([xshift=-2ex]\tikztotarget.west)
																		-- (\tikztotarget)
																		}]& \\
				& H_0(C_*,d_*) \rar["i_*"] & H_0(C_*',d_*') \rar["p_*"] & H_0(C_*'',d_*'') \rar & 0
			\end{tikzcd}
		\]
		ist eine lange exakte Folge.
	\end{enumerate}
\end{satz}
\begin{beweis}
	\[
		\begin{tikzcd}[row sep=2em, column sep=3em]
			C_{n+1} \rar["i_{n+1}"] \dar["d_{n+1}"] & C_{n+1}' \rar["p_{n+1}"] \dar["d_{n+1}'"] & C_{n+1}'' \dar["d_{n+1}''"] \\
			C_n \rar["i_n"] \dar["d_{n}"]& C_n' \rar["p_n"] \dar["d_{n}'"]& C_n'' \dar["d_{n}''"]\\
			C_{n-1} \rar["i_{n-1}"] \dar["d_{n-1}"]& C_{n-1}'  \rar["p_{n-1}"] \dar["d_{n-1}'"]& C_{n-1}'' \dar["d_{n-1}''"]\\
			C_{n-2} \rar["i_{n-2}"]& C_{n-2}' \rar["p_{n-2}"] & C_{n-2}''
		\end{tikzcd}
	\]
	\begin{enumerate}[a)]
		\item Sei $x'' \in H_n(C_*'',d_*'')$, also $x'' = [v'']$ mit $v'' \in \ker d_n''$. Da $p_n$ surjektiv ist, existiert $v' \in C_n'$ mit $p_n(v')=v''$. Es gilt
		\[
			p_{n-1} \circ d_n'(v') = d_n'' \circ p_n(v') = d_n''(v'') = 0
		\]
		Also ist $d_n'(v') \in \ker p_{n-1}$. Weiter gilt $\ker p_{n-1}= \im i_{n-1}$, also existiert ein $ v\in C_{n-1}$ mit $i_{n-1}(v)= d_n'(v')$. Dann gilt
		\[
			i_{n-2} \enbrace[\big]{d_{n-1}(v)} = d'_{n-1} \enbrace[\big]{i_{n-1}(v)} = d_{n-1}' \enbrace[\big]{d_n'(v')} = 0   
		\]
		Da $i_{n-2}$ injektiv ist, folgt somit $d_{n-1}(v)=0$. Also setzen wir $\partial_n(x'') := [v] \in H_{n-1}(C_*,d_*)$. 
	
		Es bleibt zu zeigen, dass $[v] \in H_{n-1}(C_*,d_*)$ unabhängig von der Wahl von $v''$ und $v'$ ist.
		Seien $w'' \in \ker d_n''$, $w' \in C_n'$ und  $w \in C_{n-1}$ mit $x''=[w'']$, $p_n(w')=w''$ und $i_{n-1}(w)= d_n'(w')$.
		Es ist $[v'']=[w'']$, also existiert ein $a'' \in C_{n+1}''$ mit $d_{n+1}'' (a'')=v'' - w''$. $p_{n+1}$ ist surjektiv, also existiert $a' \in C_{n+1}'$ mit
		$p_{n+1}(a')=a''$. Es gilt
		\[
			p_n \enbrace[\big]{v'-w'- d_{n+1}'(a')} = v'' - w'' - \underbrace{d_{n+1}'' \enbrace[\big]{p_{n+1}(a')}}_{=v''- w''} = 0
		\]
		Mit $\ker p_n = \im i_n$ folgt die Existenz von $a \in C_n$ mit $i_n(a)=v'- w' - d_{n+1}(a')$. Es gilt nun %bleibt zu zeigen: $d_n(-a)= v-w$.
		\begin{align*}
			i_{n-1} \enbrace[\big]{d_n(a)- (v-w)} & = i_{n-1} \enbrace[\big]{d_n(a)} - i_{n-1}(v-w) \\ &= d_n' \enbrace[\big]{i_n(a)} - \enbrace[\big]{d_n'(v')- d_n'(w')}  \\  
			&= d_n' \enbrace*{ -d_{n+1}'(a') + v'-w'} - d_n' \enbrace*{v' - w'}  \\
			&= -d_n' d_{n+1}' (a') + d_n'(v'- w') -  d_n'(v'-w') = 0
		\end{align*}
		Da $i_{n-1}$ injektiv ist, folgt $d_n(a)- (v-w) = 0$, also $d_n(a)=(v-w)$. Es folgt $[v]=[w]$.
		\item Exaktheit der langen Folge:
		\begin{description}[wide=0cm]
			\item[{$\boxed{\im (i_n)_* = \ker (p_n)_*}$ :}] 
			\enquote{$\subseteq$} folgt aus $(p_n)_* \circ (i_n)_* = \enbrace*{p_n \circ i_n}_* = (0)_* =0 $.
		
			Für \enquote{$\supseteq$} betrachte $x' \in \ker (p_n)_*$, wobei $x'= [v']$ mit $v' \in C_n'$ und $d_n'(v')=0$.
			Da nun $\benbrace*{p_n(v')}= (p_n)_* [x'] = 0$ ist, gibt es $a''\in C_n''$ mit $d_{n+1}''(a'') = p_n(v')$.
			Da $p_{n+1}$ surjektiv ist, existiert $a' \in C_{n+1}'$ mit $p_{n+1}(a')=a''$. Dann gilt
			\[
				p_n \enbrace*{v' - d_{n+1}'(a')} =p_n(v') - p_n \circ d_{n+1}'(a') = p_n(v') -
				\underbracket[0.7pt][2pt]{d_{n+1}''\enbrace[\Big]{\underbracket[0.7pt][2pt]{p_{n+1}(a')}_{=a''}}}_{=p_n(v')} = 0 
			\]
			Da $\ker p_n = \im i_n$, gibt es $a \in C_n$ mit $i_n(a)= v'- d_{n+1}'(a')$. Nun ist
			\[
				i_{n-1} \enbrace[\big]{d_n(a)} = d_n' \enbrace[\big]{i_n(a)} = d_n' \enbrace[\big]{v'- d_{n+1}'(a')}  = d_n'(v') = 0 
			\]
			Da $i_{n-1}$ injektiv ist, folgt $d_n(a)=0$. Insbesondere $[a] \in H_n \enbrace*{C_*,d_*}$. Nun ist 
			\[
				(i_n)_* [a] = [i_n(a)] = \benbrace*{v'- d_{n+1}'(a')} = [v'] = x' 
			\]
			Also ist $x' \in \im (i_n)_*$.
			\item[{$\boxed{\im (p_n)_*=\ker \partial_n}$ :}] Für \enquote{$\subseteq$} betrachte $x'' \in \im (p_n)_*$. Dann gibt es $v' \in C_n'$ mit $d_n'(v')=0$ und
			$x'' = [p_n(v')]$. Es ist $d_n'(v') \in \ker p_{n-1} = \im i_{n-1}$, also existiert ein $v \in C_{n-1}$ mit $i_{n-1}(v)=d_n'(v')$. Es gilt nun $\partial_n(x'')=[v]$.
			Wegen der Injektivität von $i_{n-1}$ ist $v=0$, also $\partial_n(x'')=[0]$.
		
			Sei nun umgekehrt $x'' \in \ker \partial_n$. Wir finden wieder $v' \in C_n'$ und $v \in C_{n-1}$ mit $x''=[p_n(v')]$, $i_{n-1}(v)=d_n'(v')$ und 
			$0=\partial_n(x'')= [v]$. Daher gibt es $a \in C_n$ mit $d_n(a)=v$. Es folgt
			\[
				d_n' \enbrace[\big]{v' - i_n(a)} = i_{n-1}(v) - i_{n-1}\enbrace[\big]{d_n(a)} = 0 
			\]
			Es gilt nun $(p_n)_* \benbrace[\big]{v' -i_n(a)}  = \benbrace[\big]{p_n(v') - p_n(i_n(a))} = \benbrace*{p_n(v')} = x''$.
			\item[{$\boxed{\im \partial_n=\ker (i_{n-1})_*}$ :}] Für \enquote{$\subseteq$} sei $x \in \im \partial_n$. Dann gibt es wieder $v' \in C_n'$ und $v \in C_{n-1}$ mit
			$d_n''\enbrace*{p_n(v')}=0$, $d_n'(v')=i_{n-1}(v)$ und $x=[v]$. Dann gilt
			\[
				(i_{n-1})_*(x) = (i_{n-1})_*[v] = [i_{n-1}(v)] = \benbrace*{d_n'(v')} = 0 
			\]
		
			Sei nun umgekehrt $x \in \ker (i_{n-1})_*$ und $v \in C_{n-1}$ mit $d_{n-1}(v)=0$ und $x=[v]$. Da nun $(i_{n-1})_*(x) = \benbrace*{i_{n-1}(v)}$ und 
			$x \in \ker (i_{n-1})_*$  ist, gibt es $v' \in C_n'$ mit $d_n'(v')=i_{n-1}(v)$. Es gilt nun
			\[
				d_n'' \enbrace*{p_n(v')} = p_{n-1} \enbrace*{d_n'(v')} = p_{n-1} \enbrace*{i_{n-1}(v)} = 0 
			\]
			Daraus folgt $\partial_n \benbrace*{p_n(v')}=x$. \qedhere
		\end{description}
	\end{enumerate}
\end{beweis}

\begin{beweis}[name={von \autoref{satz:86}},label=bew:86]
	Für jedes Paar $(X,A)$ ist die Folge  der singulären Kettenkomplexe
	\[
		\begin{tikzcd}
			\enbrace*{C_*(A;R), \partial_*^A} \rar["C_*(i;R)"] &  \enbrace*{C_*(X;R), \partial_*^X} \rar["C_*(j;R)"] & \enbrace*{C_*(X,A;R), \partial_*^{(X,A)}}
		\end{tikzcd}
	\]
	kurz exakt. Das Schlangenlemma {\ref{satz:810}} produziert $\partial_n \colon H_n(X,A;R) \to H_{n-1}(A;R)$ und die lange exakte Sequenz \eqref{eq:88}.
	Es bleibt zu zeigen, dass die Randabbildungen aus dem Schlangenlemma wie behauptet eine natürliche Transformation definieren: Sei $(f,f|_{A}) \colon (X,A) \to (Y,B)$ eine
	Abbildung von Paaren. Zu zeigen ist, dass
	\[
		\begin{tikzcd}
			H_n(X,A;R) \rar["\partial_n"]  \dar["{(f,f|_A)_*}"]& H_{n-1}(A;R) \dar["{(f|_{A})_*}"] \\
			H_n(Y,B;R) \rar["\partial_n"] & H_{n-1}(B;R)
		\end{tikzcd}
	\]
	kommutiert. Sei $x'' \in H_n(X,A;R)$. Sei $v' \in C_n(X;R)$ mit $\partial_n^X(v') \in C_n(A;R)$ und $\benbrace*{v' + C_n(A;R)}= x''$. Dann ist 
	$\partial_n(x'')= \benbrace*{\partial_n^X(v')} \in H_{n-1}(A;R)$. Dann ist 
	\[
		\enbrace*{f,f|_A}_*(x'') =  \enbrace*{f,f|_A}_* \benbrace[\big]{v' + C_n(A;R)}  = \benbrace[\big]{f_*(v')+ C_n(B;R)} \in H_n(Y,B;R) 
	\]
	Es ist 
	\(
		\partial_n^Y \enbrace[\big]{f_*(v')} = \enbrace*{f|_A}_* \enbrace[\Big]{\underbrace{\partial_n^X(v')}_{\in C_n(A;R)}} \in C_n(B;R)   
	\).
	Also
	\begin{align*}
		\partial_n^Y \enbrace[\big]{\enbrace*{f,f|_A}_*(x'') } = \benbrace[\Big]{\partial_n^Y \enbrace[\big]{f_*(v')} } = \benbrace*{f_* \circ \partial_n^X(v')} = f_* \circ 
		\partial_n(x'') \qedhere  
	\end{align*}
\end{beweis}

\begin{bemerkung}
	Für eine Abbildung $(f,f|_A) \colon (X,A) \to (Y,B)$ erhalten wir ein kommutierendes Diagramm:
	\[
		\begin{tikzcd}[row sep=2.2em]
			\cdots \rar["\partial_{n+1}"] & H_n(A;R) \rar["i^A_*"] \dar["(f|_A)_*"] & H_n(X;R) \rar["j^X_*"] \dar["f_*"]& H_n(X,A;R) \rar["\partial_n"] \dar["{(f,f|_A)_*}"]  
			& H_{n-1}(A;R) \rar \dar["(f|_A)_*"] & \cdots \\
			\cdots \rar & H_n(B;R) \rar["i^B_*"] & H_n(Y;R) \rar["j^Y_*"] & H_n(Y,B;R) \rar["\partial_n"] & H_{n-1}(B;R) \rar & \cdots
		\end{tikzcd}
	\]
\end{bemerkung}

\begin{lemma}[{name={Fünfer-Lemma}},label=lem:812]
	Seien die Zeilen in folgendem kommutativen Diagramm von $R$-Moduln exakt.
	\[
		\begin{tikzcd}
			A_4 \rar["\alpha_4"] \dar["f_4" swap, "\cong"] & A_3 \rar["\alpha_3"] \dar["f_3" swap, "\cong"] & A_2 \rar["\alpha_2"]  \dar[Firebrick3, "f_2" swap] & A_1 \rar["\alpha_1"] \dar["f_1" swap, "\cong"] & A_0  \dar["f_0" swap, "\cong"]\\
			B_4 \rar["\beta_4"] & B_3 \rar["\beta_3"] & B_2 \rar["\beta_2"] & B_1 \rar["\beta_1"] & B_0 
		\end{tikzcd}
	\]
	Sind $f_4, f_3, f_1$ und $f_0$ Isomorphismen, so ist auch $f_2$ ein Isomorphismus.
\end{lemma}
\begin{beweis}
	\begin{description}
		\item[Injektivität:] Sei $a_2 \in A_2$ mit $f_2(a_2)=0$. Wir erhalten 
		\[
			f_1 \enbrace[\big]{\alpha_2(a_2)} = \beta_2 \enbrace[\big]{f_2(a_2)} = \beta_2(0)=0. 
		\]
		Da $f_1$ ein Isomorphismus ist, folgt $\alpha_2(a_2)=0$. Damit ist $a_2 \in \ker \alpha_2 = \im \alpha_3$, also existiert $a_3 \in A_3$ mit $\alpha_3(a_3)= a_2$. Es gilt
		\[
			\beta_3 \enbrace[\big]{f_3(a_3)} = f_2 \enbrace[\big]{\alpha_3(a_3)} = f_2(a_2)=0  
		\]
		Also ist $f_3(a_3) \in \ker \beta_3 = \im \beta_4$. Da $f_4$ Isomorphismus ist, gibt es $a_4 \in A_4$ mit $\beta_4 \enbrace[\big]{f_4(a_4)} = f_3(a_3)$. Nun gilt
		\[
			f_3 \enbrace[\big]{\alpha_4(a_4)- a_3} = f_3 \enbrace[\big]{\alpha_4(a_4)} - f_3(a_3) = \beta_4 \enbrace[\big]{f_4(a_4)} - f_3(a_3) = f_3(a_3)- f_3(a_3) =0
		\]
		Da $f_3$ ein Isomorphismus ist, folgt nun $\alpha_4(a_4)=a_3$. Nun ist 
		\[
			a_2= \alpha_3(\alpha_3)= \alpha_3 \enbrace[\big]{\alpha_4(a_4)} = \underbrace{\alpha_3 \circ \alpha_4}_{=0}(a_4) =0
		\]
		\item[Surjektivität:] Sei $b_2 \in B_2$. Fall 1: $\beta_2(b_2)=0$. Da $\im \beta_3 = \ker \beta_2$ und $f_3$ Isomorphismus, existiert $a_3 \in A_3$ mit
		$\beta_3 \enbrace[\big]{f_3(a_3)} = b_2$. Also ist $b_2= \beta_3 \enbrace[\big]{f_3(a_3)} = f_2  \enbrace[\big]{\alpha_3(a_3)} \in \im f_2$. 
	
		Ist $b_2$ beliebig, so genügt es zu zeigen: Es existiert ein $a_2 \in A_2$ mit $\beta_2 \enbrace[\big]{b_2 - f_2(a_2)} =0$. $f_1$ ist ein Isomorphismus, also existiert 
		ein $ a_1 \in A_1$ mit $f_1(a_1) = \beta_2(b_2)$. Dann folgt
		\[
			f_0 \enbrace[\big]{\alpha_1(a_1)} = \beta_1 \enbrace[\big]{f_1(a_1)} = \beta_1 \enbrace[\big]{\beta_2(b_2)} = \underbrace{\beta_1 \circ \beta_2}_{=0} (b_2)=0
		\]
		Da $f_0$ ein Isomorphismus ist, folgt $\alpha_1(a_1)=0$. Da $\ker \alpha_1 = \im \alpha_2$, folgt $\alpha_2(a_2)=a_1$. Es gilt nun
		\begin{align*}
			\beta_2 \enbrace[\big]{b_2 - f_2(a_2)} = \beta_2(b_2) - \beta_2 \enbrace[\big]{f_2(a_2)} &= \beta_2(b_2) - f_1 \enbrace[\big]{\alpha_2(a_2)} \\ 
			&= \beta_2(b_2) - f_1(a_1)
			= \beta_2(b_2) - \beta_2(b_2) =0 \qedhere
		\end{align*}
	\end{description}
\end{beweis}

\begin{lemma}[{name={\enquote{2 von 3}}}]
	Sei $f \colon (X,A) \to (Y,B)$ eine Abbildung von Paaren. Seien von den drei Abbildungen
	\begin{enumerate}[(i)]
		\item $(f|_A)_* \colon  H_n(A;R) \to H_n(B;R)$
		\item $f_* \colon H_n(X;R) \to H_n(Y;R)$
		\item $(f,f|_A)_* \colon H_n(X,A;R) \to H_n(Y,B;R)$
	\end{enumerate}
	zwei für jeweils alle $n$ Isomorphismen. Dann ist auch die dritte für alle $n$ ein Isomorphismus.
\end{lemma}
\begin{beweis}
	Da die Randabbildung in der Paarfolge eine natürliche Transformation ist, erhalten wir ein kommutatives Leiterdiagramm:
	\[
		\begin{tikzcd}[column sep=1.7em, font=\small]
			H_{n+1}(A) \rar \dar["{(f|_A)_*}"] & H_{n+1}(X) \rar \dar["f_*"] & H_{n+1}(X,A) \rar["\partial_{n+1}"] \dar["{(f,f|_A)_*}"] & H_n(A) \rar \dar["(f|_A)_*"] & H_n(X) \rar \dar["f_*"] & H_n(X,A) \rar["\partial_n"] \dar["{(f,f|_A)_*}"] & H_{n-1}(A) \dar["{(f|_A)_*}"] \\
			H_{n+1}(B) \rar & H_{n+1}(Y) \rar & H_{n+1}(Y,B) \rar["\partial_{n+1}"] & H_n(B) \rar & H_n(Y) \rar & H_n(Y,B) \rar["\partial_n"] & H_{n-1}(B)
		\end{tikzcd}\marginnote{\enquote{$R$} aus Platzgründen weggelassen}
	\]
	Da die Zeilen lang exakt sind, folgt die Behauptung aus dem Fünfer-Lemma.
\end{beweis}

\newpage

\section{Ausschneidung und die Mayer-Vietoris Folge} % (fold)
\label{sec:9}

\begin{satz}[{name={Ausschneidungsisomorphismus}},label=satz:91]
	Sei $(X,A)$ ein Paar und $L \subseteq A$, sodass der Abschluss $\overline{L}$ von im Inneren $\mathring{A}$ von $A$ liegt. 
	Dann induziert die Inklusion $\enbrace*{X \setminus L, A \setminus L} \xrightarrow{j} (X,A)$ einen Isomorphismus $j_* \colon H_n(X \setminus L, A \setminus L;R) \to H_n(X,A;R)$.
\end{satz}
\begin{beweis}
	\emph{Siehe Kapitel \ref{sec:13}.}
\end{beweis}

\begin{satz}[{name={Mayer-Vietoris-Folge}}]
	Seien $A,B \subseteq X$ offen und $X=A \cup B$.\marginnote{Für simpliziale Homologie gibt es auch eine Mayer-Vietoris-Folge, siehe Aufgabe 4 von Blatt 6} 
	Seien $i_A \colon A \hookrightarrow X$, $i_B \colon B \hookrightarrow X$, $j_A \colon A \cap B \hookrightarrow A$ und 
	$j_B \colon A \cap B \hookrightarrow B$ die Inklusionen. Dann gibt es eine Randabbildung $\partial_n = \rand{X=A\cup B}_n \colon H_n(X;R) \to H_{n-1}(A \cap B;R)$, sodass
	\[
		\begin{tikzcd}[column sep=3.7em]
			\cdots \rar["\partial_{n+1}"] & H_n(A \cap B) \rar["{\begin{psmallmatrix}(j_A)_* \\ (j_B)_* \end{psmallmatrix}}"] 
			&  \bigoplus\limits_{H_n(B)}^{H_n(A)} \rar["(i_A)_* - (i_B)_*"] 
			& H_n(X) \rar["\partial_n"] & H_{n-1}(A\cap B) \rar & \cdots
		\end{tikzcd} \marginnote{\enquote{$R$} aus Platzgründen weggelassen}
	\]
	eine lange exakte Folge ist. Die Randabbildung ist dabei eine natürlicher Transformation: Seien $A',B' \subseteq X'$ offen mit $X'=A' \cup B'$ und sei $f\colon X \to X'$ mit
	$f(A) \subseteq A'$ und $f(B) \subseteq B'$. Dann kommutiert\index{natürliche Transformation}
	\[
		\begin{tikzcd}[column sep=4em, row sep=3em]
			H_n(X;R) \rar["\rand{X=A \cup B}_n"] \dar["f_*"]  & H_{n-1}(A \cap B;R) \dar["\enbrace*{f|_{A \cap B}}_*"] \\
			H_n(X';R) \rar["\rand{X'=A' \cup B'}_n"] & H_n(A' \cap B';R)
		\end{tikzcd}
	\]
\end{satz}
\begin{beweis}
	Sei $L := X \setminus B$. Dann ist $L$ abgeschlossen und $L \subseteq A=\mathring{A}$. Wir erhalten daher einen Ausschneideisomorphismus 
	\[
		H_n(B,A \cap B;R) = H_n(X \setminus L, A \setminus L;R) \xrightarrow{\enspace \cong \enspace} H_n(X,A;R) 
	\]
	Betrachte nun das folgende kommutative Leiterdiagramm, das zwei Paarfolgen vergleicht:
	\[
		\begin{tikzcd}[font=\small, column sep=1.92em, row sep=2.5em]
			\cdots \rar & H_{n+1}(B,A \cap B) \rar["\rand{B,A \cap B}_{n+1}"] \arrow{d}{(i_B,j_A)_*}[swap]{\cong} & H_n(A \cap B) \rar["(j_B)_*"] \dar["(j_A)_*"] 
			& H_n(B) \rar["(l_B)_*"] \dar["(i_B)_*"] & H_n(B,A \cap B) \rar[Firebrick3,"\rand{B,A\cap B}_n"] \arrow[Firebrick3]{d}{(i_B,j_A)_*}[swap]{\cong}
			& H_{n-1}(A \cap B) \rar \dar["(j_A)_*"] & \cdots \\
			\cdots \rar & H_{n+1}(X,A) \rar["\rand{X,A}_{n+1}"] & H_n(A) \rar["(i_A)_*"] & H_n(X) \rar[Firebrick3,"(l_X)_*"] & H_n(X,A) \rar["\partial_n^{X,A}"] & H_{n-1}(A) \rar & \cdots
		\end{tikzcd} \marginnote{\enquote{$R$} aus Platzgründen weggelassen}
	\]
	dabei sind $l_B \colon (B,\emptyset) \to (B,A \cap B)$ und $l_X \colon (X, \emptyset) \to (X,A)$ die Inklusionen von Paaren.
	Definiere nun $\partial_n \colon H_n(X;R) \to H_{n-1}(A\cap B;R)$ durch
	\[
		\partial_n := \rand{B, A\cap B}_n \circ  (i_B,j_A)_*^{-1} \circ (l_X)_* 
	\]
	Die Natürlichkeit von $\partial_{n}^{B,A \cap B}$ impliziert die Natürlichkeit von $\partial_n$ (siehe auch \autoref{satz:86}). Wir müssen noch die Exaktheit zeigen:
	\begin{itemize}
		\item Zu zeigen: $\im \begin{psmallmatrix} (j_A)_* \\ (j_B)_* \end{psmallmatrix} = \ker \enbrace[\big]{(i_A)_* - (i_B)_*} $: 
		\begin{description}
			\item[\enquote{$\supseteq$}:] Sei $(a,b) \in \ker \enbrace*{(i_A)_* - (i_B)_*}$. Also $(i_A)_*(a) = (i_B)_*(b)$. Es ist 
			\[
				(i_B,j_A)_* \enbrace[\big]{(l_B)_*(b)} = (l_X)_* \circ (i_B)_* (b) = \underbrace{(l_X)_* \circ (i_A)_*}_{=0}(a)= 0 
			\]
			$(i_B,j_A)_*$ ist ein Isomorphismus, also $(l_B)_*(b)=0$. Es gilt $\ker (l_B)_* = \im (j_B)_*$, also folgt $\exists x \in H_n(A \cap B;R) : (j_B)_*(x)=b$. Dann gilt
			\[
				(i_A)_* \enbrace[\big]{a - (j_A)_*(x)} = (i_A)_*(a) - (i_B)_* \enbrace[\big]{(j_B)_* (x)} = (i_A)_*(a) - (i_B)_* (b) = 0  
			\]
			Mit $\ker (i_A)_*  =\im \rand{X,A}_{n+1}$ und $(i_B,j_A)_*$ ist Isomorphismus, folgt $\exists y \in H_{n+1}(B,A\cap B;R)$ mit 
			\[
				\rand{X,A}_{n+1} \circ  \enbrace*{i_B,j_A}_* (y) = a - (j_A)_*(x) 
			\]
			Es folgt für $z := x + \rand{B,A \cap B}_{n+1}(y)$ nun $(j_B)_*(z) =(j_B)_*(x) = b$ und 
			\begin{align*}
				(j_A)_*(z) = (j_A)_*(x) + (j_A)_* \enbrace[\big]{\rand{B,A \cap B}_{n+1}(y)} &= (j_A)_*(x) + \rand{X,A}_{n+1} \circ  \enbrace*{i_B, j_A}_*(y) \\ 
				&= (j_A)_*(x) + a - (j_A)_*(x) = a  
			\end{align*}
			\item[\enquote{$\subseteq$}:] Sei $x \in H_n(A\cap B;R)$. Dann gilt
			\begin{align*}
				\enbrace[\big]{(i_A)_* - (i_B)_*} \begin{pmatrix}
					(j_A)_* \\ (j_B)_*
				\end{pmatrix}  (x) &= (i_A)_* \circ (j_A)_* (x) - (i_B)_* \circ (j_B)_*(x) \\
				&= (i_A)_* \circ (j_A)_* (x) - (i_A)_* \circ (j_A)_* (x) = 0
			\end{align*}
			da $i_A \circ j_A = i_B \circ j_B$. 
		\end{description}
		\item Wir zeigen nun $\im\enbrace[\big]{ (i_A)_* - (i_B)_*} = \ker \partial_n$:
		\begin{description}
			\item[\enquote{$\supseteq$}:] Sei $v \in \ker \partial_n$, also $\rand{B, A\cap B}_n \circ  (i_B,j_A)_*^{-1} \circ (l_X)_*(v)=0$. Da $\ker \rand{B, A\cap B}_n =\im (l_B)_*$,
			gibt es ein $b \in H_n(B;R)$ mit
			\[
				(l_B)_*(b) = (i_B,j_A)_*^{-1} \circ (l_X)_*(v) \iff (i_B,j_A)_* \circ(l_B)_*(b) = (l_X)_*(v)
			\]
			Nun ist
			\[
				(l_X)_* \enbrace[\big]{v- (i_B)_* b} = (l_X)_*(v) - (l_X)_* \circ (i_B)_*(b) = (l_X)_*(v) -(i_B,j_A)_* \circ(l_B)_*(b)=0
			\]
			Da $\ker (l_X)_* = \im (i_A)_*$ ist, gibt es ein $a \in H_n(A;R)$ mit $(i_A)_*(a) = v - (i_B)_* (b)$ und somit ist
			$\enbrace[\big]{(i_A)_* - (i_B)_*}(a,-b) = v$.
			\item[\enquote{$\subseteq$}:] Sei $\begin{psmallmatrix} a \\ b \end{psmallmatrix} \in \bigoplus\limits_{H_n(B)}^{H_n(A)}$. Dann gilt
			\begin{align*}
				\partial_n \enbrace[\big]{(i_A)_* - (i_B)_*}  \begin{psmallmatrix} a \\ b \end{psmallmatrix} 
				&=  \rand{B, A\cap B}_n \circ  (i_B,j_A)_*^{-1} \circ (l_X)_* \enbrace[\big]{(i_A)_*(a)- (i_B)_*(b)}  \\
				&=\hphantom{-} \rand{B, A\cap B}_n \circ  (i_B,j_A)_*^{-1} \circ \underbracket[0.7pt][2pt]{(l_X)_* \circ (i_A)_*}_{=0} (a) \\
				&\,\hphantom{=} -\underbracket[0.7pt][2pt]{\rand{B, A\cap B}_n  \circ (l_B)_*}_{=0} (b) =0
			\end{align*}
		\end{description}
		\item Noch zu zeigen: $\im \partial_n = \ker \begin{psmallmatrix} (j_A)_* \\ (j_B)_* \end{psmallmatrix}$.
		\begin{description}
			\item[\enquote{$\supseteq$}:] Sei $x \in \ker \enbrace[\big]{(j_A)_* \oplus (j_B)_*}$, also $(j_A)_*(x) = 0 = (j_B)_*(x)$. Da $\im \rand{B,A \cap B}_n=\ker (j_B)_*$ ist, 
			gibt es $v \in H_n(B,A\cap B;R)$ mit $\rand{B,A\cap B}_n(v)=x$. Nun ist
			\[
				\rand{X,A}_n \circ (i_B,J_A)_*(v) = (j_A)_* \enbrace[\big]{\rand{B,A\cap B}_n(v)} = (j_A)_*(x) = 0
			\]
			Da $\ker \rand{X,A}_n = \im (l_X)_*$, finden wir $w \in H_n(X;R)$ mit $(l_X)_*(w) = (i_B,j_A)_* (v)$. Also folgt
			\[
				\partial_n(w) = \rand{B,A\cap B}_n \circ (i_B,j_A)_*^{-1} \circ (l_X)_* (w) = \rand{B,A\cap B}_n \circ (i_B,j_A)_*^{-1}\circ (i_B,j_A)_* (v) 
				= \rand{B,A\cap B}_n(v) = x
			\]
			\item[\enquote{$\subseteq$}:] Für $v \in H_n(X;R)$ ist
			\begin{align*}
				& \hphantom{=} \enspace \enbrace[\big]{(j_A)_* \oplus (j_B)_*} \circ \partial_n(v) \\
				&= \enbrace[\Big]{(j_A)_* \circ \rand{B,A \cap B}_n \circ (i_B,j_A)_* ^{-1} \circ (l_X)_*(v) ,
				{\underbracket[0.7pt][2pt]{(j_B)_* \circ \rand{B,A\cap B}_n}_{=0}} \circ (i_B,j_A)_* ^{-1} \circ (l_X)_* (v) } \\
				&= \enbrace[\Big]{{\underbracket[0.7pt][2pt]{\rand{X,A}_n \circ (l_X)_*}_{=0}}(v), 0  } = (0,0)\qedhere
			\end{align*}
		\end{description}
	\end{itemize}
\end{beweis}

\begin{bemerkung}[{name=[Anwendbarkeit der Mayer-Vietoris-Folge]}]
	Für die Mayer-Vietoris-Folge müssen $A$ und $B$ nicht notwendig offen sein. 
	Es genügt, dass die Inklusion $(B,A \cap B) \hookrightarrow (X,A)$ einen Isomorphismus $H_*(B,A \cap B;R) \to H_*(X,A;R)$ induziert.
\end{bemerkung}

\begin{satz}[{name=[{Homologie der $n$-Sphäre}]},label=satz:94]
	Für die Homologie der $n$-Sphäre $S^n$ gilt
	\begin{equation*}
		H_k \enbrace*{S^n ;R} \cong \begin{cases}
			R \oplus R , &\text{ falls }k=n=0\\
			R, &\text{ falls } k=0, n \not= 0 \\
			R, &\text{ falls } k=n\not= 0 \\
			0, &\text{ sonst}
		\end{cases} \label{eq:94} \tag{\#}
	\end{equation*}
\end{satz}

\begin{definition}[{name=[Reduzierte Homologie]}]
	Sei $X$ ein nichtleerer topologischer Raum. 
	Sei $p \colon X  \to \set{\mathrm{pt}}$ die Projektion auf den Ein-Punkt-Raum. 
	Die \Index{reduzierte singuläre Homologie} von $X$ ist definiert durch 
	\[
		\tilde{H}_k(X;R) := \ker \enbrace[\big]{ p_* \colon H_k(X;R) \to H_k(\set{\mathrm{pt}};R )} 
	\]
\end{definition}

Einige einfachen Eigenschaften der reduzierten Homologie sind:
\begin{itemize}
	\item Es gilt $\tilde{H}_k(X;R) =H_k(X;R)$ für $k>0$, da $H_k \enbrace*{\set{\mathrm{pt}};R }=0$ ist für $k>0$ (siehe \ref{bsp:58}) und somit $p_*$ die Nullabbildung ist.
	
	Weiter gilt $H_0(X;R) = R \oplus \tilde{H}_0(X;R)$, da $H_0(\set{\mathrm{pt}};R)=R$ frei ist.
	\item Es ist nicht schwer die Eigenschaften von singulärer Homologie auf die reduzierte singuläre Homologie zu übertragen. Insbesondere ist $\tilde{H}_*$ 
	homotopieinvariant und es gibt eine Mayer-Vietoris-Folge.
	\item Es ist $\tilde{H}_k \enbrace*{\set{\mathrm{pt}}}=0$ und für jeden kontrahierbaren Raum $X$ ist somit $\tilde{H}_n(X;R)=0$ für alle $n$.
\end{itemize}

\begin{bemerkung}[{name=[{Die Aussage von \autoref{satz:94} in reduzierter Homologie}]}]
	In reduzierter singulärer Homologie wird \eqref{eq:94} zu
	\begin{equation*}
		\tilde{H}_k(S^n;R) \cong \begin{cases}
			R, &\text{ falls }k=n\\
			0 , &\text{ sonst}
		\end{cases} \label{eq:96}\tag{\#\#}
	\end{equation*}
\end{bemerkung}

\begin{beweis}[{name={von \autoref{satz:94}}}]
	Wir berechnen die reduzierte Homologie von $S^n$ per Induktion nach $n$:
	\begin{description}
		\item[Induktionsanfang: $n=0$:] Es gilt
		\begin{align*}
			H_k(S^0;R) \stackrel{\xleftarrow{\varphi}}{\cong} H_k \enbrace[\big]{\set{\mathrm{pt}};R } \oplus H_k \enbrace[\big]{\set{\mathrm{pt}};R } = \begin{cases}
				R \oplus R, &\text{ falls }k=0\\
				0 ,&\text{ sonst}
			\end{cases}
		\end{align*} 
		Seien $i_- \colon \set{\mathrm{pt}} \hookrightarrow S^0$ und $i_+ \colon \set{\mathrm{pt}} \hookrightarrow S^0$ die Inklusionen mit $\im i_{\pm} = \set{\pm 1}$.
		Der Isomorphismus $\varphi$ ist gegeben durch $\varphi = (i_-)_* + (i_+)_*$. Für
		$(a,b) \in H_k \enbrace*{\set{-1};R}\oplus H_k \enbrace*{\set{+1};R}$ gilt dann
		\[
			p_* \enbrace[\big]{i_-(a) + i_+(b)} = (p \circ i_-)_*(a) + (p \circ i_+)(b) = \id_*(a) + \id_*(b)= a+b
		\]
		Also ist $\ker p_* = \set*{(a,-a) \given a \in H_k \enbrace[\big]{\set{\mathrm{pt}};R}}$ und damit folgt
		\[
			\tilde{H}_k(S^0;R) = \begin{cases}
				0, &\text{ falls }k \not= 0\\
				R ,&\text{ falls }k=0
			\end{cases}
		\]
		\item[Induktionsschritt: $n-1  \mapsto n$:] Sei $D_+^n = S^n \setminus \set{(1,0,\ldots ,0)}$ und $D^n_- = S^n \setminus \set{-1,0,\ldots ,0}$. 
		Dann sind $D^n_+$ und $D^n_-$ offene Teilmengen von $S^n$ mit $S^n = D_+^n \cup D^n_-$. Weiter gilt $D^n_+ \simeq \set{\mathrm{pt}}$ und $D^n_- \simeq \set{\mathrm{pt}}$.
		Für den Schnitt der beiden Mengen erhalten wir
		\[
			D^n_+\cap D^n_- = \set*{x=(x_0,\ldots , x_n) \in S^n \given x_0 \not\in \set{\pm 1}\rule{0cm}{1em}} \simeq S^{n-1} = \set{x=(0, x_1, \ldots , x_n) \in S^n\rule{0cm}{1em}}.
		\]
		Betrachte nun die Mayer-Vietoris-Folge zu $S^n = D^n_+ \cup D^n_-$: 
		\[
			\begin{tikzcd}
				\underset{=0}{\tilde{H}_k(D^n_+)} \oplus \underset{=0}{\tilde{H}_k(D^n_-)} \rar & 
				\tilde{H}_k(S^n) \rar["\partial_k", "\cong"'] &
				\tilde{H}_{k-1} \enbrace*{D^n_+ \cap D^n_-} \rar & 
				\underset{=0}{\tilde{H}_{k-1}(D^n_+)} \oplus \underset{=0}{\tilde{H}_{k-1}(D^n_-)} \\
				 & & \tilde{H}_{k-1}(S^{n-1}) \uar["\cong"]& 
			\end{tikzcd} \marginnote{\enquote{$R$} aus Platzgründen weggelassen}
		\]
		Es folgt $\tilde{H}_k(S^n;R) \cong \tilde{H}_{k-1}(S^{n-1};R) \stackrel{\text{I.A.}}{=} \begin{cases}
			R, &\text{ falls }k=n\\
			0, &\text{ sonst}
		\end{cases}$ \qedhere
	\end{description}
\end{beweis}

\newpage

\section{Anwendungen} % (fold)
\label{sec:10}

\begin{satz}[{name={Invarianz der Dimension}}]
	Sind $\mathbb{R}^n$ und $\mathbb{R}^m$ homöomorph, so gilt $n=m$.
\end{satz}
\begin{beweis}
	Sei $h \colon \mathbb{R}^n \to \mathbb{R}^m$ ein Homöomorphismus. 
	Sei $x_0 \in \mathbb{R}^n$ beliebig. 
	Dann erhalten wir auch einen Homöomorphismus $k := h|_{\mathbb{R}^n\setminus \set{x_0}} \colon \mathbb{R}^n \setminus \set{x_0} \to \mathbb{R}^m  \setminus \set{h(x_0)}$. 
	Nun ist $\mathbb{R}^n \setminus \set{x_0}$ homotopieäquivalent zu $S^{n-1}$ und $\mathbb{R}^{m}\setminus\set{h(x_0)}$ ist homotopieäquivalent zu $S^{m-1}$. 
	Da $k$ ein Homöomorphismus ist, ist $k_*$ ein Isomorphismus. Wir erhalten
	\begin{align*}
		H_*(S^{n-1}) \cong H_* \enbrace[\big]{\mathbb{R}^n \setminus \set{x_0}}  \cong H_* \enbrace[\big]{\mathbb{R}^m \setminus \set{h(x_0)}} \cong H_*(S^{m-1})
	\end{align*}
	Es ist aber $H_*(S^{n-1}) \cong H_*(S^{m-1})$ genau dann, wenn $n=m$.
\end{beweis}

\begin{satz}[{name={Fixpunktsatz von Brouwer}}]
	Jede stetige Abbildung $f \colon D^n \to D^n$ besitzt einen Fixpunkt.
\end{satz}
\begin{beweis}
	Durch Widerspruch: Angenommen es gibt $f\colon D^n \to D^n$ ohne Fixpunkt. Aus dem letzten Semester wissen wir, dass es dann eine stetige Abbildung $F\colon D^n \to S^{n-1}$ 
	gibt mit $F|_{S^{n-1}} = \id_{S^{n-1}}$. In Homologie erhalten wir folgendes kommutatives Diagramm \marginnote{Œ: $n\ge 2$, $n=1$ haben wir bereits im letzen Semester 
	bewiesen}
	\[
		\begin{tikzcd}[row sep=large, column sep=1em]
			\mathbb{Z}\cong H_{n-1}(S^{n-1}) \arrow[rr,"(\id_{S^{n-1}})_* = \id" ] \arrow[dr,"i_*"'] && H_{n-1}(S^{n-1}) \cong \mathbb{Z} \\
			& H_{n-1}(D^n)=0  \urar["(F)_*"']&
		\end{tikzcd}
	\]
	wobei $i\colon S^{n-1} \hookrightarrow D^n$ die Inklusion ist. Da $F \circ i = \id_{S^{n-1}}$ ist, gilt $F_* \circ i_* = \id_{H_{n-1}(S^{n-1})}$.
\end{beweis}

\begin{definition}[{name={Abbildungsgrad}}]
	Sei $f \colon S^n \to S^n$ eine stetige Abbildung.\marginnote{reduzierte Homologie!} 
	Da $\tilde{H}_n(S^n) \cong \mathbb{Z}$ ist, gibt es eine ganze Zahl $d(f)$, so dass $f_*(x)= d(f) \cdot x$ für alle $x \in \tilde{H}_n(S^n)$. 
	Diese Zahl heißt der \Index{Abbildungsgrad} von $f$.
\end{definition}

\begin{proposition}[{name=[Einfache Eigenschaften des Abbildungsgrades]}]
	Für den Abbildungsgrad $d$ gilt
	\begin{enumerate}[a)]
		\item $d(\id_{S^n}) = 1$.
		\item Sind $f$ und $g$ homotop, so gilt $d(f)=d(g)$.
		\item $d(f \circ g) = d(f) \cdot d(g)$.
	\end{enumerate}
\end{proposition}
\begin{beweis}
	\leavevmode
	\begin{enumerate}[a),itemsep=0pt]
		\item Es ist $(\id_{S^n})_*(x)= \id_\mathbb{Z}(x)= x$.
		\item Sind $f$ und $g$ homotop, so gilt $f_*=g_*$ und damit folgt $d(f)=d(g)$.
		\item Für alle $x \in \tilde{H}_n(S^n)$ gilt: $d(f \circ g) \cdot x = (f \circ g)_* (x) = f_* \circ g_*(x)= d(f) \cdot d(g) \cdot x$. Da 
		$\tilde{H}_n(S^n) \cong \mathbb{Z}$ folgt $d(f \circ g) = d(f) \cdot d(g)$. \qedhere
	\end{enumerate}
\end{beweis}

\begin{beispiel}[{name=[{Abbildungsgrad der Spiegelung $R_i \colon S^n \to S^n$}]},label=bsp:105]
	Sei $i \in \set{0, \ldots ,n}$ und $R_i \colon S^n \to S^n$ die Spiegelung  an der $x_i$-Achse, also 
	\[
		R_i(x_0, \ldots , x_n) = (x_0, \ldots , - x_i, \ldots , x_n)
	\]
	Dann ist $d(R_i) = -1$.
\end{beispiel}
\begin{beweis}
	Sei $f \colon S^n \to S^n$ der Homöomorphismus, der die $0$-te und die $i$-te Koordinate vertauscht, also gilt $f(x_0, \ldots , x_n) = (x_i, \ldots , x_0, \ldots , x_m)$. 
	Dann ist $R_i = f \circ R_0 \circ f$. 
	Da $f$ ein Homöomorphismus ist, ist $f_*$ ein Isomorphismus und daher ist $d(f)$ invertierbar, also $d(f) \in \set{\pm 1}$. 
	Mit $d(R_i) = d(f) \cdot d(R_0) \cdot d(f)$ folgt $d(R_0)=d(R_i)$. 
	Es genügt also $R_0$ zu betrachten.
	Wir zeigen $d(R_0)= -1$ durch Induktion nach $n$.
	\begin{description}
		\item[Induktionsanfang:] Sei $i_- \colon \set{\mathrm{pt}} \to S^0$ die Abbildung mit $\im i_- = \set{-1}$ und $i_+ \colon\set{\mathrm{pt}} \to S^0 $ die Abbildung mit $\im i_+ = \set{+1}$. 
		Dann ist \marginnote{Unter $H_0(S_0)\cong \mathbb{Z} \otimes \mathbb{Z}$ induziert $R_0$ $(x,y) \mapsto (y,x)$}
		\begin{align*}
			H_0(S^0) &= \set*{(i_-)_*(a)+ (i_+)_*(b) \given a,b \in H_0(\set{\mathrm{pt}})\rule{0cm}{1em}} \qquad \text{und}\\
			\tilde{H}_0(S^0) &= \set*{(i_-)_*(a)- (i_+)_*(a) \given a \in H_0(\set{\mathrm{pt}})\rule{0cm}{1em}}
		\end{align*}
		Nun ist
		\begin{align*}
			(R_0)_* \enbrace[\Big]{(i_-)_*(a)- (i_+)_*(a)} = (R_0)_* (i_-)_* (a) - (R_0)_* (i_+)_* (a) &= (R_0 \circ i_-)_*(a) - (R_0 \circ i_+)_*(a)\\ 
			&= (i_+)_*(a) - (i_-)_*(a) \\ 
			&= - \enbrace[\big]{(i_-)_*(a) - (i_+)_*(a)} 
		\end{align*}
		Also folgt $d(R_0)=-1$, da wir den Abbildungsgrad für reduzierte Homologie definiert hatten.
		\item[Induktionsschritt:] Wir definieren zwei offene Teilmengen von $S^n$ wie folgt 
		\begin{align*}
			D_+^n &= \set*{(x_0, \ldots , x_n) \in S^n \given x_n \not= -1}\\
			D_-^n &= \set*{(x_0, \ldots , x_m) \in S^n \given x_n \not= +1}  
		\end{align*}
		Wir wissen schon aus dem Beweis von \autoref{satz:94}, dass die Randabbildung aus der dazugehörigen Mayer-Vietoris-Folge $\tilde{H}_n(S^n) \xrightarrow{\partial_n} \tilde{H}_{n-1}(D^n_+ \cap D_-^n)$ ein Isomorphismus ist. 
		Da die Randabbildung natürlich ist, erhalten wir ein kommutatives Diagramm
		\[
			\begin{tikzcd}[row sep=3em, column sep=4em]
				\tilde{H}_n(S^n) \rar["\partial_n","\cong"'] \dar["(R_0)_*"]
				& \tilde{H}_{n-1}(D^n_- \cap D^n_+) \dar["\enbrace*{R_0|_{D^n_- \cap D^n_+}}_*"]  
				& \tilde{H}_{n-1}(S^{n-1}) \lar["i_*","\cong"'] \dar["\enbrace*{R_0|_{S^{n-1}}}_*"]\\
				\tilde{H}_n(S^n) \rar["\partial_n","\cong"']
				& \tilde{H}_{n-1}(D^n_- \cap D^n_+) 
				& \tilde{H}_{n-1}(S^{n-1}) \lar["i_*","\cong"']
			\end{tikzcd}
		\]
		dabei ist $i\colon S^{n-1} \to D^n_- \cap D_+^n$ die Homotopieäquivalenz $i(x_0, \ldots , x_{n-1}) = (x_0, \ldots , x_{n-1},0)$.
		Nach Induktionsannahme ist $\enbrace*{R_0\big|_{S^{n-1}}}_*(y)= -y$ für alle $y \in \tilde{H}_{n-1}(S^{n-1})$. 
		Für $x \in \tilde{H}_n(S^n)$ folgt dann
		\begin{align*}
			(R_0)_*(x) &= (\partial_n) ^{-1} \circ  (i_*) \circ \enbrace*{R_0\big|_{S^{n-1}}}_* \enbrace[\big]{(i_* ^{-1} \circ \partial_n) (x)} \\
			&= (\partial_n) ^{-1} \circ (i_*) \enbrace*{- i_*^{-1} \enbrace*{\partial_n (x)} } \\
			&= - \partial_n ^{-1} \circ (i_*) \circ (i_*)^{-1} \circ \partial_n(x)= -x 
		\end{align*}
		Also gilt für den Abbildungsgrad der Spiegelung $d(R_0)=-1$. \qedhere
	\end{description}
\end{beweis}

\begin{beispiel}[{name=[Grad der Punktspiegelung]},label=bsp:106]
	Der Grad der Punktspiegelung $R \colon S^n \to S^n$, $x  \mapsto -x$ ist $d(R)=(-1)^{n+1}$.
\end{beispiel}
\begin{beweis}
	Nach \autoref{bsp:105} ist $d(R)= d(R_0 \circ R_1 \circ \ldots \circ R_n) = d(R_0) \cdot \ldots \cdot d(R_n)= (-1)^{n+1}$.
\end{beweis}

\begin{erinnerungA}[{name=[Vektorfelder]}]
	Sei $M$ eine $C^\infty$-Mannigfaltigkeit. 
	Ein \Index{Vektorfeld} auf $M$ ist eine stetige Abbildung $v \colon M \to \Tmap M$ mit $v(x) \in \Tmap_x M$ für alle $x \in M$.
	Für $M=S^n$ entspricht ein Vektorfeld genau einer stetigen Abbildung $v \colon S^n \to \mathbb{R}^{n+1}$ mit $\sprod*{v(x)}{x} =0$ für alle $x \in S^n$.\index{Skalarprodukt}
	\begin{figure}[th]
		\centering{
		\begin{tikzpicture}[scale=1.5]
			\foreach \r in {0,20,...,340}{
				% \filldraw (\r:1) circle[radius=0.02];
				\draw[->, thick, DodgerBlue3] (\r:1) -- ++ ($(\r+90:1)$);
			}
			\draw[ultra thick] (0,0) circle(1) node{$S^1$};
		\end{tikzpicture}
		\caption{Ein nirgends verschwindendes Vektorfeld auf $S^1$}
		}
	\end{figure}
\end{erinnerungA}

\begin{satz}[{name=[{nicht verschwindende Vektorfelder auf $S^n$}]}]
	Es gibt genau dann ein Vektorfeld ohne Nullstellen auf $S^n$, wenn $n$ ungerade ist.
\end{satz}
\begin{beweis}
	Ist $n$ ungerade, so ist $S^{n} \subseteq \mathbb{R}^{n+1} = \mathbb{C}^{k}$ mit $k=\frac{n+1}{2}$. 
	Punkte in $S^n$ sind dann genau $k$-Tupel $z=(z_1, \ldots , z_k)$ mit $\abs*{z_1}^2 + \ldots + \abs*{z_k}^2 = 1$. 
	Ein Vektorfeld entspricht dann einer stetigen Abbildung $v \colon S^n \to \mathbb{C}^k$, so dass für das kanonische komplexe Skalarprodukt $\sprod*{\cdot }{\cdot}_\mathbb{C} $ auf $\mathbb{C}^k$ gilt
	\[
		\sprod*{z}{v(z)}_\mathbb{C} \in i \mathbb{R} \marginnote{einfach nachzurechnen}
	\]
	Definiere nun $v \colon S^n \to \mathbb{C}^k$ durch $v(z) := i \cdot z$. 
	Dann ist  $i z\neq 0$ für alle $z \in S^n$ und $\sprod*{z}{i z}= -i \cdot \norm{z} \in i \mathbb{R}$. 
	Sei umgekehrt $v \colon S^n \to \mathbb{R}^{n+1}$ ein Vektorfeld ohne Nullstellen. 
	Zu zeigen: $d(R)=1$, denn dann folgt die Behauptung, da $d(R)=(-1)^{n+1}$.
	Zeige dazu: $R$ ist homotop zu $\id$. 
	Seien $H, K \colon S^n \times [0,1] \to S^n$ definiert durch 
	\[
		H(x,t) := \frac{t x + (1-t)v(x)}{\norm{tx + (1-t)v(x)} } \qquad \text{bzw.}\qquad K(x,t) :=  \frac{-t x + (1-t)v(x)}{\norm{-tx + (1-t)v(x)}}
	\]
	Da $\sprod*{v(x)}{x}=0$  und $v(x)\not=0$ folgt für alle $x \in S^n$\marginnote{für orthogonale Vektoren gilt der Satz von Pythagoras}
	\begin{align*}
		\norm{t x + (1-t)v(x)}^2 &= \norm{tx}^2 + \norm{(1-t)v(x)}^2 \not=0   \\
		\norm{-t x + (1-t)v(x)}^2 &= \norm{-tx}^2 + \norm{(1-t)v(x)}^2 \not=0
	\end{align*}
	Also sind $H$ und $K$ wohldefiniert. 
	$H$ ist Homotopie zwischen $\id$ und $x \mapsto \frac{v(x)}{\norm{v(x)}}$ und $K$ ist Homotopie zwischen $-\id=R$ und $x \mapsto \frac{v(x)}{\norm{v(x)} } $.
\end{beweis}

\begin{satz}[{name={Jordanscher Kurvensatz}}]
	Sei $f \colon S^1 \to \mathbb{R}^2$ eine stetige, injektive Abbildung. 
	Dann hat $\mathbb{R}^2 \setminus f(S^1)$ genau zwei Wegzusammenhangskomponenten. 
	Eine davon ist beschränkt, die andere nicht.\index{Jordanscher Kurvensatz}
\end{satz}
\begin{beweis}[{name={mit \ref{satz:1010}}}]
	Wir können $f$ auch als injektive Abbildung $f \colon S^1 \to \mathbb{R}^2 \subseteq S^2 = \mathbb{R}^2 \cup \set{\infty}$ auffassen. 
	Nun ist $\tilde{H}_0 \enbrace[\big]{S^2 \setminus f(S^1)} =\mathbb{Z}$ nach \autoref{satz:1010} b). 
	Also $H_0 \enbrace[\big]{S^2 \setminus f(S^1)}= \mathbb{Z} \oplus \mathbb{Z}$. 
	Daher besteht $S^2 \setminus f(S^1)$ aus zwei Wegzusammenhangskomponenten $U_1$ und $U_2$. 
	Sei ohne Einschränkungen $\infty \in U_2$. 
	Dann ist $U_2 \setminus \set{\infty}$ immer noch wegzusammenhängend: 
	Seien $x,y \in U_2$, dann gibt es eine Umgebung $V$ von $\infty \in S^2$ mit $V \setminus \set{\infty} \cong D^2 \setminus \set{0} $, $x,y \not\in V$ und $f(S^1) \cap V = \emptyset$. 
	Also sind $U_1$ und $U_2 \setminus \set{\infty}$ die Wegzusammenhangskomponenten von $\mathbb{R}^2 \setminus f(S^1)$. 
	$U_1$ ist beschränkt und $U_2$ ist unbeschränkt.
\end{beweis}

\begin{satz}[{name=[{Homologie von $S^n$ ohne Einbettung vonScheiben $D^k$ und Sphären $S^k$ mit $k<n$}]},label=satz:1010]
	\leavevmode
	\begin{enumerate}[a),itemsep=1pt]
		\item Sei $f \colon D^k \to S^n$ eine stetige, injektive Abbildung mit $0 \le k < n$. Dann gilt $\tilde{H}_i\enbrace[\big]{S^n \setminus f(D^k)} =0$ für alle $i$.
		\item Sei $f \colon S^k \to S^n$ eine stetige, injektive Abbildung mit $0 \le k < n$. Dann gilt 
		\[
			\tilde{H}_i \enbrace[\big]{S^n \setminus f(S^k)} = \begin{cases}
				\mathbb{Z}, &\text{ falls }i=n-k-1\\
				0, &\text{ sonst}
			\end{cases}
		\]
	\end{enumerate}
\end{satz}

\begin{proposition}[{name=[{Homologie von $X= \bigcup_{a \in N} U_a$ wobei $U_a \subseteq U_{a+1}$}]},label=prop:1011]
	Seien $U_1 \subseteq U_2 \subseteq \ldots  \subseteq X$ offen mit $X = \bigcup_{i=1}^\infty U_i$. 
	Seien $i_{a,b} \colon U_a \hookrightarrow U_b$ für $a < b$ und $i_a \colon U_a \hookrightarrow X$ die entsprechenden Inklusionen. Dann gilt\marginnote{Diese Proposition gilt genauso für reduzierte Homologie}
	\begin{enumerate}[(1)]
		\item Für jedes $x \in H_k(X;R)$ gibt es $a >0$ und $x_a \in H_k(U_a;R)$ mit $(i_a)_*(x_a)=x$.
		\item Ist $x_a \in H_k(U_a;R)$ mit $(i_a)_*(x_a)=0$, so gibt es $b >a$ mit $(i_{a,b})_*(x_a)=0$.
	\end{enumerate}
\end{proposition}
\begin{beweis}
	Sei $\sigma \colon \abs*{\Delta^n} \to X$ ein singulärer Simplex in $X$. Dann ist $\set{\sigma^{-1}(U_a)}_{a \in \mathbb{N}}$ eine offene Überdeckung von $\abs*{\Delta^n}$. 
	Da $\abs*{\Delta^n}$ kompakt ist, gibt es $a$ mit $\abs*{\Delta^n} \subseteq \sigma ^{-1}(U_a)$, also $\sigma \enbrace*{\abs*{\Delta^n}} \subseteq U_a$. Da Elemente  von
	$C_n(X;R)$ endliche $R$-Linearkombinationen von singulären Simplizes sind, folgt 
	\[
		C_n(X;R) = \bigcup_a C_n(U_a;R).
	\]
	Damit ergeben sich direkt (1) und (2).
\end{beweis}

\begin{beweis}[{name={von \autoref{satz:1010}}}]
	\leavevmode
	\begin{enumerate}[a)]
		\item Durch Induktion nach $k$. Für $k=0$ ist $S^n \setminus f(D^0)$ homöomorph zu $\mathbb{R}^n$ und die Behauptung folgt.
	
		Induktionsschritt $k \mapsto k+1$: Da $D^{k+1}$ und $[0,1]^{k+1}$ homöomorph sind, können wir $D^{k+1}$ durch $[0,1]^{k+1}$ ersetzen. Sei 
		$x \in \tilde{H}_i \enbrace*{S^n \setminus f \enbrace*{[0,1]^{k+1}}}$. Es ist
		\[
			S^n \setminus f \enbrace*{\set{\nicefrac{1}{2}} \times [0,1]^{k}} = \underbracket[0.8pt][1pt]{S^n \setminus f \enbrace*{\benbrace*{0,\nicefrac{1}{2}}
			\times [0,1]^k}}_{=: A} \cup \underbracket[0.8pt][1pt]{S^n \setminus f \enbrace*{\benbrace*{\nicefrac{1}{2},1} \times [0,1]^k}}_{=: B} = A \cup B
		\]
		eine Vereinigung von offenen Mengen. Weiter ist 
		\begin{align*}
			A \cap B = \enbrace*{S^n \setminus f \enbrace*{\benbrace*{0,\nicefrac{1}{2} } \times [0,1]^k} } \cap \enbrace*{S^n \setminus f 
			\enbrace*{\benbrace*{\nicefrac{1}{2} ,1} \times [0,1]^k}} =   S^n \setminus f \enbrace*{[0,1]^{k+1}}
		\end{align*}
		Die zugehörige Mayer-Vietoris-Folge liefert nun einen Isomorphismus:
		\[
			\begin{tikzcd}
				\underset{=0 \text{ nach I.A.}}{\tilde{H}_{i+1}(A \cup B)} \rar & \tilde{H}_i(A \cap B) \rar["\cong"] & \tilde{H}_i(A) \oplus \tilde{H}_i(B) \rar & 
				\underset{=0 \text{ nach I.A.}}{\tilde{H}_i(A \cup B)}
			\end{tikzcd}
		\]
		% \[
		% 	\begin{tikzcd}[column sep=1.3em, font=\small]
		% 		\underset{=0 \text{ nach I.A.}}{\tilde{H}_{i+1} \enbrace[\Big]{S^n \setminus f \enbrace*{\set{\nicefrac{1}{2}}  \times [0,1]^k}}} \rar &
		% 		\tilde{H}_i \enbrace[\Big]{S^n \setminus f \enbrace*{[0,1]^{k+1}} }  \dar["\cong"] & \\
		% 		& {\displaystyle\bigoplus\limits_{\textstyle\tilde{H}_i \enbrace*{S^n \setminus f \enbrace[\big]{[\nicefrac{1}{2},1] \times [0,1]^k} } }
		% 		^{\textstyle\tilde{H}_i \enbrace*{S^n \setminus f \enbrace[\big]{[0,\nicefrac{1}{2}] \times [0,1]^k} } }} \rar &
		% 		\underset{=0 \text{ nach I.A.}}{\tilde{H}_i \enbrace*{S^n \setminus f \enbrace*{\set{\nicefrac{1}{2}} \times [0,1]^k } } }
		% 	\end{tikzcd}
		% \]
		Angenommen $x \neq 0$. Dann ist das Bild von $x$ in $\tilde{H}_i \enbrace*{A}$ ungleich Null oder in $\tilde{H}_i \enbrace*{B}$ ungleich Null. Indem wir dieses Argument 
		iterieren, erhalten wir eine Folge von Intervallen
		\[
			[0,1] \supseteq I_1 \supseteq I_2 \supseteq \ldots \enspace\text{ sodass } \bigcap_a I_a = \set{t} 
		\]
		und das Bild von $x \in \tilde{H}_i \enbrace*{S^n \setminus f \enbrace*{[0,1]^{k+1}}}$ unter der von der Inklusion 
		$\iota_{0,a} \colon S^n \setminus f \enbrace*{[0,1]^{k+1}} \to S^n \setminus f \enbrace*{I_a \times [0,1]^k}$ induzierten Abbildung 
		$(\iota_{0,a})_*$ ungleich Null 
		ist. Nun ist aber 
		\[
			\bigcup_a \enbrace[\Big]{S^n \setminus f \enbrace*{I_a \times [0,1]^k}} = S^n \setminus f \enbrace[\big]{\set{t} \times [0,1]^k }  
		\]
		und wieder nach Induktionsannahme ist $\tilde{H}_i \enbrace*{S^n \setminus f \enbrace*{\set{t} \times [0,1]^k}} =0$.  Für die von der Inklusion
		$\iota_a \colon S^n \setminus f \enbrace*{I_a \times [0,1]^k}  \to S^n \setminus f \enbrace*{\set{t} \times [0,1]^k}$ induzierte Abbildung gilt also
		$(\iota_a)_*(x)=0$. 
		Nach \autoref{prop:1011} (2) muss $x$ dann aber schon für ein $i$ in der reduzierten Homologie $\tilde{H}_i \enbrace*{S^n \setminus f \enbrace*{I_a \times [0,1]^k} }$ trivial sein. Widerspruch!
		\item Durch Induktion nach $k$. Für $k=0$ ist $S^n \setminus f(S^0) \cong \mathbb{R}^n \setminus \set{0}$. Da $\mathbb{R}^n \setminus \set{0} \simeq S^{n-1}$
		folgt für $k=0$ die Behauptung.
	
		Induktionsschritt $k-1 \mapsto k$: Wir setzen 
		\begin{align*}
			D^k_+ &:= \set*{(x_0,\ldots ,x_k) \in S^k \given x_k \ge 0} \\
			D^k_- &:= \set*{(x_0, \ldots , x_k) \in S^k \given x_k \le 0}
		\end{align*}
		Dann ist 
		$D^k_+ \cap D^k_- = \set*{(x_0, \ldots ,x_k)\in S^k \given x_k=0} \cong S^{k-1}$. Wieder ist 
		\[
			S^n \setminus f \enbrace*{D^k_+ \cap D^k_-} =\enbrace*{S^n \setminus f(D^k_+)} \cup \enbrace*{S^n \setminus f(D^k_-)}  
		\]
		eine offene Vereinigung mit $\enbrace*{S^n \setminus f(D^k_+)} \cap (S^n \setminus f(D^k_-)) = S^n \setminus f(S^k)$. Es ist $D^k_{\pm} \cong D^{k}$, also ist
		wegen a) ist die Randabbildung in der zugehörigen Mayer-Vietoris-Folge
		\[
			\tilde{H}_i \enbrace*{S^n \setminus f \enbrace*{D^k_+ \cap D^k_-} } \xrightarrow[\cong]{\minwidthbox{\partial_i}{2em}} \tilde{H}_{i-1} \enbrace*{S^n \setminus f(S^k)}   
		\]
		ein Isomorphismus. Die Behauptung folgt dann per Induktion. \qedhere
	\end{enumerate}
\end{beweis}

\newpage

\section{CW-Komplexe} % (fold)
\label{sec:11}

\begin{definition}[{name=[Pushout]}]
	Ein kommutatives Diagramm von topologischen Räumen der Form
	\begin{equation*}
		\begin{tikzcd}[row sep=1.6em, column sep=2.1em]
			A \rar["f"] \dar["g"] & X \dar["\overline{g} "] \\
			Y \rar["\overline{f}"] & Z
		\end{tikzcd} \label{eq:111} \tag{\#}
	\end{equation*}
	\noindent\begin{minipage}{0.65\textwidth}
		heißt ein \Index{Pushout}, falls es folgende universelle Eigenschaft hat: \smallskip\\
		\centering{$\forall \hat{f} \colon X \to \hat{Z}$, $\hat{g}\colon Y \to \hat{Z}$ mit 
		$\hat{f} \circ f = \hat{g} \circ g$ gilt: 
		\[
			\exists ! \varphi \colon Z \to \hat{Z} \text{ mit } \hat{f}= \varphi \circ \overline{g} , \hat{g} = \varphi \circ \overline{f}.
		\]}
	\end{minipage}
	\hfill \begin{minipage}{0.3\textwidth}
		\hspace{2em}\(
			\begin{tikzcd}[sep=1.9em]
				A \rar["f"] \dar["g"] & X \dar["\overline{g}"] \arrow[ddr, "\hat{f}", bend left, Firebrick3] \\
				Y \rar["\overline{f}"] \arrow[drr, "\hat{g}"',bend right, Firebrick3] & Z  \drar["\varphi", DodgerBlue3] \\
				& & \hat{Z}
			\end{tikzcd}
		\)
	\end{minipage}\\[-1.2em]
	Wir sagen dann auch: \eqref{eq:111} ist der Pushout von 
	\(
		\begin{tikzcd}[sep=small]
			Y & A \rar["f"] \lar["g"']& X 
		\end{tikzcd}
	\).
\end{definition}

\begin{bemerkung}[{name=[Existenz von Pushouts]},label=bem:112]
	Jedes Diagramm $\begin{tikzcd}[sep=small]
			Y & A \rar["f"] \lar["g"']& X 
		\end{tikzcd}$ lässt sich zu einem Pushout vervollständigen:
	Betrachte dazu den Raum $X \cup_A Y := \sfrac{X \coprod Y}{f(a) \sim g(a) \forall a \in A}$ mit der Quotiententopologie. 
	Sind $\overline{f}\colon Y \to X \cup_A Y$ und $\overline{g}\colon X \to X \cup_A Y$ die von den Inklusionen $X \hookrightarrow X \coprod Y$ und $Y \hookrightarrow X \coprod Y$ induzierten Abbildungen, so ist 
	\[
		\begin{tikzcd}[row sep=1.6em]
			A \rar["f"] \dar["g"'] & X \dar["\overline{g}"] \\
			Y \rar["\overline{f}"] & X \cup_A Y
		\end{tikzcd}
	\]
	ein Pushout. 
	Existenz und Eindeutigkeit aus der universellen Eigenschaft können benutzt werden um zu zeigen, dass der Pushout eindeutig bis auf kanonischen Homöomorphismus ist.
\end{bemerkung}



\begin{definition}[{name=[CW-Komplex]}]
	Ein \Index{CW-Komplex} ist ein topologischer Raum $X$ zusammen mit einer Filtrierung durch Unterräume von $X$:\marginnote{$X^\ssbrace{i}$ bezeichnet man auch als $i$-Gerüst oder $i$-Skelett} 
	\[
		\emptyset = X^{(-1)} \subseteq X^{\ssbrace{0}} \subseteq X^{(1)} \subseteq X^{(2)} \subseteq  \ldots \subseteq X^\ssbrace{n} \subseteq X^\ssbrace{n+1} 
		\subseteq \ldots \subseteq X
	\]
	sodass die folgenden zwei Eigenschaften erfüllt sind:
	\begin{enumerate}[(i),itemsep=0pt]
		\item \Index{Zellstruktur}: Zu jedem $n \in \mathbb{N}$ gibt einen Pushout von topologischen Räumen 
		\[
			\begin{tikzcd}[row sep = 2.3em, column sep = 5.2em]
				\coprod S^{n-1} \rar["\coprod q_i^\ssbrace{n}"] \dar["\coprod j_i"', hook] 
				& X^\ssbrace{n-1} \dar[hook, "k^{(n)}"] \\
				\coprod D^n \rar["\coprod Q_i^\ssbrace{n}"] & X^\ssbrace{n}
			\end{tikzcd}
			\marginnote{Koprodukte indiziert über $i \in I^\ssbrace{n}$}
		\]
		wobei $j_i \colon S^{n-1} \hookrightarrow D^n$ und $k_n \colon X^\ssbrace{n-1} \hookrightarrow X^\ssbrace{n}$ die Inklusionen sind.
		\item \Index{Schwache Topologie}: Es ist $X= \bigcup_{n \in \mathbb{N}} X^\ssbrace{n}$ und $U \subseteq X$ ist genau dann offen, wenn $U \cap X^\ssbrace{n} \subseteq X^\ssbrace{n}$ für 
		alle $n$ offen ist.
	\end{enumerate}
	Eine Abbildung $f \colon X \to Y$ zwischen CW-Komplexen heißt \bet{zellulär}\index{zelluläre Abbildung}, falls $f\enbrace*{X^\ssbrace{n}}\subseteq Y^{(n)}$ für alle $n \in \mathbb{N}$ gilt.
\end{definition}

Zitat von \hrefsymX{https://en.wikipedia.org/wiki/CW_complex}{Wikipedia} zum Namen: \textenglish{\emph{The C stands for \enquote{closure-finite}, and the W for \enquote{weak topology}.}} C für \emph{\enquote{cell}} scheint mir allerdings auch sinnvoll.

\begin{bemerkung}
	\begin{enumerate}[(i)]
		\item Die Abbildungen $q_i^\ssbrace{n}$ und $Q_i^\ssbrace{n}$ sind \emph{nicht} Teil der Struktur eines CW-Komplex. 
		Nur die Existenz von $q_i^\ssbrace{n}$ und $Q_i^\ssbrace{n}$ wird gefordert -- sie ist nicht notwendig eindeutig. 
		Hat man $q_i^\ssbrace{n}$ und $Q_i^\ssbrace{n}$ \emph{gewählt}, so heißt $q_i^\ssbrace{n}$ die \Index{anklebende Abbildung} der $i$-ten $n$-Zelle und $Q_i^\ssbrace{n}$ die \Index{charakteristische Abbildung} der $i$-ten $n$-Zelle.
		\item Sei $\mathring{D}^n = D^n \setminus S^{n-1}$ das Innere von $D^n$. 
		Die Abbildungen $Q_i^\ssbrace{n}$ schränken sich zu einem Homöomorphismus $\coprod_{i \in I^{(n)}} \mathring{D}^n \to X^\ssbrace{n} \setminus X^\ssbrace{n-1}$ ein.\footnote{Dass dies ein Homöomorphismus ist, folgt aus der Pushout-Eigenschaft und dem letzen Teil der \autoref{bem:112}, denn es gilt demnach $X^\ssbrace{n}\cong X^\ssbrace{n-1}\cup_{\coprod S^{n-1}} \coprod D^n$.} 
		Insbesondere lässt ich $I^\ssbrace{n}$ mit der Menge der Wegzusammenhangskomponenten von $X^\ssbrace{n} \setminus X^\ssbrace{n-1}$ identifizieren. 
		Die Wegzusammenhangskomponenten heißen die \bet{offenen Zellen}\index{offene Zelle} von $X$. 
		Damit ist jeder CW-Komplex die disjunkte Vereinigung seiner offenen Zellen.
	
		\faWarning\,Achtung: Jede offene $n$-Zelle ist offen in $X^\ssbrace{n}$, aber nicht notwendig in $X$.
		\item Der Abschluss einer offenen Zelle $Q_i^\ssbrace{n}(\mathring{D}^n)$ ist $Q_i^\ssbrace{n}(D^n)$ und insbesondere kompakt. 
		Die $Q_i^\ssbrace{n}(D^n)$ heißen die \bet{abgeschlossenen Zellen}\index{abgeschlossene Zelle} und sind als Abschluss von offenen Zellen unabhängig von der Wahl der $Q_i^\ssbrace{n}$.
	
		\faWarning\,Achtung: Abgeschlossene Zellen sind Bilder von $D^n$ unter stetigen Abbildungen, aber nicht notwendig homöomorph zu $D^n$.
		\item Ein CW-Komplex $X$ heißt \bet{endlich}\index{CW-Komplex!endlicher}, wenn er nur aus endlich vielen Zellen besteht, also wenn 
		$I = \bigcup_{n \in \mathbb{N} } I^\ssbrace{n}$ endlich ist. 
		Insbesondere gibt es dann ein $n$ mit $X=X^\ssbrace{n}$.
	
		Die Dimension von $X$ ist die maximale Dimension von Zellen von $X$, also $\dim X := N$, falls $X^\ssbrace{N}=X$, $X^\ssbrace{N-1} \subsetneq X$. Gibt es kein solches $N$ so setzen wir $\dim X := \infty$.
	\end{enumerate}
\end{bemerkung}

\begin{beispiel}[label=bsp:115]
	Viele uns bereits bekannte Räume haben eine CW-Komplex-Stuktur:
	\begin{enumerate}[(i)]
		\item $S^n$ mit 
		\[
			(S^n)^{(k)} = \begin{cases}
				\emptyset, &\text{ falls }k=-1\\
				\set{(1,0,\ldots ,0)}, &\text{ falls } k=0, \ldots n-1 \\
				S^n , &\text{ falls } k \ge n 
			\end{cases}
		\]
		ist ein CW-Komplex via $S^n \cong \nicefrac{D^n}{S^{n-1}}$
		\[
			\begin{tikzcd}
				S^{n-1} \rar \dar[hook] & \enbrace*{S^n}^{(n-1)} \rar[equal] \dar[hook] & \set{(1,0,\ldots ,0)}   \\
				D^n \rar & S^n  
			\end{tikzcd}
		\]
		\item $S^n$ mit $(S^n)^{(k)} = \set{(x_0, \ldots ,x_k,0, \ldots ,0) \in S^n} $ ist ebenfalls eine CW-Struktur auf $S^n$.
		\item Sei $K=(V,\Sigma)$ ein simplizialer Komplex. 
		Es ist $\abs*{\Delta^n} \cong D^n$ und $\abs*{\partial \Delta^n} \cong S^{n-1}$. 
		Sei $\Sigma^{(n)} := \Sigma_0 \cup \Sigma_1 \cup \ldots \cup \Sigma_n$. 
		Dann wird $\abs*{K}$ durch $\abs*{K}^{(n)} := \abs*{\enbrace*{V,\Sigma^{(n)}}}$ für $n \ge 0$ zu einem CW-Komplex. 
		Ist $f \colon K \to L$ eine \hyperref[def:25]{simpliziale Abbildung}, so ist $\abs*{f} \colon \abs*{K} \to \abs*{L}$ eine zelluläre Abbildung.
		\item \label{enum:115:4}Der $n$-dimensionale reelle projektive Raum $\mathbb{R}P^n$ kann definiert werden durch\index{reell projektiver Raum} 
		\[
			\mathbb{R}P^n := \sfrac{\mathbb{R}^{n+1} \setminus \set{0}}{x \sim \lambda x} = \sfrac{S^n}{x \sim -x}
		\]
		Punkte in $\mathbb{R}P^n$ sind Äquivalenzklassen von $(n+1)$-Tupeln reeller Zahlen und werden als die sogenannten homogenen Koordinaten\index{homogene Koordinaten} $[x_0 : \ldots : x_n]$ geschrieben. Es ist 
		\[
			\mathbb{R}P^n = \set*{[x_0 : \ldots : x_n] \given (x_0, \ldots ,x_n) \in S^n \rule{0cm}{1em}} 
			= \set*{[x_0 : \ldots : x_n] \given (x_0, \ldots ,x_n) \in S^n, x_n \ge 0 \rule{0cm}{1em}}
		\]
		Wir erhalten einen Homöomorphismus $f^{(n)} \colon \nicefrac{D^n}{x \sim -x, x \in S^{n-1}} \to \mathbb{R}P^n$ mit
		\[
			f^{(n)} \enbrace*{x_1, \ldots ,x_n} = \benbrace*{x_1 : \ldots : x_n : \sqrt{1- {\textstyle\sum\nolimits_{i=1}^{n}} x_i^2}} 
		\]
		Ist $Q^\ssbrace{n} \colon D^n \to \mathbb{R}P^n$ die Komposition von $f^{(n)}$ mit der Projektion $D^n \twoheadrightarrow \nicefrac{D^n}{x \sim -x, x \in S^{n-1}}$ und 
		$q^\ssbrace{n} \colon S^{n-1} \twoheadrightarrow \mathbb{R}P^{n-1}$ die Projektion, so erhalten wir einen Pushout
		\[
			\begin{tikzcd}
				S^{n-1} \rar["q^\ssbrace{n}"] \dar[hook]& \mathbb{R}P^{n-1} \dar[hook] \\
				D^n \rar["Q^\ssbrace{n}"] & \mathbb{R}P^n
			\end{tikzcd}
		\]
		Induktiv erhalten wir einen CW-Struktur auf $\mathbb{R}P^n$ mit genau einer $k$-Zelle für $k=0,\ldots ,n$ und 
		\[
			\enbrace*{\mathbb{R}P^n}^{(k)} \cong \mathbb{R}P^k 
		\]
		\item \label{enum:115:5} Der $n$-dimensionale komplexe projektive Raum $\mathbb{C}P^n$ kann definiert werden durch 
		\[
			\mathbb{C}P^n := \sfrac{\mathbb{C}^{n+1} \setminus \set{0}}{z \sim \lambda z} = \sfrac{S^{2n+1}}{z \sim \lambda z, \lambda \in S^1 \subseteq \mathbb{C}}
		\]
		Punkte in $\mathbb{C}P^n$ sind Äquivalenzklassen von $(n+1)$-Tupeln komplexer Zahlen und werden als homogenen Koordinaten $[z_0 : \ldots : z_n]$ 
		geschrieben. Es ist wieder
		\[
			\mathbb{C}P^n = \set*{[z_0 : \ldots : z_n] \given (z_0, \ldots ,z_n) \in S^{2n+1} \rule{0cm}{1em}} 
			= \set*{[z_0 : \ldots : z_n] \given (z_0, \ldots ,z_n) \in S^{2n+1}, \abs{z_n} \ge 0 \rule{0cm}{1em}}
		\]
		Wir erhalten wieder einen Homöomorphismus $f^{(n)} \colon \nicefrac{D^{2n+1}}{z \sim \lambda z, z \in S^{2n-1},\lambda \in S^{1}} \to \mathbb{C}P^n$ durch
		\[
			f^{(n)} \enbrace*{z_1, \ldots ,z_n} = \benbrace*{z_1 : \ldots : z_n : \sqrt{1- {\textstyle\sum\nolimits_{i=1}^{n}} \abs{z_i}^2}} 
		\]
		Ist $Q^{(2n)} \colon D^{2n} \to \mathbb{C}P^n$ die Komposition von $f^{(n)}$ mit der Projektion $D^{2n} \twoheadrightarrow \nicefrac{D^{2n}}{z \sim \lambda z, z \in S^{2n}}$ und 
		$q^\ssbrace{n} \colon S^{2n+1} \to \mathbb{C}P^{n-1}$ die Projektion, so erhalten wir auch wieder einen Pushout
		\[
			\begin{tikzcd}
				S^{2n-1} \rar["q^\ssbrace{n}"] \dar[hook]& \mathbb{C}P^{n-1} \dar[hook] \\
				D^{2n} \rar["Q^\ssbrace{n}"] & \mathbb{C}P^n
			\end{tikzcd}
		\]
		Induktiv erhalten wir einen CW-Struktur auf $\mathbb{C}P^n$ mit genau einer $k$-Zelle für $k=0,2,\ldots ,2n$ und 
		\[
			\enbrace*{\mathbb{C}P^n}^{(k)} \cong \mathbb{C}P^{\floor*{\frac{k}{2}}} 
		\]
	\end{enumerate}
\end{beispiel}

\begin{lemma}[{name=[Kompaktheit von Teilmengen eines CW-Komplexes]},label=lem:116]
	Sei $X$ ein CW-Komplex und $K \subseteq X$. 
	Dann ist $K$ genau dann kompakt, wenn $K$ abgeschlossen ist und $K$ nur endlich viele offene Zellen von $X$ schneidet.
\end{lemma}
\begin{beweis}
	Zunächst: Ist $S \subseteq X$ ein Unterraum der jede offene Zelle höchstens in einem Punkt schneidet, so ist $S$ diskret. 
	Für $S_0 \subseteq S$ ist $S_0 \cap (X^\ssbrace{n} \setminus X^\ssbrace{n-1})$ abgeschlossen in $X^\ssbrace{n}$, da $S_0$ jede offene $n$-Zelle von $X$ höchstens in einem Punkt schneidet.  
	Daher ist $S_0 \cap X^\ssbrace{n} \subseteq X^\ssbrace{n}$ abgeschlossen für alle $n$. 
	Damit ist also jede Teilmenge $S_0 \subseteq S$ abgeschlossen in $X$ und $S$ somit diskret.

	Sei nun $K \subseteq X$ kompakt. Dann ist $K$ sicher abgeschlossen. 
	Sei $S \subseteq K$ ein Teilraum, der aus jeder offenen Zelle von $X$, die $K$ schneidet, genau einen Punkt enthält. 
	Wegen der Vorüberlegung ist $S$ diskret. 
	Da $S \subseteq K$ abgeschlossen und $K$ kompakt ist, ist $S$ auch kompakt. 
	Damit ist $S$ endlich und $K$ schneidet nur endlich viele offene Zellen.

	Ist nun umgekehrt $K$ abgeschlossen und schneidet nur endlich viele offene Zellen, so ist $K$ enthalten in einer endlichen Vereinigung von abgeschlossen Zellen von $X$.
	Da diese abgeschlossen Zellen kompakt sind, ist auch diese Vereinigung kompakt, K ist also ein abgeschlossener Teilraum eines kompakten Teilraums und damit selbst kompakt.
\end{beweis}

\begin{korollarB}[{name=[{ein CW-Komplex ist genau dann kompakt, wenn er endlich ist}]}]
	Ein CW-Komplex ist genau dann kompakt, wenn er endlich ist.
\end{korollarB}

\newpage

\section{Zelluläre Homologie} % (fold)
\label{sec:12}

\begin{definition}[{name=[{$n$-ter zellulärer Kettenmodul und zelluläre Randabbildung}]}]
	Sei $X$ ein CW-Komplex. Dann heißt 
	\[
		C_n^{\cell}(X;R) := H_n \enbrace*{X^\ssbrace{n}, X^\ssbrace{n-1};R} 
	\]
	der \bet{$\mathbf{n}$-te zelluläre Kettenmodul}\index{n-ter zellulärer Kettenmodul@$n$-ter zelluläre Kettenmodul} von $X$ mit Koeffizienten in $R$. Wir definieren die 
	\bet{$\mathbf{n}$-te Randabbildung}\index{n-te Randabbildung@$n$-te Randabbildung!zellulär}
	$\partial^{\cell}_n \colon C_n^{\cell}(X;R) \to C_{n-1}^{\cell}(X;R)$ durch die Komposition 
	\[
		\begin{tikzcd}[column sep=huge, row sep=large]
			H_n \enbrace*{X^\ssbrace{n}, X^\ssbrace{n-1};R} \rar["\partial_n^{\enbrace{X^\ssbrace{n}, X^\ssbrace{n-1}}}"] & H_{n-1} \enbrace*{X^\ssbrace{n-1};R} \dar["\enbrace*{j^{(n-1)}}_*"] \\
			& H_{n-1} \enbrace*{X^\ssbrace{n-1}, X^\ssbrace{n-2};R} 
		\end{tikzcd}
	\]
	Dabei ist $\partial_n^{\enbrace{X^\ssbrace{n}, X^\ssbrace{n-1}}}$ die Randabbildung aus der Paarfolge für $\enbrace*{X^\ssbrace{n},X^\ssbrace{n-1}}$ und $j^{(n-1)}$ die Inklusion
	$\enbrace*{X^\ssbrace{n-1},\emptyset} \hookrightarrow \enbrace*{X^\ssbrace{n-1}, X^\ssbrace{n-2}}$.
\end{definition}

\begin{lemma}[{name=[{es gilt $\partial_{n-1}^{\cell} \circ \partial_n^{\cell} =0$}]}]
	Es gilt
	\[
		\partial_{n-1}^{\cell} \circ \partial_n^{\cell} =0
	\]
\end{lemma}
\begin{beweis}
	Wir schreiben die Paarsequenz für $\enbrace*{X^\ssbrace{n},X^\ssbrace{n-1}}$ dreimal übereinander und erhalten:
	\[
		\tikzexternaldisable
		\begin{tikzpicture}[baseline= (a).base]
		\node[scale=0.81] (a) at (0,0){
		\begin{tikzcd}[font=\normalsize,column sep=2.4em,arrows={line width=1.23*rule_thickness}]
			H_n \enbrace*{X^\ssbrace{n}, X^\ssbrace{n-1}} \rar["\partial_n"] \arrow[rrd, bend right,"\partial_n^{\cell}"',Firebrick3]& H_{n-1}(X^\ssbrace{n-1}) \dar[equal] \rar & \cdots\\
			& H_{n-1}\enbrace*{X^\ssbrace{n-1}} \rar["j^{(n-1)}_*"] & H_{n-1} \enbrace*{X^\ssbrace{n-1}, X^\ssbrace{n-2}} \rar["\partial_{n-1}"] 
			\arrow[rrd, bend right,"\partial_{n-1}^{\cell}"',Firebrick3] & H_{n-2} \enbrace*{X^\ssbrace{n-2}} \dar[equal] \rar & \cdots \\ 
			& & & H_{n-2} \enbrace*{X^\ssbrace{n-2}}  \rar["j^{(n-2)}_*"] & H_{n-2} \enbrace*{X^\ssbrace{n-2}, X^{(n-3)}} 
		\end{tikzcd}
		};
		\end{tikzpicture}
		\tikzexternalenable\marginnote{\enquote{$R$} aus Platzgründen weggelassen}
	\]
	Also $\partial^{\cell}_{n-1} \circ \partial^{\cell}_n =0$, da wegen der Exaktheit der Paarfolge schon $\partial_{n-1} \circ j^{(n-1)}_* =0$ ist.
\end{beweis}

\begin{definition}[{name=[{Zellulärer Kettenkomplex und zelluläre Homologie}]}]
	$\enbrace*{C_*^{\cell}(X;R),\partial_*^{\cell}}$ heißt der \bet{zelluläre Kettenkomplex}\index{zellulärer Kettenkomplex} von $X$ über $R$. 
	Seine Homologie $H_*^{\cell}(X;R)$ heißt die \Index{zelluläre Homologie} von $X$.
\end{definition}

Der zelluläre Kettenkomplex und die zelluläre Homologie sind Funktoren auf der Kategorie der CW-Komplexe und zellulären Abbildungen.

\begin{lemma}[{name=[{über $\enbrace{\coprod Q_i^\ssbrace{n}, \coprod q_i^\ssbrace{n}}_* $ und $H_k \enbrace{\coprod_i D^n, \coprod_i S^{n-1};R}$}]},label=lem:124]
	Sei $X$ ein CW-Komplex. Wähle $q_i^\ssbrace{n}$ und $Q_i^\ssbrace{n}$ für die Zellstruktur auf $X$. Dann gilt
	\begin{enumerate}[a)]
		\item Die Abbildung
		\[
			\enbrace*{{\textstyle\coprod\limits_{i \in I^{(n)}}} Q_i^\ssbrace{n}, {\textstyle\coprod\limits_{i \in I^{(n)}}} q_i^\ssbrace{n}} 
			\colon \enbrace*{{\textstyle\coprod\limits_i} D^n, {\textstyle\coprod\limits_i} S^{n-1};R}  \longrightarrow
			\enbrace*{X^\ssbrace{n}, X^\ssbrace{n-1};R} 
		\]
		induziert einen Isomorphismus in Homologie.
		\item Es gilt
		\[
			H_k \enbrace*{{\textstyle\coprod\limits_{i\in I^\ssbrace{n}}} D^n, {\textstyle\coprod\limits_{i\in I^\ssbrace{n}}} S^{n-1};R} = \begin{cases}
				R[I^\ssbrace{n}], &\text{ falls }k=n\\
				0, & \text{ sonst}
			\end{cases}
		\] 
	\end{enumerate}
\end{lemma}
\begin{beweis}
	\leavevmode
	\begin{enumerate}[a)]
		\item Sei $D^n_0 := D^n \setminus \set{0}$, $\frac{1}{2} D^n := \set{\frac{x}{2} \given x \in D^n}$ und $\frac{1}{2} D^n_0 := \frac{1}{2}D^n \setminus \set{0}$. Wir definieren
		\[
			X^\ssbrace{n-1}_{++} := X^\ssbrace{n} \setminus \bigcup_{\mathclap{i \in I^{(n)}}} Q_i^\ssbrace{n}(\set{0}) \quad \text{ und } \quad 
			X^\ssbrace{n-1}_+ := X^\ssbrace{n} \setminus \bigcup_{\mathclap{i \in I^{(n)}}} Q_i^\ssbrace{n}\enbrace*{\frac{1}{2} D^n}
		\]
		Dann sind 
		$X^\ssbrace{n-1} \hookrightarrow X_+^{(n-1)} \hookrightarrow X_{++}^{(n-1)}$ Homotopieäquivalenzen, da $\coprod_{i \in I^{(n)}} Q^\ssbrace{n}_i$ ein Homöomorphismus auf den offenen 
		$n$-Zellen ist. Betrachte:
		\[
			\begin{tikzcd}[sep=large]
				\enbrace*{\coprod D^n, \coprod S^{n-1}} \dar["\enbrace*{\coprod Q_i^\ssbrace{n}, \coprod q_i^\ssbrace{n}}"] \rar["(1)"] & \enbrace*{\coprod D^n, \coprod D^n_0} \dar &
				\enbrace*{\coprod \frac{1}{2} D^n, \coprod \frac{1}{2} D_0^n} \lar["(2)"'] \dar["(3)"] \\
				\enbrace*{X^\ssbrace{n}, X^\ssbrace{n-1}} \rar["(5)"] & \enbrace*{X^\ssbrace{n}, X^\ssbrace{n-1}_{++}}    
				& \enbrace*{X^\ssbrace{n} \setminus X^\ssbrace{n-1}_+ , X^\ssbrace{n-1}_{++} \setminus X^\ssbrace{n-1}_+} \lar["(4)"'] 
			\end{tikzcd}
		\]
		$(1), (2)$ und $(5)$ sind Homotopieäquivalenzen und induzieren Isomorphismen in Homologie. $(3)$ ist ein Homöomorphismus und induziert einen Isomorphismus in Homologie.
		$(4)$ erfüllt die Vorraussetzungen für den \hyperref[satz:91]{Ausschneideisomorphismus} und induziert daher auch einen Isomorphismus in Homologie. Damit folgt a).
		\item Da die Randabbildung aus der Paarsequenz in diesem Fall ein Isomorphismus ist, gilt 
		\begin{align*}
			H_k \enbrace*{\coprod_i D^n, \coprod_i S^{n-1};R} \xleftarrow{\cong} \bigoplus_i H_k (D^n,S^{n-1};R) &\xrightarrow{\cong} \bigoplus_i \tilde{H}_{k-1}(S^{n-1};R)\\
			&\cong \bigoplus_i \begin{cases}
				R, &\text{ falls }k=n\\
				0, &\text{ sonst}
			\end{cases} \\
			&\cong \begin{cases}
				R[I^{(n)}], &\text{ falls }k=n\\
				0, &\text{ sonst}
			\end{cases} \qedhere
		\end{align*}
	\end{enumerate}
\end{beweis}

\begin{bemerkung}[{name=[{Matrixform der zellulären Randabbildung}]}]
	Wir haben soeben gezeigt, dass die folgenden Isomorphismen existieren:
	\[
		\begin{tikzcd}[row sep=tiny]
			C_n^{\cell}(X;R) \arrow[draw=none]{d}[sloped,auto=false]{=} \rar["\partial_n^{\cell}"] & C_{n-1}^{\cell}(X;R) \arrow[draw=none]{d}[sloped,auto=false]{=}\\
			H_n \enbrace*{X^\ssbrace{n},X^\ssbrace{n-1};R} \rar \arrow[draw=none]{d}[sloped,auto=false]{\cong}
			& H_{n-1} \enbrace*{X^\ssbrace{n-1},X^\ssbrace{n-2};R} \arrow[draw=none]{d}[sloped,auto=false]{\cong}\\
			R[I^{(n)}] \rar & R[I^{(n-1)}]  
		\end{tikzcd}
	\]
	Die Randabbildung $\partial_n^{\cell}$ wird unter diesen Isomorphismen zu einer $I^{(n-1)} \times I^{(n)}$-Matrix. Für $j \in I^{(n-1)}$ und $k \in I^{(n)}$ ist der 
	$(j,k)$-Eintrag dieser Matrix genau der Abbildungsgrad einer Selbstabbildung der $(n-1)$-Sphäre: 

	Die $\enbrace{Q_i^\ssbrace{n-1}, q_i^\ssbrace{n-1}}$ induzieren einen Homöomorphismus
	$\nicefrac{\coprod_i D^{n-1}}{\coprod_i S^{n-2}} \cong \nicefrac{X^\ssbrace{n}}{X^\ssbrace{n-1}}$. Durch Komposition mit der Projektion auf die $j$-te Zelle erhalten 
	wir die folgenden Abbildung
	\[
		\nicefrac{\coprod_i D^{n-1}}{\coprod_i S^{n-2}} \xrightarrow{p_j^{(n-1)}}  \nicefrac{D^{n-1}}{S^{n-2}} \cong S^{n-1}.
	\] 
	Der $(j,k)$-te Eintrag ist nun der 
	Abbildungsgrad von 
	\[
		\begin{tikzcd}
			S^{n-1} \rar["q^{(n-2)}_k"] & X^\ssbrace{n-1} \rar & \nicefrac{X^\ssbrace{n-1}}{X^\ssbrace{n-2}} \rar["\cong"] & \nicefrac{\coprod_i D^{n-1}}{\coprod_i S^{n-2}} \rar["p_j^{(n-1)}"]
			& S^{n-1}
		\end{tikzcd}
	\]
\end{bemerkung}

\begin{definition}[{name=[natürlicher Isomorphismus]}]
	Eine \hyperref[def:16]{natürliche Transformation} $\tau$ zwischen Funktoren $F,G\colon \mathcal{C} \to \mathcal{D}$ heißt ein \Index{natürlicher Isomorphismus}, wenn $\tau_C \colon F(C) \to G(C)$ für alle Objekte $C$ von $\mathcal{C}$ ein Isomorphismus in $\mathcal{D}$ ist.
\end{definition}

\begin{satz}[{name=[Natürlicher Isomorphismus zwischen zellulärer und singulärer Homologie]}]
	Für CW-Komplexe gibt es einen natürlichen Isomorphismus $\tau_X \colon H_*^{\cell}(X;R) \xrightarrow{\cong} H_*(X;R)$.
\end{satz}
\begin{beweis}[{name={mit \autoref{lem:128}}}]
	Betrachte folgendes Diagramm; die diagonalen Folgen sind Ausschnitte aus den Paarsequenzen: 
	\[
		{\begin{tikzcd}[row sep=large, column sep=tiny, arrows=thick]
			& & & \underset{=0}{H_n \enbrace*{X^\ssbrace{n+1}, X^\ssbrace{n}}} \\ 
			0 \stackrel{\mbox{\scriptsize\ref{lem:128} a)}}{=} H_n \enbrace*{X^\ssbrace{n-1}} \arrow[dr,Firebrick3] & 
			& H_n \enbrace*{X^\ssbrace{n+1}} \rar["\cong","\mbox{\scriptsize\ref{lem:128} b)}"']  \arrow[ur, Goldenrod2] & H_n(X)\\  
			& H_n \enbrace*{X^\ssbrace{n}} \arrow[ur, two heads, Goldenrod2] \arrow[dr, hook,"j^{(n)}_*", Firebrick3]\\
			H_{n+1} \enbrace*{X^\ssbrace{n+1},X^\ssbrace{n}} \arrow[ur, Goldenrod2,"\partial_{n+1}^{X^\ssbrace{n+1},X^\ssbrace{n}}"]\arrow[rr,"\partial_{n+1}^{\cell}"] 
			& & H_n \enbrace*{X^\ssbrace{n},X^\ssbrace{n-1}} \arrow[rr,"\partial_n^{\cell}"] \arrow[dr,"\partial_n^{X^\ssbrace{n},X^\ssbrace{n-1}}", Firebrick3]
			& & H_{n-1}  \enbrace*{X^\ssbrace{n-1},X^\ssbrace{n-2}}  \\
			& & & H_{n-1} \enbrace*{X^\ssbrace{n-1}} \arrow[ur, hook, DodgerBlue3, "j^{(n-1)}_*"'] \\
			& & \underset{=0}{H_{n-1}\enbrace*{X^\ssbrace{n-2}} } \arrow[ur, DodgerBlue3]
		\end{tikzcd}}
		\marginnote{\enquote{$R$} aus Platzgründen weggelassen}
	\]
	Es folgt, dass $j_*^{(n)}$ wegen Exaktheit einen Isomorphismus $H_n(X^\ssbrace{n};R) \xrightarrow{\cong} \ker\partial_n^{\cell}$ induziert. Da $j_*^{(n)}$ injektiv ist,
	induziert $j_*^{(n)}$ außerdem einen Isomorphismus $\im \partial_{n+1}^{X^\ssbrace{n+1}, X^\ssbrace{n}} \xrightarrow{\cong} \im \partial_{n+1}^{\cell}$.
	% Es folgt, dass $j_*^{(n)}$ Isomorphismen $H_n(X^\ssbrace{n};R) \xrightarrow{\cong} \ker \partial_n^{\cell}$ und
	% $\im \partial_{n+1}^{X^\ssbrace{n+1}, X^\ssbrace{n}} \xrightarrow{\cong} \im \partial_{n+1}^{\cell}$ induziert.
	Insgesamt erhalten wir natürliche Isomorphismen
	\[
		H_n(X;R) \cong H_n \enbrace*{X^\ssbrace{n+1};R} \cong \nicefrac{H_n(X^\ssbrace{n};R)}{ \im \partial_{n+1}^{X^\ssbrace{n+1}, X^\ssbrace{n}}}
		\cong \nicefrac{\ker \partial_n^{\cell}}{\im \partial_{n+1}^{\cell}} \cong H_n^{\cell}(X;R) \qedhere
	\]
\end{beweis}

\begin{lemma}[{name=[{Die Homologie des $n$-Gerüsts und die von $X^\ssbrace{n} \hookrightarrow X$ induzierte Abbildung}]},label=lem:128]
	Sei $X$ ein CW-Komplex.
	\begin{enumerate}[a)]
		\item Für $k >n$ ist $H_k(X^\ssbrace{n};R)=0$.
		\item Für $k <n$ induziert die Inklusion $l_n\colon X^\ssbrace{n} \hookrightarrow X$ einen Isomorphismus $(l_n)_* \colon H_k(X^\ssbrace{n};R) \to H_k(X;R)$.
	\end{enumerate}
\end{lemma}
\begin{beweis}
	\leavevmode
	\begin{enumerate}[a)]
		\item Ist $k>n$, so erhalten wir aus der Paarfolge einen Isomorphismus
		\[
			\begin{tikzcd}[column sep=small]
				\underset{\textstyle\StackText{\ref{lem:124}}{=} 0}{H_{k+1}\enbrace*{X^\ssbrace{n}, X^\ssbrace{n-1};R}} \rar &  H_k(X^\ssbrace{n-1};R) \rar["\cong"] & 
				H_k(X^\ssbrace{n};R) \rar 
				& H_k \underset{\textstyle\StackText{\ref{lem:124}}{=}0}{\enbrace*{X^\ssbrace{n}, X^\ssbrace{n-1};R} }
			\end{tikzcd}
		\]
		Da $H_k(X^{\ssbrace{0}};R)=0$ für $k>0$, folgt a) nun durch endliche Induktion für $n=0,1, \ldots , k-1$.
		\item Ist $n >k$, so erhalten wir aus der Paarfolge
		\[
			\begin{tikzcd}[column sep=small]
				\underset{\textstyle\StackText{\ref{lem:124}}{=}0}{H_{k+1} \enbrace*{X^\ssbrace{n+1},X^\ssbrace{n};R}} \rar & H_k \enbrace*{X^\ssbrace{n};R} \rar["\cong"] 
				& H_k \enbrace*{X^\ssbrace{n+1};R} \rar 
				& \underset{\textstyle\StackText{\ref{lem:124}}{=}0}{H_k \enbrace*{X^\ssbrace{n+1},X^\ssbrace{n};R}}
			\end{tikzcd}
		\]
		Für $m \ge n >k$ induziert die Inklusion $X^\ssbrace{n} \hookrightarrow X^\ssbrace{m}$ daher einen Isomorphismus in Homologie $H_k(X^\ssbrace{n};R) \to H_k(X^{(m)};R)$. 
		Ist $\dim X$ endlich, so folgt die Behauptung. 
	
		Im allgemeinen Fall beobachten wir zunächst, dass es zu jeder stetigen Abbildung $\sigma\colon\abs{\Delta^j}\to X$ ein $n$ gibt mit 
		$\im \sigma \subseteq X^\ssbrace{n}$, da $\im \sigma \subseteq X$ kompakt ist und daher nur endlich viele offene Zellen von $X$ trifft (siehe \autoref{lem:116}).
		Also gilt
		\(
			C_*(X;R) = \bigcup_n C_*(X^\ssbrace{n};R)
		\).
		Damit folgt \begin{enumerate}[(1)]
			\item $\forall x \in H_k(X;R) $ existiert ein $n$, sodass $x$ im Bild von $ H_k(X^\ssbrace{n};R) \to H_k(X;R)$.
			\item $\forall x\in H_k(X^\ssbrace{n};R)$ mit trivialem Bild von $x$ in $H_k(X;R)$, gibt es ein $m \ge n$, sodass das Bild von $x$ in $H_k(X^\ssbrace{m};R)$
			trivial ist.
		\end{enumerate}
		Da für $m \ge n >k$ die Abbildung $H_k(X^\ssbrace{n};R) \to H_k(X^\ssbrace{m};R)$ ein Isomorphismus ist, ergibt sich damit die Behauptung. \qedhere 
	\end{enumerate}
\end{beweis}

\begin{korollar}[{name=[Die Eulercharakteristik ist eine topologische Invariante]},label=kor:129]
	Die Eulercharakteristik eines endlichen simplizialen Komplexes ist eine topologische Invariante.
\end{korollar}
\begin{beweis}
	Sei $a_n$ die Anzahl der $n$-Simplizes von $K$. 
	Dann gibt es eine CW-Struktur auf $\abs*{K}$, die für jedes $n$ genau $a_n$ Zellen hat.\marginnote{Man beachte, dass \ref{satz:410} für beliebige $\mathbb{Z}$-Kettenkomplexe und deren Homologie gilt}
	\[
		\chi(K) = \sum_{i=0}^{\infty} (-1)^i a_i \StackText{\ref{lem:124}}{=} \sum_{i=0}^{\infty} (-1)^i \Rg C_i^{\cell} \enbrace[\big]{\abs*{K}} 
		\StackText{\ref{satz:410}}{=}
		\sum_{i=0}^{\infty} (-1)^i \Rg H_i^{\cell} \enbrace[\big]{\abs*{K}} = \sum_{i=0}^{\infty} (-1)^i \Rg H_i \enbrace[\big]{\abs*{K} }    
	\]
	Da die $H_i \enbrace[\big]{\abs*{K}}$ eine topologische Invariante sind, ist auch $\chi(K)$ eine topologische Invariante.
\end{beweis}

\begin{bemerkung}[{name=[{n-te Bettizahl und Eulercharakteristik eines CW-Komplexes}]}]
	Ist $X$ ein CW-Komplex mit endlich vielen $n$- und $n+1$-Zellen, so ist $H_n(X)$ ein endlich erzeugter $\mathbb{Z}$-Modul. 
	Man definiert dann die \bet{$\mathbf{n}$-te Bettizahl}\index{n-te Bettizahl@$n$-te Bettizahl} von $X$ als $b_n := \Rg H_n(X)$. 
	Ist $X$ ein endlicher CW-Komplex, so definiert man die Euler-Charakteristik von $X$ als\index{Euler-Charakteristik!eines CW-Komplex}
	\[
		\chi(X) := \sum_{n=0}^{\infty} (-1)^{n} b_n.
	\]
	Ist $a_n$ die Anzahl der $n$-Zellen von $X$, so gilt $\chi(X)=\sum_{n=0}^{\infty} (-1)^n a_n$.
\end{bemerkung}

\begin{beispiel}[{name=[{Homologie des komplex projektiven Raumes}]}]
	Es gilt 
	\[
		H_k \enbrace*{\mathbb{C}P^n;R} \cong  \begin{cases}
			R, &\text{ falls }k=0,2,4,\ldots ,2n\\
			0 ,&\text{ sonst}
		\end{cases}
	\]
\end{beispiel}
\begin{beweis}
	Aus \hyperref[enum:115:5]{Beispiel \ref*{bsp:115} \ref*{enum:115:5}} wissen wir, dass es auf $\mathbb{C}P^n$ eine CW-Struktur gibt mit genau einer Zelle in der Dimension $0,2,4, \ldots ,2n$ (und keine weiteren 
	Zellen). Der zelluläre Kettenkomplex hat daher folgende Form:
	\[
		\begin{tikzcd}[row sep=small,column sep=2.0em]
			0 & 1 & 2 & 3 & 4 & \cdots & 2n-1 & 2n & 2n +1 & \cdots \\
			R & 0 \lar & R \lar& 0 \lar & R \lar & \cdots \lar & 0 \lar& R \lar & 0 \lar & \cdots \lar
		\end{tikzcd}
	\]
	Alle Randabbildungen sind trivial und die Homologie von $\mathbb{C}P^n$ stimmt mit den zellulären Kettenkomplexen überein.
\end{beweis}

\begin{beispiel}[{name=[{ Homologie des reell projektiven Raumes}]}]
	Es gilt 
	\[
		H_k(\mathbb{R}P^n;\mathbb{Z})\cong \begin{cases}
			\mathbb{Z}, &\text{ falls }k=0 \text{ oder }(k=n \text{ und $k$ ungerade})\\
			\nicefrac{\mathbb{Z}}{2 \mathbb{Z}} ,&\text{ falls } 0<k<n \text{ und $k$ ungerade} \\
			0 ,&\text{ sonst}
		\end{cases}
	\]
\end{beispiel}
\begin{beweis}
	Aus \hyperref[enum:115:4]{Beispiel \ref*{bsp:115} \ref*{enum:115:4}} wissen wir, dass es eine CW-Struktur auf $\mathbb{R}P^n$ gibt, die für $k=0,\ldots ,n$ genau eine $k$-Zelle und keine weiteren Zellen besitzt.
	Der zelluläre Kettenkomplex von $\mathbb{R}P^n$ (über $\mathbb{Z}$) hat also die Form:
	\[
		\begin{tikzcd}[column sep=large]
			\mathbb{Z} & \mathbb{Z} \lar["\partial_1^{\cell}"'] & \mathbb{Z} \lar["\partial_2^{\cell}"'] & \cdots \lar & \mathbb{Z} \lar & \mathbb{Z} 
			\lar["\partial_n^{\cell}"'] & 0 \lar
		\end{tikzcd}
	\]
	Wir müssen die Randabbildung verstehen. 
	In der CW-Struktur können wir $(\mathbb{R}P^n)^{(k)}$ mit $\mathbb{R}P^k$ identifizieren. 
	Die anklebende Abbildung $q^\ssbrace{k+1} \colon S^k \to \mathbb{R}P^k$ der $(k+1)$-Zelle ist die Projektion $(x_0, \ldots ,x_k) \mapsto [x_0 : \ldots : x_k]$. 
	Die charakteristische Abbildung $Q^\ssbrace{k+1} \colon D^{k+1} \to \mathbb{R}P^{k+1}$ der $(k+1)$-Zelle ist gegeben durch 
	\[
		(x_0, \ldots ,x_k) \mapsto \benbrace*{x_0 : \ldots : x_k : {\textstyle\sqrt{\cramped 1- \norm{(x_1,\ldots, x_k)}^2}}} 
	\]
	Betrachte folgendes kommutatives Diagramm:\marginnote{Das linke untere Quadrat kommutiert, da die Randabbildung aus der Paarfolge eine natürliche Transformation ist. Das rechte, da die zugrunde liegenden Abbildungen kommutieren.}
	\[
		\begin{tikzcd}[column sep=huge, row sep=large]
			C_{k+1}^{\cell}(\mathbb{R}P^n) \arrow[rr,"\partial_{k+1}^{\cell}"] \dar[equal] & & C_k^{\cell}(\mathbb{R}P^n) \dar[equal]\\
			H_{k+1} \enbrace*{\mathbb{R}P^{k+1},\mathbb{R}P^k} \rar["\partial^{\mathbb{R}P^{k+1},\mathbb{R}P^k}"] & H_k(\mathbb{R}P^k) \rar
			& H_k \enbrace*{\mathbb{R}P^k,\mathbb{R}P^{k-1}} \\
			H_{k+1}(D^{k+1},S^k) \uar["\enbrace*{Q^\ssbrace{k+1},q^\ssbrace{k+1}}_*","\cong"'] \rar["\partial^{D^{k+1},S^k}","\cong"'] & \tilde{H}_k(S^k) \uar["\enbrace*{q^\ssbrace{k+1}}_*"] \rar["j_*"] 
			& H_k(S^k,S^{k-1}) \uar["\enbrace*{q^\ssbrace{k+1},q^\ssbrace{k}}_*"']
		\end{tikzcd}
	\]
	Das Bild von $\partial_{k+1}^{\cell}$ stimmt also mit dem Bild der Komposition 
	\[
		\begin{tikzcd}[column sep=5.5em]
			H_k(S^k) \rar["j_*"] & H_k(S^k,S^{k-1}) \rar["\enbrace*{q^\ssbrace{k+1},q^\ssbrace{k}}_*"] & H_k \enbrace*{\mathbb{R}P^k, \mathbb{R}P^{k-1}} 
		\end{tikzcd}
	\]
	überein. 
	Sei $i_{\pm} \colon D^k \to S^k$ definiert durch $i_{\pm}(x)= \enbrace[\big]{\pm x,\pm\sqrt{\cramped 1-{\scriptscriptstyle\norm{\textstyle x}^{2}}}}$. 
	Sei $l\colon S^{k-1} \hookrightarrow i_-(D^k)$ die Inklusion und $R \colon S^k \to S^k$ die Punktspiegelung $x \mapsto -x$. 
	Es ist $i_+ = R \circ i_-$. Betrachte
	\marginnote{faktorisieren über 0}
	\begin{equation*}
		\begin{tikzcd}[column sep=huge, row sep=large]
			& \tilde{H}_k(S^k) \dar["j_*"] \\
			\bigoplus\limits_{\displaystyle{\color{light_gray}y_- \in } H_k(D^k,S^{k-1})}^{\displaystyle {\color{light_gray}y_+ \in }H_k(D^k,S^{k-1})} 
			\rar["\bigoplus\limits^{\scriptstyle(i_+,i_+|_{S^{k-1}})_*}_{\scriptstyle(i_-,i_-|_{S^{k-1}})_*}","\cong"'] 
			& H_k(S^k,S^{k-1}) \dar["\enbrace*{q^\ssbrace{k+1},q^\ssbrace{k}}_*"] \rar["\cong"',"{\begin{pmatrix}
				\enbrace*{\id_{S^k},l}_* \\ \enbrace{R,R|_{S^{k-1}}}_*
			\end{pmatrix}}"] 
			& \bigoplus\limits_{\displaystyle H_k \enbrace*{S^k,i_-(D^k)} {\color{light_gray} \ni x_-}}^{\displaystyle H_k \enbrace*{S^k,i_-(D^k)} {\color{light_gray} \ni x_+}} \\
			& H_k \enbrace*{\mathbb{R}P^k,\mathbb{R}P^{k-1}} 
		\end{tikzcd}
		\label{eq:1213} \tag{\#}
	\end{equation*}
	In der Übungen haben wir gesehen, dass die horizontalen Abbildungen Isomorphismen sind.
	Sei $x \in \tilde{H}_k(S^k)$ ein Erzeuger.
	Seien $(x_+,x_-)$ und $(y_+,y_-)$ die Elemente, die $j_*(x)$ unter den horizontalen Isomorphismen in \eqref{eq:1213} entsprechen.
	Mit Hilfe der Paarfolge zu $\enbrace{S^k,i_-(D^k)}$ sehen wir, dass die Inklusion $j_- \colon (S^k, \emptyset) \hookrightarrow \enbrace{S^k,i_-(D^k)}$ einen Isomorphismus $(j_-)_* \colon \tilde{H}_k(S^k) \to \tilde{H}_k \enbrace{S^k, i_-(D^k)}$ induziert. 
	Es ist nun $(\id_{S^k},l) \circ j = j_-$ und $(R,R|_{S^{k-1}}) \circ j= j_- \circ R$. 
	Damit ist $x_+$ ein Erzeuger von $H_k(S^k,i_+(D^k))$. 
	Weiter entspricht $R_*$ der Multiplikation mit $(-1)^{k+1}$ nach \autoref{bsp:106}. Damit folgt
	\begin{align*}
		 x_- = \enbrace*{R,R\big|_{S^{k+1}}}_* \circ j_* (x) = \enbrace*{\id_{S^k},l} \circ j_* \circ R_*(x) = \enbrace*{\id_{S^k},l}_* \circ j_* \enbrace*{(-1)^{k+1}x} = 
		 (-1)^{k+1}\cdot x_+    
	\end{align*}
	% Da $(\id_{S^k},l) \circ (i_+,i_+|_{S^{k-1}}) = (R,R|_{S^{k-1}}) \circ (i_-,i_-|_{S^{k-1}})$, folgt auch $y_- = (-1)^{k+1} y_+$.
	Die Komposition der horizontalen Isomorphismen in \eqref{eq:1213} ist gegeben durch\marginnote{Ich sehe noch immer nicht ganz ein, warum dies so ist. Über genauere Ausführungen wäre ich sehr dankbar}
	\[
		\begin{pmatrix}
			\enbrace*{i_+,l \circ i_+|_{S^{k-1}}}_* & 0 \\
			0 & \enbrace*{R,R|_{S^{k-1}}}_*  \circ \enbrace*{i_-, i_-|_{S^{k-1}}}_*
		\end{pmatrix}=
		\begin{pmatrix}
			\enbrace*{i_+, l \circ i_+|_{S^{k-1}}}_* & \\
			& \enbrace*{i_+,l \circ i_+|_{S^{k-1}}}_*
		\end{pmatrix}
	\]
	Es folgt, dass $y_+$ ein Erzeuger von $H_k(D^k,S^{k-1})$ ist und $y_- = (-1)^{k+1} y_+$
	% \todo{Für die Gleichung unten brauchen wir doch $\scriptscriptstyle\enbrace*{i_-,i_-|_{S^{k-1}}}_*(y_-) = (-1)^{k+1} \enbrace*{i_+,i_+|_{S^{k-1}}}_*(y_+)$ ??}
	ist. 
	In \eqref{eq:1213} wird $y_+$ auf $(q^\ssbrace{k+1},q^\ssbrace{k})_* \circ (i_+,i_+|_{S^{k-1}})_*(y_+) = \enbrace*{Q^\ssbrace{k},q^\ssbrace{k}}_*(y_+)$ abgebildet. 
	Insbesondere wird $y_+$ auf einen Erzeuger $z_k \in H_k(\mathbb{R}P^k,\mathbb{R}P^{k-1})$ abgebildet, da diese Abbildung nach \ref{lem:124} ein Isomorphismus ist. Es folgt 
	\begin{align*}
		\enbrace*{q^\ssbrace{k+1},q^\ssbrace{k}}_* \enbrace[\big]{j_*(x)} &= \enbrace*{q^\ssbrace{k+1},q^\ssbrace{k}}_* \enbrace*{\enbrace*{i_+,i_+\big|_{S^{k-1}}}_*(y_+) +
		\enbrace*{i_-,i_-\big|_{S^{k-1}}}_*(y_-) } \\
		&= \enbrace*{Q^\ssbrace{k},q^\ssbrace{k}}_* \enbrace*{y_+ + (-1)^{k+1} y_+} \\
		& = \begin{cases}
			2 z_k, &\text{ falls }k \text{ ungerade}\\
			0 , &\text{ sonst}
		\end{cases}  
	\end{align*}
	Es folgt 
	\[
		\im  \partial_{k+1}^{\cell} = \begin{cases}
		2 \cdot C^{\cell}_k(X), &\text{ falls } k \text{ ungerade}\\
		0 , &\text{ sonst}
	\end{cases}
	\]
	Der zelluläre Kettenkomplex ist daher isomorph zu:
	\[
		\begin{tikzcd}
			\mathbb{Z} & \mathbb{Z} \lar["0"'] & \mathbb{Z} \lar["2"'] & \mathbb{Z} \lar["0"'] & \cdots \lar & \overset{n}{\mathbb{Z}} \lar & 0 \lar & 0 \lar
		\end{tikzcd}
	\]
	Womit die Behauptung folgt.
\end{beweis}

Für $R=\nicefrac{\mathbb{Z}}{2 \mathbb{Z}}$ ist
\[
	H_k \enbrace*{\mathbb{R}P^n; \nicefrac{\mathbb{Z}}{2 \mathbb{Z}}} = \begin{cases}
		\nicefrac{\mathbb{Z}}{2 \mathbb{Z}}, &\text{ falls }k=0,\ldots ,n\\
		0 ,&\text{ sonst}
	\end{cases}
\]

\begin{satz}[{name=[{Isomorphie zwischen $C_*(K;R)$ und $C^{\cell}_n \enbrace*{\abs*{K};R}$}]}]
	Sei $K$ ein geordneter simplizialer Komplex. 
	Für $\sigma = \set{v_0, \ldots ,v_n} \in \Sigma_n$ mit $v_0 < v_1 < \ldots  < v_n$ sei $f_\sigma \colon \abs*{\Delta^n} \to \abs*{K}$ gegeben durch $f_\sigma \enbrace*{\sum_{i=0}^{n} t_i e_i} = \sum_{i=0}^{n} t_i v_i$. 
	Wir erhalten 
	\[
		[f_\sigma] \in H_n \enbrace*{\abs*{K}^{(n)}, \abs*{K}^{(n-1)};R} = C^{\cell}_n \enbrace*{\abs*{K};R}. 
	\]
	Diese Zuordnung definiert einen Isomorphismus zwischen dem simplizialen Kettenkomplex von $K$ und dem zellulären Kettenkomplex von $\abs*{K}$
	\[
		\tau \colon C_*(K;R) \longrightarrow C_*^{\cell}\enbrace*{\abs*{K},R } 
	\]
	Insbesondere ist $H_*(K;R) \cong H_*(\abs*{K};R)$ eine topologische Invariante von $K$.
\end{satz}
\begin{beweis}
	Betrachte 
	\[
		\begin{tikzcd}
			\smashoperator[r]{\bigoplus\limits_{\sigma \in \Sigma_n}} R \dar["\cong"'] \rar["\cong"]
			& \smashoperator[r]{\bigoplus\limits_{\sigma \in \Sigma_n}} H_n \enbrace*{\abs*{\Delta^\sigma}, \abs*{\partial \Delta^\sigma};R} \dar["\cong"'] \\
			C_n(K;R) \rar["\tau"] & H_n \enbrace*{\abs*{K}^{(n)}, \abs*{K}^{(n-1)};R} 
		\end{tikzcd}
	\]
	Nun ist auch $\tau$ ein Isomorphismus, da $[f_\sigma]$ ein Erzeuger von $H_n \enbrace*{\abs*{\Delta^\sigma}, \abs*{\partial \Delta^\sigma};R}$ ist, siehe Blatt 8. 
	Es bleibt zu zeigen, dass $\sigma \mapsto  [f_\sigma]$ mit den Randabbildungen vertauscht. 
	Dies ergibt sich sofort aus $f_{\partial_j \sigma} = \partial_j f_\sigma$.
\end{beweis}

\newpage

\section{Kleine Simplizes} % (fold)
\label{sec:13}

\emph{In diesem Abschnitt holen wir den Beweis des Ausschneidungssatzes nach, der vorher schon benutzt wurde.}

\begin{satz}[{name={Ausschneidung}}]
	Sei $(X,A)$ ein Paar von topologische Räumen. Sei $L \subseteq A$ mit $\overline{L} \subseteq \mathring{A}$. Dann induziert die Inklusion 
	$(X \setminus L, A \setminus L) \hookrightarrow (X,A)$ den Ausschneidungs-Isomorphismus $H_* \enbrace*{X \setminus L, A \setminus L;R} \grenzw{\cong} H_*(X,A;R)$.
\end{satz}

\begin{figure}[ht]
	\centering{
	\tikzsetnextfilename{skizze_kleine_simplizes_ausschneidung}
	\begin{tikzpicture}[scale=0.8, line join=bevel]
		% \draw[help lines] (-4,-3) grid (4,3);
		\draw[very thick] (0,0) ellipse(4cm and 2.5cm);
		\node at (-3,2) {$X$};

		\draw[ultra thick] (1,0.4) ellipse(2cm and 1.5cm);
		\node at (-0.5,1.8) {$A$};

		\draw[ultra thick, Firebrick3, line join=miter] (0,-0.5) rectangle (2,1) node[midway]{$L$};

		% großer Simplex
		\coordinate (a1) at (-2,-1.7);
		\coordinate (b1) at (0.3,-0.3);
		\coordinate (c1) at (-3,0.7);
		\draw[very thick,dashed,DodgerBlue3] (a1) .. controls +(1,0) and +(170:1) .. (b1) .. controls +(120:1.5) and +(-50:1.5) .. (c1) .. controls +(-80:1.7) and +(110:1.5) .. cycle;
		\foreach \x in {a1,b1,c1}{
			\draw[fill=DodgerBlue3, thick] (\x) circle[radius=0.04];
		}

		% kleiner Simplex in X\A
		\coordinate (a2) at (-2.3,1);
		\coordinate (b2) at (-1.9,0.8);
		\coordinate (c2) at (-2,1.5);
		\draw[very thick, SeaGreen3] (a2) .. controls +(-40:0.2) and +(180:0.3) .. (b2) .. controls +(70:0.4) and +(-40:0.5) .. (c2) .. controls +(-80:0.7) and +(110:0.5) .. cycle;
		\foreach \x in {a2,b2,c2}{
			\draw[fill=SeaGreen3,thick] (\x) circle[radius=0.04];
		}

		% kleiner Simplex in X und A
		\coordinate (a3) at (0.5,-1.5);
		\coordinate (b3) at (1.3,-1.6);
		\coordinate (c3) at (0.9,-0.8);
		\draw[very thick, SeaGreen3] (a3) .. controls +(20:0.2) and +(160:0.3) .. (b3) .. controls +(70:0.4) and +(-40:0.5) .. (c3) .. controls +(-80:0.5) and +(75:0.5) .. cycle;
		\foreach \x in {a3,b3,c3}{
			\draw[fill=SeaGreen3,thick] (\x) circle[radius=0.04];
		}

		% kleiner Simplex in L und A
		\coordinate (a4) at (1.75,-0.1);
		\coordinate (b4) at (2.3,0.4);
		\coordinate (c4) at (1.8,0.75);
		\draw[very thick, SeaGreen3] (a4) .. controls +(20:0.2) and +(160:0.3) .. (b4) .. controls +(85:0.5) and +(10:0.2) .. (c4) .. controls +(-95:0.7) and +(75:0.5) .. cycle;
		\foreach \x in {a4,b4,c4}{
			\draw[fill=SeaGreen3,thick] (\x) circle[radius=0.04];
		}

		% kleiner Simplex in L und A
		\coordinate (a5) at (-0.2,0.75);
		\coordinate (b5) at (0.3,0.6);
		\coordinate (c5) at (0.1,1.3);
		\draw[very thick, SeaGreen3] (a5) .. controls +(20:0.2) and +(160:0.3) .. (b5) .. controls +(85:0.5) and +(10:0.2) .. (c5) .. controls +(220:0.3) and +(80:0.5) .. cycle;
		\foreach \x in {a5,b5,c5}{
			\draw[fill=SeaGreen3,thick] (\x) circle[radius=0.04];
		}

		% kleiner Simplex in X und A
		\coordinate (a6) at (0.7,1.55);
		\coordinate (b6) at (1.3,1.6);
		\coordinate (c6) at (1,2.3);
		\draw[very thick, SeaGreen3] (a6) .. controls +(20:0.2) and +(160:0.3) .. (b6) .. controls +(85:0.5) and +(10:0.2) .. (c6) .. controls +(-80:0.3) and +(60:0.5) .. cycle;
		\foreach \x in {a6,b6,c6}{
			\draw[fill=SeaGreen3,thick] (\x) circle[radius=0.04];
		}
	\end{tikzpicture}
	\caption{Skizze zur Ausschneidung mit kleinen Simplizes}}
\end{figure}

\begin{definition}[{name=[{kleine Simplizes}]}]
	Sei $\mathcal{U}$ eine offene Überdeckung von $X$. 
	Ein singulärer Simplex $\sigma \colon \abs*{\Delta^n} \to X$ heißt \bet{$\mathcal{U}$-klein}\index{U-klein@$\mathcal{U}$-klein}, falls es ein $U \in \mathcal{U}$ gibt mit $\im \sigma  \subseteq U$. 
	Die $\mathcal{U}$-kleinen Simplizes erzeugen den Unterkettenkomplex $\enbrace*{C_*^{\mathcal{U}}(X;R),\partial_*}$ des simplizialen Kettenkomplexes $\enbrace*{C_*(X;R),\partial_*}$.
\end{definition}

\begin{satz}[{name={Satz über kleine Simplizes}}]
	Die Inklusion $C_*^\mathcal{U}(X;R) \hookrightarrow C_*(X;R)$ induziert einen Isomorphismus in Homologie.
\end{satz}

\begin{bemerkung}[{name=[{Satz über kleine Simplizes für Paare von topologischen Räumen}]}]
	Der Satz über kleine Simplizes lässt sich leicht auf Paare von topologischen Räumen verallgemeinern. 
	Für eine Überdeckung $\mathcal{U}$ von $X$ setzen wir 
	\[
		C_*^\mathcal{U} \enbrace*{X,A;R} := \frac{C_*^\mathcal{U}(X;R)}{C_*(A;R) \cap C_*^\mathcal{U}(X;R)} 
	\]
	Dann induziert die Inklusion $C_*^\mathcal{U}(X,A;R) \hookrightarrow C_*(X,A;R)$ einen Isomorphismus in Homologie. 
	Der Beweis benutzt die Paarfolge und das Fünfer-Lemma.
\end{bemerkung}

\begin{beweis}[{name={des Ausschneidungssatzes}}]
	Seien $\mathcal{U} := \set{\mathring{A},X \setminus \overline{L}}$ und $\mathcal{U}_L := \set{\mathring{A} \setminus L, X \setminus \overline{L}}$. 
	Nach dem Satz über kleine Simplizes induzieren die beiden vertikalen Inklusionen in folgendem Diagramm Isomorphismen in Homologie:
	\[
		\begin{tikzcd}[row sep=large]
			C_*^{\mathcal{U}_L} \enbrace*{X \setminus L, A \setminus L;R} \rar["\cong"] \dar[hook]
			& C_*^\mathcal{U}(X,A;R) \dar[hook] \\
			C_*(X \setminus L, A \setminus L;R) \rar["(\star)"] & C_*(X,A;R)
		\end{tikzcd}
	\]
	Wie man sich leicht überlegt, ist außerdem die obere Abbildung ein Isomorphismus.
	Daher induziert auch die Inklusion $(\star)$ einen Isomorphismus in $H_*$.
\end{beweis}

\begin{wrapfigure}[5]{r}[2cm]{0.33\textwidth}
	\centering{
	\tikzsetnextfilename{zerteilung_simplex}
	\begin{tikzpicture}[scale=4.0, rotate=200]
		\coordinate (a) at (0,0);
		\coordinate (b) at (1,0);
		\coordinate (c) at (60:1);
	
		\draw[thick] (a) -- (b) -- (c) -- cycle;
	
		\coordinate (ab) at ($(a)!.5!(b)$);
		\coordinate (ac) at ($(a)!.5!(c)$);
		\coordinate (bc) at ($(b)!.5!(c)$);
	
		\draw[darkgray,thick, name path=erster] (a) -- (bc);
		\draw[darkgray,thick, name path=zweiter] (b) -- (ac);
		\draw[darkgray, thick] (c) -- (ab);
		\path[name intersections={of=erster and zweiter, name=i}] (i-1) coordinate (m);
	
		\coordinate (aab) at ($(a)!.5!(ab)$);
		\coordinate (am) at ($(a)!.5!(m)$);
		\coordinate (abm) at ($(ab)!.5!(m)$);
		\coordinate (abb) at ($(ab)!.5!(b)$);
		\coordinate (bm) at ($(b)!.5!(m)$);
	
		\draw[gray, thick, name path=dritter] (a) -- (abm);
		\draw[gray, thick, name path=vierter] (ab) -- (am);
		\draw[gray, thick] (m) -- (aab);
		\path[name intersections={of=dritter and vierter, name=i}] (i-1) coordinate (m1);
	
		\draw[gray, thick, name path=fuenfter] (b) -- (abm);
		\draw[gray, thick, name path=sechster] (m) -- (abb);
		\draw[gray, thick] (ab) -- (bm);
		\path[name intersections={of=fuenfter and sechster, name=i}] (i-1) coordinate (m2);
	
	
		% drawing circles
		\foreach \x in {am,bm,abm,aab,abb}{
			\draw[thick,fill=IndianRed2] (\x) circle[radius=0.01];
		}
		\foreach \x in {a,b,c}{
			\draw[thick,fill=DodgerBlue2] (\x) circle[radius=0.015];
		}
		\foreach \x in {ab,bc,ac}{
			\draw[thick,fill=SeaGreen3] (\x) circle[radius=0.015];
		}
		\draw[thick,fill=Yellow2] (m) circle[radius=0.015];
		\foreach \x in {m1,m2}{
			\draw[thick,fill=PaleTurquoise2] (\x) circle[radius=0.01];
		}
	\end{tikzpicture}
	\caption{Zerteilung eines $2$-Simplex}\label{fig:zerteilung_simplex}
	}
\end{wrapfigure}
Um den Satz über kleine Simplizes beweisen zu können, müssen wir uns überlegen, wie wir einen Simplex systematisch in kleinere zerlegen können, siehe dazu \cref{fig:zerteilung_simplex}. Dies führt uns zu folgender Definition


\begin{definition}[{name=[{Schwerpunkt und \enquote{Abkegeln}}]}]
	$s_n := \enbrace*{\frac{1}{n+1}, \ldots , \frac{1}{n+1}} \in \abs*{\Delta^n}$ heißt der \Index{Schwerpunkt} von $\abs*{\Delta^n}$. 
	Ist $\sigma \colon \abs*{\Delta^p} \to \abs*{\Delta^n}$ ein singulärer $p$-Simplex, so definieren wir $(K_{\Delta^n})(\sigma) \colon \abs*{\Delta^{p+1}} \to \abs*{\Delta^n}$ durch \marginnote{\enquote{Abkegeln}}
	\[
		\enbrace*{K_{\Delta^n}} (\sigma) \enbrace*{t_0, \ldots , t_{p+1}} = t_0 \cdot s_n + (1-t_0) \cdot \sigma \enbrace*{\frac{(t_1,\ldots ,t_{p+1})}{1-t_0}} 
	\]
	Wir erhalten eine Abbildung $\enbrace*{K_{\Delta^n}}_p \colon C_p \enbrace*{\abs*{\Delta^n};R} \to C_{p+1} \enbrace*{\abs*{\Delta^n};R}$.
\end{definition}

\begin{figure}[ht]
	\centering{
	\tikzsetnextfilename{abkegeln}
	\begin{tikzpicture}[scale=3.5, rotate=230]
		\coordinate (a) at (0,0);
		\coordinate (b) at (1,0);
		\coordinate (c) at (60:1);
	
		\draw[thick] (a) -- (b) -- (c) -- cycle;
	
		\coordinate (ab) at ($(a)!.5!(b)$);
		\coordinate (ac) at ($(a)!.5!(c)$);
		\coordinate (bc) at ($(b)!.5!(c)$);
	
		\path[name path=erster] (a) -- (bc);
		\path[name path=zweiter] (b) -- (ac);
		\path[name intersections={of=erster and zweiter, name=i}] (i-1) coordinate (m);
		
		\draw[DarkOrange2, thick] (m) -- (ab) node[midway, above right,xshift=-0.2cm]{$K_{\Delta^n}(\sigma)$};
		\foreach \x in {0.1,0.2,...,0.9,1.0}{
			\draw[DarkOrange2, thick] (m) -- ($(ab)!\x!(b)$);
		}
		\draw[very thick, Firebrick2] (ab) -- (b) node[midway, above left]{\normalsize$\sigma$};
		
		% drawing circles
		\foreach \x in {a,b,c}{
			\draw[thick,fill=DodgerBlue2] (\x) circle[radius=0.015];
		}
		\foreach \x in {ab,bc,ac}{
			\draw[thick,fill=SeaGreen3] (\x) circle[radius=0.015];
		}
		\draw[thick,fill=Yellow2] (m) circle[radius=0.015];
	\end{tikzpicture}
	\caption{Skizze des \enquote{Abkegelns}}
	}
\end{figure}

\begin{lemma}[{name=[{Kettenhomotopie zwischen $\const_{s_n}$ und $\id$ auf $\abs*{\Delta^n}$}]},label=lem:136]
	Sei $\const_{s_n} \!\colon \abs*{\Delta^n} \to \abs*{\Delta^n}$ die konstante Abbildung mit $\im (\const_{s_n}) = \set{s_n} $. 
	Dann ist $\enbrace*{K_{\Delta^n}}_*$ eine \hyperref[def:76]{Kettenhomotopie} zwischen $\id$ und $(\const_{s_n})_*$.
\end{lemma}
\begin{beweis}
	Sei $\sigma$ ein singulärer $p$-Simplex in $\abs*{\Delta^n}$. Dann ist 
	\begin{align*}
		\partial_{p+1} \circ \enbrace*{K_{\Delta^n}}_p (\sigma)  = \sum_{j=0}^{p+1} (-1)^{j} \enbrace*{K_{\Delta^n}}_p(\sigma) \circ \iota_{p+1,j} &= \sigma + 
		\sum_{j=1}^{p+1} (-1)^j (K_{\Delta^n})_p(\sigma) \circ \iota_{p+1,j} \\
		&= \sigma + \begin{cases}
			\sum_{j=0}^{p} (-1)^{j+1}  (K_{\Delta^n})_{p-1} \enbrace*{\sigma \circ \iota_{p,j}} , & p \ge 1\\
			- \const_{s_n} \colon \abs*{\Delta^0} \to \set{s_n} , & p=0  
		\end{cases} \\
		&= \sigma - \begin{cases}
			(K_{\Delta^n})_{p-1} (\partial_p \circ \sigma) , &\text{ falls }p \ge 1\\
			\const_{s_n}  , &\text{ falls } p=0  
		\end{cases}
	\end{align*}
	Also ist $\partial_{p+1} \circ (K_{\Delta^n})_p + (K_{\Delta^n})_{p-1} \circ \partial_p= \id - (\const_{s_n})_*$.
\end{beweis}

\begin{definitionUnshaded}[{name=[{Zerteilungsoperatoren}]}]
	Wir definieren induktiv natürliche Transformationen $B_n \colon C_n(-;R) \to C_n(-;R)$ und $H_n \colon C_n(-;R) \to C_{n+1}(-;R)$ wie folgt: 
	Für $n=0$ setzen wir $(B_0)_X = \id$. Sei $B_k$ nun schon für $k=0, \ldots ,n-1$ definiert. 
	Betrachte zunächst $\id_{\abs*{\Delta^n}} \in C_n \enbrace*{\abs*{\Delta^n};R}$. Dafür setzen wir 
	\[
		(B_n)_{\abs*{\Delta^n}} \enbrace*{\id_{\abs*{\Delta^n}}} :=  (K_{\Delta^n})_{n-1} \circ (B_{n-1})_{\abs*{\Delta^n}}\circ \partial_n \enbrace*{\id_{\abs*{\Delta^n}}} 
	\]
	Ist $\sigma \colon \abs*{\Delta^n} \to X$ ein beliebiger singulärer $n$-Simplex in $X$, so ist $\sigma = \sigma_*\enbrace*{\id_{\abs*{\Delta^n}}}$, also setzen wir\footnote{$(B_n)_X(\sigma) = (B_n)_X\enbrace*{\sigma_*(\id_{\abs*{\Delta^n}})}= \sigma_*\enbrace[\big]{(B_n)_{\abs*{\Delta^n}(\id_{\abs*{\Delta^n}})}}$ muss wegen Natürlichkeit gelten}
	\[
		(B_n)_X(\sigma) := \sigma_* \enbrace*{ (B_n)_{\abs*{\Delta^n}} \enbrace*{\id_{\abs*{\Delta^n}}}} 
	\]
	Für $n=0$ und $\sigma \colon\abs*{\Delta^0} \to X$ setzen wir $(H_0)_X(\sigma)$ als den konstanten $1$-Simplex, dessen Bild der Punkt $\sigma\enbrace*{\abs*{\Delta^0}}$ ist.
	% $\tau$ mit $\im \tau = \im \sigma$.
	Ist $H_{n-1}$ schon definiert, so setzen wir für einen singulären $n$-Simplex $\sigma \colon \abs*{\Delta^n} \to X$
	\[
		(H_n)_X(\sigma) := \sigma_* \enbrace[\Big]{ (K_{\Delta^n})_n \enbrace[\big]{ \id_{\abs*{\Delta^n}} - (H_{n-1})_{\abs*{\Delta^n}}\circ\partial_n 
		\enbrace{\id_{\abs*{\Delta^n}}}}} 
	\]
\end{definitionUnshaded}

\begin{figure}[ht]
	\centering{
	\tikzsetnextfilename{Zerteilungsoperatoren}
	\begin{tikzpicture}[scale=1.1, thick]
		% \draw[help lines] (0,0) grid (8,4);
		
		% oberer Teil
		\node at (-1,2.7) {$B_n$};
		\coordinate (c) at (0,2.7);
		\draw (2,2.7) coordinate (d) -- ++(2,0) coordinate (e);
		\draw ($(d)!0.5!(e) + (0,0.15)$) -- ($(d)!0.5!(e) + (0,-0.15)$);
		\coordinate (mid) at (7,2.5);
		\coordinate (f) at ($(mid) + (90:1)$);
		\coordinate (g) at ($(mid) + (210:1)$);
		\coordinate (h) at ($(mid) + (-30:1)$);
		\draw (f) -- (g) -- (h) -- cycle;
		\draw (f) -- ($(g)!.5!(h) + (0,-0.1)$);
		\draw (g) -- ($(f)!.5!(h) + (30:0.1)$);
		\draw (h) -- ($(g)!.5!(f) + (150:0.1)$);
		
		% unterer Teil
		\node at (-1,0.5) {$H_n$};
		
		\draw (0,0) coordinate (a) -- (0,1) coordinate (b);
		
		\draw (2,-0.5) coordinate (i) -- ++(2,0) coordinate (j) -- ++(0,2) coordinate (k) -- ++ (-2,0) coordinate (l) -- cycle;
		\draw ($(i)!.5!(j)$) -- (k);
		\draw ($(i)!.5!(j)$) -- (l);
		
		\draw ($(5.5,-0.5) + (0,0.125)$) coordinate (m) -- ++(3,0) coordinate (n) -- ($(m)!.5!(n) + (0,0.75)$) coordinate (o) -- cycle;
		\draw (m) -- ++ (0,1) coordinate (p);
		\draw (n) -- ++ (0,1) coordinate (q);
		\draw[dashed] (o) -- ++ (0,1) coordinate (r);
		\draw (p) -- (q) -- (r) -- cycle;
		\coordinate (mitte) at ($(o) + (0,-0.45)$);
		
		\coordinate (mitte2) at ($(m)!.5!(n)$);
		\draw[Chartreuse3, dashed] (m) -- (n) --(q) -- (mitte2) -- (p) -- cycle;
		\draw[Chartreuse3, dashed] (m) -- (mitte) --(n);
		\draw[Chartreuse3, dashed] (mitte2) -- (mitte) -- (p);
		\draw[Chartreuse3, dashed] (mitte) -- (q);
		
		\foreach \y in {p,q,r}{
			\draw[IndianRed2, densely dashed] (mitte) -- (\y);
		}
		\draw[IndianRed2, densely dashed] (p) -- (q) -- (r) -- cycle;
		
		
		% punkte
		\foreach \x in {a,b,...,r}{
			\draw[fill=DodgerBlue2] (\x) circle[radius=0.045];
		}
		
	\end{tikzpicture}
	\caption{Die natürlichen Transformationen $B_n$ und $H_n$ für $n=0,1,2$}}
\end{figure}

\begin{lemma}[{name=[{Kettenhomotopie zwischen $(H_*)_X$ und der Identität}]},label=lem:138]
	Für jeden topologischen Raum $X$ ist $(B_*)_X$ eine Kettenabbildung und $(H_*)_X$ eine \hyperref[def:76]{Kettenhomotopie} zwischen $(B_*)_X$ und der Identität.
\end{lemma}
\begin{beweis}
	Wir müssen zeigen, dass folgende Gleichung gilt
	\begin{equation}
		\partial_{n+1} \circ (H_n)_X + (H_{n-1})_X \circ \partial_n = (B_n)_X - \id \tag{\#} \label{eq:1310}
	\end{equation}
	Da $\id$ eine Kettenabbildung ist, ist dann auch $(B_*)_X$ eine Kettenabbildung und $(H_*)_X$ eine Kettenhomotopie zwischen $\id$ und $(B_*)_X$. Wir beweisen \eqref{eq:1310}
	durch Induktion nach $n$. Für $n=0$ und $\sigma \colon \abs*{\Delta^0} \to X$ ist
	\[
		\partial_1 \enbrace[\big]{ (H_0)_X (\sigma)} = 0 \qquad \text{ und } \qquad  (B_0)_X(\sigma) = \sigma 
	\]
	Also gilt \eqref{eq:1310} für $n=0$. Für den Induktionsschritt $n-1 \mapsto n$ betrachten wir $\sigma \colon \abs*{\Delta^n} \to X$. Es gilt
	\begin{align*}
		\partial_{n+1} \circ (H_n)_X (\sigma) &= \sigma_* \circ \partial_{n+1} \circ (H_n)_{\abs*{\Delta^n}} \enbrace*{\id_{\abs*{\Delta^n}}} \\
		&= \sigma_* \circ \partial_{n+1} \enbrace[\Big]{(K_{\abs*{\Delta^n} })_n \enbrace*{ \id_{\abs*{\Delta^n}}- (H_{n-1})_{\abs*{\Delta^n}} \circ \partial_n 
		\enbrace{\id_{\abs*{\Delta^n}}}}} \\
		&\StackTextClap{vgl. \ref{lem:136}}{=} \sigma_* \enbrace[\Big]{- (K_{\abs*{\Delta^n}})_{n-1} \circ \partial_n \enbrace*{\id_{\abs{\Delta^n}} 
		- (H_{n-1})_{\abs*{\Delta^n}} \circ \partial_n (\id_{\abs*{\Delta^n} }) } } \\
		&\phantom{ = } + \sigma_* \enbrace*{\id_{\abs*{\Delta^n}} - (H_{n-1})_{\abs*{\Delta^n}} \circ \partial_n (\id_{\abs*{\Delta^n}})} \\
		&\stackrel{\mathclap{\text{I.A.}}}{=} \sigma_*  \enbrace*{- (K_{\abs*{\Delta^n}})_{n-1} \circ \partial_n (\id_{\abs*{\Delta^n}})}  
		- \sigma_* \enbrace[\Big]{- (K_{\abs*{\Delta^n}})_{n-1}  \enbrace*{\id - (B_{n-1})_{\abs*{\Delta^n}}}  \partial_n (\id_{\abs*{\Delta^n} }) } \\
		&\phantom{ = } + \sigma- (H_{n-1})_X \circ \partial_n(\sigma) \\
		&= - \sigma_* \enbrace*{(K_{\abs*{\Delta^n}})_{n-1} \circ (B_{n-1})_{\abs*{\Delta^n}} \circ \partial_n(\id_{\abs*{\Delta^n}})}
		+ \sigma - (H_{n-1})_X \circ \partial_n(\sigma) \\
		&= - \sigma_* \enbrace[\big]{(B_n)_{\abs*{\Delta^n}} \enbrace{\id_{\abs*{\Delta^n}}}} + \sigma - (H_{n-1})_X \circ \partial_n (\sigma) \\
		&= -(B_n)_X(\sigma) + \sigma - (H_{n-1})_X \circ \partial_n(\sigma)
	\end{align*}
	Also ist $\partial_{n+1} \circ H_n(\sigma) + H_{n-1} \circ \partial_{n}(\sigma) = \sigma - (B_n)_X(\sigma)$.
\end{beweis}

\begin{lemma}[{name=[{Eigenschaft von $(B_n)_X$ in Bezug auf eine offene Überdeckung}]},label=lem:139]
	Sei $\mathcal{U}$ eine offene Überdeckung von $X$. Sei $\alpha\in C_n(X;R)$. 
	Dann existiert ein $k_0\ge 0$ so, dass\marginnote{$(B_n)^{k}_X$ meint die $k$-fache Komposition von $(B_n)_X$} $(B_n)^{k}_X(\alpha)\in C^\mathcal{U}_n(X;R)$ für alle $k\ge k_0$.
\end{lemma}
\begin{beweis}
	Es genügt einen singulären Simplex $\sigma \colon \abs*{\Delta^n} \to X$ zu betrachten. Wir setzen $\sigma^* \mathcal{U}:= \set*{\sigma ^{-1}(U) \given U \in \mathcal{U}}$. Dann gilt
	\[
		\sigma_* \enbrace[\big]{C_*^{\sigma^*\mathcal{U}} \enbrace*{\abs*{\Delta^n};R}} \subseteq C_*^\mathcal{U}(X;R). 
	\]
	Da $(B_n)_X^k(\sigma) = \sigma_* \enbrace[\big]{(B_n)_{\abs*{\Delta^n}}^k (\id_{\abs*{\Delta^n} })}$ genügt es zu zeigen, dass 
	$(B_n)^k_{\abs*{\Delta^n}} (\id_{\abs*{\Delta^n} }) \in C_*^{\sigma^* \mathcal{U}}\enbrace*{\abs*{\Delta^n};R}$ für ein geeignetes $k>0$. Sei $\varepsilon>0$ eine 
	Lebesguezahl\footnote{siehe \url{https://de.wikipedia.org/wiki/Lebesguezahl}} für $\sigma^*\mathcal{U}$, das heißt 
	\[
		\forall x \in \abs*{\Delta^n} : \exists V \in \sigma^*\mathcal{U} : B_\varepsilon(x) \subseteq V 
	\]
	Wähle nun $k>0$ mit $\enbrace[\big]{\frac{n}{n+1}}^k \cdot \diam(\id_{\abs*{\Delta^n}})< \varepsilon$. Dann gilt nach \autoref{lem:1313} (iii)
	\[
		(B_n)_{\abs*{\Delta^n}}^k (\id_{\abs*{\Delta^n}}) \in C_*^{\sigma^*\mathcal{U}} \enbrace*{\abs*{\Delta^n};R} \qedhere
	\]
\end{beweis}

\begin{beweis}[{name={des Satzes über kleine Simplizes}}]
	Sei $x \in H_n(X;R)$. Wähle $\alpha \in C_n(X;R)$ mit $[\alpha]=x$. 
	Nach \autoref{lem:139} gibt es ein $k$ mit $(B_n)_X^k(\alpha) \in C_n^\mathcal{U}(X;R)$.
	Mit \autoref{lem:138} folgt $x=[\alpha]= \benbrace*{(B_n)_X^k(\alpha)}$. 
	Daher ist die von $C_*^\mathcal{U}(X;R) \hookrightarrow C_*(X;R)$ induzierte Abbildung in $H_*$ surjektiv.

	Sei nun $x$ im Kern von $H_n \enbrace*{C_*^\mathcal{U}(X;R)} \to H_n \enbrace*{C_*(X;R)}$. 
	Sei $\alpha \in C_n^\mathcal{U}(X;R)$ mit $x=[\alpha]$. 
	Da $x$ im Kern liegt, gibt es $\beta \in C_{n+1}(X;R)$ mit $\partial_{n+1}(\beta)=\alpha$. 
	Nach \autoref{lem:139} gibt es ein $k$ mit $(B_n)_X^k(\beta) \in C_n^\mathcal{U}(X;R)$. 
	Mit \autoref{lem:138} folgt nun
	\[
		x = [\alpha] = \benbrace*{(B_n)_X^k(\alpha)} = \benbrace*{(B_n)_X^k(\partial_{n+1} \beta)} = \benbrace*{\partial_n \enbrace*{(B_{n+1})_X^k (\beta)} }   
	\]
	Also ist $x=0$ in $H_* \enbrace*{C_*^\mathcal{U}(X;R)}$.
\end{beweis}

\begin{bemerkung}[{name=[{Abbildungen $s_n \colon C_n \to C_{n+1}'$, die den Rang erhöhen}]}]
	Seien $(C_*,\partial_*)$ und $(C_*',\partial_*')$ Kettenkomplexe und $s_n \colon C_n \to C_{n+1}'$ eine Abbildung, die den Grad um $1$ erhöht. 
	Dann ist $s_{n-1} \circ \partial_n + \partial_{n+1}' \circ s_n$ eine Kettenabbildung, da
	\[
		\enbrace*{s_{n-1} \circ \partial_n + \partial_{n+1}' \circ s_n} \circ \partial_{n+1} = \partial_{n+1}' \circ  \enbrace*{s_n \circ \partial_{n+1} 
		+ \partial_{n+2}' \circ s_{n+1}}  
	\]
	Insbesondere gilt: Wenn $s \circ \partial + \partial' \circ s = f -g$ für eine Kettenabbildung $f \colon C_* \to C'_*$, dann ist auch $g$ eine Kettenabbildung.
\end{bemerkung}

\begin{definition}[{name=[{Affiner singulärer Simplex}]}]
	Seien $v_0, \ldots ,v_p \in \abs*{\Delta^n}$. 
	Dann bezeichnen wir mit $[v_0,\ldots ,v_p]$ den singulären $p$-Simplex mit $\sum_{i=0}^{p} t_i e_i \to \sum_{i=1}^{p} t_i v_i$.
	Singuläre Simplizes dieser Form heißen \bet{affin}\index{affiner singulärer Simplex}.
\end{definition}

\begin{definition}[{name=[{Durchmesser eines singulären Simplizes $\abs*{\Delta^p} \to \abs*{\Delta^n}$}]}]
	Sei $\norm{\cdot}$ eine Norm auf $\mathbb{R}^{n+1}$. 
	Für einen singulären Simplex $\sigma \colon \abs*{\Delta^p} \to \abs*{\Delta^n}$ definieren wir den \Index{Durchmesser} von $\sigma$ durch
	\[
		\diam(\sigma) := \max \set[\big]{\norm{\sigma(x)- \sigma(y)}  \given x,y \in \abs*{\Delta^p}} 
	\]
\end{definition}

\begin{lemma}[{name=[{Eigenschaften von affinen singulären Simplizes}]},label=lem:1313]
	Sei $[v_0, \ldots , v_p]$ ein affiner singulärer $p$-Simplex in $\abs*{\Delta^n}$. Dann gilt:
	\begin{enumerate}[(i)]
		\item Es gilt $\diam \enbrace*{[v_0, \ldots ,v_p]}= \max \set*{\norm{v_i - v_j} \given 0 \le i,j\le p\rule{0cm}{0.9em}}$
		\item Es ist 
		\[
			(B_p)_{\abs*{\Delta^n}} \enbrace[\big]{[v_0, \ldots ,v_p]} =
			\smashoperator{\sum_{\hspace{3em}\sigma \in \sum_{\set{v_0,\ldots ,v_p}}}}
			\enspace\sgn(\sigma)  \cdot  
			\underbrace{\benbrace*{\frac{\sum_{i=0}^{p} \sigma(v_i)}{p+1}, \frac{\sum_{i=0}^{p-1} \sigma(v_i)}{p} ,\ldots , \sigma(v_0)  } }_{=: \tau_\sigma}
		\]
		\item Für $\sigma \in \sum_{\set{v_0, \ldots ,v_p}}$ ist $\diam \tau_\sigma \le \frac{n}{n+1} \diam [v_0,\ldots ,v_p]$.
	\end{enumerate}
\end{lemma}
\begin{beweis}
	\leavevmode
	\begin{enumerate}[(i)]
		\item Seien $v,v' \in \im \enbrace*{[v_0, \ldots ,v_p]}$. Dann ist $v = \sum_{i=0}^{p} t_i v_i$ und $v' = \sum_{i=0}^{p} t_i' v_i$ mit $t_i, t_i' \in [0,1]$, sowie
		$\sum_{i=0}^{p} t_i =1 = \sum_{i=0}^{p} t_i'$. Es gilt nun
		\begin{align*}
			\norm*{v-v'} = \norm*{\sum_{i=0}^{p} t_i v_i - v'} &= \norm*{\sum_{i=0}^{p} t_i v_i - \sum_{i=0}^{p} t_i v'} = \norm*{\sum_{i=0}^{p} t_i(v_i -v')}
			\le \sum_{i=0}^{p} t_i \norm{v_i -v'}   \\
			&\le \enbrace*{\sum_{i=0}^{p} t_i} \max_i \norm*{v_i -v'} = \max_i \norm*{v_i -v} \\
			&= \max_i \norm*{v_i - \sum_{j=0}^{p} t'_j v_j} \\
			&= \max_i \norm*{ \sum_{j=0}^{p} t'_j v_i + \sum_{j=0}^p t'_j v_j} 
			\le \ldots \le \max_{i,j} \norm{v_i - v_j}.   
		\end{align*}
		\item Per Induktion nach $p$. Für $p=0$ ist $(B_0)_{\abs{\Delta^n}}=\id$ und die Aussage wahr.
	
		Für $p-1 \mapsto p$ genügt es $n=p$ und $[v_0,\ldots ,v_p] =[e_0, \ldots ,e_p]$ zu betrachten. Es gilt
		\begin{align*}\marginnote{Aus Platzgründen wird der Index $\abs*{\Delta^p}$ bei den Abbildungen $(K_{\abs*{\Delta^p}})_*$ und $(B_*)_{\abs*{\Delta^p}}$ weggelassen}
			B_p \enbrace[\big]{[v_0,\ldots ,v_p]} &= K_{p-1} \circ B_{p-1} \circ \partial_p \enbrace[\big]{[e_0, \ldots ,e_p]} \\
			&= \sum_{i=0}^{p} (-1)^i \cdot K_{p-1} \circ B_{p-1} \enbrace[\big]{[e_0, \ldots , e_{i-1}, e_{i+1}, \ldots ,e_p]} \\
			&\stackrel{\mathclap{\text{I.V.}}}{=} \sum_{i=0}^{p} (-1)^i \cdot
			\smashoperator{\sum_{\hspace{4em}\sigma \in \sum_{\set{e_0, \ldots, e_p}\setminus \set{e_i}}} }
			\enspace\sgn(\sigma) \cdot K_{p-1} 
			\enbrace*{\textstyle\benbrace*{\frac{\sigma(e_0) + \ldots +\sigma(e_{i-1}) + \sigma(e_{i+1}) + \ldots  + \sigma(e_p)}{p},\ldots  } } \\
			&= \sum_{i=0}^{p} (-1)^i \cdot 
			\smashoperator{\sum_{\hspace{4em}\sigma \in \sum_{\set{e_0, \ldots, e_p}\setminus \set{e_i}}} }
			\enspace \sgn(\sigma)\cdot 
			\benbrace*{\textstyle\frac{e_0 + \ldots + e_p}{p+1}, \frac{\sigma(e_0) + \ldots + \sigma(e_{i-1}) + \sigma(e_{i+1}) + \ldots  + \sigma(e_p)}{p}, \ldots} \\
			&=\ldots = \smashoperator{\sum_{\hspace{3em}\sigma \in \sum_{\set{e_0, \ldots ,e_n}}}} \enspace\sgn(\sigma) \cdot \benbrace*{\textstyle \frac{\sum_{i=0}^{p} \sigma(e_i)}{p+1}, \ldots , \sigma(e_0)} 
		\end{align*}
		\item Es genügt $\sigma=\id$ zu betrachten. Es ist 
		\begin{align*}
			\hspace{-1em}\diam \benbrace*{\textstyle\frac{v_0+ \ldots +v_p}{p+1}, \frac{v_0 + \ldots + v_{p-1}}{p}, \ldots ,v_0} &\stackrel{\text{(i)}}{=} \max_{i <j}
			\norm*{\textstyle \frac{v_0+ \ldots +v_i}{i+1} - \frac{v_0 + \ldots + v_j}{j+1}}  \\
			&= \max_{i <j} \norm*{\textstyle\frac{v_0+ \ldots +v_i}{i+1} - \frac{v_0 + \ldots + v_i}{j+1} - \frac{v_{i+1}+ \ldots + v_j}{j+1}} \\
			&= \max_{i<j} \norm*{\textstyle\frac{(j+1)- (i+1)}{(i+1)(j+1)}(v_0 + \ldots + v_i) - \frac{1}{j+1} (v_{i+1} + \ldots +v_j)}  \\
			&= \max_{i<j} \frac{j-i}{j+1} \cdot \norm[\bigg]{\textstyle{\underbrace{\frac{v_0 + \ldots +v_i}{i+1}}_{\in \im [v_0, \ldots ,v_n]}} - 
			\underbrace{\frac{v_{i+1} + \ldots  + v_j}{j-i}}_{\in \im [v_0, \ldots , v_n]}} \\
			&\le  \max_{i<j} \frac{j-i}{j+1} \diam [v_0, \ldots, v_n] \\
			&\le \max_j \frac{j}{j+1} \diam [v_0, \ldots ,v_n] \le \frac{n}{n+1} \diam [v_0, \ldots , v_n]  \qedhere
		\end{align*}
	\end{enumerate}
\end{beweis}

\newpage

\section{Axiome für Homologie} % (fold)
\label{sec:14}

Mit $V \colon \TOP^2 \to \TOP^2$ bezeichnen wir in diesem Kapitel den Funktor $V(X,A) = (A,\emptyset)=A$.

\begin{definition}[{name=[Homologietheorie]},label=def:141]
	Eine \Index{Homologietheorie} mit Werten in $R$-Moduln ist ein Funktor 
	\[
		E_* \colon \HTOP^2 \longrightarrow \textsc{Gr}\text{-}R\text{-}\MOD
	\]
	zusammen mit einer natürlichen Transformation $\partial_{*+1} \colon E_{*+1}  \to E_{*} \circ V$,\footnote{Genauer: Sei 
	$S \colon \textsc{Gr}\text{-}R\text{-}\MOD \to \textsc{Gr}\text{-}R\text{-}\MOD$ der Funktor $(S M)_n = M_{n+1}$. Dann ist $\partial_*$ eine natürliche Transformation $S \circ E_* \to E_* \circ V$, für $(X,A)$ also $S \enbrace[\big]{E_*(X,A)} \to E_* \enbrace[\big]{V(X,A)}$, $E_{n+1}(X,A) \xrightarrow{\partial_*}E_n(A)$.} sodass folgende Axiome gelten:
	\begin{enumerate}[(i)]
		\item \Index{Paarfolge}: Sei $(X,A)$ ein Paar von topologischen Räumen. 
		Seien $X= (X,\emptyset) \xrightarrow{\minwidthbox{j}{1em}} (X,A)$ und $A= (A,\emptyset) \xrightarrow{\minwidthbox{i}{1em}} X=(X,\emptyset)$ die Inklusionen. Dann ist 
		\[
			\begin{tikzcd}[font=\small]
				\cdots \rar & E_{n+1}(X,A) \rar["\partial_{n+1}"] & E_n(A) \rar["E_n(i)"] & E_n(X) \rar["E_n(j)"] & E_n(X,A) \rar["\partial_{n}"] & E_{n-1}(A) \rar & \cdots
			\end{tikzcd}
		\] 
		eine lange exakte Folge.
		\item \Index{Ausschneidung}: Sei $(X,A)$ ein Paar von topologischen Räumen und $L \subseteq A$ mit $\overline{L} \subseteq \mathring{A}$. 
		Dann induziert die Inklusion
		$i \colon \enbrace*{X \setminus L, A \setminus L} \hookrightarrow (X,A)$ einen Isomorphismus
		\[
			E_*(i) \colon E_* \enbrace*{X \setminus L, A \setminus L} \xrightarrow{\enspace \cong \enspace} E_*(X,A)  
		\] 
	\end{enumerate}
\end{definition}

Homologietheorien sind homotopieinvariant. Oft wird dies als eigenes Axiom formuliert.

\begin{bemerkung}[{name=[Mayer-Vietoris-Folge und disjunkten Vereinigungen]}]
	Für jede Homologietheorie gibt es eine Mayer-Vietoris-Folge. 
	Dies ist eine formale Konsequenz aus der Paarfolge und dem Ausschneidungsisomorphismus. 
	Eine einfache Folgerung aus der Mayer-Vietoris-Folge ist, dass jede Homologietheorie mit \emph{endlichen} disjunkten Vereinigungen verträglich ist: 
	Sei dazu $X= X_1 \amalg \cdots \amalg X_n$ und $j_i \colon X_i \to X$ die Inklusion, dann ist 
	\[
		E_*(X_1) \oplus \cdots \oplus E_*(X_n) \xrightarrow{E_*(j_1) \oplus \cdots \oplus E_*(j_n)} E_*(X)
	\]
	ein Isomorphismus.
\end{bemerkung}

\begin{definition}[{name=[{Dimensionsaxiom und Axiom über disjunkte Vereinigungen}]},label=def:143]
	Zwei weitere Axiome, die Homologietheorien mit Werten in $R$-Moduln erfüllen können, sind:
	\begin{description}
		\item[Dimensionsaxiom:] \index{Dimensionsaxiom}Für den Einpunktraum $\set{x_0}$ ist $E_0(\set{x_0})\cong R$ und $E_n(\set{x_0})=0$ für $n \not=0$.
		\item[Disjunkte Vereinigung:] \index{disjunkte Vereinigung!Axiom} Sei $X = \coprod_{i \in I} X_i$ und $j_i \colon X_i \to X$ die Inklusion. Dann ist 
		\[
			\bigoplus_{i \in I} E_*(j_i) \colon \bigoplus_{i \in I} E_*(X_i) \longrightarrow E_*(X)
		\]
		ein Isomorphismus.
	\end{description}
\end{definition}

\begin{bemerkung}[{name=[{über Z-graduierte Moduln}]}]
	In Kapitel \ref{sec:6} haben wir $\textsc{Gr}\text{-}R\text{-}\MOD$ als die Kategorie der $\mathbb{N}$-graduierten $R$-Moduln definiert. 
	Für viele Homologietheorien muss man diese Kategorie etwas vergrößern und $\mathbb{Z}$"=graduierte $R$-Moduln zulassen. 
	Objekte sind dann Folgen $(M_n)_{n \in \mathbb{Z}}$ über $\mathbb{Z}$ von $R$-Moduln $M_n$.
\end{bemerkung}

\begin{bemerkung}[{name=[{Homologie der Sphäre und zur zellulären Homologie}]}]
	\leavevmode
	\begin{itemize}
		\item Für jede Homologietheorie $E_*$ gilt: 
		\(
			E_*(S^n) \cong E_* \enbrace*{\set{\mathrm{pt}}} \oplus E_{*-n} \enbrace*{\set{\mathrm{pt}} }  
		\)\marginnote{Aufgabe 1 Blatt 12}
		\item Die Konstruktion des zellulären Kettenkomplexes in singulärer Homologie benutzte nur die Axiome aus \autoref{def:141}. 
		Für die Identifikation der zellulären Homologie mit der singulären Homologie waren zusätzlich die Axiome aus \autoref{def:143} notwendig. 
		Ohne das Axiom über disjunkte Vereinigungen kann man nur endliche CW-Komplexe behandeln.
		Ohne das Dimensionsaxiom erhält man an der Stelle der zellulären Kettenkomplex eine sogenannte \Index{Spektralfolge}, die Atiyah-Hirzebruch Spektralfolge. 
	\end{itemize}
\end{bemerkung}

\begin{definition}[{name=[Morphismus von Homologietheorien]}]
	Seien $(E_*,\partial_*^E)$ und $(F_*,\partial_*^F)$ Homologietheorien. 
	Ein \Index{Morphismus von Homologietheorien} ist eine natürliche Transformation $\tau_* \colon E_* \to F_*$, sodass $\tau_* \circ \partial_{*+1}^E = \partial_*^F \circ \tau_{*+1}$. 
	Also kommutiert für alle $(X,A)$ das folgende Diagramm
	\[
		\begin{tikzcd}[sep=large]
			E_{n+1}(X,A) \rar["\partial_{n+1}^E"] \dar["\tau_{n+1,(X,A)}"] & E_n(A) \dar["\tau_{n,A}"] \\
			F_{n+1}(X,A) \rar["\partial_{n+1}^F"] & F_n(A)
		\end{tikzcd}
	\]
\end{definition}

\begin{satz}[{name=[{Isomorphie der Homologie von CW-Komplexen von zwei Homologietheorien}]}]
	Sei $\tau_* \colon (E_*,\partial_*^E) \to (F_*,\partial_*^F)$ ein Morphismus von Homologietheorien. Ist 
	\[
		(\tau_*)_{\set{\mathrm{pt}}} \colon E_*\enbrace[\big]{\set{\mathrm{pt}}} \to F_* \enbrace[\big]{\set{\mathrm{pt}}}
	\]
	ein Isomorphismus, dann ist $(\tau_*)_X$ für jeden endlichen CW-Komplex ein Isomorphismus. 
	Sind $E_*$ und $F_*$ mit beliebigen disjunkten Vereinigungen verträglich, so ist $(\tau_*)_X$ für alle CW-Komplexe ein Isomorphismus.
\end{satz}
\begin{beweis}
	\leavevmode
	\begin{description}
		\item[Schritt 1:] Sei $X$ ein $0$-dimensionaler endlicher CW-Komplex, also eine endliche Menge mit der diskreten Topologie. 
		Für $x \in X$ sei $i_x \colon \set{x} \to X$ die Inklusion. Wir erhalten folgendes Diagramm:
		\[
			\begin{tikzcd}[column sep=5.5em, row sep=2.5em]
				E_*(X) \rar["(\tau_*)_X"] & F_*(X) \\
				\bigoplus\limits_{\mathclap{x \in X}} E_*(\set{x}) \uar["\cong"] \rar["\bigoplus_{x \in X} (\tau_*)_{\set{x}}","\cong"'] 
				& \bigoplus\limits_{\mathclap{x \in X}} F_*(\set{x}) \uar["\cong"]
			\end{tikzcd}
		\]
		Das Diagramm kommutiert, da $\tau_*$ eine natürliche Transformation ist.
		\item[Schritt 2:] Angenommen $(\tau_*)_X$ ist ein Isomorphismus für jeden endlichen CW-Komplex von Dimension echt kleiner als $n$. 
		Sei $X$ ein $n$-dimensionaler endlicher CW-Komplex:
		\[
			\begin{tikzcd}
				\coprod\limits_{\mathclap{i \in I}} S^{n-1} \rar["\coprod q_i^\ssbrace{n}"] \dar & X^\ssbrace{n-1} \dar\\
				\coprod\limits_{\mathclap{i \in I}} D^{n} \rar["\coprod Q_i^\ssbrace{n}"] & X
			\end{tikzcd}
		\]
		Sei $X_+^\ssbrace{n-1} := X^\ssbrace{n-1} \cup \coprod_{i \in I} Q^\ssbrace{n}_i \enbrace*{D^n \setminus \set{0}}$ und $Z := \coprod_{i \in I} Q^\ssbrace{n}_i(\mathring{D^n})$. 
		Dann sind $X^\ssbrace{n-1}_+$, $Z$ und $Z\cap X_+^\ssbrace{n-1}$ homotopieäquivalent zu endlichen CW-Komplexen mit Dimension echt kleiner $n$.
		Insbesondere sind $(\tau_*)_{X_+^\ssbrace{n-1}}$, $(\tau_*)_Z$ und $(\tau_*)_{X^\ssbrace{n-1}_+ \cap Z}$ Isomorphismen. 
		Da $X^\ssbrace{n-1}_+ \subseteq X$ und $Z \subseteq X$ offen sind, erhalten wir in $E_*$ und $F_*$ zugehörigen Mayer-Vietoris-Folgen:
		\[
			\!\!\begin{tikzcd}[column sep=1.8em, row sep=2.5em,font=\small]
				\cdots \rar & E_* \enbrace{X_+^\ssbrace{n-1} \cap Z} \rar \dar["(\tau_*)_{X_+^\ssbrace{n-1} \cap Z}","\cong"']
				& E_*\enbrace{X_+^\ssbrace{n-1}} \oplus E_*(Z) \rar \dar["(\tau_*)_{X_+^\ssbrace{n-1}} \oplus (\tau_*)_{Z}","\cong"']
				& E_*(X) \rar["\partial_*^{\mathrm{MV},E}"] \dar["(\tau_*)_X"]
				& E_{*-1} \enbrace{X_+^\ssbrace{n-1} \cap Z} \rar \dar["(\tau_{*-1})_{X_+^\ssbrace{n-1} \cap Z}", "\cong"'] & \cdots  \\
				\cdots \rar & F_* \enbrace{X_+^\ssbrace{n-1} \cap Z} \rar & F_*\enbrace{X_+^\ssbrace{n-1}} \oplus F_*(Z) \rar 
				& F_*(X) \rar["\partial_*^{\mathrm{MV},F}"] & F_{*-1} \enbrace{X_+^\ssbrace{n-1} \cap Z} \rar & \cdots
			\end{tikzcd}
		\]
		Da $\tau_*$ eine natürliche Transformation ist, kommutierten die beiden linken Vierecke. 
		Da $\tau_*$ mit der Randabbildung der Paarfolgen von $E_*$ und $F_*$ kommutiert, kommutiert auch $\tau_*$ mit den Randabbildungen der Mayer-Vietoris-Folgen in $E_*$ und $F_*$. A
		lso kommutiert auch das das dritte Viereck. 
		Mit dem 5er-Lemma (\ref{lem:812}) folgt, dass auch $(\tau_*)_X$ ein Isomorphismus ist.
	
		Erfüllen $E_*$ und $F_*$ auch das Dimensionsaxiom, so zeigen die Argumente aus Schritt 1 und 2, dass $(\tau_*)_X$ ein Isomorphismus ist für jeden endlich dimensionalen CW-Komplex.
		\item[Schritt 3:] Sei $X$ ein CW-Komplex. Sei
		\[
			X^+ := \set*{(x,t) \given x \in X^\ssbrace{n}, t \in \enbrace*{n+ 1/2,\infty}\rule{0cm}{1em}} \subseteq X \times \enbrace*{1/2 ,\infty}
		\]
		Sei weiter
		\[
			X_0 := X^+ \cap \enbrace[\Big]{X \times \enbrace[\big]{(1/2,\infty) \setminus \set{1,3,4,\ldots } } } \quad \,\,
			X_1 := X^+ \cap \enbrace[\Big]{X \times \enbrace[\big]{(1/2,\infty) \setminus \set{2,4,6,\ldots } } } 
		\]
		Dann sind $X_0,X_1$ und $X_0 \cap X_1$ homotopieäquivalent zu disjunkten Vereinigungen von endlichen dimensionalen CW-Komplexen. 
		Sind $E_*$ und $F_*$ mit beliebigen Vereinigungen verträglich, so folgt dass $(\tau_*)_{X_0}$, $(\tau_*)_{X_1}$ und $(\tau_*)_{X_0\cap X_1}$ Isomorphismen sind. 
		Da $X$ homotopieäquivalent zu $X^+$ ist, $X^+ =X_0 \cup X_1$ und $X_0$ und $X_1$ offen in $X^+$ sind, folgt wie in Schritt 2 mit der 5er-Lemma und der Mayer-Vietoris-Folge, dass auch $(\tau_*)_X$ ein Isomorphismus ist. \qedhere
	\end{description}
\end{beweis}

\newpage
\section{\protect{Zusammenhang zwischen Fundamentalgruppen und Homologiegruppen}} % (fold)
\label{sec:15}

Zur Vorbereitung müssen wir kurz wiederholen, was eine \emph{Abelisierung} ist:

\begin{definition}[{name=[Kommutatorgruppe]}]
	Sei $G$ eine Gruppe. Dann heißt die Untergruppe von $G$, die von Elementen der Form $[g,h]:= g h g^{-1} h ^{-1}$ erzeugt wird \Index{Kommutatorgruppe} von $G$. 
	Wir schreiben hierfür auch $[G,G] \subset G$.
\end{definition}

\begin{definition}[{name=[{Abelisierung}]}]
	Der Quotient $G^{\ab} := \sfrac{G}{[G,G]}$ heißt \Index{Abelisierung}\marginnote{Für eine abelsche Gruppe $G$ ist sinnigerweise $G^{\ab}=G$.} (manchmal auch \emph{Abelianisierung}) von G und ist offensichtlich stets abelsch.
\end{definition}

\begin{bemerkung}
	\leavevmode
	\begin{itemize}
		\item $\GRUPPEN \grenzw{\ab} \ABELGRUPPEN , \quad G \longmapsto G^{\ab} $
		definiert einen Funktor.
		\item Dieser Funktor hat die folgende universelle Eigenschaft: Zu jeder abelschen Gruppe $A$ und einem Gruppenhomomorphismus $\varphi \colon G \to A$ existiert ein eindeutiger
		Gruppenhomomorphismus $\varphi^{\ab} \colon G^{\ab} \to A$, sodass folgendes Diagramm kommutiert
		\[
			\begin{tikzcd}
				G \rar["\varphi"] \dar[two heads]& A \\
				G^{\ab} \urar["\varphi^{\ab}"']
			\end{tikzcd}
		\]
		\item Sei $X$ ein topologischer Raum. 
		Da $\abs*{\Delta^1} \cong [0,1]$, liefert jeder Pfad $\omega \colon [0,1] \to X$ einen $1$-Simplex in $C_1(X;\mathbb{Z})$.
	\end{itemize}
\end{bemerkung}

\begin{lemma}[{name=[{Für Pfade $\omega,\omega'$ ist $\omega * \omega' - \omega- \omega'$ der Rand eines $2$-Simplex}]},label=lem:154]
	Seien $\omega,\omega' \colon [0,1] \to X$ Pfade in $X$ mit $\omega(1)=\omega'(0)$,\marginnote{\enquote{$*$} bezeichnet die Verkettung von Pfaden und wird noch links nach rechts gelesen} dann ist $\omega * \omega' - \omega- \omega' \in C_1(X;\mathbb{Z})$ der Rand eines $2$-Simplex.
\end{lemma}
\begin{beweis}
	Zur Erinnerung: Die geometrische Realisierung des $2$-Simplex ist gegeben durch 
	\[
		\abs*{\Delta^2} = \set*{(x_0,x_1,x_2) \in \mathbb{R}^3 \given \sum_{i=0}^{2}x_i=1, x_j \ge 0 \,\forall j \in \set{0,1,2}}.
	\] 
	Betrachte nun folgenden singulären Simplex $\sigma \colon \abs*{\Delta^2} \to X$ definiert durch 
	\[
		(x_0,x_1,x_2) \longmapsto(\omega*\omega') \enbrace*{\frac{1}{2} (1+x_0-x_2)} = \omega * \omega' \enbrace*{\frac{1}{2}(2x_0 + x_1) }  
	\]
	Die Ränder sind
	\begin{align*}
		(\delta_0 \sigma)(s) &= \sigma(0,s,1-s) = \omega * \omega' \enbrace*{\frac{1}{2} s}=\omega(s) \\
		(\delta_1 \sigma)(s) &= \sigma(s, 0, 1-s) = \omega * \omega'(s) \\
		(\delta_2 \sigma)(s) &= \sigma(s,1-s,0) = (\omega * \omega') \enbrace*{\frac{1}{2} (1+s)} = \omega'(s)
	\end{align*}
	Dann folgt $\partial \sigma = \delta_0 \sigma - \delta_1 \sigma + \delta_2 \sigma = \omega - \omega * \omega' + \omega'$.
\end{beweis}
\begin{figure}[th]
	\centering{
	$\vcenter{\hbox{\begin{tikzpicture}[scale=2]
		\coordinate (a) at (0,0);
		\coordinate (b) at (1,0);
		\coordinate (c) at (60:1);
		\draw (c) -- ($ (a)!0.5!(b) $) node[midway, below left]{$\downarrow$} node[midway, below right]{$\downarrow$};
		\draw[thick] (a) -- (b) node[midway, below]{$\omega * \omega'$} -- (c) node[midway, above right]{$\omega'$} -- cycle node[midway, above left]{$\omega$};
	\end{tikzpicture}}} \quad \xrightarrow{\hspace{1.2cm}} \quad 
	\begin{tikzpicture}[scale=2, my brace/.style={decorate, decoration={brace, amplitude=5pt, mirror, raise=3pt}, gray},baseline=-0.8em] % baseline=-\heightof{$\to$}/2
		\draw[thick] (1,0.09)--(1,0);
		\draw[thick] (0,0) -- (2,0);
		\draw[my brace] (0,0) -- (1,0) node[midway, below=6pt]{$\omega$};
		\draw[my brace] (1,0) -- (2,0) node[midway, below=6pt]{$\omega'$};
	\end{tikzpicture}
	$ \caption{Skizze zum Beweis von \autoref{lem:154}}}
\end{figure}

\begin{lemma}[{name=[{Summe inverser Wege und der konstante Wege}]},label=lem:155]
	\leavevmode
	\begin{enumerate}[a)]
		\item Sei $\omega \colon [0,1] \to X$ ein Pfad und $\omega ^{-1}(s) = \omega(1-s)$ der umgekehrte Pfad. 
		Dann ist $\omega + \omega ^{-1}$ der Rand einer $2$-Kette, also von einem Element aus $C_2(X;\mathbb{Z})$.
		\item Der konstante Pfad ist der Rand einer $2$-Kette.
	\end{enumerate}
\end{lemma}
\begin{beweis}
	\leavevmode
	\begin{enumerate}[a)]
		\item Sei $\sigma \colon \abs*{\Delta^2} \to X$ gegeben durch $\sigma(x_0,x_1,x_2) = \omega(x_1)$. Dann ist 
		\[
			\partial \sigma = \omega- {\underbrace{\const_{\omega(0)}}_{=: \partial \tau}} + \omega^{-1}
		\]
		Also gilt nach b) $\omega + \omega ^{-1} = \partial(\sigma+\tau)$.
		\item Der Rand des konstanten $2$-Simplex ist der konstante Pfad. \qedhere
	\end{enumerate}
\end{beweis}

\begin{figure}[ht]
	\centering{
	$\vcenter{\hbox{\begin{tikzpicture}[scale=2]
		\coordinate (a) at (0,0);
		\coordinate (b) at (-30:1);
		\coordinate (c) at (30:1);
		\draw (a) -- ($ (c)!0.5!(b) $) node[pos=0.6, below]{$\uparrow$} node[pos=0.85, below=3pt]{$\uparrow$} node[pos=0.6, above]{$\downarrow$} 
		node[pos=0.85, above=3pt]{$\downarrow$};
		\draw[thick] (a) -- (b) node[midway, below]{$\omega$} -- (c) node[midway, right]{$\const_{\omega(0)}$} -- cycle node[midway, above]{$\omega^{-1}$};
	\end{tikzpicture}}} \quad \xrightarrow{\hspace{1.2cm}} \quad 
	\vcenter{\hbox{
	\begin{tikzpicture}[scale=2]
		\draw[thick] (1,0.06)--(1,-0.06);
		\draw[thick] (0,0.06)--(0,-0.06);
		\draw[thick] (0,0) -- (1,0);
	\end{tikzpicture}
	}}
	$ \caption{Skizze zum Beweis von \autoref{lem:155}}}
\end{figure}

\begin{lemma}[{name=[{Für punktiert homotope Pfade $\omega$ und $\omega'$ ist $\omega-\omega'$ im Rand einer $2$-Kette}]},label=lem:156]
	Sind $\omega,\omega' \colon [0,1] \to X$ Pfade in $X$ mit $\omega(0)=\omega'(0)$ und $\omega(1)= \omega'(1)$, dann gilt: 
	Ist $\omega$ homotop zu $\omega'$ relativ zu den Endpunkten, dann ist $\omega- \omega'$ der Rand einer $2$-Kette.
\end{lemma}
\begin{beweis}
	Sei $H \colon [0,1] \times [0,1] \to X$ eine Homotopie zwischen $\omega$ und $\omega'$ mit $H(0,s) = \omega(s)$, $H(1,s)=\omega'(s)$, $H(t,0)=\const_{\omega(0)}$ und $H(t,1)= \const_{\omega(1)}$.
	Wähle einen Homöomorphismus durch
	\[
		\kappa \colon \nicefrac{[0,1] \times [0,1]}{[0,1] \times \set{0}} \to \abs*{\Delta^2} , \qquad [t,s] \longmapsto \enbrace[\big]{st, s(1-t), (1-s)}
	\]
	Betrachte nun $\sigma := H \circ \kappa ^{-1} \colon \abs*{\Delta^2} \to X$. 
	Die Ränder sind $\delta_0 \sigma = \omega$ und $\delta_1 \sigma = \omega'$ und $\delta_2 \sigma = \const_{\omega(1)}$. 
	Also folgt, dass $\omega- \omega'$ der Rand einer $2$-Kette ist.
\end{beweis}
\begin{figure}[th]
	\centering{
	$\vcenter{\hbox{\begin{tikzpicture}[scale=1.7]
		\coordinate (a) at (0,0);
		\coordinate (b) at (1,0);
		\coordinate (c) at (1,1);
		\coordinate (d) at (0,1);
		\draw[thick] (a) -- (b) node[midway, below]{$\omega$} -- (c) node[midway, right]{$\const_{\omega(1)}$} -- (d) node[midway, above]{$\omega'$}
		-- cycle node[midway, left]{$\const_{\omega(0)}$};
	\end{tikzpicture}}} \qquad 
	\begin{tikzpicture}
		\draw[-{Classical TikZ Rightarrow[width=2mm, length=1mm]}, decorate, decoration={snake, amplitude=0.4mm, segment length=2.5mm}] (0,0) -- (1,0);
	\end{tikzpicture}
	\qquad 
	\vcenter{\hbox{
	\begin{tikzpicture}[scale=1.7]
		\coordinate (a) at (0,0);
		\coordinate (b) at (1,0);
		\coordinate (c) at (1,1);
		\draw[thick] (a) -- (b) node[midway, below]{$\omega$} -- (c) node[midway, right]{$\const_{\omega(1)}$} node[above]{$\phantom{\omega'}$}-- cycle node[midway, above left]{$\omega'$};
	\end{tikzpicture}
	}}
	$ \caption{Skizze zum Beweis von \autoref{lem:156}}}
\end{figure}

Falls $\omega \colon [0,1] \to X$ eine Schleife ist, dann ist $\partial \omega = \omega(1)- \omega(0)=0$. 
Also repräsentiert $\omega$ ein Element\marginnote{$\bbenbrace{\cdot}$ für die Homologieklassen und $[\cdot]$ für Homotopieklassen} $\bbenbrace{\omega} \in H_1(X)$. 
Nach \autoref{lem:156} erhalten wir eine wohldefinierte Abbildung $\theta \colon \pi_1(X,x_0) \to H_1(X)$. 
Nach \autoref{lem:154} bildet diese Abbildung die Verknüpfung von Schleifen auf Summen von $1$-Ketten ab.
\begin{align*}
	\bbenbrace*{\omega * \omega'} = \bbenbrace*{\omega} + \bbenbrace*{\omega'} \quad ,\quad  \bbenbrace*{\omega ^{-1}} = - \bbenbrace*{\omega}
\end{align*}
Zusammen mit \autoref{lem:155} erhalten wir einen Gruppenhomomorphismus. 
Aus der universellen Eigenschaft der Abelisierung erhalten wir 
\[
	\theta^{\ab} \colon \pi_1(X,x_0)^{\ab} \longrightarrow H_1(X;\mathbb{Z})
\]
Diese Abbildung heißt \Index{Hurewicz-Abbildung}.

\begin{satz}[{name={Hurewicz}}]
	Für wegzusammenhängende Räume ist $\theta^{\ab}$ ein Isomorphismus.
\end{satz} 
\begin{beweis}
	Sei $X$ wegzusammenhängend. Wähle für jedes $x \in X$ einen Pfad $\lambda_x$ von $x_0$ nach $x$. Wähle $\lambda_{x_0} = \const_{x_0}$. Betrachte 
	\[
		\overline{\rho} \colon C_1(X;\mathbb{Z}) \longrightarrow \pi_1(X,x_0)^{\ab} \quad , 
		\qquad  \omega \longmapsto \benbrace*{\lambda_{\omega(0)} * \omega * \lambda_{\omega(1)}^{-1}}
	\]
	Dies ist ein Gruppenhomomorphismus, den wir auf den Generatoren angegeben haben. 
	Hier benutzen wir, dass $\pi_1(X,x_0)^{\ab}$ abelsch ist.
	Nach \autoref{lem:158} induziert $\overline{\rho}$ einen Gruppenhomomorphismus $\rho \colon H_1(X;\mathbb{Z}) \to \pi_1(X,x_0)^{\ab}$. 
	Für eine Schleife $\omega$ in $x_0$ gilt
	\begin{align*}
		(\rho \circ \theta)[\omega] = \benbrace*{\lambda_{\omega(0)} * \omega * \lambda_{\omega(0)} ^{-1}} = \benbrace*{\lambda_{x_0} * \omega * \lambda_{x_0}^{-1}}=[\omega] 
	\end{align*}
	Die Zuordnung $x \mapsto \lambda_x$ liefert einen Gruppenhomomorphismus $\lambda \colon C_0(X;\mathbb{Z}) \to C_1(X;\mathbb{Z})$  
	\[
		\sum\nolimits_{i} n_i \cdot x_i  \longmapsto \sum\nolimits_i n_i \cdot \lambda_{x_i}
	\]
	Mit \autoref{lem:159} folgt $(\theta \circ \rho)\enbrace[\big]{\bbenbrace{c}} = \bbenbrace{c}$ für alle $\bbenbrace{c} \in H_1(X)$. 
	Damit ist $\theta^{\ab}$ ein Isomorphismus für wegzusammenhängende Räume $X$.
\end{beweis}

\begin{lemma}[label=lem:158,{name=[Ränder von 2-Simplizes]}]
	$\overline{\rho}$ bildet die Ränder von $2$-Simplizes auf $1 \in \pi_1(X,x_0)^{\ab}$ ab.
\end{lemma}
\begin{beweis}
	Sei $\sigma \in C_2(X;\mathbb{Z})$ ein $2$-Simplex. Es gilt mit $\omega_i := \delta_i \sigma$
	\begin{align*}
		\overline{\rho}(\partial \sigma) &= \overline{\rho}(\delta_0 \sigma - \delta_1 \sigma + \delta_2 \sigma) \\
		&= \overline{\rho}(\omega_0) \cdot \overline{\rho}(\omega_1) ^{-1}   \cdot \overline{\rho}(\omega_2) \\
		&= \benbrace[\Big]{\lambda_{\omega_0(0)} * \omega_0 * \lambda_{\omega_0(1)} ^{-1} * 
		{\underbrace{\enbrace*{\lambda_{\omega_1(0)} * \omega_1 * \lambda_{\omega_1(1)}^{-1}} ^{-1}}_{=\lambda_{\omega_0(1)} * \omega_1 ^{-1} * \lambda_{\omega_1(0)} ^{-1}}}
		 * \lambda_{\omega_2(0)} * \omega_2 * \lambda_{\omega_2(1)}} \\
		 &= \benbrace[\Big]{\lambda_{\omega_0(0)} * {\underbrace{\omega_0 * \omega_1^{-1} * \omega_2}_{=: \gamma}} * \lambda_{\omega_2(1)=\omega_0(0)} ^{-1}} \\
		 \intertext{Der Pfad $\gamma$ ist die Schleife, die einmal auf dem Rand des $2$-Simplex herumläuft. Also 
		$\gamma \colon [0,1] \to \abs*{\partial \Delta^2} \hookrightarrow \abs*{\Delta^2} \to X$. Aber $\abs*{\partial \Delta^2} \hookrightarrow \abs*{\Delta^2}$ ist homotop zur 
		konstanten Abbildung auf $\omega_0(0)$ mit einer Homotopie, die den Punkt $\omega_0(0)$ fixiert.}
		&= \benbrace*{\lambda_{\omega_0(0)} * \lambda_{\omega_0(0)} ^{-1}} = 1 \qedhere
	\end{align*}
\end{beweis}

\begin{figure}[th]
	\centering{
	\begin{tikzpicture}[scale=1.8, font=\small]
		\clip (-1.67,-0.65) rectangle (1.5,1.51);
		% \draw (-2,-2) grid[step=5mm] (2,2);
		\coordinate (a) at (0,0);
		\coordinate (b) at (1,0);
		\coordinate (c) at (60:1);
		\coordinate (d) at (-1.4,1);
		\draw[thick,
			postaction={decorate},
			decoration={markings,
				mark=at position 1/6 with {\arrow{to reversed}},
				mark=at position 0.5 with {\arrow{to}},
				mark=at position 5/6 with {\arrow{to}}}] 
		(a) -- (b) node[midway, below]{$\omega_1$} node[below right]{$\omega_1(0)$} 
		-- (c) node[midway, above right]{$\omega_2$}  node[above right]{$\omega_0(0)$}
		-- cycle node[midway, above left]{$\omega_0$} node[below]{$\omega_0(1)$};
		\draw[dashed, thick] (d) node[left]{$x_0$}.. controls +(-70:0.8) and +(125:0.6) .. (a) node[midway, above,sloped]{$\lambda_{\omega_0(1)}$};
		\draw[dashed, thick] (d) .. controls +(-95:1.0) and +(-108:1.3) .. (b) node[midway, below,sloped]{$\lambda_{\omega_1(0)}$};
		\draw[dashed, thick] (d) .. controls +(20:0.7) and +(125:0.8) .. (c) node[midway, above,sloped]{$\lambda_{\omega_0(0)}$};
		
		\foreach \x in {a,b,c,d}{
			\draw[fill=DodgerBlue3, thick] (\x) circle[radius=0.03];
		}
	\end{tikzpicture}
	\caption{Skizze zum Beweis von \autoref{lem:158}}}
\end{figure}

\begin{lemma}[label=lem:159,{name=[Bild von 1-Simplizes unter der Komposition]}]
	Sei $\omega \colon \abs*{\Delta^1} \to X$ ein $1$-Simplex, dann gilt $\theta \circ \overline{\rho}(\omega) = \bbenbrace{\omega + \lambda_{\partial \omega}}$
\end{lemma}
\begin{beweis}
	Es gilt
	\[
		(\theta \circ \overline{\rho}) (\omega) = \theta \enbrace*{ \benbrace*{\lambda_{\omega(0)} * \omega * \lambda_{\omega(1)}^{-1}} } =
		\left[\!\!\left[\lambda_{\omega(0)} * \omega * \lambda_{\omega(1) ^{-1}} \right]\!\!\right] = \bbenbrace{\omega}  
		\underbrace{-\bbenbrace{\lambda_{\omega(1)}} + \bbenbrace{\lambda_{\omega(0)}}}_{\lambda_{\omega(0) -\omega(1)} = \lambda_{\partial \omega}} \qedhere
	\] 
\end{beweis}

\begin{korollar}[{name=[Erste Homologiegruppe des Produkts wegzusammenhängender Räume]}]
	Sind $X_1, \ldots , X_n$ wegzusammenhängende Räume, dann gilt
	\[
		H_1(X_1 \times \ldots \times X_n) \cong H_1(X_1) \oplus \ldots \oplus H_1(X_n)
	\]
\end{korollar}
\begin{beweis}
	\begin{align*}
		H_1(X_1 \times \ldots \times X_n) &\cong \pi_1 \enbrace*{X_1 \times \ldots \times X_n, (x_0^\ssbrace{1}, \ldots , x_0^\ssbrace{n})}^{\ab} \\
		&\cong \enbrace*{\pi_1 \enbrace*{X_1,x_0^\ssbrace{1}} \times \ldots \times \pi_1 \enbrace*{X_n, x_0^\ssbrace{n}}}^{\ab} \\
		&\cong \pi_1 \enbrace*{X_1, x_0^\ssbrace{1}}^{\ab} \oplus \ldots \oplus \pi_1 \enbrace*{X_n, x_0^\ssbrace{n}}^{\ab} \tag*{(Übungsaufgabe)} \\
		&\cong H_1(X_1) \oplus \ldots \oplus H_1(X_n) \qedhere
	\end{align*}
\end{beweis}

\begin{definition}[{name=[{n-te absolute Homotopiegruppe}]}]
	Sei $X$ ein topologischer Raum mit Basispunkt $x_0 \in X$, sei $I=[0,1]$ und $n \in \mathbb{N}_0$. 
	Die \bet{$\mathbf{n}$-te absolute Homotopiegruppe}\index{Homotopiegruppe} $\pi_n(X,x_0)$ ist die Menge der Äquivalenzklassen von stetigen Abbildungen $\omega\colon I^n \to X$ mit $\omega(\partial I^n)=\set{x_0}$. 
	Die Äquivalenzrelation ist Homotopie relativ zu $\partial I^n$, d.h. $H \colon I^n \times [0,1] \to X$ erfüllt $H_t(\partial I^n)= \set{x_0}$ für alle $t \in [0,1]$. 
	Die Gruppenverknüpfung in $\pi_n(X,x_0)$ ist gegeben durch 
	\[
		(\omega * \omega')(s_1,\ldots ,s_n) = \begin{cases}
			\omega(2s_1, s_2, \ldots ,s_n), &\text{ falls }0\le s_1 \le \frac{1}{2} \\
			\omega'(2s_1 -1, s_2, \ldots , s_n) ,&\text{ falls } \frac{1}{2} \le s_1 \le 1
		\end{cases}
	\]
\end{definition}

Bild für $n=2$:
\[	
	\vcenter{\hbox{\begin{tikzpicture}[scale=1.2]
		\draw (0,0) rectangle (1,1);
		\node at (0.5,0.5) {$\omega\vphantom{\omega^\prime}$};
	\end{tikzpicture}}}
	 * 
 	\vcenter{\hbox{\begin{tikzpicture}[,scale=1.2]
 		\draw (0,0) rectangle (1,1);
		\node at (0.5,0.5){$\omega^\prime$};
 	\end{tikzpicture}}}
	= 
	\vcenter{\hbox{\begin{tikzpicture}[scale=1.2]
		\draw (0,0) rectangle (1,1);
		\draw (0.5,0) -- (0.5,1);
		\node at (0.25,0.5) {$\omega\vphantom{\omega^\prime}$};
		\node at (0.75,0.5) {$\omega^\prime$};
	\end{tikzpicture}}}
\]

\begin{bemerkung}[{name=[{Eigenschaften von $\pi_n(X,x_0)$}]}]
	Die Verknüpfung hat die folgenden Eigenschaften:
	\begin{itemize}
		\item Die Verknüpfung \enquote{$*$} ist assoziativ, hat $\const_{x_0}$, also die konstante Abbildung auf $x_0$, als neutrales Element und 
		$\omega ^{-1}(s_1,\ldots ,s_n)= \omega(1-s_1,s_2, \ldots ,s_n)$ als inverses Element.
		\item Die Verknüpfung \enquote{$*$} ist abelsch für $n \ge 2$ (Eckmann-Hilton Argument\footnote{siehe auch \enquote{Analysis, Topologie und Geometrie} 14.5 und 14.6})
		\[
			\vcenter{\hbox{\begin{tikzpicture}[scale=1.4]
				\draw (0,0) rectangle (1,1);
				\draw (0.5,0) -- (0.5,1);
				\node at (0.25,0.5) {$\omega\vphantom{\omega^\prime}$};
				\node at (0.75,0.5) {$\omega^\prime$};
			\end{tikzpicture}}}
			\enspace\leadsto\enspace	
			\vcenter{\hbox{\begin{tikzpicture}[scale=1.4, every node/.style={draw,minimum size=5mm,inner sep=2pt,fill=white}]
				\draw[fill=DodgerBlue3!30] (0,0) rectangle (1,1);
				\node at (0.25,0.75) {$\omega$};
				\node at (0.75,0.25) {$\omega^\prime$};
			\end{tikzpicture}}}
			\enspace\leadsto\enspace
			\vcenter{\hbox{\begin{tikzpicture}[scale=1.4, every node/.style={draw,minimum size=5mm,inner sep=2pt,fill=white}]
				\draw[fill=DodgerBlue3!30] (0,0) rectangle (1,1);
				\node at (0.75,0.75) {$\omega$};
				\node at (0.25,0.25) {$\omega^\prime$};
			\end{tikzpicture}}}
			\enspace\leadsto\enspace
			\vcenter{\hbox{\begin{tikzpicture}[scale=1.4]
				\draw (0,0) rectangle (1,1);
				\draw (0.5,0) -- (0.5,1);
				\node at (0.25,0.5) {$\omega^\prime$};
				\node at (0.75,0.5) {$\omega\vphantom{\omega^\prime}$};
			\end{tikzpicture}}}	
		\]
		\item $\pi_n\colon \TOP  \to \ABELGRUPPEN$ ist ein Funktor, das heißt eine stetige Abbildung $f \colon (X,x_0) \to (Y,y_0)$ induziert einen Gruppenhomomorphismus 
		$\pi_n(f) \colon \pi_n(X,x_0) \to \pi_n(Y,y_0)$.
		\item Da $\nicefrac{I^n}{\partial I^n} \cong S^n$ ist, ergibt sich $\pi_n(X,x_0)$ auch aus den punktierten Homotopieklassen von Abbildungen $(S^n,*) \to (X,x_0)$.
		\item Falls $x_0$ und $x_0'$ in der gleichen Wegzusammenhangskomponente von $X$ liegen, dann induziert ein Weg von $x_0$ nach $x_0'$ einen Isomorphismus $\pi_n(X,x_0) \cong \pi_n(X,x_0')$.
	\end{itemize}
\end{bemerkung}

\begin{definition}[{name=[{n-zusammenhängender Raum}]}]
	Sei $n \in \mathbb{N}_0$. 
	Ein topologischer Raum $X$ heißt \bet{$\mathbf{n}$-zusammenhängend}\index{n-zusammenhängend@$n$-zusammenhängend}, falls $\pi_0(X)=\set{*}$ und $\pi_k(X,x_0)$ trivial ist für alle $k \le n$ und einen Basispunkt $x_0 \in X$.
\end{definition}

\begin{bemerkung}[{name=[Abhängigkeit von der Wahl eines Basispunktes]}]
	Da ein $n$-zusammenhängender Raum $X$ für $n \ge 0$ mindestens wegzusammenhängend ist, gilt $\pi_k(X,x_0) \cong \pi_k(X,x_0')$. 
	Somit ist die Definition unabhängig von der Wahl des Basispunktes.
\end{bemerkung}

Aus der Definition erhalten wir die folgenden Sprechweisen: $0$-zusammenhängend entspricht wegzusammenhängend und $1$"=zusammenhängend entspricht wegzusammenhängend und einfach zusammenhängend.

\begin{definition}[{name=[Relative Homotopiegruppen]}]
	Sei $X$ ein topologischer Raum, $A \subset X$ ein Unterraum und $x_0 \in A$ ein Basispunkt von $A$ und $X$. 
	Wir identifizieren $I^{n-1} \subset I^n$ mit der Seitenfläche von $I^n$, für die $s_n=0$ gilt. Außerdem setzen wir 
	\[
		J^{n-1} := \overline{\partial I^n \setminus I^{n-1}}
	\]
	Die \Index{relative Homotopiegruppe} $\pi_n(X,A,x_0)$ besteht aus den Homotopieklassen von stetigen Abbildungen $\omega\colon I^n  \to X$ mit $\omega(I^{n-1}) \subset A$ und $\omega(J^{n-1})=\set{x_0}$. 
	Die Homotopien müssen ebenfalls $H_t(I^{n-1}) \subset A$ und $H_t(J^{n-1}) =\set{x_0}$ für alle $t \in [0,1]$ erfüllen.
	$\pi_n(X,A,x_0)$ ist eine Menge für $n \ge 1$, eine Gruppe für $n \ge 2$ und eine abelsche Gruppe für $n > 2$. 
	Alternativ können wir relative Homotopiegruppen auch wie folgt beschreiben: $\pi_n(X,A,x_0)$ sind Homotopieklassen von Abbildungen $(D^n,S^{n-1},*) \to (X,A,x_0)$.
\end{definition}

\begin{figure}[hbt]
	\centering{
	$\enbrace*{\vcenter{\hbox{\begin{tikzpicture}[scale=1.4,my brace/.style={decorate, decoration={brace, amplitude=5pt, mirror, raise=3pt}},font=\small]
		\draw (0,0) -- (1,0) -- (1,1) node[midway,right]{$\const_{x_0}$} -- (0,1) node[midway, above]{$\const_{x_0}$} -- cycle node[midway, left]{$\const_{x_0}$};
		\draw[my brace] (0,0) -- (1,0) node[midway, below=6pt]{$ \to A$};
		\node at (0.5,0.5) {$\omega\vphantom{\omega^\prime}$};
	\end{tikzpicture}}}}
	 * 
 	\enbrace*{\vcenter{\hbox{\begin{tikzpicture}[,scale=1.4, my brace/.style={decorate, decoration={brace, amplitude=5pt, mirror, raise=3pt}},font=\small]
 		\draw (0,0) -- (1,0) -- (1,1) node[midway,right]{$\const_{x_0}$} -- (0,1) node[midway, above]{$\const_{x_0}$} -- cycle node[midway, left]{$\const_{x_0}$};
		\draw[my brace] (0,0) -- (1,0) node[midway, below=6pt]{$ \to A$};
		\node at (0.5,0.5){$\omega^\prime$};
 	\end{tikzpicture}}}}
	= 
	\vcenter{\hbox{\begin{tikzpicture}[scale=1.4,font=\small, my brace/.style={decorate, decoration={brace, amplitude=5pt, mirror, raise=3pt}}]
		\draw (0,0) -- (1,0) -- (1,1) node[midway,right]{$\const_{x_0}$} -- (0,1) node[midway, above]{$\const_{x_0}$} -- cycle node[midway, left]{$\const_{x_0}$};
		\draw (0.5,0) -- (0.5,1);
		\draw[my brace] (0,0) -- (1,0) node[midway, below=6pt]{$ \to A$};
		\node at (0.25,0.5) {$\omega\vphantom{\omega^\prime}$};
		\node at (0.75,0.5) {$\omega^\prime$};
	\end{tikzpicture}}}$
	\caption{Verknüpfung in der relativen Homotopiegruppe für $n=2$}
	}
\end{figure}

\begin{bemerkung}[{name=[{Homotopiegruppen bilden keine verallgemeinerte Homologietheorie}]}]
	Homotopiegruppen bilden \emph{keine} verallgemeinerte Homologietheorie, da sie nicht das Ausschneidungsaxiom erfüllen.
	Es gibt eine lange exakte Sequenz von Homotopiegruppen:
	\[
		\begin{tikzcd}
			\cdots \rar & \pi_n(A,x_0) \rar & \pi_n(X,x_0) \rar & \pi_n(X,A,x_0) \rar["\partial"] & \pi_{n-1}(A,x_0) \rar & \cdots
		\end{tikzcd}
	\]
\end{bemerkung}

\begin{definition}[{name={Hurewicz-Abbildung}}]
	Sei $\vartheta \colon \abs*{\Delta^n} \xrightarrow{\cong} D^n$ ein Homöomorphismus. 
	Dann gilt $\vartheta \enbrace*{\abs*{\partial \Delta^n}}= S^{n-1}$ und $\vartheta$ liefert eine Homologieklasse $\bbenbrace{\vartheta} \in H_n(D^n,S^{n-1}) \cong \mathbb{Z}$. 
	Diese Klasse ist ein Generator der Gruppe $H_n(D^n,S^{n-1})$. 
	Sei nun $[\omega] \in \pi_n(X,A,x_0)$ repräsentiert durch $\omega \colon (D^n,S^{n-1},*) \to (X,A,x_0)$. Setze 
	\[
		\theta_n \colon \pi_n(X,A,x_0) \longrightarrow H_n(X,A) \enspace, \quad [\omega] \longmapsto \omega_*\bbenbrace{\vartheta} 
	\]
	für $n \ge 2$. Dies ist wohldefiniert, denn falls $\omega$ homotop zu $\omega'$ ist (bezüglich $(A,x_0)$), dann ist $\omega_*=\omega_*'$.
\end{definition}

\begin{lemma}[{name=[Die Hurewicz-Abbildung ist ein Gruppenhomomorphismus]},label=lem:1518]
	$\theta_n \colon \pi_n(X,A,x_0) \longrightarrow H_n(X,A)$ ist ein Gruppenhomomorphismus.
\end{lemma}
\begin{beweis}
	Wir brauchen eine alternative Beschreibung der Gruppenverknüpfung: 
	Sei dazu $\gamma$ die Projektion $D^n \to \nicefrac{D^n}{S^{n-2}} \cong D^n \vee D^n$. Dann gilt $\gamma(S^{n-1})= S^{n-1} \vee S^{n-1}$ und folglich ist $\gamma$ eine Abbildung von Raumpaaren
	$\gamma \colon (D^n,S^{n-1}) \to (D^n \vee D^n, S^{n-1} \vee S^{n-1})$. Seien nun $\omega, \omega' \colon (D^n, S^{n-1},*) \to (X,A,x_0)$ 
	Für die Verknüpfung $[\omega * \omega'] \in \pi_n(X,A,x_0)$ gilt
	\[
		\benbrace[\big]{(\omega \vee \omega') \circ \gamma} = [\omega * \omega'] 
	\]
	\begin{figure}[bht]
		\centering{
		\begin{tikzpicture}[scale=0.8,>={Classical TikZ Rightarrow[width=2mm,length=1mm]}]
			\coordinate (a) at (5,-1.2);
			\coordinate (b) at (7,-1.2);
			\coordinate (x1) at (5.7,-1.8);
			\coordinate (x2) at (6.3,-1.8);
			\draw (-1,0) arc (180:360:1cm and 0.5cm);
		    \draw[dashed, black!60] (-1,0) arc (180:0:1cm and 0.5cm);
		    \draw (0,0) circle (1cm);
		    \shade[ball color=blue!20!white,opacity=0.20] (0,0) circle (1cm);
		
			\draw[->] (1.5,0) -- (3.5,0) node[midway, above]{$\gamma$};
		
			\draw (4,0) arc (180:360:1cm and 0.5cm);
		    \draw[dashed, black!60] (4,0) arc (180:0:1cm and 0.5cm);
		    \draw (5,0) circle (1cm);
		    \shade[ball color=blue!20!white,opacity=0.20] (5,0) circle (1cm);
		
			\draw (6,0) arc (180:360:1cm and 0.5cm);
		    \draw[dashed, black!60] (6,0) arc (180:0:1cm and 0.5cm);
		    \draw (7,0) circle (1cm);
		    \shade[ball color=blue!20!white,opacity=0.20] (7,0) circle (1cm);
			\draw[->] (a) .. controls +(-90:0.4) and +(180:0.4).. (x1) node[midway, below left,yshift=2mm]{$\omega\vphantom{\omega^\prime}$};
			\draw[->] (b) .. controls +(-90:0.4) and +(0:0.4).. (x2) node[midway, below right,yshift=2mm]{$\omega'$};
			\node at (6,-1.8){$X$};
		\end{tikzpicture}
		\caption{Die Abbildung $\gamma$ aus \autoref{lem:1518}}}
	\end{figure}%
	Sei $p_j \colon (D^n \vee D^n, S^{n-1} \vee S^{n-1}) \to (D^n,S^{n-1})$ die Abbildung, die die $j$-te Kopie von $D^n$ auf den Basispunkt abbildet und auf der anderen Kopie die 
	Identität ist. Sei weiter $i_k \colon (D^n,S^{n-1}) \to (D^n \vee D^n, S^{n-1} \vee S^{n-1})$ die Inklusion auf die $k$-te Kopie von $(D^n, S^{n-1})$. Betrachte nun die Abbildung
	\[
		\enbrace[\big]{(p_1)_*, (p_2)_*} \colon H_n(D^n \vee D^n, S^{n-1} \vee S^{n-1}) \longrightarrow H_n(D^n,S^{n-1}) \oplus H_n(D^n,S^{n-1}) 
	\]
	Dies ist ein Isomorphismus mit inverser Abbildung
	\[
		(i_{1})_* + (i_2)_* \colon H_n(D^n,S^{n-1}) \oplus H_n(D^n,S^{n-1}) \longrightarrow H_n(D^n \vee D^n, S^{n-1} \vee S^{n-1})
	\]
	denn es gilt
	\begin{align*}
		\enbrace[\big]{(p_{1})_*, (p_{2})_*} \circ \enbrace[\big]{(i_{1})_* + (i_{2})_*} = \enbrace*{
		\enbrace[\Big]{\underbracket[0.8pt][1.2pt]{p_1 \circ i_1}_{\id}}_* + \enbrace[\Big]{\underbracket[0.8pt][1.2pt]{p_1 \circ i_2}_{\const}}_* , 
		\enbrace[\Big]{\underbracket[0.8pt][1.2pt]{p_2 \circ i_1}_{\const}}_* + \enbrace[\Big]{\underbracket[0.8pt][1.2pt]{p_2 \circ i_2}_{\id}} _*} = \id
	\end{align*}
	Also ist $\enbrace[\big]{(p_1)_*, (p_2)_*}$ eine surjektive Abbildung und Quelle und Ziel sind freie abelsche Gruppen von gleichem Rang. Die Abbildung $p_j \circ \gamma$ ist homotop 
	zur Identität (Übung). Also folgt 
	$\enbrace[\big]{(p_1)_*, (p_2)_*} \circ \gamma \enbrace*{\bbenbrace{\vartheta}} = \enbrace*{\bbenbrace{\vartheta}, \bbenbrace{\vartheta}} $ und weiter
	\begin{align*}
		\theta_n \enbrace*{\benbrace*{\omega * \omega'}} = (\omega \vee \omega')_* \circ \gamma_* \enbrace*{\bbenbrace{\vartheta}} &= 
		(\omega \vee \omega')_* \circ \enbrace[\big]{(i_1)_* + (i_2)_*} \circ \enbrace[\big]{(p_1)_*, (p_2)_*}   \circ \gamma_* \enbrace*{\bbenbrace{\vartheta}} \\
		&= (\omega \vee \omega')_* \circ \enbrace*{(i_1)_* \enbrace*{\bbenbrace{\vartheta}} + (i_2)_* \enbrace*{\bbenbrace{\vartheta}} } \\
		&= \omega_* \enbrace*{\bbenbrace{\vartheta}}  + \omega'_* \enbrace*{\bbenbrace{\vartheta}} \qedhere
	\end{align*}
\end{beweis}

\begin{satz}[{name={Hurewicz}}]
	Seien $(X,A)$ mit $A \subset X$ beide wegzusammenhängend. 
	Es gelte weiter $\pi_1(A,x_0)=1$ und $\pi_k(X,A,x_0)=0$ für $1 \le k \le n-1$, $n \ge 2$. Dann ist 
	\[
		\theta_n \colon \pi_n(X,A,x_0) \longrightarrow H_n(X,A)
	\]
	ein Isomorphismus.
\end{satz}

\begin{korollar}[{name=[{Isomorphie von Homotopie- und Homologiegruppen für $(n-1)$-zsmhg. Räume}]},label=korr:1520]
	Ist $X$ ein $(n-1)$-zusammenhängender Raum für ein $n \ge 2$, dann ist 
	\[
		\theta_n \colon \Underbrace{\pi_n(X,x_0)}{=\pi_n(X,\set{x_0},x_0)} \longrightarrow \enspace H_n(X,x_0)\cong H_n(X)
	\]
	ein Isomorphismus.
\end{korollar}

\begin{satz}[{name=[{Homotopiegruppen der Sphäre}]}]
	Sei $n \ge 1$. Dann ist $\pi_n(S^n,*)\cong \mathbb{Z}$ und $\pi_k(S^n,*)=0$ für $k <n$.
\end{satz}
\begin{beweis}
	Für $n=1$ gilt $\pi_0(S^n,*)= \set{*}$ und $\pi_1(S^1,*)\cong \mathbb{Z}$. Für $n>1$ wissen wir aus dem letzten Semester, dass $\pi_1(S^n,*)=1$ und $S^n$ wegzusammenhängend ist, also 
	$1$-zusammenhängend ist.

	Angenommen $\pi_i(S^n,*)=0$ für $i=0,\ldots ,k$ und $0 \le k < n$. Dann ist $S^n$ $k$-zusammenhängend und wir erhalten aus dem vorigen \autoref{korr:1520}
	\[
		\pi_{k+1}(S^n,*) \xrightarrow{\minwidthbox{\cong}{2em}} H_{k+1}(S^n)= \begin{cases}
			\mathbb{Z}, &\text{ falls }k+1=n\\
			0 , &\text{ sonst}
		\end{cases}
	\]
	Falls $k+1<n$ ist, so ist auch $\pi_{k+1}(S^n,*)=0$, sonst sind wir fertig.
\end{beweis}

\newpage
\section{Die Fundamentalklasse} % (fold)
\label{sec:16}
\todo{das könnte man nochmal genau durchgehen}
\begin{beispiel}[{name=[Homologiegruppen verschiedener Mannigfaltigkeiten]}]
	Wir betrachten zu einigen uns gut bekannten Mannigfaltigkeiten die Homologiegruppen in der Dimension der Mannigfaltigkeit:
	\begin{multicols}{2}
		\begin{itemize}
			\item $H_n(\mathbb{R}^n)=0$
			\item $H_n(S^n)=\mathbb{Z}$
			\item $H_2(T^2)=\mathbb{Z}$
			\item $H_{2n}(\mathbb{C}P^n)=\mathbb{Z}$
			\item $H_{2n}(S^n \times \mathbb{R}^n)=0$
			\item $H_n(\mathbb{R}P^n)=\mathbb{Z}$, falls $n$ ungerade, $0$ sonst.
		\end{itemize}
	\end{multicols}
\end{beispiel}

\begin{bemerkung}[{name=[Homologie von Mannigfaltigkeiten]}]
	Sei $M$ ein $n$-Mannigfaltigkeit und $x \in M$. 
	Da $x$ in $M$ eine Umgebung hat,\marginnote{Ausgeschnitten wird das Komplement dieser Umgebung!} die homöomorph zu $\mathbb{R}^n$ ist, folgt mit dem Ausschneidungsisomorphismus und der Paarfolge
	\[
		H_k\enbrace[\big]{M, M\setminus \set{x};R} \cong H_k \enbrace[\big]{\mathbb{R}^n,  \mathbb{R}^n \setminus \set{0};R} \cong H_{k-1}\enbrace[\big]{\mathbb{R}^n \setminus 
		\set{0};R} \cong H_{k-1}(S^{n-1};R)\cong \begin{cases}
			R &\text{ falls }k=n\\
			0 &\text{ sonst}
		\end{cases}
	\]
\end{bemerkung}

\begin{definition}[{name=[Homologische Orientierung in einem Punkt]}]
	Sei $M$ eine $n$-Mannigfaltigkeit und $x \in M$. 
	Eine (homologische) \bet{Orientierung}\index{homologische Orientierun!im Punkt} (über $R$) im Punkt $x$ ist ein Erzeuger $\mu_x$ von $H_n(M,M\setminus \set{x};R)$.
\end{definition}

Über $R=\mathbb{Z}$ gibt es in jedem Punkt genau zwei Orientierungen, da $\mathbb{Z}$ genau zwei Erzeuger hat. 
Über $R=\mathbb{F}_2$ gibt es in jedem Punkt eine eindeutige Orientierung, da $\mathbb{F}_2$ genau einen Erzeuger hat.

\begin{notation}[{name=[{Einschränkung eines Elementes $v\in H_n(X,X\setminus B;R)$}]}]
	Seien $A \subseteq B \subseteq X$ Unterräume. Sei $i\colon (X,X\setminus B) \hookrightarrow \enbrace*{X,X\setminus A}$ die Inklusion. 
	Ist $v\in H_n(X,X\setminus B;R)$, so nennen wir 
	\[
		v\big|_A := i_*(v)
	\]
	die Einschränkung von $v$ auf $A$. Für $A=\set{x}$ schreiben wir $v|_x := v|_{\set{x}}$.
\end{notation}

\begin{definition}[{name=[{Homologische Orientierung der gesamten Mannigfaltigkeit}]}]
	Sei $M$ eine $n$-Mannigfaltigkeit. 
	Eine (homologische) \bet{Orientierung}\index{homologische Orientierun!global} von $M$ über $R$ ist eine Zuordnung $M \ni x \mapsto \mu_x \in H_n(M,M \setminus\set{x};R)$, die für jeden Punkt $x\in M$ eine Orientierung über $R$ im Punkt $x$ auswählt, sodass folgende Bedingung erfüllt ist:
	Für alle $x \in M$ gibt es eine Umgebung $U$ von $x \in M$ und $\mu_U \in H_n(M,M \setminus U;R)$ mit $\mu_y =\mu_U|_y$ für alle $y \in U$.

	Gibt es eine solche Orientierung, so heißt $M$ $R$-\bet{orientierbar}\index{orientierbare Mannigfaltigkeit}. 
	Eine Mannigfaltigkeit zusammen mit einer Orientierung heißt eine $R$-\bet{orientierte Mannigfaltigkeit}\index{orientierte Mannigfaltigkeit}. 
	Lassen wir $R$ in den Bezeichnungen weg, so ist immer $R=\mathbb{Z}$ gemeint.
\end{definition}

\begin{beispiel}[{name=[Orientierung der Sphäre]}]
	Sei $\mu_{S^n} \in \tilde{H}_n(S^n;R)$ ein Erzeuger. 
	Für $x \in S^n$ ist $\tilde{H}_*(S^n \setminus \set{x};R)= 0$ und daher die Einschränkungsabbildung $\tilde{H}_n(S^n;R) \to \tilde{H}_*(S^n,S^n \setminus \set{x};R)= 0$ ein Isomorphismus und $\mu_{S^n}|_x$ ist ein Erzeuger von $\tilde{H}_n(S^n, S^n\setminus \set{x};R)$. 
	Daher ist $x \mapsto \mu_{S^n}|_x$ eine Orientierung von $M$ über $R$.
\end{beispiel}

Eine \Index{geschlossene Mannigfaltigkeit} ist eine kompakte Mannigfaltigkeit (ohne Rand).

\begin{satz}[{name=[{$R$-Orientierbarkeit geschlossener Mannigfaltigkeiten}]}]
	Sei $M$ eine geschlossene $n$-Mannigfaltigkeit. 
	Dann ist $M$ genau dann $R$-orientierbar, wenn es $\mu \in H_n(M;R)$ gibt, sodass $\mu|_x \in H_n(M,M\setminus \set{x};R)$ für jedes $x$ eine Orientierung von $M$ in $x$ ist. 
	Ist $M$ zusammenhängend, so ist $\mu$ ein Erzeuger von $H_n(M;R)$ und $H_n(M;R)\cong R$.
\end{satz}
\begin{beweis}[{name={mit \autoref{lem:1611}}}]
	$H_n(M;R)$ und $\Gamma(M,\mathcal{H})$ sind beide mit $\coprod$ verträglich. 
	Daher können wir auch annehmen, dass $M$ zusammenhängend ist.
	Da $M$ kompakt ist, können wir in dem \autoref{lem:1611} $A=M$ setzen und wir erhalten für jede Orientierung $s \in \Gamma(M,\mathcal{H})$ eine Klasse $\mu \in H_n(M;R)$ mit $\mu|_x=s(x)$. 
	Nach der Überlagerungstheorie aus dem letzten Semester ist jeder Schnitt von $\mathcal{H}$ schon durch seinen Wert an einem Punkt festgelegt. 
	Insbesondere ist $\ev_x \colon\Gamma(M,\mathcal{H}) \to H_n(M,M \setminus \set{x};R ) \cong R$ injektiv. 
	Ist $M$ $R$-orientierbar, so ist $\ev_x$ ein Isomorphismus und es folgt
	\[
		H_n(M;R) \cong \Gamma(M,\mathcal{H}) \cong H_n(M,M \setminus \set{x};R)\cong R \qedhere
	\]
\end{beweis}

\begin{definition}[{name=[{Disjunkte Vereinigung der Homologieklassen}]}]
	Sei $M$ eine geschlossene $n$-Mannigfaltigkeit. 
	Wir definieren $\mathcal{H}=\mathcal{H}_R$ als die disjunkte Vereinigung aller $H_n(M,M\setminus \set{x};R)$ mit $x \in M$. 
	Auf $\mathcal{H}$ definieren wir nun eine Topologie  durch:
	\begin{align*}
		V \subseteq \mathcal{H} \text{ offen } :\iff \enspace& \forall v \in V \text{ mit } v \in H_n(M,M \setminus \set{x};R) \text{ gibt es eine Umgebung $U$ von $x$ in $M$}\\
		&\text{und $v_U \in H_n(M,M \setminus U,R)$ mit } v_U\big|_x=v \text{ und } v_U\big|_y \in V \text{ für alle $y \in U$.}
	\end{align*}
	Wir definieren $\pi \colon \mathcal{H} \to M$ durch $\pi(v)=x$ für $v \in H_n(M,M \setminus \set{x};R)$. 
	Es gilt $\pi ^{-1}(\set{x})= H_n(M,M \setminus \set{x};R)$.
	Mit $\Gamma(M,\mathcal{H})$ bezeichnen wir die Menge aller stetigen Schnitte $s\colon M \to \mathcal{H}$ für $\pi$, das heißt 
	\[
		\Gamma(M,\mathcal{H})= \set[\big]{s \colon M \to \mathcal{H} \text{ stetig} \given \pi  \circ s =\id_M} 
	\]
\end{definition}

\begin{bemerkung}[{name=[{Eigenschaften der Vereinigung der Homologieklassen}]},label=bem:1612]
	Einige grundlegende Eigenschaften:
	\begin{enumerate}[1),itemsep=1pt]
		\item Für jedes $x \in M$ erbt $H_n(M,M \setminus \set{x};R) \subseteq \mathcal{H}$ die diskrete Topologie.
		\item $\pi \colon \mathcal{H} \to M$ ist eine Überlagerung.
		\item $\Gamma(M,\mathcal{H})$ ist ein $R$-Modul. Sind $s,s' \in \Gamma(M,\mathcal{H})$ und $r \in R$, so können wir $s+s'$ und $r \cdot s$ punktweise erklären:
		\begin{align*}
			(s+s') (x) &:= s(x) + s'(x) \in H_n(M,M \setminus \set{x};R )\\
			(r \cdot s) (x) &:= r \cdot s(x) \in H_n(M,M \setminus \set{x};R )
		\end{align*}
		\item $R$-Orientierungen von $M$ entsprechen genau den Schnitten $s \in \Gamma(M,\mathcal{H})$, für die gilt: Für alle $x \in M$ ist 
		$s(x) \in H_n(M,m \setminus \set{x};R)$ ein Erzeuger.
		\item Eine Folgerung aus der Eindeutigkeit im Hebungssatz aus der Überlagerungstheorie ist folgende Aussage:
	
		Sei $M$ zusammenhängend. Dann ist für jedes $x \in M$ die offensichtliche Auswertungsabbildung $\ev_x\colon \Gamma(M,\mathcal{H}) \to H_n(M,M \setminus \set{x};R)$, $\ev_x(s)=s(x)$ injektiv.
		\item Sei $v \in H_n(M;R)$. Dann erhalten wir durch $(s_v)(x) := v|_x$ einen Schnitt $s_v\in \Gamma(M,\mathcal{H})$. Dies definiert einen $R$-Modul-Homomorphismus
		$H_n(M;R) \to \Gamma(M,\mathcal{H})$, $v \mapsto s_v$.
		\item Es gibt immer den Nullschnitt in $\Gamma(M,\mathcal{H})$, $x \mapsto 0\in H_n(M,M\setminus \set{x};R)$.
		\item Ist $M$ zusammenhängend, so ist $\Gamma(M,\mathcal{H})\not=0$ genau dann, wenn $M$ $R$-orientierbar ist.
		\item Für $M$ zusammenhängend gilt: $M$ ist $R$-orientierbar $\iff H_n(M;R) \cong R$. 
	\end{enumerate}
\end{bemerkung}

\begin{beispiel}[{name=[Für den Restklassenkörper $\mod 2$ erhält man eine 2-fache Überlagerung]}]
	Für $R=\mathbb{F}_2$ ist $\pi\colon \mathcal{H} \to M$ eine $2$-fache Überlagerung. 
	Da $\mathcal{H}$ immer den Nullschnitt zulässt, ist $\mathcal{H}$ dann trivial. 
	Insbesondere ist jede $n$-Mannigfaltigkeit $\mathbb{F}_2$-orientierbar.
\end{beispiel}

\begin{lemma}[{name=[{über kompakten Teilraum einer $n$-Mannigfaltigkeit}]},label=lem:1611]
	Sei $M$ eine $n$-Mannigfaltigkeit und $A \subseteq M$ kompakt. Dann gilt
	\begin{enumerate}[a)]
		\item \label{enum:16:1} Zu jedem $s \in \Gamma(M,\mathcal{H})$ gibt es ein eindeutiges $v_A \in H_n(M, M \setminus A;R)$ mit $v_A|_x=s(x)$ für alle $x \in A$.
		\item \label{enum:16:2}Für $i>n$ ist $H_i(M,M \setminus A;R)=0$.
	\end{enumerate}
\end{lemma}
\begin{beweis}
	Wir beweisen die Aussage in vier Schritten:
	\begin{enumerate}[label=\textbf{Schritt \arabic*:},wide=0cm]
		\item Seien $A_1,A_2 \subseteq M$ kompakt. Wir zeigen: Gelten \ref{enum:16:1} und \ref{enum:16:2} für $A_1$, $A_2$ und $A_1 \cap A_2$, so auch für $A_1 \cup A_2$.
		Mit der Mayer-Vietoris-Sequenz aus \autoref{bem:1612} folgt aus 
		\[
			H_i(M,M \setminus A_1;R) = H_i(M,M \setminus A_2;R) = H_{i+1}(M,M \setminus A_1 \cap A_2;R) = 0
		\]
		direkt $H_i(M, M \setminus A_1 \cup A_2;R)=0$. Daher gilt \ref{enum:16:2} für $A_1 \cup A_2$.
		Seien $v_1 \in H_n(M,M \setminus A_1;R)$ und $v_2 \in H_n(M,M \setminus A_2;R)$ mit $v_1\big|_x= s(x)$ für alle $x \in A_1$ und $v_2\big|_x=s(x)$ für alle $x \in A_2$.
		Dann gilt 
		\[
			(v_1|_{A_1 \cap A_2})\big|_x = s(x) = (v_2|_{A_1\cap A_2})\big|_x
		\]
		für alle $x \in A_1 \cap A_2$. Da $A_1 \cap A_2$ die Eindeutigkeitsaussage aus \ref{enum:16:1} erfüllt, folgt $v_1\big|_{A_1 \cap A_2} = v_2\big|_{A_1 \cap A_2}$.
		Betrachte nun die Mayer-Vietoris-Folge
		\[
			\begin{tikzcd}[font=\small,column sep=1.5em, row sep=0.1em]
				 H_{n+1}(M,M \setminus A_1 \cap A_2) \rar["\partial"] & H_n(M,M \setminus A_1 \cup A_2) \rar & 
				{\displaystyle\bigoplus_{H_n(M,M \setminus A_2)}^{H_n(M,M\setminus A_1)}}
				\rar & H_n(M,M\setminus A_1 \cap A_2) \\
				0 \ar[draw=none]{u}[sloped,auto=false]{=} & & (v_1,v_2) \arrow[draw=none]{u}[sloped,auto=false]{\in} & v_1|_{A_1 \cap A_2} = v_2|_{A_1 \cap A_2}
				\arrow[draw=none]{u}[sloped,auto=false]{\in}
			\end{tikzcd}
			\marginnote{\enquote{$R$} aus Platzgründen weggelassen}
		\]
		Daher gibt es ein eindeutiges $v \in H_n(M,M \setminus A_1 \cup A_2;R)$ mit $v|_{A_1} =v_1$ und $v|_{A_2}=v_2$. Für $x \in A_1$ folgt $v|_{x}=v_1|_x=s(x)$ und für 
		$x \in A_2$ folgt genauso $v|_x = v_2|_x =s(x)$. Ist nun $v' \in H_n(M,M \setminus A_1 \cup A_2;R)$ eine weitere Homologieklasse mit $v'|_x=s(x)$ für alle 
		$x \in A_1 \cup A_2$, so folgt $v'|_{A_1}=v_1$ und $v'|_{A_2}=v_2$ aus der Eindeutigkeit von $v_1$ und $v_2$. Daher gilt auch $v=v'$.
		\item Das Lemma ist richtig, wenn $A$ eine abgeschlossene konvexe Teilmenge von $M=\mathbb{R}^n$ ist. Dann ist 
		$(\mathbb{R}^n, \mathbb{R}^n \setminus A) \to (\mathbb{R}^n, \mathbb{R}^n \setminus \set{x})$ eine Homotopieäquivalenz von Paaren und daher 
		$H_n(\mathbb{R}^n, \mathbb{R}^n \setminus A;R) \to H_n(\mathbb{R}^n,\mathbb{R}^n \setminus \set{x};R)$ ein Isomorphismus. Folglich gibt es 
		$v \in H_n(\mathbb{R}^n,\mathbb{R}^n \setminus A;R)$ mit $v|_x=s(x)$. Da $s|_A$ und $s_v$ beides Schnitte der Überlagerung $\pi ^{-1}(A) \to A$ sind, folgt mit der
		Eindeutigkeit im Hebungssatz für alle $y \in A$
		\[
			v\big|_y = s_v(y) = s(y)
		\]
		Ist $v' \in H_n(\mathbb{R}^n, \mathbb{R}^n \setminus A;R)$ eine zweite Homologieklasse mit $v'|_x=s(x)$, so folgt $v=v'$, da $v'|_x =s(x)=v|_x$ ist und $w \mapsto w|_x$
		ein Isomorphismus $H_n(\mathbb{R}^n,\mathbb{R}^n \setminus A;R) \to H_n(\mathbb{R}^n, \mathbb{R}^n\setminus \set{x};R )$ ist. Offenbar ist 
		$\mathbb{R}^n \setminus A \simeq S^{n-1}$ und daher $H_i(\mathbb{R}^n \setminus A;R)=0$ für alle $i >n$.
		\item Das Lemma gilt für $M=\mathbb{R}^n$ und eine beliebige kompakte Teilmenge $A \subseteq M$: Ist $A$ die Vereinigung von endlichen vielen konvexen und kompakten 
		Teilmengen von $\mathbb{R}^n$, so folgt die Behauptung für $A$ aus Schritt 1 und 2. Für ein beliebiges kompaktes $A \subseteq \mathbb{R}^n$ folgt die Existenzaussage in 
		\ref{enum:16:1} aus Schritt 2. Wähle dazu $B \subseteq \mathbb{R}^n$ konvex und kompakt mit $A \subseteq B$. Nach Schritt 2 gibt es 
		$v \in H_n(\mathbb{R}^n, \mathbb{R}^n \setminus B;R)$ mit $v|_x=s(x)$ für alle $x \in B$. Für $v|_A \in H_n(\mathbb{R}^n,\mathbb{R}^n \setminus A;R)$ gilt dann 
		$(v|_A)\big|_x = v|_x=s(x)$ für alle $x \in A$.
	
		Sei nun $v \in H_d(\mathbb{R}^n,\mathbb{R}^n \setminus A;R)$ mit $d \ge n$. Falls $d=n$ ist, so sei $v|_x=0$ für alle $x \in A$. Wir müssen zeigen: $v=0$. Sei 
		$z \in C_d(\mathbb{R}^n;R)$ ein Repräsentant von $v$. Insbesondere ist $\partial z \in C_{d-1}(\mathbb{R}^n \setminus A;R)$. $\partial z$ ist eine 
		$R$-Linearkombination von endlich vielen singulären $d-1$-Simplizes in $\mathbb{R}^n \setminus A$. Da jeder dieser singulären $d-1$-Simplizes ein kompaktes Bild in 
		$\mathbb{R}^n \setminus A$ hat, finden wir endlich viele Bälle $B_i \subseteq \mathbb{R}^n$ mit
	
		\begin{enumerate*}[(1),itemjoin={,\hspace{1em}},itemjoin*={{\hspace{1ex}und\hspace{1ex}}}]
			\item $\forall i : B_i \cap A \not= \emptyset$
			\item $A \subseteq B_1 \cup \ldots \cup B_n$
			\item $\partial z \in C_{d-1}\enbrace[\big]{\mathbb{R}^n \setminus (B_1 \cup \ldots \cup B_n);R}$
		\end{enumerate*}
	
		Es folgt, dass $z$ auch eine Homologieklasse $v' \in H_d \enbrace[\big]{\mathbb{R}^n,\mathbb{R}^n \setminus (B_1 \cup \ldots \cup B_n);R}$ repräsentiert mit $v'|_A=v$.
		Ist $d>n$, so folgt $v'=0$ und daher auch $v=v'|_A=0$. Ist $d=n$, so wählen wir $x_i \in B_i \cap A$. Da $v'|_{x_i}=v|_{x_i}= 0$ ist, folgt $v'|_{B_i}=0$, weil die 
		Einschränkungsabbildung $w \mapsto w|_{x_i}$ einen Isomorphismus 
		$H_n(\mathbb{R}^n, \mathbb{R}^n \setminus B_i;R) \to H_n(\mathbb{R}^n,\mathbb{R}^n \setminus \set{x_i};R )$ definiert. Daher gilt $v'|_y=0$ für alle 
		$y \in B_1 \cup \ldots \cup B_n$. Es folgt $v'=0$ und damit auch $v=v'|_A=0$.
		\item Für den allgemeinen Fall schreibt man $A=A_1 \cup \ldots \cup A_n$, sodass jedes $A_i$ kompakt ist und in einem Kartengebiet $U_i$ von $M$ enthalten ist. Da wir 
		Ausschneidungsisomorphismen $H_*(U_i,U_i\setminus A;R) \to H_*(M,M \setminus A;R)$ erhalten und $U_i \cong \mathbb{R}^n$ ist, gilt die Behauptung nach Schritt 3 für 
		jedes $A_i$. Mit Schritt 1 folgt die Behauptung auch für $A$. \qedhere
	\end{enumerate}
\end{beweis}

\begin{bemerkung}[{name=[{Existenz einer Mayer-Vietoris-Folge für $A,B \subseteq X$ abgeschlossen}]}]
	\leavevmode
	\begin{enumerate}[1)]
		\item Zu Unterräumen $A \subseteq B \subseteq X$ gibt es eine kurze exakte Folge von Kettenkomplexen 
		\[
			\begin{tikzcd}
				C_*(B,A;R) \rar[hook] & C_*(X,A;R) \rar[two heads] & C_*(X,B;R) 
			\end{tikzcd}
		\]
		Damit erhalten wir eine lange exakte Folge in Homologie, die man als lange exakte Folge des 
		Tripels bezeichnet:
		\[
			\begin{tikzcd}[sep=small]
				\cdots \rar & H_{n+1}(X,B;R) \rar["\partial"] & H_n(B,A;R) \rar & H_n(X,A;R) \rar & H_n(X,B;R) \rar & \cdots
			\end{tikzcd}
		\]
		\item Zu abgeschlossenen Unterräumen $A,B \subseteq X$ ist $A \setminus B \subseteq X \setminus A \cap B$ abgeschlossen und $X \setminus B\subseteq X \setminus A \cap B$
		offen. Da $A \setminus B  \subseteq X \setminus B$ und 
		\begin{align*}
			\enbrace*{X \setminus A \cap B} \setminus (A \setminus B) &= X \setminus A \qquad \text{ sowie} \\
			\enbrace*{X \setminus B} \setminus \enbrace*{A \setminus B} &= X  \setminus (A \cup B)
		\end{align*}
		erhalten wir einen Ausschneideisomorphismus 
		\[
			H_*\enbrace[\big]{X \setminus A, X \setminus (A \cup B)} \cong H_* \enbrace*{X \setminus A \cap B, X \setminus B;R}
		\]
		Zusammen mit der Paarfolge aus 1) erhalten wir eine Mayer-Vietoris-Folge:
		\[
			\!\begin{tikzcd}[font=\small,column sep=1.1em]
				\cdots \rar & H_{n+1}(X,X \setminus A \cap B) \rar["\partial"] & H_n(X,X \setminus A \cup B) \rar & 
				{\displaystyle\bigoplus_{H_n(X,X \setminus B)}^{H_n(X,X\setminus A)}}
				\rar & H_n(X,X\setminus A \cap B) \rar & \cdots
			\end{tikzcd}
			\marginnote{\enquote{$R$} aus Platzgründen weggelassen}
		\]
	\end{enumerate}
\end{bemerkung}

\cleardoubleoddemptypage
\pagenumbering{Alph}
\setcounter{page}{1}
\printindex
\listoffigures
% \todototoc
% \listoftodos[To-do's und andere Baustellen]
\end{document}
